\documentclass[11pt,a4paper]{article}
\usepackage{fontspec}
\defaultfontfeatures{Ligatures=TeX}
\setmainfont{lmroman10-regular.otf}[BoldFont=lmroman10-bold.otf,ItalicFont=lmroman10-italic.otf,BoldItalicFont=lmroman10-bolditalic.otf]
\setsansfont{lmsans10-regular.otf}[BoldFont=lmsans10-bold.otf,ItalicFont=lmsans10-oblique.otf,BoldItalicFont=lmsans10-boldoblique.otf]
\setmonofont{lmmono10-regular.otf}[Scale=0.88,Ligatures=TeXOff]
\usepackage{amsmath,amssymb,amsthm,mathtools,mathrsfs}
\DeclareMathSizes{8}{8}{6}{6}
\DeclareMathSizes{9}{9}{7}{6}
\DeclareMathSizes{10}{10}{7}{6}
\usepackage[margin=25mm,headheight=16pt]{geometry}
\usepackage{booktabs,array,longtable,enumitem,xcolor}
\usepackage{fancyhdr}
\usepackage{xurl}
\usepackage[unicode,colorlinks=true,linkcolor=blue!38!black,
 citecolor=blue!38!black,urlcolor=blue!38!black]{hyperref}
\usepackage{bookmark}
\hypersetup{pdftitle={Missing-digit primes in bases six through eleven},
 pdfauthor={Tseng},pdfsubject={Analytic number theory; missing digits; computer-assisted proofs},
 pdfkeywords={missing-digit primes, multiple omitted digits, Harman sieve, Fourier moments}}
\pdfstringdefDisableCommands{\def\leqslant{<=}\def\geqslant{>=}\def\asymp{asymp}\def\sim{sim}}
\newtheoremstyle{paperplain}{8pt}{8pt}{\normalfont}{}{\bfseries}{.}{.6em}{}
\theoremstyle{paperplain}
\newtheorem{theorem}{Theorem}[section]
\newtheorem{lemma}[theorem]{Lemma}
\newtheorem{proposition}[theorem]{Proposition}
\newtheorem{corollary}[theorem]{Corollary}
\newtheorem{definition}[theorem]{Definition}
\newtheorem{remark}[theorem]{Remark}
\numberwithin{equation}{section}
\newcommand{\A}{\mathcal A}
\newcommand{\e}{\mathrm e}
\newcommand{\ind}{\mathbf1}
\newcommand{\dd}{\,\mathrm d}
\setlist[enumerate]{leftmargin=2.2em,itemsep=4pt,topsep=4pt}
\setlist[itemize]{leftmargin=2.2em,itemsep=3pt,topsep=4pt}
\setlength{\parindent}{1.4em}
\setlength{\parskip}{2.5pt}
\setlength{\emergencystretch}{2em}
\makeatletter
\renewcommand*\l@subsection{\@dottedtocline{2}{1.8em}{3.4em}}
\makeatother
\linespread{1.04}
\allowdisplaybreaks[2]
\pagestyle{fancy}
\fancyhf{}
\fancyhead[L]{\small Missing-digit primes in bases six through eleven}
\fancyhead[R]{\small Tseng}
\fancyfoot[C]{\thepage}
\renewcommand{\headrulewidth}{.3pt}
\title{\bfseries Missing-digit primes in bases six through eleven}
\author{Tseng}
\date{}

\begin{document}
\maketitle
\begin{abstract}
Let $A(x)$ count the positive integers below $x$ whose ordinary
base-$b$ expansions use only the allowed digits. When one digit is
forbidden in base $6$, $7$, $8$ or $9$, or two digits are forbidden
in base $10$, the number of primes counted by $A(x)$ is of order
$A(x)/\log x$. When one digit is forbidden in base $10$ or $11$, we
obtain an asymptotic equivalent with the local factor determined by
the allowed units digits. All estimates hold at arbitrary real endpoints.

For the order estimates a common sieve separates the arithmetic of a
digit alphabet from the geometry of the prime factors. Distribution in
residue classes supplies the upper bound. For the lower bound we compare
products of primes when a group of factors lies in a specified size
range; the remaining loss is an integral on ordered simplices that
depends only on this range. It reduces to a single triangle in base
eight. The asymptotic results instead require a common model for the
linear and bilinear sums. The Fourier moment inequalities and the
bounds for the sieve losses have finite arithmetic certificates.
\end{abstract}
\noindent\textbf{Keywords:} missing-digit primes; digital Fourier moments;
Type I/II estimates; Harman sieve; computer-assisted proofs.\\
\textbf{2020 Mathematics Subject Classification:} 11N05, 11N35, 11N36, 11A63.
\setcounter{tocdepth}{1}
\begingroup
\small
\linespread{.88}\selectfont
\setlength{\parskip}{0pt}
\tableofcontents
\endgroup
\setcounter{tocdepth}{1}
\newpage

\section{Introduction}
\label{sec:introduction}

A restriction on every digit removes a power of the integers up to $x$,
but does not favor divisibility by a fixed prime away from the radix.
This suggests a simple prediction: after correcting for the units digit,
the surviving integers should be prime with frequency $1/\log x$.
The correction is elementary. The difficulty is to distinguish primes
from products of several large primes inside a set which has large
Fourier coefficients at many rational frequencies.

For a base $b\geqslant3$ and a digit $d\in\{0,\ldots,b-1\}$, let
\[
 \mathscr A_{b,d}
 =\{n\geqslant1:\text{the ordinary base-}b
                    \text{ expansion of }n\text{ omits }d\},
\]
and write
\[
 A_{b,d}(x)=\#(\mathscr A_{b,d}\cap[1,x)),\qquad
 \pi_{b,d}(x)=\#\{p<x:p\text{ is prime and }p\in\mathscr A_{b,d}\}.
\]
Ordinary expansions have no leading zeros. In the one-digit
notation above, only the specified digit is forbidden; an integer need not use every other digit. In base eleven,
$d=10$ denotes the eleventh digit symbol. Put
\begin{equation}\label{intro:constants}
 c_b=\frac{\log(b-1)}{\log b},\qquad
 h_{b,d}=\varphi(b)-\mathbf1_{(d,b)=1},\qquad
 \kappa_{b,d}=\frac{b\,h_{b,d}}{(b-1)\varphi(b)}.
\end{equation}
Here $h_{b,d}$ counts the allowed last digits coprime to $b$, and
$\varphi$ is Euler's totient function. Logarithms without an indicated
base are natural.

\begin{theorem}[Orders of magnitude in bases six through nine]
\label{main:order}
For $b\in\{6,7,8,9\}$ and every $d\in\{0,\ldots,b-1\}$, one has
\[
 \pi_{b,d}(x)\asymp\frac{A_{b,d}(x)}{\log x}
 \asymp\frac{x^{c_b}}{\log x}
 \qquad(x\to\infty)
\]
through arbitrary real values.
\end{theorem}

\begin{theorem}[Exact asymptotics in bases ten and eleven]
\label{main:exact}
For $b\in\{10,11\}$ and every $d\in\{0,\ldots,b-1\}$,
\begin{equation}\label{main:asymptotic}
 \pi_{b,d}(x)\sim
 \kappa_{b,d}\frac{A_{b,d}(x)}{\log x}
 \qquad(x\to\infty),
\end{equation}
where $x$ tends to infinity through arbitrary real values.
\end{theorem}

For a two-element set \(S\subset\{0,\ldots,9\}\), let \(A_{10,S}(x)\)
count the positive integers below \(x\) whose ordinary decimal expansions
omit both digits in \(S\). Let \(\pi_{10,S}(x)\) count the primes among
these integers.

\begin{theorem}[Two omitted decimal digits]\label{main:two-decimal}
For every two-element \(S\subset\{0,\ldots,9\}\),
\[
 \pi_{10,S}(x)\asymp\frac{A_{10,S}(x)}{\log x}
                 \asymp\frac{x^{\log_{10}8}}{\log x}
 \qquad(x\to\infty).
\]
The comparison constants can be chosen independently of \(S\), and
the limit is through arbitrary real values.
\end{theorem}

The implied constants may be taken uniformly over the alphabets in
each statement. The lower bounds in bases six through nine do not
identify a leading constant. In contrast, the main term in bases ten
and eleven depends on the actual number of admitted integers below $x$.

The local factor in \eqref{intro:constants} compares the proportion
$h_{b,d}/(b-1)$ of allowed unit last digits with the proportion
$\varphi(b)/b$ of units among all integers. The use of $A_{b,d}(x)$
also matters. Although $A_{b,d}(x)\asymp x^{c_b}$, this counting
function has persistent fluctuations. For example, choose an allowed
nonzero digit $a$ and put $C=ab+d$. Every integer in
$[Cb^k,(C+1)b^k)$ begins with the forbidden prefix $a,d$, so
\[
 A_{b,d}(Cb^k)=A_{b,d}((C+1)b^k).
\]
The corresponding powers of the endpoints have the fixed ratio
$((C+1)/C)^{c_b}>1$. Thus the main term in
\eqref{main:asymptotic} cannot in general be replaced by a fixed
constant times $x^{c_b}/\log x$.

\subsection{Relation to previous work}

Distribution in residue classes is an arithmetic property of the digit
set itself \cite{ErdosMauduitSarkozy1998}. It underlies the almost-prime
results of Dartyge and Mauduit \cite{DartygeMauduit2000,DartygeMauduit2001};
even a bound of two on the number of prime factors does not separate
primes from semiprimes. For one omitted decimal digit Maynard
\cite[Theorem~1.1]{Maynard2019} proved the two-sided order at arbitrary
endpoints. His equation~(1.1) gives the predicted decimal multiplier;
the discussion following Theorem~1.2 gives the general multiplier and
records the one-omission asymptotic for $b\ge12$.

For every $b\ge10$ and every excluded digit, Karwatowski
\cite[Theorem~1]{Karwatowski2022} proved two-sided bounds at $X=b^k$.
His 2019 announcement \cite{KarwatowskiDMV2019} also includes the two
outer digits in base nine; his base-nine paper proves the bounds for
these complete words \cite[Theorem~1]{KarwatowskiBaseNine}.
The lower bound in Theorem~\ref{main:order} covers every digit in bases
six through eight and the seven interior digits in base nine. In bases
ten and eleven Theorem~\ref{main:exact} identifies the leading term,
rather than just its order.

Excluding two decimal digits leaves $8^k$ complete words instead of
$9^k$. Maynard's multiple-omission theorem assumes a sufficiently large
radix \cite[Theorem~1.2]{Maynard2019}; in particular it gives no lower
bound for two decimal omissions. Theorem~\ref{main:two-decimal} gives
the expected order, not an asymptotic equivalent, for each of the
$45$ unordered pairs.

The proofs use the digital circle method and the analysis of
exceptional frequency pairs developed by Maynard
\cite[Sections~10 and~13--15]{Maynard2019}. The further analytic and sieve inputs depend on the radix.

For two decimal omissions we compare smooth boxes of actual prime factors. The filters are applied to the full
digital weight before the row correlations are estimated. In base six,
Theorem~\ref{six:rough} gives an independent weighted bound for
rough primes and semiprimes using divisor distribution alone; the
prime lower bound uses the stronger comparison for groups of primes.

\subsection{Strategy of proof}\label{sec:outline}

\paragraph{The comparison required by a sieve.}
Write $t_i=\log p_i/\log X$ for the prime factors of a number of size
$X$. Up to endpoint errors these exponents lie on $\sum t_i=1$.
A bilinear estimate on $[\alpha,\beta]$ evaluates a marked sum when some
proper subset of the exponents has sum in that interval. A comparison
restricted to actual products of primes can evaluate the same sum once
its rough cofactor has been factored, but does not give an estimate for
arbitrary bilinear coefficients. In both cases the arithmetic task is
to establish the comparison; the remaining sieve task is geometric:
which marked regions of the simplex are left, and with what sign?

\paragraph{One sieve, five applications.}
The positive sieve of Theorem~\ref{common:positive-sieve} starts by
marking two primes. It either stops with a positive term, or makes
two more Buchstab steps whose intermediate terms can be compared.
Only at the end do we factor the entire cofactor and ask whether
\emph{any} subproduct is in the comparison window. Distinct markings
of the same integer must be charged separately. For instance, at
the octal window a triple with exponents $.29,.34,.37$ has no
computable subproduct; either of the two pairs containing $.29$
occurs in the expansion. Its flat mass must therefore be counted
twice, not once (Section~\ref{cost:octal}). The arithmetic inputs
and the extent of the resulting loss differ as follows:
\begin{center}
\begin{tabular}{@{}lcccl@{}}\toprule
Radix & Digits omitted & Divisor level $T$ & Sieve window & Upper bound for loss\\\midrule
Six   & one & $.550$ & $[.423,.480]$ & $.986$\\
Seven & one & $.560$ & $[.420,.480]$ & $.934$\\
Eight & one & $.640$ & $[.400,.500]$ & $.186012$\\
Nine  & one & $.633$ & $[.3155,.3691]$ & $.834204$\\
Ten   & two & $.560$ & $[.420,.480]$ & $.934$\\\bottomrule
\end{tabular}
\end{center}
The divisor level gives an upper limit $2\kappa/T$ for $\pi(x)\log x/A(x)$
(Proposition~\ref{common:upper}), where $\kappa$ is the ratio of the
unit-digit density to $\varphi(b)/b$. The last column measures the loss against the prime mass of the
local model; a loss below one guarantees a positive lower bound. In base eight
the unmatched terminal tuples are exactly triples on a triangle, and
the bound is analytic. In base nine the comparison window crosses
$1/3$. A bound on the size of the next marked product, using the two
marks already present, lets the same stopping rule work there too.

\paragraph{How the comparisons are obtained.}
Part~\ref{part:geom} develops the two approaches for bases nine and
eight. In base nine an estimate for frequency pairs near fractions
with a common denominator supplies a Type~II interval; factoring the
rough cofactor turns it into a prime-shape comparison. In base eight
the calculation is made directly on products of primes. Projections
onto dyadic periods separate the range accessible by characters from
the range controlled by short row differences.

Part~\ref{part:cuts} treats the other three families of alphabets. At a cut in the
digit word the smaller periods are compared by characters and Mellin
inversion. For the remaining frequencies a signed row estimate exploits
affine overlaps between lower words. One cut suffices in base seven.
In the decimal and senary arguments different ranges of Fourier
amplitudes require different cuts; the unequal powers of two and
three in a senary period require a joint moment on the two grids.

Part~\ref{part:exact} concerns one omission in base ten or eleven.
Reversing a digit decomposition and bounding an affine correlation
extend the bilinear range below one quarter. The decisive extra
ingredient is a single nonnegative comparison model valid for both
Type~I and Type~II; the asymptotic sieve then yields an equivalent
for every fixed prefix. Lemma~\ref{prelim:endpoints} passes from
these equivalents to all real endpoints.

\clearpage
\part{Digital distribution and the geometry of the sieve}\label{part:framework}
\section{Digital words, local factors, and Fourier products}
\label{sec:preliminaries}

Digital moment estimates are expressed in terms of suffix words of a fixed
length. Their product structure governs the Fourier analysis, while the
allowed unit last digits determine the arithmetic density.

\subsection{Ordinary expansions and fixed-length words}

Fix $b,d$, write $D=\{0,\ldots,b-1\}\setminus\{d\}$, and set
\[
 \mathcal W_k=\left\{\sum_{j=0}^{k-1}\varepsilon_jb^j:
                 \varepsilon_j\in D\right\},\qquad X=b^k,
 \qquad A=(b-1)^k=X^{c_b}.
\]
The $k$ positions here define a suffix word. If $d\ne0$, this set includes
zero and integers with shorter ordinary expansions; if $d=0$, it contains
only words of actual length $k$. For a fixed allowed positive prefix $v$,
the set $v b^k+\mathcal W_k$ always consists of ordinary expansions.

\begin{lemma}[Digital counting]\label{prelim:count}
For fixed $b,d$, one has $A_{b,d}(x)\asymp_b x^{c_b}$. More precisely,
\[
 A_{b,d}(b^k)=
 \begin{cases}
 \displaystyle\frac{(b-1)((b-1)^k-1)}{b-2},&d=0,\\[3pt]
 (b-1)^k-1,&d\ne0.
 \end{cases}
\]
Every interval $I$ of length $Y\geqslant1$ satisfies
\begin{equation}\label{prelim:interval}
 \#(\mathscr A_{b,d}\cap I)\ll_b Y^{c_b}.
\end{equation}
The same estimate holds for fixed-length words and arbitrary translates.
\end{lemma}
\begin{proof}
If $0$ is forbidden, there are $(b-1)^j$ words of actual length $j$; if a
nonzero digit is forbidden, there are $(b-2)(b-1)^{j-1}$ such words.
Summation gives the endpoint formulas, and comparison with adjacent powers
of $b$ proves the first assertion. Choose
$b^{\ell-1}\leqslant Y<b^\ell$. The interval $I$ meets at most two standard
digital cylinders of length $b^\ell$. A cylinder with a positive prefix
contains at most $(b-1)^\ell$ allowed suffixes. The first cylinder, which
contains zero, has $O_b((b-1)^\ell)$ allowed ordinary expansions by the
endpoint formula just proved. Adding these bounds gives the result.
For fixed-length words the same cylinder bound holds when $\ell\leqslant k$;
when $\ell>k$, their total size already gives the assertion. Translation
of the interval proves the translated version.
\end{proof}

\begin{lemma}[The last-digit local density]\label{prelim:local}
Among the allowed words of length $k\geqslant1$, exactly
$h_{b,d}(b-1)^{k-1}$ are coprime to $b$. For ordinary expansions,
\[
 \#\{n<x:n\in\mathscr A_{b,d},\ (n,b)=1\}
   =\frac{h_{b,d}}{b-1}A_{b,d}(x)+O_b(1).
\]
Relative to the unit density $\varphi(b)/b$ among ordinary integers,
the corresponding factor is $\kappa_{b,d}$.
\end{lemma}
\begin{proof}
There are $h_{b,d}$ allowed unit last digits, and all other positions are
free, proving the fixed-length formula. For ordinary expansions, group
integers by the positive prefix $\lfloor n/b\rfloor$. Each complete allowed
group has $b-1$ members and exactly $h_{b,d}$ units. There is at most one
incomplete group at the endpoint; all single-digit integers contribute
$O_b(1)$. Thus the discrepancy is bounded independently of the number of
digit positions. Dividing by the unit density $\varphi(b)/b$ gives the factor.
\end{proof}

\subsection{Fourier products and moments}

Write $\mathrm e(t)=\exp(2\pi i t)$ and
$\|t\|=\min_{n\in\mathbb Z}|t-n|$, and define
\[
 P_{b,d}(z)=\sum_{a\in D}z^a,\qquad
 f_{b,d}(t)=\frac{|P_{b,d}(\mathrm e(t))|}{b-1},
\]
\[
 F_X(t)=\frac1A\left|\sum_{n\in\mathcal W_k}\mathrm e(nt)\right|
       =\prod_{j=0}^{k-1}f_{b,d}(b^jt).
\]
In particular, $0\leqslant F_X\leqslant1$ and
$f_{b,b-1-d}=f_{b,d}$, since
$P_{b,b-1-d}(z)=z^{b-1}P_{b,d}(z^{-1})$. The terminal-digit density is
given by Lemma~\ref{prelim:local}.

For $p>0$, let
\[
 H_p(X;x)=\sum_{a=0}^{X-1}F_X((a+x)/X)^p,
 \qquad
 (T_ph)(x)=\sum_{j=0}^{b-1}f_{b,d}((j+x)/b)^p h((j+x)/b).
\]
Expansion one digit at a time gives the exact identity
\begin{equation}\label{prelim:transfer}
 T_p^k1(x)=H_p(b^k;x).
\end{equation}
Some later certificates use the normalization $\mathcal L_p=b^{-1}T_p$;
their supersolution constants are correspondingly divided by $b$.

\begin{lemma}[Iteration of a positive supersolution]\label{prelim:supersolution}
If a continuous positive $1$-periodic function $h$ satisfies
$T_ph\leqslant\Lambda h$, then for every $k\geqslant0$ and $x\in\mathbb R$,
\[
 H_p(b^k;x)\leqslant\frac{\max h}{\min h}\Lambda^k.
\]
If $\Lambda\leqslant b^g$, then $H_p(X;x)\ll X^g$ and
$\int_0^1F_X(t)^p\,dt\ll X^{g-1}$.
\end{lemma}
\begin{proof}
Iterate using $1\leqslant h/\min h$ and positivity of the operator.
Integrating \eqref{prelim:transfer} over $0\leqslant x\leqslant1$ and
changing variables gives $X\int_0^1F_X(t)^p\,dt$ on the left.
\end{proof}

The following two even moments need no computational certificate:
\begin{equation}\label{prelim:even}
 H_2(X;x)=X^{1-c_b},\qquad
 H_4(X;x)\leqslant
 X^{\log_b(1+(b-1)^{-3})}.
\end{equation}
The first identity is discrete Parseval. For the second, the one-digit
polynomial $f_{b,d}^4$ has nonnegative Fourier coefficients, so its sum over
$b$ equally spaced sample points is largest at $x=0$. At each nontrivial
$b$th root of unity, $|P_{b,d}|=1$, and this maximum is therefore
$1+(b-1)^{-3}$. Iteration of the positive transfer operator proves the bound.

The continuous inequalities needed for the other moment bounds are obtained
from rational values on finite grids. Appendix~\ref{app:continuous}
gives the interpolation and curvature estimates that justify this passage
from nodes to the whole circle.

\subsection{From fixed prefixes to arbitrary endpoints}

A prime asymptotic for each fixed admissible prefix determines the answer
at arbitrary real endpoints. The order of limits in the following argument
requires uniformity only over finitely many prefixes at a time.

\begin{lemma}[Completion from fixed-prefix cylinders]\label{prelim:endpoints}
Suppose that a constant $\kappa>0$ has the following property: for every
fixed admissible positive prefix $v$,
\begin{equation}\label{prelim:prefix-hypothesis}
 \#\{p\text{ prime}:p\in vb^k+\mathcal W_k\}
 =(\kappa+o_v(1))\frac{(b-1)^k}{k\log b}
 \qquad(k\to\infty).
\end{equation}
Then, with ordinary expansions and arbitrary real $x\to\infty$,
\[
 \pi_{b,d}(x)\sim\kappa\,\frac{A_{b,d}(x)}{\log x}.
\]
\end{lemma}
\begin{proof}
Apply \eqref{prelim:prefix-hypothesis} to the finitely many allowed
nonzero one-digit prefixes and sum over suffix lengths. Geometric
summation gives
\begin{equation}\label{prelim:initial-prime-bound}
 \pi_{b,d}(b^s)\ll_{b,d,\kappa}\frac{(b-1)^s}{s}.
\end{equation}
The finitely many small suffix lengths are absorbed in the constant.

Take $b^K\leqslant x<b^{K+1}$, fix $L_0\geqslant2$, and set
$k=K-L_0+1$, $q=b^k$, $V=\lfloor x/q\rfloor<b^{L_0}$.
Partition $[1,x)$ into $[1,q)$, the complete cylinders
$[vq,(v+1)q)$ for $1\leqslant v<V$, and the terminal partial cylinder.
For fixed $L_0$, the hypothesis is uniform over this finite collection
of prefixes. An inadmissible prefix contributes nothing. In every complete
admissible cylinder, $\log x/(k\log b)\to1$ and there are exactly
$(b-1)^k$ allowed integers, so their total discrepancy is
$o_{L_0}((b-1)^K)$.

A terminal admissible partial cylinder has at most $(b-1)^k$ allowed
integers and at most the prime mass of its complete cylinder.
The initial segment is bounded by
\eqref{prelim:initial-prime-bound} and Lemma~\ref{prelim:count}.
Thus
\[
 \bigl|\pi_{b,d}(x)\log x-\kappa A_{b,d}(x)\bigr|
 \leqslant o_{L_0}((b-1)^K)
       +O_{b,d,\kappa}((b-1)^{K-L_0+1}).
\]
The constant in the second term is independent of $L_0$: for each fixed
depth the complete-cylinder prime bound is eventually at most twice
its main term, and $K\geqslant2L_0$ ensures
$\log x/(k\log b)\leqslant2$. Since
$A_{b,d}(x)\geqslant(b-2)(b-1)^{K-1}$, let first $K\to\infty$ and
then $L_0\to\infty$.
\end{proof}


\section{From Fourier mass to divisor distribution}\label{common:divisors}
The first moment alone already gives more than the upper bound for
primes. It distributes the digits among all residue classes of
moduli coprime to the radix up to a fixed power of $X$. This is the part
of the argument which does not distinguish one omitted digit from two.

Let $b$ be fixed and $\mathcal D\subset\{0,\ldots,b-1\}$ contain
two consecutive digits. Write $s=|\mathcal D|$, $c=\log_b s$,
$X=b^k$, $A=s^k$, $L=\log X$; $\mathcal A_X$ denotes its $k$-place
words, and $F_X$ their Fourier magnitude divided by $A$. Write
$u_d$ for the number of allowed unit digits.

\begin{lemma}[Sampling a block of digits]\label{common:sampling}
Let $P$ have integer frequencies in $[-Y,Y]$, $Y\ge1$. For
$r\ge1$, $C\ge0$ fixed and a $\Delta$-spaced set $\Theta$ on the circle,
\begin{equation}\label{arith:sampling}
 \sum_{\theta\in\Theta}\ \sup_{\|v\|\le C/Y}|P(\theta+v)|^r
 \ll_{C,r}(Y+\Delta^{-1})\int_0^1 |P(t)|^r\,dt .
\end{equation}
Suppose $F_Y$ is the modulus of a normalized radix-$b$ word polynomial,
$Y=b^k$. Assume each contiguous subword of scale $D=b^j$, with
positions reindexed from zero, has $r$th integral moment $O(D^{g-1})$.
Then for any $1\le H\le q\le Y$, any $H$ consecutive integers $i$,
and any shift $u$,
\[
 \sum_i\sup_{|v|\le C/q} F_Y((i+u)/q+v)^r\ll_{b,C,r}H^g .
\]
In particular $q$ need not divide $Y$.
\end{lemma}
\begin{proof}
Write $K_M$ for the Fej\'er kernel with Fourier coefficients
$(1-|h|/M)_+$. For $M=\lceil Y\rceil$ the kernel $2K_{2M}-K_M$
reproduces $P$ and is bounded in absolute value by
$B_M(t)\ll M/(1+M\|t\|)^2$. Its majorant has bounded integral,
so H\"older bounds $|P|^r$ by a constant times $|P|^r*B_M$.
Moving a centre by $C/Y$ enlarges the majorant by a bounded factor;
ordering the $\Delta$-separated centres by distance gives a sum of
translated majorants $O_C(M+\Delta^{-1})$. Integrate to obtain
\eqref{arith:sampling}.

For the second assertion choose the largest radix powers $U\le q/H$
and $D\le H$. Retain the $\log_b D$ digits starting at position
$\log_b U$ and bound every other one-digit modulus by one. The
retained polynomial has scale $D$; its sampling centres have step
$U/q\asymp_b 1/D$, and their unwrapped span is less than one.
Split them into two arcs of span at most $1/2$; on each the circular
spacing is at least $U/q$. The independent displacements are at most
$CU/q\ll C/D$. Apply \eqref{arith:sampling} to each arc and use the
assumed integral bound. If $D=1$, there are fewer than $b$ summands.
\end{proof}

\begin{lemma}[Occupancy of digital intervals and residues]\label{common:occupancy}
For an alphabet of size $s$ and $Y=b^k$, set $c=\log_b s$.
An interval of length $H$ meets $O_b((1+H)^c)$ complete words.
Any class modulo $1\le q\le Y$ contains $O_b((Y/q)^c)$ such words.
Both bounds hold after imposing any further digit restrictions.
\end{lemma}
\begin{proof}
For $1\le H\le Y$ cover the interval by a bounded number of aligned
blocks of the least radix-power length at least $H$; only the suffix
is free on each. The remaining lengths are immediate. For a residue
class take the greatest radix power $D\le q$ and fix the prefix
above its $\log_b D$ suffix positions. There are $(Y/D)^c$ choices,
and a suffix in $[0,D)$ belongs to the specified class in at most one way.
\end{proof}


\begin{proposition}[Type I estimate]\label{arith:typeI}
Suppose that $1/2<\sigma<c$ and that
$\int_0^1F_Z(t)\,dt\ll Z^{-\sigma}$ for every power $Z$ of the base.
Define
\[
 D_X(q)=\sup_{I,a}\left|
 \#\{n\in\mathcal A_X\cap I:n\equiv a\pmod q\}
 -\frac{\#(\mathcal A_X\cap I)}q\right|,
\]
where $I$ ranges over real intervals. For fixed $K\ge0$ and
$\epsilon>0$, there is a constant $c_1>0$ such that
\begin{equation}\label{arith:typeI-max}
 \sum_{\substack{q\le Q\\(q,b)=1}}\tau(q)^K D_X(q)
 \ll A\left(e^{-c_1\sqrt{\log X}}+
                      Q^{1+\epsilon}X^{-\sigma}\right).
\end{equation}
In particular, for every fixed $\theta<\sigma$ and $J>0$,
\begin{equation}\label{arith:typeI-unit}
 \sum_{\substack{q\le X^\theta\\(q,b)=1}}\tau(q)^K
 \left|\#\{n\in\mathcal A_X:(n,b)=1,\ q\mid n\}
                -\frac{u_d A}{sq}\right|
 \ll A L^{-J}.
\end{equation}
In \eqref{arith:typeI-max} one may restrict the leading and final
digits independently to nonempty subsets of $\mathcal D$, or take
a convex average of cyclic translates of such a word, using its
own interval mass as the mean.
\end{proposition}
\begin{proof}
Decompose $I\cap[0,X)$ into aligned radix intervals, with at most
$2(b-1)$ intervals of each length $Z=b^j$. On an interval with an
admissible fixed prefix, the free lower digits form a translate of
$\mathcal A_Z$. Thus
\[
 U_X(t)=2b\sum_{Z\le X,\ Z=b^j}Z^cF_Z(t)
\]
simultaneously bounds the absolute values of the Fourier sums over all
intervals.
The reduced fractions $h/q$ with $R\le q<2R$ are separated by at least
$1/(4R^2)$. When $Z\ge R^2$, discard the higher-digit factors and replace
$F_Z$ by the lower-digit word whose length is comparable to $R^2$;
otherwise use length $Z$ directly. By \eqref{arith:sampling},
\[
 \sum_{q\sim R}\ \sum_{\substack{0<h<q\\(h,q)=1}}F_Z(h/q)
 \ll
 \begin{cases}R^{2-2\sigma},&Z\ge R^2,\\
 R^2Z^{-\sigma},&Z<R^2.
 \end{cases}
\]
Put $E_q=q^{-1}\sum_{0<h<q,(h,q)=1}U_X(h/q)$.
Since $c>\sigma$, summing geometrically over $Z$ gives
\begin{equation}\label{arith:primitiveI}
 \sum_{q\sim R}E_q\ll A\{R^{1-2\sigma}+RX^{-\sigma}\}.
\end{equation}

For small denominators we also use contraction along digital orbits.
The alphabet contains two consecutive digits, so there is a constant
$c_2>0$ for which the one-digit factor
$f(t)=s^{-1}|\sum_{a\in\mathcal D}e(at)|$
satisfies $f(t)\le\exp(-c_2\|t\|^2)$.
If $(h,q)=(q,b)=1$, every $b^jh/q$ is nonintegral and is at distance
at least $1/q$ from the nearest integer. When that distance is less
than $1/(2b)$, the next step multiplies it by exactly $b$.
Hence a fixed amount of contraction occurs every $O(\log(2q))$ steps,
and therefore
\[
 F_Z(h/q)\ll\exp\left(-c_3\frac{\log Z}{\log(2q)}\right).
\]
Take $R_0=\exp(t_0\sqrt{\log X})$, with $t_0$ sufficiently small.
The terms with $Z<Xe^{-\sqrt{\log X}}$ in $U_X$ already have total
mass $O(Ae^{-c_4\sqrt{\log X}})$; apply the orbit bound to the remaining
lengths. Decreasing $t_0$ if necessary gives
\[
 \sum_{2\le q\le R_0,(q,b)=1}\tau(q)^K E_q
       \ll A e^{-c_5\sqrt{\log X}}.
\]
Orthogonality and reduction of fractions yield
\[
 D_X(q)\le\sum_{r\mid q,\ r>1}\frac rq E_r.
\]
Interchange the finite nonnegative sums, using
$\tau(rj)\le\tau(r)\tau(j)$ and
$\sum_{j\le T}\tau(j)^K/j\ll(\log(2T))^{C_K}$.
For small $r$ use the preceding bound; for large $r$ use
\eqref{arith:primitiveI} and $\tau(r)^K\ll r^{\epsilon_0}$, where
$0<\epsilon_0<\min(\epsilon/4,(2\sigma-1)/2)$.
Summation over dyadic ranges, with the logarithmic factors absorbed into
the stretched-exponential decay or into $Q^{\epsilon-\epsilon_0}$,
proves \eqref{arith:typeI-max}.

Finally, put $Z=X/b$. Every admissible word coprime to $b$ has a unique
representation $n=bm+a$, where $m\in\mathcal A_Z$, $a\in\mathcal D$ and
$(a,b)=1$. For $(q,b)=1$, the condition $q\mid n$ specifies one residue
class for $m$ modulo $q$ for each final digit.
Thus the error is at most $u_dD_Z(q)$, and the main term is
$u_dZ^c/q=u_dA/(sq)$.
Choose $(1+\epsilon)\theta<\sigma$ in \eqref{arith:typeI-max} to obtain
\eqref{arith:typeI-unit}. The same proof allows an interval restriction,
retaining the actual digital count in that interval as the mean.
\end{proof}

Indeed, a prescribed leading digit restricts the interval; for a
prescribed final digit $a$, pull the interval and the residue class
back under $n=bm+a$. Sum over the bounded number of choices. A
cyclic translation splits an interval into at most two intervals,
with a corresponding shift of the residue on each. The bounds thus
pass to convex averages by the triangle inequality.

\begin{proposition}[An upper bound at every endpoint]\label{common:upper}
Let $\mathcal D$ be an alphabet of size $s\ge2$ in radix $B$, with
$u>0$ unit digits. Write $c=\log_Bs$, $\gamma=Bu/(s\varphi(B))$.
At $X=B^k$ let $f$ be the indicator of the $k$-place words ending in
a unit digit. Suppose that for some $0<T<\min(1,2c)$ and every fixed $J$
one has, with $I$ ranging over subintervals of $[0,X)$,
\[
 \sum_{q\le X^T,\ (q,B)=1}\tau(q)^2
 \sup_I\left|\sum_{n\in I}f(n)(\mathbf1_{q\mid n}-1/q)\right|
 \ll_J X^c(\log X)^{-J}.
\]
For ordinary expansions let $A(x)$ count the
allowed positive integers below $x$ and $\pi(x)$ their primes. Then
\[
 \limsup_{x\to\infty}\frac{\pi(x)\log x}{A(x)}\le \frac{2\gamma}{T}.
\]
For complete $k$-position words the same constant bounds the upper
limit with denominator $s^k$ and $x=B^k$.
\end{proposition}
\begin{proof}
Set $z=X^{T/2}$ and let $\mathcal S$ consist of the squarefree
$m\le z$ coprime to $B$. Selberg's upper sieve
\cite{FriedlanderIwaniec2010} in this case is just a quadratic form.
Put $G=\sum_{m\in\mathcal S}1/\varphi(m)$,
$y_m=\mu(m)/(G\varphi(m))$ and
$\lambda_a=a\sum_{a\mid m,\ m\in\mathcal S}\mu(m/a)y_m$ for $a\in\mathcal S$.
Then $\lambda_1=1$ and $|\lambda_a|\le1$: in $G$ retain the distinct
terms $hv$ for $h\mid a$, $v\in\mathcal S$, $(v,a)=1$, $v\le z/a$,
and use $\sum_{h\mid a}1/\varphi(h)=a/\varphi(a)$.
Expanding $(a,b)=\sum_{m\mid a,b}\varphi(m)$ and inverting gives
\[
 \sum_{a,b\in\mathcal S}\frac{\lambda_a\lambda_b}{[a,b]}
 =\sum_{m\in\mathcal S}\varphi(m)y_m^2=G^{-1}.
\]
Every unit prime above $z$ survives with weight
$(\sum_{a\in\mathcal S,\ a\mid p}\lambda_a)^2=1$.
The other integers have nonnegative weight. Grouping the expansion by
$[a,b]\le X^T$ costs at most $\tau([a,b])^2$ per modulus. Thus for
every interval $I\subset[0,X)$, writing $M(I)=\sum_{I} f$, we have
\[
 \sum_{p\in I} f(p)\le z+M(I)/G+O_J(X^c(\log X)^{-J}).
\]
The lower bound on $G$ needs no asymptotic for squarefree integers.
Grouping all positive integers with radical $m$ gives harmonic mass
$\prod_{p\mid m}(p-1)^{-1}=1/\varphi(m)$. Consequently
$G\ge\sum_{n\le z,(n,B)=1}1/n=(\varphi(B)/B)\log z+O_B(1)$.
The unit mass of any interval of words is $u/s$ times its total
word count, up to $O_B(1)$: partition by deleting the last digit.
This proves the complete-word estimate, including the $O_B(1)$
primes dividing $B$.

For the ordinary count choose $B^{K-1}\le x<B^K$. For length $j\le K$
apply the interval bound at $X=B^j$ on
$[B^{j-1},\min(x,B^j))$, interpreted as empty when necessary.
These intervals count each admissible positive integer once, whether
or not zero is permitted. Lengths $j<K-\sqrt K$ contribute
$O_B(s^{K-\sqrt K})$. On the remaining lengths
$\log(B^j)\sim\log x$ uniformly. The errors sum to
$o(s^K/\log x)$, since $T/2<c$ and $\sum_{j\le K}s^j\ll_s s^K$.
Finally $A(x)\asymp_{B,s}s^K$: at least $(s-1)s^{K-3}$ words of
ordinary length $K-2$ are available, and the geometric sum bounds
all lengths from above. Summing the main terms proves the claim.
\end{proof}


\section{A positive sieve for grouped prime products}
\label{dt:sieve}

Our starting point is Buchstab's identity in Harman's sieve
\cite{Harman2007}; compare Maynard's marked-prime decomposition
\cite[Section~6]{Maynard2019}. We use an even-step stopping rule and
charge the loss only after completing each cofactor to a full tuple.

A comparison for products of primes does not give a comparison for all
bilinear sums. The distinction matters when the digit alphabet is small:
the coefficients of a product of primes have arithmetic structure even
when coefficients bounded by a divisor function need not. We give the
sieve in a form that uses only the comparison available in this situation.
The loss is a function of a real interval, not of the radix or alphabet.

Let $B$ be a fixed radix, $s\ge2$ the size of its digit alphabet and
$u>0$ the number of its allowed digits coprime to $B$. Put
\[
 X=B^k,\quad Q=s^k=X^c,\quad L=\log X,\quad
 \gamma=\frac{Bu}{\varphi(B)s},\quad \mathcal M=\frac{\gamma Q}{L}.
\]
On $[1,X)$ let $f$ indicate the allowed $k$-position words with unit
last digit, and let $f_0(n)=(\gamma Q/X)\mathbf1_{(n,B)=1}$; extend
both by zero. Their masses on this interval coincide, since
$\sum f=uQ/s$ and $B\mid X$. Fix $D>3/c$ and set
$I_X=[XL^{-D},X)$. Write $b=f-f_0$ and $m_g=\sum_{I_X}g$ for
$g=f,f_0$. An interval of length $H\ge1$ meets at most two aligned
blocks of the least radix-power length $B^j\ge H$, each containing at
most $s^j$ words (or at most $Q$ when $j>k$). Consequently
\begin{equation}\label{dt:sieve-mass}
 m_f-m_{f_0}=O(QL^{-Dc}+Q/X).
\end{equation}

\begin{theorem}[The grouped-product sieve]\label{common:positive-sieve}
Let $1/4<\alpha<\beta\le1/2$ and $1-\beta<T<1$, with $1-T/2<c$.
Assume the following two comparisons.
\begin{enumerate}
\item\emph{Divisors.} For each fixed $J,K>0$,
\begin{equation}\label{dt:sieve-type-i}
 \sum_{\substack{d\le X^T\\(d,B)=1}}\tau(d)^K
       (|R_f(d)|+|R_{f_0}(d)|)\ll_{J,K}QL^{-J},
 \qquad R_g(d)=\sum_{\substack{n\in I_X\\d\mid n}}g(n)-m_g/d .
\end{equation}
\item\emph{Products of primes.} For fixed rank and roughness,
the affine comparison \eqref{dt:sieve-flat-comparison} holds whenever
one prescribed group of factors has limiting exponent in a compact
subinterval of $(\alpha,\beta)$. The precise one-sided condition on
affine regions is given immediately below.
\end{enumerate}
Then, for the nonnegative geometric loss in \eqref{dt:sieve-loss},
\[
 \liminf_{k\to\infty}\frac{1}{\mathcal M}\sum_{p\in I_X}f(p)
 \ge 1-\mathscr L(\alpha,\beta).
\]
\end{theorem}

The statement separates arithmetic from combinatorics. In fact the
same prime lower bound holds along any sequence $X\to\infty$, with
$Q=X^c$, $0<c\le1$, and fixed $B,\gamma>0$ ($B$ an integer radix),
for a nonnegative $f=f_X$ uniformly bounded in $X$ and supported on units in $[1,X)$.
In place of the word identities assume that its mass is
$(\gamma\varphi(B)/B+o(1))Q$ and that its mass on each of
$[1,2XL^{-D}]$ and $[X(1-L^{-D}),X)$ is $o(Q/L)$.
Indeed the proof below uses only these mass estimates, the two
comparisons, and boundedness to remove large squares. The conversion
to ordinary digit counts, by contrast, uses the word structure.

The flat sequence satisfies $R_{f_0}(d)=O_B(Q/X)$ for $(d,B)=1$:
divide the multiples by $d$ and count unit classes. The condition
$T<1$ absorbs the resulting errors even with a fixed divisor weight.

\subsection{From smooth boxes to affine prime shapes}
\label{dt:sieve-smooth-shapes}
First work with an interior window, with room for a small enlargement. Fix the
number \(r\) of prime factors and a positive lower exponent
\(\epsilon_0\). Suppose a family of tuples has
\(p_i\ge X^{\epsilon_0}\), product in \(I_X\), and is defined by a
finite Boolean combination of affine conditions in
\(t_i=\log_X p_i\). Suppose a prescribed proper set of factors
forms a group with exponent in a closed subinterval of the buffered window at every point where the closure of the
affine region from $\sum t_i<1$ meets $\sum t_i=1$. This is a
condition on limiting exponent vectors, not just on points satisfying
the strict inequalities on the simplex itself.
Its complementary factors form the other group. Before smoothing, resolve any affine form constant on
the simplex. If its constant is nonzero its sign is eventually fixed
on $I_X$; if zero it is $a(\sum t_i-1)$, whose sign (including the
case $a=0$) is exact since the product is strictly less than $X$.
The other affine walls have simplex measure zero. The comparison required is
\begin{equation}
 \sum_{\boldsymbol p\ {\rm in\ the\ family}} b(p_1\cdots p_r)
       =o(\mathcal M).
 \label{dt:sieve-flat-comparison}
\end{equation}
Finite unions of such statements may be used after assigning each
tuple to its first admissible choice of group.

\begin{proposition}[A sufficient smooth comparison]\label{common:box-transfer}
In place of the second hypothesis of Theorem~\ref{common:positive-sieve},
it suffices to know the following estimate on boxes of primes. Fix
$r,\epsilon>0$ with $r$ a positive integer. On each fixed-multiple box
at scales $P_i\ge X^\epsilon$, let $w$ be smooth with product supported
in $I_X$ and with prescribed polylogarithmic bounds on derivatives in
the logarithmic variables. Write $D_w(v)\,dv$ for the image of
$w(\boldsymbol y)\prod_i dy_i/\log y_i$ under multiplication.
Whenever a proper group of scales has product $X^m$, assume
\[
 \sum_{\boldsymbol p} w(\boldsymbol p) f(\textstyle\prod p_i)
 =\frac B{\varphi(B)}\sum_n f(n)D_w(n)+O_J(QL^{-J})
\]
uniformly for $m$ in each compact subinterval of $(\alpha,\beta)$,
and for every fixed $J$. Then the affine comparison follows.
\end{proposition}
\begin{proof}
First take a fixed smooth weight \(W(t_1,\ldots,t_r)\), with each
exponent bounded below by a positive constant and with the group
within the buffered range on a neighborhood of the simplex.
Choose a smooth product cutoff \(\chi_X\), supported on \(I_X\)
and equal to one on
\([2XL^{-D},X(1-L^{-D})]\), with derivatives in \(\log v\)
bounded by powers of \(L\). The endpoint losses for \(f\) are
\(o(\mathcal M)\) by the interval word bound. The losses for \(f_0\)
are \(o(\mathcal M)\) by counting integers in the two endpoint
intervals. In both estimates the number of tuples representing
an integer is bounded in terms of the fixed positive roughness:
such an integer has only a bounded number of prime factors with
multiplicity. In particular this smoothing does not use a
uniform divisor bound by a fixed power of \(L\).

A smooth multiplicative partition of unity decomposes the
weighted tuple sum into \(O(L^r)\) fixed-multiple boxes. On any
box meeting the support of the test, the exponents of the
scales differ from the actual exponents by \(O(1/L)\) and
the product of the scales is \(X L^{O(1)}\). The buffer
permits the application of the prime-box hypothesis in every
such box. The derivative bounds of the tests are polylogarithmic
in the coordinate scaled to each box. Summing its errors with a
sufficiently large logarithmic saving gives \(o(\mathcal M)\).
The continuous densities add by linearity. On writing \(v\) for
the product, their sum is
\[
 \mathfrak D_X(v)=\frac{\chi_X(v)}{L}
      \int W(t_1,\ldots,t_{r-1},t_r)
          \prod_{i=1}^r\frac1{t_i}\,dt_1\cdots dt_{r-1},
 \qquad t_r=\log_Xv-\sum_{i<r}t_i .
\]
This follows by setting the last prime variable equal to
\(v/\prod_{i<r}v_i\) and then taking logarithms in the first
\(r-1\) variables. The integral is continuous in \(\log_X v\)
near \(1\), and \(\log_X v=1+O_D(\log L/L)\) on \(I_X\). If
\(C_W\) is its value at \(1\), the box comparison therefore gives
\[
 \frac{B}{\varphi(B)}\sum_{v\in\mathbb N} f(v)\mathfrak D_X(v)
       =\mathcal M C_W+o(\mathcal M).
\]
On the flat side the prime number theorem in the same smooth
boxes gives
\[
 \sum_{\boldsymbol p}f_0(p_1\cdots p_r)
      W(\boldsymbol t)\chi_X(p_1\cdots p_r)
       = \frac{\gamma Q}{X}\int_0^\infty\mathfrak D_X(v)\,dv
           +o(\mathcal M)
       =\mathcal M C_W+o(\mathcal M).
\]
All primes in the box are coprime to $B$ for large \(X\).
For the first equality, replace the prime sum in one variable
at a time by its integral. Partial summation with the prime
number theorem gives a relative error bounded by a fixed power
of \(L\) times \(\exp(-c_0\sqrt L)\) against the product of the
box scales. The scales have product \(O(X)\), and the number
of boxes is polylogarithmic.

For a sharp affine family, sandwich its indicator between smooth
nonnegative weights whose difference is supported near the
affine walls. The supports can be kept in the buffered group
range. This follows from the fixed margin on the simplex,
compactness on the positive exponent region, and the estimate
\(\log_X v=1+o(1)\) on the product support. The limiting
\(\prod_i t_i^{-1}\) density is bounded on this region, and
the walls have measure zero. Apply the smooth estimate
to the upper and lower weights, first let \(X\) tend to
infinity, and then shrink the neighborhoods of the walls.
Positivity of \(f,f_0\) gives \eqref{dt:sieve-flat-comparison}.
The same prime-number-theorem calculation and sandwich yield
the flat asymptotic for every such affine family without a
group hypothesis. In particular the flat asymptotic for ordered
tuples has density
\begin{equation}
 \mathcal M\,
 \frac{dt_1\cdots dt_{r-1}}{\prod_{i=1}^r t_i},
 \qquad \sum_i t_i=1,\quad t_1\le\cdots\le t_r .
 \label{dt:sieve-flat-density}
\end{equation}
The coordinate measure here is \(dt_1\cdots dt_{r-1}\).
\end{proof}

\subsection{Comparison of marked rough sums}
\label{dt:sieve-rough}

Set
\[
 \begin{gathered}
 \eta=\beta-\alpha,\qquad z_*=X^\eta ,\\
 d_*=1-\beta,\qquad c_*=1-\alpha .
 \end{gathered}
\]
For \(P^-(1)=+\infty\) write
\[
 S_e(z)=\sum_{\substack{em\in I_X\\P^-(m)>z}} b(em).
\]
For a bounded number of distinct ordered prime marks comprising
\(e\), each greater than \(z_*\), and any fixed affine mark family
with \(e\le X^{c_*}\), we claim
\begin{equation}
             \sum_e S_e(z_*)=o(\mathcal M).
 \label{dt:sieve-rough-comparison}
\end{equation}
The sum is over marked tuples and includes the empty-tuple case.
For \(X^{d_*}\le e\le X^{c_*}\), factor the rough cofactor into
its uniquely ordered prime tuple. It is nontrivial since
\(e<X L^{-D}\) for large \(X\), and the resulting tuple has
bounded rank. The cofactor, complementary to \(e\), has limiting
exponent in \([\alpha,\beta]\), so \eqref{dt:sieve-flat-comparison}
applies. The lower bound on each factor and the order restrictions
are affine conditions of the required kind.

For smaller \(e\), fix \(h>0\) with \(d_*+h<T\); take a
fixed \(0<\epsilon<\min(h,\eta)\) and put \(z_0=X^\epsilon\).
Use Buchstab's exact identity for any \(z\ge z_0\),
\begin{equation}
 S_e(z)=S_e(z_0)-
    \sum_{z_0<q\le z}
      \sum_{\substack{eqm\in I_X\\P^-(m)\ge q}} b(eqm),
 \label{dt:sieve-buchstab-exact}
\end{equation}
where \(q\) runs over primes. For an inner threshold \(P^-(m)\ge q\)
the next inserted prime is strictly smaller than \(q\).
Expand only while the current marked product is below
\(X^{d_*}\), and stop at its first crossing of this threshold.
Each newly inserted prime is at most \(X^\eta\), hence a
crossing lies in \([X^{d_*},X^{c_*}]\). Each terminal term
compares by factoring its nontrivial cofactor, whose least
prime exceeds \(z_0\). Repeated prime factors, if present at
this stage, may be removed as described below. All paths have
bounded depth at fixed \(\epsilon\). At a nonterminal term
the contribution with cutoff \(z_0\) has its marks above
\(z_0\). For any fixed value of the current product, the
original marks are its prime factors above \(z_*\); the
inserted ones occur in uniquely decreasing order. Thus
the multiplicity at the cutoff \(z_0\) is bounded by the
factorial of the number of original marks, independently of
\(\epsilon\).

To bound these terms, let
\(\Pi(z_0)=\prod_{p\le z_0,\ p\nmid B}p\).
The fundamental lemma of the dimension-one sieve
\cite{Harman2007}, with density \(1/p\), gives upper and lower
weights \(\lambda_d^\pm\) of absolute value at most one, on
squarefree \(d\mid\Pi(z_0)\), \(d\le X^h\). They bound the
sifted indicator by their divisor sums and satisfy
\[
 \sum_d\frac{\lambda_d^\pm}{d}
    =V(z_0)\bigl(1+O(e^{-c_1h/\epsilon})\bigr),\qquad
 V(z_0)=\prod_{\substack{p\le z_0\\p\nmid B}}(1-1/p).
\]
Apply the sieve inequalities separately to \(f\) and \(f_0\).
Both sequences are supported on units modulo $B$, and a
marked product with \(P^-(e)>z_0\) is coprime to every sieve
divisor. Thus the main term uses \(m_g/(ed)\), and
\[
 |S_e(z_0)|\ll \frac{V(z_0)}{e}
      \left(|m_f-m_{f_0}|+(m_f+m_{f_0})e^{-c_1h/\epsilon}\right)
     +\sum_{\substack{d\mid \Pi(z_0)\\d\le X^h}}
         (|R_f(ed)|+|R_{f_0}(ed)|).
\]
In summing over \(e<X^{d_*}\) with \(P^-(e)>z_0\), the
factorization of \(ed\) into its large and small prime parts
is unique. The combined modulus is at most \(X^{d_*+h}\),
so \eqref{dt:sieve-type-i} bounds the total remainder.
Mertens' theorem gives \(V(z_0)\ll(\epsilon L)^{-1}\) and
\[
 \sum_{\substack{e<X^{d_*}\\P^-(e)>z_0}}\frac1e
   \le\prod_{X^\epsilon<p\le X^{d_*}}(1-1/p)^{-1}
   \ll\epsilon^{-1}.
\]
Since $|m_f-m_{f_0}|=o(Q)$ by the mass hypothesis (for words, \eqref{dt:sieve-mass}), the normalized total of the
nonterminal terms is
\(O(\epsilon^{-2}e^{-c_1h/\epsilon})+o_\epsilon(1)\).
The terminal terms give \(o_\epsilon(1)\) after
normalization by \(\mathcal M\). Letting \(X\) tend to
infinity and then \(\epsilon\) decrease to zero proves
\eqref{dt:sieve-rough-comparison}.

\subsection{Positive leaves of the Buchstab expansion}
\label{dt:sieve-leaves}

For every fixed \(v>0\), the \(X^v\)-rough integers having a square
factor have combined \(f+f_0\) mass \(o(\mathcal M)\). For a
prime square \(p^2\) with \(X^v<p\le X^{T/2}\), use
\eqref{dt:sieve-type-i} at modulus \(p^2\) and sum \(O(Q/p^2)\).
For larger \(p\) the total integer count is
\(O(X\sum_{j>X^{T/2}}j^{-2})=O(X^{1-T/2})=o(Q/L)\).
The same bounds hold in any fixed-rank marked sum, since
each rough integer has bounded prime-factor multiplicity and
admits only a bounded number of markings. This justifies
removing squareful products in the terminal uses of
\eqref{dt:sieve-flat-comparison} above, at fixed \(\epsilon\).
It also permits using \eqref{dt:sieve-rough-comparison} for
squarefree marked sums at the fixed cutoff \(z_*\).
In the remainder of the expansion, write \(S_e^{\rm sf}\)
for \(S_e\) restricted to squarefree products.

For large \(X\), the interval \(I_X\) lies above \(\sqrt X\).
Every integer in \(I_X\) with least prime greater than
\(\sqrt X\) is prime. The squarefree version of
\eqref{dt:sieve-buchstab-exact} gives
\begin{equation}
 \sum_{p\in I_X}b(p)
     =S_1^{\rm sf}(z_*)-
       \sum_{z_*<p\le\sqrt X}S_p^{\rm sf}(p).
 \label{dt:sieve-first-step}
\end{equation}
Here and below a least-factor condition \(\ge p\) on the
cofactor becomes \(>p\) when \(p\) is already a mark,
because the product is squarefree. The first term compares
by \eqref{dt:sieve-rough-comparison}.

We may replace the threshold \(p\) in the second term by
\(p'=\min(p,\sqrt{X/p})\) at cost \(o(\mathcal M)\).
Indeed, a newly admitted cofactor has every prime factor
greater than \(\sqrt{X/p}\) and is at most \(X/p\);
it must be a prime at most \(p\). For any fixed
\(\delta>0\), this is impossible on \(I_X\) if
\(p\le X^{1/2-\delta}\). For the remaining \(p\) with
\(X^{1/2-\delta}<p\le\sqrt X\), the total of the two
nonnegative sources on the marked \(z_*\)-rough integers
is \(O_\eta(\delta\mathcal M)+o(\mathcal M)\).
This follows from \eqref{dt:sieve-rough-comparison} for the
single mark \(p\) and from the reference density
\eqref{dt:sieve-flat-density}. The flat strip has
bounded density at fixed \(\eta\) and width \(\delta\).
Letting \(\delta\) decrease to zero proves the replacement.

Expanding \(-S_p^{\rm sf}(p')\) down to \(z_*\) removes a
compared term and produces the positive states
\begin{equation}
 S_{pq}^{\rm sf}(q),\qquad
 z_*<q<p,\quad p\le\sqrt X,\quad q\le\sqrt{X/p}.
 \label{dt:sieve-initial-states}
\end{equation}
For a positive state \(S_e^{\rm sf}(z)\), let
\(s=\log_X e\) and \(\zeta=\log_X z\). The threshold \(z\)
is the least marked prime. If \(s+\zeta<c_*\), apply
Buchstab's identity twice down to \(z_*\):
\begin{equation}
 S_e^{\rm sf}(z)
   =S_e^{\rm sf}(z_*)-
        \sum_{z_*<v<z} S_{ev}^{\rm sf}(z_*)+
        \sum_{z_*<w<v<z} S_{evw}^{\rm sf}(w).
 \label{dt:sieve-two-step}
\end{equation}
The first two terms compare since their marked products
have exponent less than \(c_*\). The last terms are
positive states. If \(s+\zeta\ge c_*\), retain the state
as a leaf. The expansions have bounded depth because each
appended prime exceeds \(z_*\). Initially \(s\le3/4\) and
\(\zeta\le1/3\) by \eqref{dt:sieve-initial-states}. An expandable state contains at least two marks, so $2\zeta\le s$ and
$s+\zeta<c_*$ imply $\zeta<c_*/3$. After two insertions the new
marked exponent is less than $c_*+\zeta<4c_*/3<1$, since $\alpha>1/4$. Thus a leaf
that meets the product support has a nontrivial cofactor.

Factor each leaf's cofactor into its increasing prime tuple
and keep all markings with their multiplicities. A completed
leaf is called matched if a proper set of its factors has
exponent in \([\alpha,\beta]\). For each matched tuple,
choose the first such subset in a fixed order. The marking,
the successive decisions of the expansion, the ordering
conditions, and this subset choice are finite affine conditions
at each fixed rank. Order walls and the tests on a proper
mark set or proper subset give measure zero on \(\sum_i t_i=1\).
The sum over matched leaves is
\(o(\mathcal M)\) by \eqref{dt:sieve-flat-comparison}.
The remaining leaf terms retain a positive sign in
\eqref{dt:sieve-first-step}. Their \(f\)-contributions
are nonnegative, so their \(b\)-contribution is bounded
below by minus their flat mass.

Let \(\ell(\boldsymbol t)\) count the terminal markings
of an ordered tuple \(\eta\le t_1\le\cdots\le t_r\),
\(\sum_i t_i=1\), according to the inequalities in
\eqref{dt:sieve-initial-states} and \eqref{dt:sieve-two-step}.
Its value on the boundary walls is immaterial. Define
\begin{equation}
 \mathscr L
  =\sum_{3\le r\le 1/\eta}
     \int_{\substack{\eta\le t_1\le\cdots\le t_r,\
             \sum_i t_i=1\\\text{no proper subset sum in }[\alpha,\beta]}}
      \frac{\ell(\boldsymbol t)}{\prod_{i=1}^r t_i}
          \,dt_1\cdots dt_{r-1}.
 \label{dt:sieve-loss}
\end{equation}
The rank bound follows from $r\eta\le1$.
The flat asymptotic \eqref{dt:sieve-flat-density} applies
to each marking, with the finite-\(X\) exponent sum
\(\log_X n=1+o(1)\). In particular the sum over unmatched
leaves has flat mass \((\mathscr L+o(1))\mathcal M\).
The prime number theorem gives
\(\sum_{p\in I_X}f_0(p)=(1+o(1))\mathcal M\), and hence
\begin{equation}
     \sum_{p\in I_X} f(p)
           \ge (1-\mathscr L+o(1))\mathcal M .
 \label{dt:sieve-lower}
\end{equation}


This proves Theorem~\ref{common:positive-sieve} in the buffered case.
To remove the buffer, approach $[\alpha,\beta]$ from within.
Locally the rank is bounded and all exponents are bounded away from
zero. The subset and marking tests are locally constant off finitely
many affine walls of measure zero: the subset is proper, a stopping
test counts the least mark twice and an unmarked coordinate zero
times, and the initial and ordering tests also have nonconstant
coefficient vectors. At a rank transition $r\eta=1$ the allowed
ordered simplex collapses to a point, null for $r\ge3$. Dominated convergence in
\eqref{dt:sieve-loss} gives continuity of the loss; take the limit
after the lower limit in $k$. If $B^k\le x<B^{k+1}$, each prime so
counted has an allowed ordinary expansion, and $s^k>x^c/s$.
In particular the lower limit for $\pi(x)\log x/x^c$ is at least
$(\gamma/s)(1-\mathscr L)$ whenever the loss is less than one.


\clearpage
\section{The cost of an unmatched product}\label{dt:sieve-certificate}
The arithmetic information has now been reduced to a pair of
endpoints. Four windows occur in this paper, rather than five:
base seven and the two decimal omissions use exactly the same
measure, including its marking multiplicities.

\subsection{A triangular loss with two markings}\label{cost:octal}
At $[\alpha,\beta]=[2/5,1/2]$ put $\mathscr L_8=\mathscr L(2/5,1/2)$.
The roughness threshold is $1/10$, and the
stopping test is $s+\zeta\ge3/5$ for marked sum $s$ and least mark
$\zeta$. Outside the null boundary walls, complementation shows that an unmatched tuple
has no proper subset sum in $[2/5,3/5]$. An exponent above $3/5$
cannot be marked, since the initial larger mark is at most $1/2$ and
marks subsequently decrease. This forces $s<2/5$ and, since there
are at least two marks, $\zeta\le s/2$; thus no such leaf stops.


Otherwise every coordinate is below $2/5$. Call one large if it exceeds
$1/5$. If there are at most two large coordinates, start with one
if present, omit the other if present, and add small coordinates until
first reaching $2/5$. The available mass exceeds $3/5$ and the
overshoot is at most $1/5$, contradicting nonmatching. Hence there
are at least three large coordinates. Each pair of these sums to more
than $3/5$, so any three sum to more than $9/10$, leaving no room for
a fourth rough coordinate. 

Conversely, triples with all coordinates
below $2/5$ are rough and unmatched. For $t_1<t_2<t_3$ just the
pairs $\{1,2\},\{1,3\}$ occur: $t_j<1/2$, $t_j+2t_1<1$, their
cofactor exceeds $t_1$, and both stop already at $s>3/5$. The pair
$\{2,3\}$ fails the cofactor test. This counts the positive markings,
not merely the ordered triples.

With $z=1-x-y$, the exact loss is therefore $2/3!$ times the integral
of $1/(xyz)$ over the triangle $x,y,z\le2/5$. Symmetry and
$1/(xyz)=1/(xy)+1/(xz)+1/(yz)$ reduce this to the integral of
$1/(xy)$ there. Put $x=(2-u)/5$, $y=(2-v)/5$. Then on
$u,v\ge0$, $u+v\le1$ the integrand with Jacobian is
$1/((2-u)(2-v))$. Expanding both geometric series gives an exact expression:
\[
 \mathscr L_8=\sum_{h\ge0}
 \frac{\sum_{i=0}^h i!(h-i)!}{2^{h+2}(h+2)!}.
\]
Indeed $\int u^iv^{h-i}\,du\,dv=i!(h-i)!/(h+2)!$ on the unit
triangle. Before integration the numerator of the homogeneous term
is at most $(u+v)^h\le1$, so truncation at $h=N$ costs at most
$2^{-N-3}$. At $N=30$ the rational sum gives
$.18601165999<\mathscr L_8<.18601166012$.


\subsection{An ordered-simplex bound}
Fix a rank $r$, a mesh $N$ with $N\eta\ge1$, and put $j_i=\lfloor Nt_i\rfloor$
in increasing order. Only $j_1\ge\lfloor N\eta\rfloor$ and
$1\le H=N-\sum j_i<r$ can contribute positive volume.
For a subset \(J\) of cardinality \(a\), write \(s_J=\sum_{i\in J}j_i\).
Its scaled exponent sum lies in
\begin{equation}
    \left[s_J+\max(0,H-r+a),\ s_J+\min(H,a)\right].
 \label{dt:sieve-subset-cell}
\end{equation}
Discard the cell if this interval lies in
\([N\alpha,N\beta]\) for some proper subset.

For a remaining cell and a proper even mark set \(J\),
put \(i=\min J\). Bounds on its scaled stopping test
\(N(s+\zeta)\) are
\[
  \underline u_J=s_J+j_i+\max(0,H-r+|J|),\qquad
  \overline u_J=s_J+j_i+\min(H,|J|)+1 .
\]
Define \(U_J\) recursively to be the maximum of the two
numbers below:
\[
 \begin{aligned}
 &1_{\{i=1,\ \overline u_J\ge Nc_*\}},\\
 &1_{\{\underline u_J<Nc_*,\ |J|+2<r\}}
       \sum_{1\le v<w<i} U_{J\cup\{v,w\}} .
 \end{aligned}
\]
The second number is defined to be zero if its indicator
vanishes. The first number bounds stopping and the second
bounds expanding. An actual state does one or the other
at each tuple, so the maximum is an upper bound even when
the cell straddles the stopping wall. The value of
\(\ell\) throughout the cell is at most the sum of
\(U_{\{i,j\}}\) over \(1\le i<j\le r\) with \(2j_j<N\).
This initial test drops \(t_j+2t_i\le1\), thereby
enlarging the upper bound. The possible equality at \(2t_j=1\)
has zero volume.

The coordinate slice volume of the cell before imposing
order is
\[
 \frac{N^{1-r}}{(r-1)!}\operatorname{Eul}(r-1,H-1),
\]
where \(\operatorname{Eul}\) denotes the Eulerian numbers. In the first
\(r-1\) fractional coordinates, the sum belongs to
\([H-1,H]\), giving this cube-slice formula by
inclusion-exclusion. Imposing order divides the volume by
\(\prod_b m_b!\), where \(m_b\) are the multiplicities of
the equal \(j_i\) blocks. The slice measure is invariant
under coordinate permutations, and orders only within
equal-index blocks remain to be imposed. Since
\(\prod_i t_i^{-1}\le N^r/\prod_i j_i\), the cell's
integral contribution is bounded by its marking upper
count times
\[
  \frac{N\operatorname{Eul}(r-1,H-1)}
       {(r-1)!\prod_i j_i\,\prod_b m_b!}.
\]
Round this last factor upward to a multiple of \(10^{-15}\)
before multiplying by the marking upper count. All tests
and rounds are integer operations, including the Eulerian
recurrence
\[
 \operatorname{Eul}(n,h)
  =(n-h)\operatorname{Eul}(n-1,h-1)
   +(h+1)\operatorname{Eul}(n-1,h),
 \qquad \operatorname{Eul}(0,0)=1 .
\]
with out-of-range entries zero.

\subsection{Three directed upper sums}
\begin{proposition}[Loss outside the computable factor ranges]\label{nine:certificate}
The losses at $[.420,.480]$, $[.423,.480]$, and $[.3155,.3691]$
are respectively less than $.933364$, at most $.985573022923228$,
and less than $.834204$.
\end{proposition}
\begin{proof}
The following are integer upper sums in units of $10^{-15}$.
Each row adds the cell bounds just proved over the indicated ranks;
the individual rank totals are checked by the accompanying program
\path{supplement/common_sieve/verify.py}.
\[
\begin{array}{c@{\qquad}r@{\qquad}c@{\qquad}r}
\hline
[\alpha,\beta]&N&r&\text{upper numerator}\\ \hline
{[.420,.480]}&640&3\text{--}9&922064374609841\\
{[.423,.480]}&1000&3,4&946517051229487\\
 &400&5,6&38812654325222\\
 &200&7\text{--}9&243305312404\\
 &160&10\text{--}17&12056115\\
{[.3155,.3691]}&512&3\text{--}12&834189512933385\\ \hline
\end{array}
\]
To see how to check the table, start with the empty prefix of bin
indices, appending nondecreasing $j_i\ge\lfloor N\eta\rfloor$
so that $1\le H<r$ remains possible. At each step store the
attainable sums of $a$ indices: the new sums are the old ones and
the old $(a-1)$-sums translated by the index just added. If an
$a$-sum $v$ already satisfies $[v,v+a]\subset[N\alpha,N\beta]$
with $0<a<r$, every completion is matched. For each surviving full
cell apply \eqref{dt:sieve-subset-cell} and the marking recursion.
These operations, the Eulerian recurrence and upward rounding are
implemented in \path{supplement/common_sieve/simplex.cpp}.

For the analytic tail note that a terminal marking is an even subset
containing the least prime; reading the marks in decreasing order
recovers its history. There are at most $2^{r-2}$ such subsets.
Writing $t_i=\eta+v_i$ shows that $\prod_i t_i$ is at least
$\eta^r+\eta^{r-1}\sum_i v_i=\eta^{r-1}(1-(r-1)\eta)$.
The ordered simplex has coordinate volume
$(1-r\eta)^{r-1}/(r!(r-1)!)$. At $\eta=.060$ the remaining
ranks therefore cost at most
\begin{equation}\label{dt:sieve-tail}
 \sum_{r=10}^{16}
 \frac{2^{r-2}(1-r\eta)^{r-1}}{r!(r-1)!\eta^{r-1}(1-(r-1)\eta)}
 <11298778668719\cdot10^{-15}.
\end{equation}
This gives the sharper total $\mathscr L<.933363153278560$.
At $\eta=.057$ the four rows exhaust the ranks and sum to
$985573022923228$. At $\eta=.0536$ even the weaker product bound
$\eta^r$ gives a contribution less than $.000013819$ for ranks
$13$ through $18$, the last possible rank. Adding the finite sum
proves the last assertion.
\end{proof}

\clearpage
\part{Two routes to comparison: bases nine and eight}\label{part:geom}
For nonary words the Fourier argument controls a bilinear sum before
the coefficients are specialized to prime factors. For octal words
the order is reversed: keeping the prime factors visible makes
character estimates available. Both arguments end at the affine
comparison of Section~\ref{dt:sieve-smooth-shapes}.

\section{A weighted estimate for common denominators}\label{geo:section}

A large bilinear kernel can force two frequencies close to rationals with
the same denominator \cite[Lemmas~13.2 and~14.1]{Maynard2019}. Squaring
the Fourier mass of a denominator row counts exactly the pairs in that
row. The problem is to sum these squares without a power loss at small
denominators; there even a small factor $X^\varepsilon$ would exceed
the logarithmic saving sought in Theorem~\ref{geo:typeii}.

Fix a prime-power base $b=p^r$ and a forbidden digit $d$, and put
\[
 \mathcal D=\{0,\ldots,b-1\}\setminus\{d\},\qquad
 c=c_b=\frac{\log(b-1)}{\log b},\qquad Z=b^j.
\]
Let $\mathcal W(Z)$ be the set of all allowed words of length $j$, and write
$P_Z(t)=Z^{-c}\sum_{n\in\mathcal W(Z)}\mathrm e(nt)$ and $F_Z=|P_Z|$.
Thus $0\le F_Z\le1$, and for powers $U,V$ of the base,
\begin{equation}\label{geo:product}
 F_{UV}(t)=F_U(t)F_V(Ut),\qquad
 \int_0^1 F_U(t)^2\,dt=U^{-c}.
\end{equation}
For the moment, suppose that the following estimates hold uniformly over all
powers of the base and all real shifts $x$:
\begin{equation}\label{geo:moments}
 \sum_{a<Z}F_Z\left(\frac{a+x}{Z}\right)\ll Z^{1-\sigma},
 \qquad \int_0^1F_Z(t)\,dt\ll Z^{-\sigma},
 \qquad \sum_{a<Z}F_Z(a/Z)^{4/3}\ll Z^{g_{4/3}},
\end{equation}
where $1/2<\sigma<c$. Proposition~\ref{nine:fourier} supplies these
moment bounds for the base-nine application.

There are two complementary scales. When the rational windows are sparse,
direct row estimates retain their decay in the denominator and window width.
When they are dense, four-cycles control the off-diagonal part of an incidence
matrix, producing the $4/3$ moment in the weighted estimate. Keeping these
scales separate is what allows the final bound to save arbitrary powers of
the logarithm even for short denominators. All implicit constants may depend
on the fixed base, alphabet, moment exponents, and stated fixed parameters,
but never on the word length or the denominators being summed.

\subsection{The weighted energy estimate}

For a fixed denominator, the row mass below is the total frequency weight
near its reduced rational points. The bilinear argument requires the sum
of the squares of these masses, so pairs sharing many denominators must
be counted with their full weight.

For $D\ge1,\Delta\ge1$ and fixed $C>0$, take integer rows
$D\le q<2D$ and columns $a\in\mathbb Z/Z\mathbb Z$, and define
\begin{equation}\label{geo:incidence-matrix}
 M_{q,a}=\mathbf1_{\{\exists h\pmod q:\ (h,q)=1,
       \ \|a/Z-h/q\|\le C\Delta/Z\}},\qquad D\le q<2D.
\end{equation}
The row and column sets may be restricted arbitrarily in what follows.
Write $H=D^2\Delta/Z$.
We shall first treat the range
\begin{equation}\label{geo:det-hyp}
 H\ge1,\qquad 4CD\Delta<Z,\qquad
 48C^2D^3\Delta^2<Z^2.
\end{equation}

\begin{proposition}[Weighted common-denominator second moment]\label{geo:energy}
Assume~\eqref{geo:det-hyp}, let $0\le w_a\le1$, and put
\[
 H_u=\sum_aw_a^u,\qquad S_1=\max_q\sum_aw_aM_{q,a}.
\]
Then, for every $\varepsilon>0$,
\begin{equation}\label{geo:energy-bound}
 \sum_q\left(\sum_aw_aM_{q,a}\right)^2
 \ll_\varepsilon Z^\varepsilon\left\{
 S_1H_1+\frac{\Delta^{3/2}D^2}{Z}H_{4/3}^{3/2}
 +\frac{\Delta^2D^3}{Z^2}H_1^2\right\}+O(DZ^{-18}).
\end{equation}
\end{proposition}

The first term is the diagonal contribution of the row Gram matrix.
The other two terms control its off-diagonal part through four-cycles. For a set
of $R$ columns, the first four-cycle contribution has size proportional
to $R^{3/2}$. A weight layer of height $\mu$ therefore contributes
\[
 \mu^2R^{3/2}=(\mu^{4/3}R)^{3/2}.
\]
This is the source of the $4/3$ moment in~\eqref{geo:energy-bound}.
We prove the estimate by counting short affine relations and then showing
that many shared denominators force an especially short relation.

\subsection{Short affine relations}

\begin{lemma}[Counting pairs admitting an affine relation]\label{geo:incidence}
Let $E\subset\{0,\ldots,Z-1\}$, $R=|E|$, and $1\le V\le W$.
Let $I(E;V,W)$ count the ordered pairs $(a,a')\in E^2$ for which there are
integers $v_i$ satisfying
\[
 v_1a+v_2a'+v_3Z+v_4=0,\qquad (v_1,v_2,v_3)\ne(0,0,0),
\]
with $|v_1|,|v_2|,|v_3|\le V$ and $|v_4|\le W$.
Then, for every $\varepsilon>0$,
\begin{equation}\label{geo:incidence-bound}
 I(E;V,W)\ll_\varepsilon
 Z^\varepsilon\left(RVW+\frac{R^2V^2W}{Z}\right).
\end{equation}
\end{lemma}

\begin{proof}
If $V^2W\ge Z$, the second term already bounds $R^2$, so we may assume that
$V^2W<Z$. In particular, $W<Z$. If $v_1=0$ and $v_2\ne0$, then
$v_2a'=-v_3Z-v_4$.
There are $O(VW)$ choices for the right-hand side. When it is nonzero, the
divisor bound gives at most $O_\varepsilon(Z^\varepsilon)$ factorizations;
when it vanishes, we must have $v_3=v_4=0$ and $a'=0$.
Multiplying by the $R$ choices for $a$ gives
$O_\varepsilon(RVWZ^\varepsilon)$. The case $v_2=0$ is identical,
and these two coefficients cannot both vanish.

For the remaining relations, interchange $a,a'$ and change signs as needed,
then assign each pair a witness on a dyadic scale $S\le V$ satisfying
\[
 S\le v_2<2S,\qquad 0<|v_1|\le2S.
\]
Let $I_S$ count the pairs assigned to this scale.
Set $J=\lfloor SZ/(100W)\rfloor$.
If $SZ/W$ is bounded, then $S^2W/Z\gg1$ and the desired bound is immediate.
Otherwise, $1\le J<Z$ and $J\asymp SZ/W$. Use the Fej\'er kernel of height one,
\[
 B_J(t)=J^{-2}\left|\sum_{0\le m<J}\mathrm e(mt)\right|^2,
 \qquad f(t)=\sum_{a\in E}B_J(t-a/Z).
\]
Its Fourier support lies in $|n|<J$, and orthogonality on the full grid of
$Z$ points gives
\begin{equation}\label{geo:fejer-norms}
 \int f=R/J,\qquad \|f\|_\infty\le Z/J,
 \qquad \|f\|_2^2\le ZR/J^2.
\end{equation}
Consider the periodic operator
\begin{equation}\label{geo:periodic-operator}
 (T_Sf)(t)=
 \sum_{0<|m_1|\le2S}\ \sum_{S\le m_2<2S}
 \ \sum_{j=0}^{m_2-1}f\left(\frac{m_1t+j}{m_2}\right).
\end{equation}
For a given relation, take $m_1=-v_1,m_2=v_2,j\equiv-v_3\pmod {m_2}$.
At $t=a/Z$, the image lies within circular distance $W/(SZ)$ of $a'/Z$.
Since $J\le SZ/(100W)$, the corresponding two Fej\'er atoms are both bounded
below by an absolute positive constant on an interval of length $\gg1/J$.
Expanding $fT_Sf$ and retaining one witness for each ordered pair therefore gives
\begin{equation}\label{geo:positive-witness}
 I_S/J\ll\int_0^1 f(t)(T_Sf)(t)\,dt.
\end{equation}
The short intervals need not be disjoint: distinct pairs carry distinct
pairs of atom labels in the expansion, and all terms are nonnegative.
At the circular seam one may use a one-sided interval; the complete sum over
$j\bmod m_2$ ensures that the operator remains periodic.

For an input frequency $n$, summing over the intercept gives
\[
 \sum_{j\bmod m_2}\mathrm e\left(n\frac{m_1t+j}{m_2}\right)
 =\begin{cases}m_2\mathrm e(\ell m_1t),&n=\ell m_2,\\0,&m_2\nmid n.
 \end{cases}
\]
Thus the nonzero-frequency matrix is generated by
$n=\ell m_2$, $k=\ell m_1$, $\ell\ne0$, with weight $m_2$ for each contribution.
For fixed $n$ we have $m_2\mid n$, and for fixed $k$ we have $m_1\mid k$.
The other multiplier has $O(S)$ choices, each with weight $O(S)$.
Inputs with $|n|<J$ produce outputs with $|k|<2J$, so both row sums and column
sums are $O_\varepsilon(S^2Z^\varepsilon)$. Schur's test yields
\[
 \|T_S^{\ne0}\|_{2\to2}\ll_\varepsilon S^2Z^\varepsilon.
\]
The constant mode has eigenvalue $O(S^3)$, and, because $m_1\ne0$, no nonzero
input frequency contributes to the constant mode. Combining
\eqref{geo:fejer-norms} and~\eqref{geo:positive-witness}, we obtain
\[
 I_S\ll_\varepsilon
 J\left\{S^3(R/J)^2+S^2Z^\varepsilon ZR/J^2\right\}
 \ll_\varepsilon RSWZ^\varepsilon+R^2S^2W/Z.
\]
Summing the geometric series over $S$ and including the cases with a zero
coordinate proves the lemma.
\end{proof}

\subsection{Shared denominators and four-cycles}


For two distinct columns, their codegree is the number of denominators shared
by the two frequencies. If that number is $t$, the integer lifts of their
rational approximations lie in one plane, and the product of its successive
minima is $O(1/t)$. Their cross product consequently gives an affine relation
whose height decreases by the same factor. Summing this gain with weight
$t(t-1)$ gives the four-cycle estimate needed for the weighted energy.

\begin{lemma}[Anisotropic codegree bound]\label{geo:codegree}
Assume~\eqref{geo:det-hyp}.
For any $E\subset\{0,\ldots,Z-1\}$, let $R=|E|$ and
$t_{a,a'}=\sum_qM_{q,a}M_{q,a'}$.
Then
\begin{equation}\label{geo:fourcycles}
 \sum_{\substack{a,a'\in E\\a\ne a'}}t_{a,a'}(t_{a,a'}-1)
 \ll_\varepsilon Z^\varepsilon
 \left(R\Delta H^2+\frac{R^2\Delta H^3}{Z}\right).
\end{equation}
\end{lemma}

\begin{proof}
Fix a pair with $t=t_{a,a'}\ge2$ common denominators.
Each denominator determines a unique integer lift $(h_1,h_2,q)$ such that
\[
 e_1=h_1-qa/Z,\quad e_2=h_2-qa'/Z,
 \quad |q|\le2D,\quad |e_i|\le2CD\Delta/Z.
\]
The transformation $(h_1,h_2,q)\mapsto(q,e_1,e_2)$ has determinant of absolute
value one. By~\eqref{geo:det-hyp}, the integer determinant of any three lifted
vectors has absolute value less than one, so the vectors are coplanar.
Lifts with distinct positive denominators cannot be proportional, since both
fractions are reduced. They therefore span a lattice of rank two in this plane.

Let $\lambda_1\le\lambda_2\le1$ be the successive minima of this section lattice
with respect to the symmetric box above. The two-dimensional lattice-point
bound gives $t\ll(\lambda_1\lambda_2)^{-1}$, and hence
$\lambda_1\lambda_2\ll1/t$.
To see this counting bound, transform the box linearly to a unit cube and take
a reduced basis in the resulting plane. The two basis coordinates have
$O(1+\lambda_i^{-1})$ possible values.
Choose corresponding independent integer vectors $w,w'$, and set
$(v_1,v_2,v_3)=w\times w'$.
Writing $w=(h_1,h_2,q)$ and $w'=(h'_1,h'_2,q')$, we have
\[
 v_1=e_2q'-qe'_2,\qquad v_2=qe'_1-e_1q',
\]
and
\[
 v_1a+v_2a'+v_3Z=Z(e_1e'_2-e_2e'_1).
\]
The left-hand side is an integer. Let $v_4$ be its negative; then
\begin{equation}\label{geo:height-codegree}
 |v_1|,|v_2|,|v_3|\ll_C H/t,\qquad
 |v_4|\ll_C\Delta H/t,
 \qquad(v_1,v_2,v_3)\ne0.
\end{equation}
The product of error terms in $v_3$ is
$O((CD\Delta/Z)^2\lambda_1\lambda_2)$, which is
$O_C(H\lambda_1\lambda_2)$ by $4CD\Delta<Z$ and $D\ge1$.
This also gives $t\ll_CH$: if the first height bound is smaller than one,
there is no such integer configuration.

On a dyadic range $2\le T\le t<2T$, apply Lemma~\ref{geo:incidence} with
$V\asymp H/T$ and $W\asymp\Delta H/T$. This gives
\[
 \#\{(a,a')\in E^2:a\ne a',\ T\le t_{a,a'}<2T\}
 \ll_\varepsilon Z^\varepsilon
 \left(\frac{R\Delta H^2}{T^2}
       +\frac{R^2\Delta H^3}{ZT^3}\right).
\]
Multiply by $T^2$ and sum over dyadic scales. The extra logarithm in the first
term is absorbed into $Z^\varepsilon$, while the second term is a geometric
series. This proves~\eqref{geo:fourcycles}.
\end{proof}

\begin{proof}[Proof of Proposition~\ref{geo:energy}]
First take $w=\mathbf1_E$, and denote the row degrees by $\ell_q$ and the column
degrees by $d_a$. Farey separation gives $d_a\ll1+H\ll H$.
For the actual $0$--$1$ incidence matrix $M$, we have the exact identity
\begin{equation}\label{geo:trace-identity}
 \sum_{q\ne q'}(MM^*)_{q,q'}^2
 =\sum_a d_a(d_a-1)
  +\sum_{a\ne a'}t_{a,a'}(t_{a,a'}-1).
\end{equation}
Split $MM^*$ into its diagonal and off-diagonal parts. The former has operator
norm $\max_q\ell_q$, while the operator norm of the latter is at most its
Hilbert--Schmidt norm. Using
$\|M\mathbf1_E\|_2^2\le R\|MM^*\|$ and Lemma~\ref{geo:codegree}, we obtain
\begin{equation}\label{geo:unweighted-energy}
 \sum_q\ell_q^2\ll_\varepsilon Z^\varepsilon
 \left\{R\max_q\ell_q+R^{3/2}\Delta^{1/2}H
             +R^2\Delta^{1/2}H^{3/2}Z^{-1/2}\right\}.
\end{equation}
The repeated-column contribution
$\sum_a d_a(d_a-1)\ll RH^2$ is absorbed by the first four-cycle term,
since $\Delta\ge1$.

For general weights, decompose into layers
$E_\mu=\{a:\mu/2<w_a\le\mu\}$, with $\mu$ dyadic, and truncate at
$w_a\ge Z^{-10}$. There are $O(\log Z)$ layers. Write
$R_\mu=|E_\mu|$ and $\ell_{q,\mu}=\sum_{a\in E_\mu}M_{q,a}$.
Cauchy's inequality in the layer index gives
\[
 \sum_q\left(\sum_\mu\mu\ell_{q,\mu}\right)^2
 \ll(\log Z)\sum_\mu\mu^2\sum_q\ell_{q,\mu}^2.
\]
In the first term of~\eqref{geo:unweighted-energy}, use
$\mu\ell_{q,\mu}\le2S_1$ and
$\sum_\mu\mu R_\mu\le2H_1$. For the other two terms, use
\[
 \sum_\mu\mu^2R_\mu^{3/2}
 \le\left(\sum_\mu\mu^{4/3}R_\mu\right)^{3/2}
 \le4H_{4/3}^{3/2},\qquad
 \sum_\mu\mu^2R_\mu^2\le4H_1^2.
\]
This exhibits the origin of the $4/3$ moment directly.
Each row sum of the omitted tail is at most $Z^{-9}$, so its squared norm is
$O(DZ^{-18})$; the cross term is absorbed by $(u+v)^2\le2u^2+2v^2$.
Substitute $H=D^2\Delta/Z$ and absorb the fixed logarithmic loss by increasing
$\varepsilon$ slightly.
\end{proof}

\subsection{Digital row bounds on rational grids}

\begin{lemma}[Positive envelopes and separated sampling]\label{geo:sampling}
Let $Q$ be a trigonometric polynomial of degree at most $U\ge1$,
let $j=1,2$, and fix $C_0>0$. There is a nonnegative trigonometric
polynomial $E$ such that
\begin{equation}\label{geo:envelope}
 \sup_{|v|\le C_0/U}|Q(t+v)|^j\le E(t),\qquad
 \deg E=O_{C_0}(U),\qquad \int E\ll_{C_0}\int|Q|^j.
\end{equation}
If $\Theta\subset\mathbb R/\mathbb Z$ has circular spacing at least
$0<\delta\le1$, then
\begin{equation}\label{geo:separated-sampling}
 \sum_{t\in\Theta}\sup_{|v|\le C_0/U}|Q(t+v)|^j
 \ll_{C_0}(U+\delta^{-1})\int_0^1|Q(t)|^j\,dt.
\end{equation}
Moreover, every nonnegative trigonometric polynomial of degree at most
$U_0$ satisfies
\begin{equation}\label{geo:grid-positive}
 \sum_{h\bmod q}E(h/q)
 =q\sum_{\ell\in\mathbb Z}\widehat E(q\ell)
 \le(q+2U_0)\widehat E(0).
\end{equation}
\end{lemma}

\begin{proof}
Choose a smooth Fourier cutoff equal to one on $[-U,U]$ and supported
in $[-2U,2U]$. Its periodic kernel $\Psi_U$ reproduces $Q$ and satisfies
$|\Psi_U(t)|\ll_A U(1+U\|t\|)^{-2A-2}$ for every fixed $A$.
For a sufficiently large fixed integer $A$, use the positive kernel
\[
 K_U(t)=U\sum_{n\in\mathbb Z}\left\{
 \left(\frac{\sin\pi U(t+n)}{\pi U(t+n)}\right)^{2A}
 +\left(\frac{\sin\pi(U(t+n)-1/2)}{\pi(U(t+n)-1/2)}\right)^{2A}
 \right\}.
\]
It has degree $O_A(U)$ and bounded integral. The two real-line summands
have no common zero and their sum is comparable to
$(1+|Ut|)^{-2A}$. Consequently,
\[
 \sup_{|v|\le C_0/U}|\Psi_U(t+v)|\ll_{A,C_0}K_U(t).
\]
The integral form of H\"older's inequality, applied to
$Q=\Psi_U*Q$, now gives~\eqref{geo:envelope} with
$E=C_{A,C_0}K_U*|Q|^j$.
For fixed $x$, separation and the decay of $K_U$ give
\[
 \sum_{t\in\Theta}K_U(t-x)
 \ll U\sum_{m\ge0}(1+Um\delta)^{-2A}
 \ll U+\delta^{-1}.
\]
Integrating this bound against $|Q(x)|^j$ proves
\eqref{geo:separated-sampling}.
Finally, orthogonality on a complete grid gives the equality in
\eqref{geo:grid-positive}; its inequality follows from
$|\widehat E(n)|\le\widehat E(0)$ and the Fourier support of $E$.
\end{proof}

\begin{lemma}[Coprime row sums and window averages]\label{geo:rows}
Define
\[
 A_q^Z(\Delta)=\sum_{a<Z}F_Z(a/Z)M_{q,a}\qquad(q\ge1).
\]
Under~\eqref{geo:moments}, the following estimates hold uniformly.
Every sum and maximum in these three estimates is restricted to $(q,b)=1$:
\begin{align}
 \max_{q\sim D}A_q^Z(\Delta)&\ll
  \Delta^{1-\sigma}D^{1-c},&D^2\Delta&\le Z,
    \label{geo:row-small}\\
 \sum_{q\sim D}A_q^Z(\Delta)&\ll
  (D^2\Delta)^{1-\sigma},&D^2\Delta&\le Z,
    \label{geo:sum-small}\\
 \max_{q\sim D}A_q^Z(\Delta)&\ll
  Z^{\sigma-c}(D\Delta)^{1+c-2\sigma},&D^2\Delta&\ge Z,
    \label{geo:row-large}
\end{align}
where the last estimate only requires $D\Delta\ll Z$; changing the boundary
between the ranges by a fixed constant does not affect the conclusion.
The first two estimates incur no small-power loss in the ambient length $Z$.
\end{lemma}

\begin{proof}
We first prove a row-sum bound at the natural scale.
Write $Z=NM$, where $M,N$ are powers of the base, $M\le N$, and $D_0=N/M$.
By~\eqref{geo:product},
\[
 F_Z(t)=F_M(t)F_{D_0}(Mt)F_M(Nt).
\]
Apply $uv\le(u^2+v^2)/2$ to the two outer blocks and then use
\eqref{geo:envelope}, obtaining
\[
 \sup_{|v|\le C/Z}F_Z(t+v)
 \ll E_{D_0}(Mt)\{E_M(t)+E_M(Nt)\},
\]
where $\deg E_{D_0}=O(D_0)$, $\int E_{D_0}\ll D_0^{-\sigma}$,
$\deg E_M=O(M)$, and $\int E_M\ll M^{-c}$.
The first product has degree $O(N)$ and integral
$O(M^{-c}D_0^{-\sigma})$: average first conditional on $Mt\pmod1$.
Only $O(1)$ Fourier coefficients of $E_M$ at frequencies divisible by $M$
remain, and each is bounded by its mean.
For the second term, $(M,q)=1$ allows the exact permutation of $Mh\bmod q$:
\[
 \sum_{h\bmod q}E_{D_0}(Mh/q)E_M(Nh/q)
 =\sum_{h\bmod q}E_{D_0}(h/q)E_M(D_0h/q).
\]
The product on the right still has degree $O(N)$, and averaging first
conditional on $D_0t$ gives the same integral bound.
Thus~\eqref{geo:grid-positive} proves
\begin{equation}\label{geo:outer-blocks}
 \sum_{h\bmod q}\sup_{|v|\le C/Z}F_Z(h/q+v)
 \ll(q+N)M^{-c}D_0^{-\sigma}.
\end{equation}
For $\sqrt Z\ll q\ll Z$, choose a neighboring power $N\asymp q$ of the base
with $N\mid Z$ and $N\ge\sqrt Z$, to obtain
\begin{equation}\label{geo:outer-large}
 \sum_{h\bmod q}\sup_{|v|\le C/Z}F_Z(h/q+v)
 \ll Z^{\sigma-c}q^{1+c-2\sigma}.
\end{equation}
Rounding to powers of the base, and the parity of the word length, change
only constants depending on $b$.

For a wider window, choose a power $J\asymp\Delta$ of the base with $J\mid Z$,
and put $Z'=Z/J$. On the integer grid we have the exact factorization
\[
 F_Z(a/Z)=F_{Z'}(a/Z)F_J(a/J).
\]
Each rational window contains $O(J)$ consecutive integers, hence a bounded
number of complete grids modulo $J$ and at most two partial grids.
Sum the second factor over those integers first; positivity and
\eqref{geo:moments} give $O(J^{1-\sigma})$. Then take a local supremum of the first
factor:
\begin{equation}\label{geo:window-average}
 A_q^Z(\Delta)\ll\Delta^{1-\sigma}
 \sum_{\substack{h\bmod q\\(h,q)=1}}
 \sup_{|v|\le C'/Z'}F_{Z'}(h/q+v).
\end{equation}
If $D^2\Delta\ge Z$, apply~\eqref{geo:outer-large} at length $Z'$ to obtain
\eqref{geo:row-large}.

If $D^2\Delta\le Z$, choose a power $M\asymp D$ of the base such that
$M^2\mid Z'$. Discarding the middle digit block gives
\[
 F_{Z'}(t)\le F_M(t)F_M((Z'/M)t).
\]
Apply the arithmetic mean of the squares, permute the coprime residue
classes, and use~\eqref{geo:envelope}. Each envelope for $F_M^2$ has mean
$O(M^{-c})$, and the sampling cost is $O(q+M)$.
Equation~\eqref{geo:window-average} now gives~\eqref{geo:row-small}.
For the sum over all $q\sim D$, shorten the remaining word further to a length
$U\asymp D^2$. The reduced fractions are separated by $\gg D^{-2}$,
and the positive envelope has degree $O(D^2)$ and mean $O(D^{-2\sigma})$.
Lemma~\ref{geo:sampling} together with~\eqref{geo:window-average} gives
$\Delta^{1-\sigma}D^{2-2\sigma}$, which is~\eqref{geo:sum-small}.
Every loss in this argument is an absolute constant once the base is fixed;
bounded $D$ is handled by increasing that constant.
\end{proof}

\subsection{From reduced fractions to the full minor-arc sum}

The two fractions in a simultaneous approximation may have different
reduced denominators. We therefore retain both divisor sums until after
squaring. Weighted Cauchy places the divisor loss on the reduced
denominators, where the preceding estimates supply decay. The original
minor-arc restriction first removes the small rational rows; the digital
fibres can then be enlarged without restoring those rows.

\begin{proposition}[Simultaneous rational approximations with all denominators]\label{geo:all-denom}
Suppose that~\eqref{geo:moments} holds and that
\begin{equation}\label{geo:range-y}
 \sigma>1-c/2,\qquad
 0<y<\min\{2/3,\sigma,1-3g_{4/3}/2\}.
\end{equation}
For $X=b^k$ and $L=\log X$, define
\[
 \mathfrak M_B=
 \bigcup_{\substack{1\le q\le L^B\\(h,q)=1}}
 \{t\in\mathbb R/\mathbb Z:\|t-h/q\|\le L^B/X\},
 \qquad
 \mathcal R_B=\{0\le a<X:a/X\notin\mathfrak M_B\}.
\]
For every $J_0>0$ and fixed $C>0$, there is a choice of $B$ such that,
uniformly for $1\le Y\le X^y$,
\begin{equation}\label{geo:all-denom-bound}
 \sum_{\substack{a_1,a_2\in\mathcal R_B\\
 \exists,1\le q\le CY, h_1,h_2\in\mathbb Z:\
 |qa_i/X-h_i|\le CY/X\ (i=1,2)}}
 F_X(a_1/X)F_X(a_2/X)
 \ll_{J_0,C}Y L^{-J_0}.
\end{equation}
Neither fraction $h_i/q$ is required to be reduced.
\end{proposition}

\begin{proof}
We first control primitive denominators coprime to the base.
For $D\Delta\ll Z^y$, Lemma~\ref{geo:rows} gives
\begin{equation}\label{geo:small-energy}
 \frac1{D\Delta}\sum_{q\sim D,(q,b)=1}A_q^Z(\Delta)^2
 \ll\Delta^{1-2\sigma}D^{2-c-2\sigma},
 \qquad D^2\Delta\le Z.
\end{equation}
If $D^2\Delta>Z$, then $H>1$, $D\Delta/Z\ll Z^{y-1}$ and
$D^3\Delta^2/Z^2\le (D\Delta)^3/Z^2\ll Z^{3y-2}$. Thus both strict
inequalities in \eqref{geo:det-hyp} hold for large $Z$; bounded lengths
are absorbed into the constant.
Insert~\eqref{geo:moments} and~\eqref{geo:row-large}, and divide by $D\Delta$.
The resulting three terms are
\[
 Z^{1-c}(D\Delta)^{c-2\sigma},\qquad
 Z^{-1+3g_{4/3}/2}D\Delta^{1/2},\qquad
 Z^{-2\sigma}D^2\Delta.
\]
Since $D\Delta\ge\sqrt Z$ and $c<2\sigma$, these are at most, respectively,
\begin{equation}\label{geo:three-powers}
 Z^{1-\sigma-c/2},\qquad Z^{y-1+3g_{4/3}/2},\qquad
 Z^{2y-2\sigma}.
\end{equation}
Choose a fixed
\[
 0<\delta_0<
 \tfrac12\min\{\sigma+c/2-1,\ 1-y-3g_{4/3}/2,\ 2\sigma-2y\}.
\]
All three entries are positive by~\eqref{geo:range-y}.
Taking the incidence exponent $\varepsilon$ sufficiently small absorbs
both $Z^\varepsilon$ and the fixed logarithmic losses, giving
\begin{equation}\label{geo:primitive-energy}
 \frac1{D\Delta}\sum_{q\sim D,(q,b)=1}A_q^Z(\Delta)^2
 \ll
 \begin{cases}
 \Delta^{-\rho}D^{-\tau},&D^2\Delta\le Z,\\
 Z^{-\delta_0},&D^2\Delta>Z,
 \end{cases}
 \quad
 \rho=2\sigma-1>0,\quad\tau=c+2\sigma-2>0.
\end{equation}

Now write a reduced denominator as $d=gq$, where $g$ is a power of $p$ and
$(q,p)=1$. Let $g_0$ be the smallest power of $b$ not less than $g$.
Then $1\le g_0/g\le p^{r-1}$, and in the present range $g_0\mid X$ for all
sufficiently large $X$.
Put $Z=X/g_0$ and write uniquely
$a=t+Z\ell$, $0\le t<Z$, $0\le\ell<g_0$.
The digital product gives
\begin{equation}\label{geo:smooth-fiber}
 F_X(a/X)=F_{g_0}\left(\frac{t/Z+\ell}{g_0}\right)F_Z(t/Z).
\end{equation}
If $a/X$ is within $C\Delta/X$ of the reduced fraction $h/(gq)$, then $t/Z$
is within $C\Delta/Z$ of the reduced fraction $(g_0/g)h/q$.
Relaxing the restriction on $\ell$ and applying the first moment on shifted
grids yields
\begin{equation}\label{geo:fiber-row}
 A_{gq}^X(\Delta)\ll g_0^{1-\sigma}A_q^Z(\Delta).
\end{equation}
If $gq\Delta\ll Y$, then
$q\Delta\ll Y/g_0\le Z^y$ and $Z\gg X^{1-y}$.
Consequently, the constants in~\eqref{geo:primitive-energy} remain uniform.

To handle unreduced common denominators, we keep the divisor loss on the
reduced denominator. For nonnegative $B_d$ and $\omega>0$, weighted Cauchy gives
\[
 \left(\sum_{d\mid q}B_d\right)^2
 \le\sigma_{-\omega/2}(q)\sum_{d\mid q}d^{\omega/2}B_d^2.
\]
Use $\sigma_{-a}(dm)\le\sigma_{-a}(d)\sigma_{-a}(m)$,
$\sum_{m\le U}\sigma_{-a}(m)\le U\zeta(1+a)$, and
$\sigma_{-a}(d)\ll_\omega d^{\omega/2}$, and interchange nonnegative sums.
This gives
\begin{equation}\label{geo:divisor-cauchy}
 \sum_{q\sim Q}\left(\sum_{d\mid q}B_d\right)^2
 \ll_\omega Q\sum_{d\le2Q}d^{\omega-1}B_d^2.
\end{equation}

Decompose the original common denominators into $q\sim Q$, and take
$\Delta\asymp\max(1,Y/Q)$, so that $Q\Delta\asymp Y$.
Let $B_d$ be the mass of row $d$ in~\eqref{geo:incidence-matrix}, still summed
only over the original set $a\in\mathcal R_B$.
The reduced denominator of every unreduced fraction divides $q$, so the
pair mass $\mathcal P_Q$ supplied by this block satisfies
\[
 \mathcal P_Q
 \le\sum_{q\sim Q}\left(\sum_{d\mid q}B_d\right)^2.
\]
The two fractions in a pair may have different reduced denominators;
the weighted divisor inequality~\eqref{geo:divisor-cauchy} accounts for
their cross terms.
First group by $d=gq'$ and $q'\sim D$.
Choose $T$ later, and take $B>T+2$.
If $gD\Delta\le L^T$, every reduced denominator in the block is at most
$2L^T$, and every window has radius at most $CL^T/X$.
Thus the original mass $B_d$ vanishes identically.
This deletion takes place before removing the original major-arc restriction
and before enlarging the sum over the digital fibers.
Divide each remaining block by $Y$. Its bound from
\eqref{geo:divisor-cauchy} and~\eqref{geo:fiber-row} is
\[
 \frac QY(gD)^{\omega-1}g_0^{2-2\sigma}
       \sum_{\substack{q'\sim D\\(q',b)=1}}A_{q'}^Z(\Delta)^2
 \ll g^{-\rho+\omega}D^\omega
       \frac1{D\Delta}
       \sum_{\substack{q'\sim D\\(q',b)=1}}A_{q'}^Z(\Delta)^2.
\]
Here $Q/Y\asymp\Delta^{-1}$ and $g_0/g$ is bounded in terms of the base.
Thus~\eqref{geo:primitive-energy} bounds the block by
\begin{equation}\label{geo:block-final}
 g^{-\rho+\omega}D^\omega
 \begin{cases}
 \Delta^{-\rho}D^{-\tau},&D^2\Delta\le Z,\\
 Z^{-\delta_0},&D^2\Delta>Z,
 \end{cases}
 \qquad Z=X/g_0.
\end{equation}
Choose
\[
 \omega<\min\{\rho/2,\tau/2,\delta_0(1-y)/(2y)\}.
\]
The large-scale part is
$O(X^{-\delta_0(1-y)/2})$, since $D\ll X^y$ and $Z\gg X^{1-y}$.
The small-scale part is at most
$(gD\Delta)^{-\lambda}$, where
$\lambda=\min\{\rho-\omega,\tau-\omega\}>0$.

Every retained small-scale block satisfies $gD\Delta>L^T$, and hence
contributes at most $L^{-\lambda T}$.
Take $\lambda T>J_0+10$ and sum over the $O(L^3)$ blocks to obtain
\eqref{geo:all-denom-bound}.
This order of operations also handles the rational point zero:
its repeated occurrence among the full denominator rows is removed by the
original major-arc restriction before the enlargement.
\end{proof}

\section{The Type II interval from moment bounds}\label{geo:typeii-section}

The two fractional moments have different roles. An estimate
$\sum_{a<X}F_X(a/X)^s\ll X^{g_s}$ controls the small-amplitude tail
and determines $\kappa=g_s/(2-s)$, while the
$4/3$ moment controls common denominators through the preceding section.
Together they give the interval $(\kappa,1-2\kappa)$ below. We use
Maynard's bilinear Cauchy reduction \cite[\S13]{Maynard2019}, inserting
the weighted common-denominator estimate in its simultaneous-approximation
alternative.

\begin{theorem}[Type II minor-arc estimate from moment bounds]\label{geo:typeii}
Assume the hypotheses of Proposition~\ref{geo:all-denom}, and suppose that,
for some $4/3\le s<2$,
\begin{equation}\label{geo:fractional}
 \sum_{a<X}F_X(a/X)^s\ll X^{(2-s)\kappa},\qquad
 \kappa<1/3,\qquad 2\kappa<y.
\end{equation}
Then every compact interval
$[\alpha_0,\beta_0]\subset(\kappa,1-2\kappa)$
is a Type II window: for each $J>0,C_0>0$, one can choose $B$ such that
\begin{equation}\label{geo:typeii-bound}
 \left|\sum_{h\in\mathcal R_B}F_X(h/X)\gamma_h
 \sum_{n\sim N}\sum_{m\sim M}u_nv_m\mathrm e(-hnm/X)\right|
 \ll_{J,C_0}X(\log X)^{-J},
\end{equation}
uniformly for $X^{\alpha_0}\le N\le X^{\beta_0}$, $NM\le C_0X$,
and all $|u_n|,|v_m|,|\gamma_h|\le1$.
The intervals of support may be shortened, and the coefficients may be
multiplied by any fixed power of the logarithm.
\end{theorem}

Stretch the thin direction of the slab to the scale of its other two
directions. An excess of lattice points over volume excludes rank three.
A line then gives a short integer relation, while a plane gives a common
denominator through its integer normal.

\begin{lemma}[The lattice dichotomy for an integer slab]\label{geo:slab-dichotomy}
Fix $C\ge1$. There is a constant $K_0=K_0(C)$ with the following property.
Let $X\ge2$ be an integer, let $N\ge1$, and suppose that
\[
 C^{-1}N/X\le\delta\le1,\qquad K\ge K_0,\qquad
 \delta N^2K\ge K_0.
\]
For $a_1,a_2\in\{0,\ldots,X-1\}$, put
$\mathbf A=(a_1,a_2,X)$ and
\[
 \mathcal S=\{v\in\mathbb Z^3:|v|\le CN,\
                         |\mathbf A\cdot v|\le C\delta X\}.
\]
If $|\mathcal S|\ge\delta N^2K$, then at least one of the following holds:
\begin{enumerate}
 \item there are integers $v_1,v_2,v_3,v_4$, with
 $(v_1,v_2,v_3)\ne0$, such that
 \[
  v_1a_1+v_2a_2+v_3X+v_4=0,\qquad
  \max_{1\le i\le4}|v_i|\ll_C X/(N^2K);
 \]
 \item there are integers $q,h_1,h_2$ such that
 \begin{equation}\label{geo:simultaneous}
  1\le q\ll_C X/(NK),\qquad
  |qa_i/X-h_i|\ll_C1/(NK)\quad(i=1,2).
 \end{equation}
\end{enumerate}
No coprimality is asserted in the second alternative.
\end{lemma}

\begin{proof}
This is the combination of the geometric reductions in
\cite[Lemmas 13.2 and 14.1]{Maynard2019}, with the lower bound on the
number of points made explicit.
Stretch the direction $\mathbf A/|\mathbf A|$ by $N/\delta$ and leave its
orthogonal complement fixed; denote this map by $\Phi$.
The lattice $\Phi(\mathbb Z^3)$ has determinant $N/\delta$, and
$\Phi(\mathcal S)$ lies in a ball of radius $O_C(N)$.
By the fixed-dimensional basis and successive-minima estimates
\cite[Section~5]{Maynard2019}, a reduced lattice basis $b_i$, with
$s_1\le s_2\le s_3$ and $s_i=|b_i|$, satisfies
\[
 \left|\sum_i z_i b_i\right|\asymp\sum_i|z_i|s_i,\qquad
 s_1s_2s_3\asymp\det\Phi(\mathbb Z^3)=N/\delta.
\]
Thus the coordinates of any point in that ball satisfy
$|z_i|\ll_C N/s_i$.
If $s_3$ is at most a sufficiently large constant times $N$, the ball
contains only
$O_C(N^3/(s_1s_2s_3))=O_C(\delta N^2)$ lattice points, contradicting
$K\ge K_0$ when $K_0$ is large enough.
Otherwise every point in the ball has $z_3=0$.
Consequently the linear span of $\mathcal S$ has dimension at most two.
It has positive dimension since $|\mathcal S|\ge K_0>1$.

If this span is a line, it passes through the origin, and its integer
points are multiples of a primitive vector $v$.
Because $\mathcal S$ is symmetric and contains the origin, its largest
multiple has coefficient $|j|\gg|\mathcal S|$.
The two defining inequalities of $\mathcal S$ therefore give
\[
 |v|\ll_C\frac1{\delta NK}\ll_C\frac X{N^2K},
 \qquad
 |\mathbf A\cdot v|\ll_C\frac X{N^2K}.
\]
Taking $v_4=-\mathbf A\cdot v\in\mathbb Z$ proves the first alternative.

If the span has dimension two, take a reduced basis of the transformed
section lattice, of lengths $s_1,s_2$.
Since $\Phi(\mathcal S)$ contains two independent vectors of length
$O_C(N)$, both basis lengths are $O_C(N)$.
The planar lattice-point bound gives
\[
 \delta N^2K\le|\mathcal S|\ll_C\frac{N^2}{s_1s_2},
 \qquad s_1s_2\ll_C\frac1{\delta K}.
\]
Pull the basis back to integer vectors $w_1,w_2$. Then
\[
 |w_i|\le s_i,\qquad
 |\mathbf A\cdot w_i|\ll_C\delta Xs_i/N.
\]
For $v=w_1\times w_2\ne0$, the projection of $\mathbf A$ onto
$\operatorname{span}(w_1,w_2)$ is
\[
 E=\frac{(\mathbf A\cdot w_1)(w_2\times v)
       +(\mathbf A\cdot w_2)(v\times w_1)}{|v|^2}.
\]
Hence
$|E|\ll_C\delta Xs_1s_2/(N|v|)\ll_C X/(NK|v|)$ and
$\mathbf A=\lambda v+E$.
Since $|v|\ge1$, $|\mathbf A|\asymp X$, and $A_3=X$, taking $K_0$
large enough gives $|v_3|\asymp|v|$.
Change the sign of $v$ so that $q=v_3>0$ and put $h_i=v_i$.
Then
\[
 q\le|v|\le s_1s_2\ll_C\frac1{\delta K}\ll_C\frac X{NK},
 \qquad
 |qa_i-Xh_i|=|v_3E_i-v_iE_3|
 \ll_C|v|\,|E|\ll_C\frac X{NK}.
\]
This is the second alternative.
\end{proof}

\begin{proof}[Proof of Theorem~\ref{geo:typeii}]
Choose $u_0$ sufficiently close to $1-2\kappa$ that
\begin{equation}\label{geo:u-choice}
 \beta_0<u_0<1-2\kappa,\qquad 1-u_0<y,\qquad
 2\alpha_0+u_0>1.
\end{equation}
These three requirements follow from $\beta_0<1-2\kappa$,
$2\kappa<y$, and $\alpha_0>\kappa$, respectively.
Choose fixed $\varepsilon_0,\xi>0$ sufficiently small that
\begin{align*}
 u_0+2\kappa+2(s-1)\varepsilon_0&<1,\\
 2\kappa-2\alpha_0+(3s-4)\varepsilon_0+\xi&<0,\\
 6\kappa-1-3\alpha_0+6(s-1)\varepsilon_0+2\xi&<0.
\end{align*}
The first inequality follows from the choice of $u_0$, the second from
$\alpha_0>\kappa$, and the third from
$6\kappa-1-3\alpha_0<3\kappa-1<0$.
These choices are made before choosing the requested logarithmic saving.
Let
$E=\{a:F_X(a/X)\ge X^{-\kappa-\varepsilon_0}\}$.
Writing $g=g_s=(2-s)\kappa$, we have
\begin{equation}\label{geo:tail}
 \sum_{a\notin E}F_X(a/X)^2\ll X^{-(2-s)\varepsilon_0},
 \qquad H_E:=\sum_{a\in E}F_X(a/X)
 \ll X^{\kappa+(s-1)\varepsilon_0}.
\end{equation}
Let $S(h)$ denote the bilinear exponential sum in~\eqref{geo:typeii-bound}.
Group the products by $nm$; their coefficients are bounded by the divisor
function and supported in $O_{C_0}(X)$.
Folding modulo $X$ introduces only $O_{C_0}(1)$ aliases.
Parseval's identity and
$\sum_{n\le CX}\tau(n)^2\ll_CX(\log X)^3$ give
\[
 \sum_{h<X}|S(h)|^2\ll_{C_0}X^2(\log X)^3.
\]
Cauchy's inequality therefore bounds the cost of discarding $a\notin E$ by
$O(X^{1-(2-s)\varepsilon_0/2}(\log X)^{3/2})$.

Split the retained frequencies into $O(\log X)$ amplitude layers
$E_T=\{a:F_X(a/X)\asymp T^{-1}\}$,
$1\le T\ll X^{\kappa+\varepsilon_0}$.
Fix one layer, retaining the restriction to $\mathcal R_B$.
Apply Cauchy's inequality in $m$, then take absolute values of the geometric
sums over $m$. The resulting square is at most $CX/N$ times the positive quantity
\begin{equation}\label{geo:cauchy-kernel}
 \mathcal K=
 \sum_{a_1,a_2\in E_T\cap\mathcal R_B}F_X(a_1/X)F_X(a_2/X)
 \sum_{1\le n_1,n_2\le C_1N}
 \min\left(\frac{C_2X}{N},
 \left\|\frac{n_1a_1-n_2a_2}{X}\right\|^{-1}\right).
\end{equation}
Our target is $\mathcal K\ll_JNX(\log X)^{-J}$.
For each pair $(a_1,a_2)$, denote the inner sum in
\eqref{geo:cauchy-kernel} by $\mathscr H_N(a_1,a_2)$.
It lies between constant multiples of $N^2$ and $XN$.
Partition the pairs according to
$\mathscr H_N(a_1,a_2)\asymp N^2K$, with $K$ dyadic and
$1\ll K\ll X/N$.
For all layers with $NK\le X^{u_0}$, equation~\eqref{geo:tail} gives directly
\begin{equation}\label{geo:small-K}
 \mathcal K_{\rm small}
 \ll N X^{u_0+2\kappa+2(s-1)\varepsilon_0}\log X
 =o\bigl(NX(\log X)^{-J}\bigr).
\end{equation}
Choose $\varepsilon_0$ so that the exponent here is strictly less than one.

The remaining layers satisfy $NK>X^{u_0}$.
Decompose distances into ranges $N/X\ll\delta\le1$:
each kernel layer is bounded by $C\delta^{-1}$ times the indicator that the
distance is at most $C\delta$.
Hence there is a choice of $\delta$ for which the integer slab
\begin{equation}\label{geo:slab}
 |v|\le C_3N,\qquad
 |v_1a_1-v_2a_2-v_3X|\le C_4\delta X
\end{equation}
contains $\gg\delta N^2K'$ lattice points, where
$K'=K/(C_5\log X)$.
These points come from the original positive integers $n_1,n_2$;
enlarging to the symmetric slab only increases their number.
Each $(v_1,v_2)$ permits only $O(1)$ choices for $v_3$, so
$\delta\ll1/K'$.
Moreover,
\begin{equation}\label{geo:slab-margins}
 K'\gg X^{u_0-\beta_0}/\log X\longrightarrow\infty,
 \qquad
 \delta N^2K'\gg X^{2\alpha_0+u_0-1}/\log X\longrightarrow\infty.
\end{equation}

The two lower bounds in~\eqref{geo:slab-margins} ensure that
Lemma~\ref{geo:slab-dichotomy} applies uniformly to every remaining
layer once $X$ is sufficiently large. Change the signs of the second
and third coordinates in~\eqref{geo:slab}; enlarge the fixed constant $C$
and choose $C_5$ so that the lemma applies with $K'$ in place of $K$.
Put $Y=X/(NK')$. The bounds on $K$ and~\eqref{geo:u-choice} give
\[
 1\ll\log X\ll Y\ll X^{1-u_0}\log X<X^y.
\]

For pairs in the simultaneous-approximation alternative,
Proposition~\ref{geo:all-denom} bounds the weighted mass by
$O_J(Y(\log X)^{-J})$.
Multiplying by the kernel bound $O(N^2K)$ costs only the ratio
$K/K'=O(\log X)$, giving a contribution
$O_J(NX(\log X)^{-J+1})$ to~\eqref{geo:cauchy-kernel}.
All constants are uniform in the amplitude, kernel-size, and distance layers.

For the line alternative, Lemma~\ref{geo:slab-dichotomy} gives a nonzero
integer relation of the form in Lemma~\ref{geo:incidence}, with all four
heights $O(V)$, where $V=X/(N^2K')$.
If $V$ is too small, integrality excludes this case; otherwise
Lemma~\ref{geo:incidence} applies with both height parameters comparable
to $V$.

Write $R=|E_T|$, $\zeta=\log_XT$, and $h=g+(s-1)\zeta$.
The $s$th moment gives
\[
 R/T\ll X^h,\qquad R\ll X^{h+\zeta}.
\]
Put $\nu=\log_XN$ and $u=\log_X(NK')$.
The weighted mass of pairs giving a line is bounded both by $(R/T)^2$ and by
\begin{equation}\label{geo:line-mass}
 X^\varepsilon T^{-2}\left(RV^2+R^2V^3/X\right).
\end{equation}
Relative to the target scale $Y=X/(NK')$, the exponent in the trivial bound is
$2h-1+u$.
Use that bound directly when $u\le1-2h-\xi$, with the $\xi$ fixed above.
Otherwise $1-u<2h+\xi$. Since $V=X^{1-\nu-u}$ and $Y=X^{1-u}$,
the two terms in \eqref{geo:line-mass}, divided by $Y$ and before the
factor $X^\varepsilon$, have exponents at most
$h-\zeta-2\nu+(1-u)$ and $2h-1-3\nu+2(1-u)$. Thus they are at most
\[
 3h-\zeta-2\nu+\xi,\qquad 6h-1-3\nu+2\xi.
\]
Since $s\ge4/3$ and
$\zeta\le\kappa+\varepsilon_0+O(1/\log X)$,
\[
 3h-\zeta=3g+(3s-4)\zeta
 \le2\kappa+(3s-4)\varepsilon_0+o(1),\qquad
 h\le\kappa+(s-1)\varepsilon_0+o(1).
\]
The two exponents are therefore at most
\begin{equation}\label{geo:line-powers}
 2\kappa-2\alpha_0+(3s-4)\varepsilon_0+\xi+o(1),\qquad
 6\kappa-1-3\alpha_0+6(s-1)\varepsilon_0+2\xi+o(1).
\end{equation}
The initial choice of $\varepsilon_0,\xi$ makes both exponents strictly
negative. Choose the incidence loss $\varepsilon$ smaller than these
fixed margins.
The line contribution thus saves a fixed power, sufficient to absorb
$K/K'$ and all logarithmic losses from the amplitude, distance, and kernel-size
decompositions.
Combining this with~\eqref{geo:small-K} and the rank-two estimate proves
$\mathcal K\ll_JNX(\log X)^{-J}$.
Return to the initial Cauchy inequality and increase $J$ to obtain
\eqref{geo:typeii-bound}.
\end{proof}


\section{Arithmetic comparison on prime-factor regions}\label{arith:section}

We verify two inputs of Theorem~\ref{common:positive-sieve}: interval
Type I distribution (Proposition~\ref{arith:typeI}) and comparison on
affine regions of prime factorizations with an admissible subproduct
(Proposition~\ref{arith:polytope}). The former also applies to the
octal words in Section~\ref{sec:octal-typeI}.

Let the base $b=\ell^\nu\ge4$ be a fixed prime power, let $d$ be the excluded
digit, and write
\[
 \mathcal D=\{0,\ldots,b-1\}\setminus\{d\},\quad
 X=b^k,\quad c=\log_b(b-1),\quad A=(b-1)^k,
\]
\[
 \mathcal A_X=\left\{\sum_{j=0}^{k-1}d_jb^j:d_j\in\mathcal D\right\},
 \qquad P_X(t)=A^{-1}\sum_{n\in\mathcal A_X}e(nt),\quad F_X=|P_X|.
\]
Leading zeros are allowed here when $d\ne0$; when $d=0$, $\mathcal A_X$
contains only admissible integers with exactly $k$ digits. Put $L=\log X$
and
\begin{equation}\label{arith:localfactor}
 u_d=\#\{a\in\mathcal D:(a,b)=1\},\qquad
 \gamma_d=\frac{b\,u_d}{\varphi(b)(b-1)}.
\end{equation}
All sieve sums and prime-factor sums are restricted to $1\le n<X$.
The digital Fourier polynomial may include zero, since the arithmetic
weight paired with it vanishes there.
The number of marked prime factors, the coefficients of every polyhedral
condition, and the sieve endpoints are fixed as $X\to\infty$.
Implied constants may depend on these data and on $b,d$.
\begin{lemma}[A uniform bound for small products]\label{arith:small}
For $1\le Y\le X$, we have
\[
 \#(\mathcal A_X\cap[0,Y])\le(b-1)Y^c.
\]
Suppose that a family of marked factorizations has total absolute weight
at most $C(\log X)^K$ at each positive integer. Its combined digital and
normalized reference mass on $n\le X^{1-\delta}$ is at most
\begin{equation}\label{arith:small-bound}
 O\left(CA(\log X)^K
       \{X^{-c\delta}+X^{-\delta}\}\right).
\end{equation}
Consequently, if $\delta\log X/\log\log X\to\infty$, this quantity is
smaller than every fixed $A(\log X)^{-J}$.
\end{lemma}
\begin{proof}
Choose $b^j\le Y<b^{j+1}$. If zero is allowed, all digits above position
$j$ must vanish, leaving at most $(b-1)^{j+1}\le(b-1)Y^c$ integers.
If zero is excluded and $k>j+1$, there are no such integers; in the
remaining cases the total number of words gives the same bound.
The reference weight is $\gamma_d A/X$, so the reference mass of integers
$n\le Y$ is $O(AY/X)$. Multiplying by the bound for the absolute weight
at each integer proves the second assertion.
A product of a fixed number $r$ of marked primes has at most $r!$
markings, and its logarithmic weight is at most $L^r$. If all prime
factors exceed a fixed $X^\epsilon$, their number is at most
$\lfloor1/\epsilon\rfloor$, so this case also satisfies the stated
hypothesis.
\end{proof}

\subsection{The major-arc comparison on prime boxes}

\begin{proposition}[Prime-factor boxes with logarithmic weights]\label{arith:major}
Fix $\epsilon>0$ and an integer $r\ge1$, and put $\delta=L_2^{-1}$,
where $L_2=\log\log X$.
Take $a_i\ge\epsilon/2$ and $\sum_{i<r}a_i<1-\epsilon/2$, and define
\[
 R(n)=\sum_{\substack{p_1\cdots p_r=n\\
 X^{a_i}<p_i\le X^{a_i+\delta}\ (i<r)\\
 p_r\ge\max(X^{\epsilon/4},X^{1-\sum_{i<r}a_i-r\delta})}}
       \prod_{i=1}^r\log p_i\qquad(1\le n<X).
\]
Write $S_R(t)=\sum_{n<X}R(n)e(nt)$ and $W_R=S_R(0)$.
Let $\mathfrak M(C)\subset\{0,\ldots,X-1\}$ consist of the indices $h$
for which $h/X$ is within distance $L^C/X$ of a reduced rational point
with denominator at most $L^C$. For fixed $C,J>0$,
\begin{equation}\label{arith:majorformula}
 \frac1X\sum_{h\in\mathfrak M(C)}
       \left(\sum_{n\in\mathcal A_X}e(nh/X)\right)S_R(-h/X)
   =\gamma_d\frac AX W_R+O(A L^{-J}).
\end{equation}
The implied constant is uniform in the admissible box endpoints.
\end{proposition}
\begin{proof}
If a reduced denominator $q\le L^C$ has a prime factor other than $\ell$,
the first $k/3$ orbit points $b^ja/q$ are nonintegral. A perturbation
$|v|\le L^C/X$ reduces their distance lower bound $1/q$ by at most
a factor of two. The orbit argument used for Type I therefore gives
\[
 F_X(a/q+v)\ll_C
       \exp\left(-c_C\frac{\log X}{\log\log X}\right).
\]
There are $O(L^{3C})$ grid frequencies on the major arcs, and
$S_R\ll X L^{O_r(1)}$, so these denominators contribute an admissible
error.

The remaining denominators are $q=\ell^j$. Denote the product of the
first $r-1$ primes by $m$, and their marked logarithmic coefficient
by $C_X(m)$. We have
\[
 \sum_m\frac{C_X(m)}m\ll L^{r-1},\qquad m\le X^{1-\epsilon/3}
\]
for sufficiently large $X$.
Temporarily remove the lower bound on the last prime, obtaining
$\widetilde S_R(t)=\sum_mC_X(m)\sum_{p<X/m}(\log p)e(mpt)$.
The added products do not exceed $X^{1-\epsilon/12}+X^{1-\delta}$.
Chebyshev's bound and the preceding reciprocal sum imply
\begin{equation}\label{arith:freeprime}
 \sup_t|\widetilde S_R(t)-S_R(t)|\ll_J XL^{-J}
\end{equation}
for every fixed $J$.

By Siegel--Walfisz (in the logarithmic-modulus form of
\cite[(5.1)]{Maynard2019}), for
$Y=X/m\ge X^{\epsilon/3}$,
$q\le L^C$, $(a,q)=1$ and $0\le z\le Y$,
\[
 \sum_{\substack{p\le z\\p\equiv a\pmod q}}\log p
       =\frac z{\varphi(q)}+O_{C,J,\epsilon}(YL^{-J}).
\]
For $z<Y^{1/2}$ the trivial bound $O(Y^{1/2}\log Y)$ suffices.
Otherwise sum (5.1) of the cited paper on dyadic intervals ending at
$z,z/2,\ldots$ until the remainder is $O(Y^{1/2})$. Their logarithmic
scales are comparable to $L$, and deleting proper prime powers costs
$O(Y^{1/2}\log^2 Y)$.
Partial summation then yields, for $|v|\le L^C$,
\begin{equation}\label{arith:APprime}
 \sum_{p<Y}(\log p)e(amp/q)e(vmp/X)
 =\frac{\mu(q)}{\varphi(q)}\int_0^Y e(vmt/X)\,dt
       +O(q(1+|v|)YL^{-J}).
\end{equation}
Here $(m,q)=1$, and the finitely many primes that are not units can be
absorbed into the error.

For sufficiently large $X$, we have $q\mid X$, so the offsets on the
major-arc integer grid are exactly $v/X$, with $v\in\mathbb Z$.
When $v\ne0$,
$\int_0^{X/m}e(vmt/X)dt=m^{-1}\int_0^Xe(vu/X)du=0$.
Thus \eqref{arith:freeprime}--\eqref{arith:APprime} handle all nonzero
offsets. At zero offset, the same calculation with $q=1$ gives
\[
 W_R=X\sum_m\frac{C_X(m)}m+O_J(XL^{-J}),
 \qquad
 S_R(a/q)=\frac{\mu(q)}{\varphi(q)}W_R+O_J(XL^{-J})
 \quad(q=\ell^j).
\]
Since $\mu(\ell^j)=0$ for $j\ge2$, only $0$ and $a/\ell$,
$1\le a<\ell$, remain. Put $S_{\mathcal A}(t)=AP_X(t)$.
Their digital coefficients satisfy
\[
 S_{\mathcal A}(0)-\frac1{\ell-1}
                  \sum_{a=1}^{\ell-1}S_{\mathcal A}(a/\ell)
 =\frac{\ell}{\ell-1}\#\{n\in\mathcal A_X:(n,b)=1\}
 =\gamma_d A.
\]
Multiply this identity by the arithmetic main term $W_R$ and divide
by $X$ to obtain \eqref{arith:majorformula}. All logarithmic powers
introduced by the denominators, offsets and preceding prime sums
are absorbed by choosing $J$ sufficiently large in Siegel--Walfisz.
\end{proof}

\subsection{From boxes to prime-factor regions}

We now assume that the bounded-coefficient minor-arc estimate of
Theorem \ref{geo:typeii} holds on a fixed short-factor window
$[\alpha_0,\beta_0]\subset(0,1/2)$.
Choose a sieve window strictly inside it:
\begin{equation}\label{arith:window}
 \alpha_0<\alpha<\beta<\beta_0,
 \qquad \mathcal W=[\alpha,\beta]\cup[1-\beta,1-\alpha].
\end{equation}

\begin{proposition}[Comparison on prime-factor polytopes]\label{arith:polytope}
Fix an integer $r\ge2$ and $\epsilon>0$, and consider marked prime tuples with
$p_1\cdots p_r<X$ and $p_i\ge X^\epsilon$.
One may impose any Boolean combination of finitely many fixed affine
inequalities, whose variables are either $\log p_i/\log X$ or
$\log p_i/\log(p_1\cdots p_r)$.
Suppose that each admissible tuple has a nonempty proper subproduct
whose logarithm to base $X$ lies in $\mathcal W$. Then
\begin{equation}\label{arith:polytopeformula}
 \sum_{\boldsymbol p}^{\mathcal R}
       1_{\mathcal A_X}(p_1\cdots p_r)
 =\gamma_d\frac AX\sum_{\boldsymbol p}^{\mathcal R}1
       +o(A/L).
\end{equation}
The same conclusion holds when the subproduct condition is normalized
by the total product. Finitely many such regions may be added, with
their respective multiplicities retained.
\end{proposition}
\begin{proof}
First convert conditions normalized by the total product exactly.
Write $n=p_1\cdots p_r$, $t_i=\log p_i/L$ and
$s_0=\sum_i t_i=\log n/L>0$. For any comparison sign $\bowtie$,
\begin{equation}\label{arith:homogenize}
 \lambda_0+\sum_i\lambda_i\frac{\log p_i}{\log n}\ \bowtie\ 0
 \quad\Longleftrightarrow\quad
 \sum_i(\lambda_i+\lambda_0)t_i\ \bowtie\ 0.
\end{equation}
This handles tests written entirely in one normalization. More generally
a mixed affine test has left side $C+U\cdot t+V\cdot t/s_0$, with fixed
vectors $U,V$. On $1-\delta\le s_0<1$ it differs from the affine form
$C+(U+V)\cdot t$ by $O(\delta)$. If the restriction of that form to
$\sum t_i=1$ vanishes identically, then $U+V=-C\mathbf1$ and the exact
left side is $(s_0-1)U\cdot t/s_0$. The first factor has fixed negative
sign, so this is an ordinary affine sign test (including equality).
Otherwise the test has fixed sign or a nonconstant affine wall on the
simplex. Only boxes in its $O(\delta)$ strip can be ambiguous.

First consider a box from Proposition \ref{arith:major}.
Choose a proper group that stays strictly inside the analytic window,
with a fixed buffer, on the retained part of the box. For a reflected
window use the complementary group: its exponent is $s_0$ minus that
of the original group. The last prime in the box ranges over a band of
width $O_r(\delta)$ in exponent after imposing its redundant upper bound
below. Thus either grouped factor varies by only $O_r(\delta)$ in
exponent even on the enclosing product box. Write $u$ for the short
group and $v$ for the other.
Each coefficient has absolute value at most $r!L^r$; repeated primes
do not affect this bound.
First insert the implicit upper bound $p_r\le X^{1-\sum_{i<r}a_i}$,
and then split $u,v$ into dyadic ranges.
Discard every dyadic pair whose support has no point with $uv<X$.
On each surviving pair $u\sim U,v\sim V$ we have $UV<X$ by its lower
endpoints; throughout its rectangle $uv\le4X$. Rounding the short scale
to a dyadic endpoint costs at most $\log_X2$, absorbed by the buffer.
Separate the sole remaining coupling $uv<X$ by the truncated Perron
formula with threshold $X-1/2$, real part $1/L$, and height $X^{10}$.
The logarithmic distance of an integer $uv$ from this threshold is
$\gg X^{-1}$, giving an error $O(X^{-9})$ per term, still negligible
after summation.
The total variation of the Perron measure is $O(L)$, and the inserted
factors $u^{-it},v^{-it}$ are bounded coefficients.
The minor-arc estimate therefore applies, with its arbitrary logarithmic
saving absorbing $r!L^r$ and the losses from these decompositions.
Together with \eqref{arith:majorformula}, this gives
\begin{equation}\label{arith:boxcomparison}
 \sum_{n\in\mathcal A_X}R(n)
 =\gamma_d\frac AX\sum_{n<X}R(n)+O_J(AL^{-J}).
\end{equation}

Put $\delta=1/\log\log X$ and write
$e(\boldsymbol a)=(a_1,\ldots,a_{r-1},1-\sum_{i<r}a_i)$.
Within the box, each prime-factor logarithm is
$L e_i(\boldsymbol a)(1+O_{r,\epsilon}(\delta))$.
The prime number theorem and partial summation give
\[
 \sum_{X^{a_i}<p\le X^{a_i+\delta}}\frac{\log p}{p}
       \ll\delta L,
 \qquad \sum_{n<X}R(n)\ll X(\delta L)^{r-1}.
\]
Hence the unweighted reference mass of a box is
$O(X\delta^{r-1}/L)$.
Positivity and \eqref{arith:boxcomparison} give the corresponding
bound for the digital box. After removing the logarithmic weights,
the comparison error is
\begin{equation}\label{arith:boxerror}
 O(A\delta^r/L)+O_J(AL^{-J}).
\end{equation}

By Lemma \ref{arith:small}, we may first discard total products
$n<X^{1-\delta}$. Partition the first $r-1$ logarithms of the remaining
tuples into boxes of width $\delta$.
There are $O(\delta^{-(r-1)})$ boxes, and their exponent vectors are
within $O_r(\delta)$ of the simplex points $e(\boldsymbol a)$ above.
The fixed buffer in \eqref{arith:window} ensures that, in every box
meeting the target region, the specified subproduct remains in the
larger analytic window. This also covers a subproduct condition
normalized by $n$, since $1-\delta\le s_0<1$ and the two
normalizations differ by $O_r(\delta)$.
If more than one subproduct is available, partition according to a
fixed order on the finitely many proper subsets, obtaining finitely
many regions of the same type.

Before counting boundary boxes, resolve every affine expression
that is constant on $\sum t_i=1$.
A nonzero constant has a fixed sign in a sufficiently small
neighborhood of that simplex. If the constant is zero, the
fixed-$X$ expression is $\lambda(\sum t_i-1)$: its truth value is
determined by $n<X$ when $\lambda\ne0$, and by exact equality when
$\lambda=0$. This includes the homogeneous expressions in
\eqref{arith:homogenize}.
Every remaining affine hyperplane meets at most
$O(\delta^{-(r-2)})$ boxes.
The positive box bounds show that the total boundary mass is
$O(A\delta/L)+O_J(AL^{-J})$; summing \eqref{arith:boxerror} gives the
same bound for the interior errors.
This handles both strict and nonstrict inequalities, as well as the
ordering of the prime factors.
Letting $X\to\infty$ proves the proposition.
\end{proof}

\section{Base nine: full tuples and a positive sieve margin}\label{nine:sieve-section}

We first verify the Type II interval for all nine alphabets. We then
apply the common two-step positive tree and bound its marked simplex loss.

\subsection{Moment bounds and the Type II interval}\label{nine:fourier-section}

\begin{proposition}\label{nine:fourier}
For $b=9$ and every forbidden digit $d=0,\ldots,8$,
equations~\eqref{geo:moments} and~\eqref{geo:fractional} hold with the common choices
\begin{equation}\label{nine:moment-values}
 \sigma=\frac{127}{200},\qquad s=\frac32,\qquad
 g_s=\frac{6307}{40000},\qquad
 \kappa=\frac{6307}{20000},\qquad
 g_{4/3}=\frac{13607}{60000}.
\end{equation}
The first and $3/2$ moments both hold uniformly on grids with an arbitrary real shift.
\end{proposition}

The continuous operator inequalities are proved in
Appendix~\ref{app:nine-moments}. For the common-denominator estimate, the
useful interpolation is directly between the first and $3/2$ moments:
\[
 \sum_a F_X(a/X)^{4/3}
 \le\left(\sum_a F_X(a/X)\right)^{1/3}
      \left(\sum_a F_X(a/X)^{3/2}\right)^{2/3}.
\]
This gives $g_{4/3}=(73/200)/3+2(6307/40000)/3=13607/60000$,
which is the exponent used below.

\begin{proposition}[Choice of parameters in base nine]\label{nine:parameters}
For every forbidden digit in base nine, the common-denominator estimate holds
with $y=633/1000$. A Type II window is
\begin{equation}\label{nine:buffer}
 \alpha_0=\frac{6309}{20000},\qquad
 \beta_0=\frac{7383}{20000}.
\end{equation}
We choose the interior sieve parameters
\begin{equation}\label{nine:sieve-params}
 \alpha=\frac{631}{2000}=.3155,\quad
 \beta=\frac{3691}{10000}=.3691,\quad
 \eta=\beta-\alpha=\frac{67}{1250}=.0536.
\end{equation}
\end{proposition}

\begin{proof}
The integer comparison $8^{500}>9^{473}$ gives $c>473/500$.
The positive margins required by~\eqref{geo:three-powers} and the determinant
condition are, respectively,
\[
 \sigma+c/2-1>\frac{27}{250},\quad
 1-y-3g_{4/3}/2=\frac{1073}{40000},\quad
 2\sigma-2y=\frac1{250},\quad 2-3y=\frac{101}{1000}.
\]
Take $\delta_0=1/2000$ and $\omega=10^{-6}$. In \eqref{geo:block-final}
the sparse exponents satisfy $\rho=27/100$ and $\tau>27/125$;
the dense exponent saves $\delta_0(1-y)-y\omega=182867/10^9$.
Also $y-2\kappa=23/10000>0$, as required by \eqref{geo:fractional}.
For Theorem~\ref{geo:typeii}, take
$u_0=1477/4000$, $\varepsilon_0=10^{-6}$, and $\xi=10^{-5}$.
The corresponding six positive margins are
\begin{align*}
 1-u_0-2\kappa-2(s-1)\varepsilon_0&=49/10^6,&
 u_0-\beta_0&=1/10000,\\
 y-(1-u_0)&=9/4000,&
 2\alpha_0+u_0-1&=3/20000,\\
 2\alpha_0-2\kappa-(3s-4)\varepsilon_0-\xi&=379/2000000,&
 3\alpha_0+1-6\kappa-6(s-1)\varepsilon_0-2\xi&=54227/10^6.
\end{align*}
Also $1-\beta=.6309<.633<\sigma=.635$, leaving room for the Type I level.
The interval $[\alpha,\beta]$ lies strictly inside~\eqref{nine:buffer}.
\end{proof}

\subsection{The common positive tree at the base-nine window}
\label{nine:common-tree}

The arithmetic inputs now determine the window in the grouped-product sieve. Here $X=9^k$, $A=8^k$, $L=\log X$,
$\mathcal M=\gamma_d A/L$, and $\gamma_d=9/8$ if $3\mid d$ and $15/16$
otherwise, by \eqref{arith:localfactor}. The parameters are
\[
 (\alpha,\beta,\eta)=(.3155,.3691,.0536),\qquad
 d_*=.6309,\qquad c_*=.6845.
\]
The interval straddles $1/3$. Apply the affine-shape version of
Theorem~\ref{common:positive-sieve} with $T=.633$:
$1-\beta<T<\sigma$ and $1-T/2=.6835<c$. On $I_X=[XL^{-D},X)$ take
$D>3/c$. For the digital Type I hypothesis, fixing the unit last digit
leaves an interval and one residue class modulo each modulus coprime
to nine; apply \eqref{arith:typeI-max} at length $X/9$ with
$(1+\epsilon)T<\sigma$ and sum over the allowed unit digits. The mean
is the actual interval mass divided by the modulus. For the flat source
counting unit residues has error $O(A/X)$ per modulus, absorbed since
$T<1$. For the second hypothesis apply Proposition~\ref{arith:polytope}:
there is a buffer of $1/20000$ on both sides. All prime factors exceed
a fixed positive power of $X$, so the digital and flat weights on their
products coincide with those of the common theorem. Removing products
below $XL^{-D}$ costs $O(A L^{-Dc}+A L^{-D})=o(\mathcal M)$ by the
interval mass bound and bounded marking multiplicity. This gives the
required affine comparison.

Thus both inputs to the common sieve hold on the indicated window,
and Proposition~\ref{nine:certificate} bounds its loss.

\subsection{Prime counts at arbitrary real endpoints}

\begin{theorem}\label{nine:theorem}
For each $d\in\{0,\ldots,8\}$, using the usual base-nine expansion without
leading zeros, we have
\[
 \pi_{9,d}(x)\asymp\frac{A_{9,d}(x)}{\log x}
 \asymp\frac{x^{\log_9 8}}{\log x}.
\]
More explicitly, for all sufficiently large real $x$, uniformly in the
forbidden digit,
\begin{equation}\label{nine:explicit}
 \frac1{52}\frac{x^{\log_9 8}}{\log x}
 \le\pi_{9,d}(x)\le33\frac{x^{\log_9 8}}{\log x}.
\end{equation}
\end{theorem}

\begin{proof}
By the common positive-tree comparison in Section~\ref{nine:common-tree}
and Proposition~\ref{nine:certificate}, $(15/16)(1-\mathscr L)>(15/16)(.165796)>2/13$
and $\min_d\gamma_d=15/16$.
Hence every sufficiently large $X=9^k$ satisfies
\[
 \#\{p<X:p\in\mathcal W(X)\}\ge\frac2{13}\frac{8^k}{\log X}.
\]
For $9^k\le x<9^{k+1}$ this yields the lower constant $1/52$.
Proposition~\ref{common:upper} gives the upper limit $2\gamma_d/.633$
relative to $A_{9,d}(x)/\log x$. Since
$A_{9,d}(x)\le(8/7)8^{k+1}\le (64/7)x^c$ and $\gamma_d\le9/8$,
the upper limit relative to $x^c/\log x$ is less than $33$.
Lemma~\ref{prelim:count} supplies the comparison of normalizations.
\end{proof}

\clearpage

\section{Octal divisor estimates}
\label{sec:octal-typeI}

For the octal argument put $D_d=\{0,\ldots,7\}\setminus\{d\}$,
$X=8^k$, $L=\log X$, $A=7^k$, $\rho=A/X$, and
$\delta=1-\log_8 7$. Let $\mathcal A_X$ be the length-$k$ words
over $D_d$, identified with integers in $[0,X)$. Put $e(t)=e^{2\pi i t}$
and write the one-digit and word amplitudes as
\[
 f_d(t)=\frac17\left|\sum_{a\in D_d}e(at)\right|,
 \qquad
 F_X(t)=\frac1A\left|\sum_{n\in\mathcal A_X}e(nt)\right|
       =\prod_{j=0}^{k-1}f_d(8^jt).
\]
We use the normalized transfer operator
$\mathcal T_d g(u)=\tfrac18\sum_{j=0}^7 f_d((u+j)/8)g((u+j)/8)$.
Fix any $d\in\{0,\ldots,7\}$. Put $K_d=D_d\cap\{1,3,5,7\}$ and
write $\gamma_d=2|K_d|/7$. Thus $\gamma_d=8/7$ if $d$ is even and
$6/7$ if $d$ is odd. We work with $f=1_{\mathcal A_X}1_{\rm odd}$,
of mass $\gamma_d A/2$ on a full period. If $d=0$, every word in
$\mathcal A_X$ has exactly $k$ ordinary digits. If $d\ne0$, padding
with zeros represents every allowed integer in $[1,X)$ in $\mathcal A_X$.
On $\mathbb Z/X\mathbb Z$, let $E_{X,Q}f$ denote the conditional mean
of $f$ on classes modulo a dyadic divisor $Q$ of $X$.

\begin{lemma}[Uniform Type I input]\label{octal:uniformTypeI}
For every missing octal digit, the hypotheses of Proposition~\ref{arith:typeI}
hold with $\sigma=247/400$.
\end{lemma}
\begin{proof}
The supplement supplies positive periodic piecewise linear functions
$g_d$ with $4096$ rational nodes each (Appendix~\ref{app:reproduce}).
They satisfy
\[
\mathcal T_dg_d\le (2769/10000)g_d\qquad (d=0,1,2,3).
\]
On the refined grid of $n=16384$ nodes the verifier encloses cosines
using Machin's identity and the degree-$48$ Taylor polynomial with
remainder bounded by $4^{50}/50!$. The identity
\[
 49 f_d(t)^2=7+2\sum_{j=1}^7 C_j\cos(2\pi jt),\qquad
 C_j=8-j-1_{d+j\le 7}-1_{d-j\ge0},
\]
has nonnegative coefficients. Upper enclosures for the cosines and an
upward-rounded square root thus give an upper residual at each node.
Writing $\lambda=2769/10000$, the residual
$\mathcal T_d g_d-\lambda g_d$ has Lipschitz constant at most
\[
 \tfrac18\big(\operatorname{Lip}(f_d)\max g_d+\operatorname{Lip}(g_d)\big)
       +\lambda\operatorname{Lip}(g_d).
\]
Here $\operatorname{Lip}(f_d)\le 32\pi/7<704/49$ by centering the
polynomial at $7/2$, and $\operatorname{Lip}(g_d)$ is computed from its
linear segments. The verifier adds this constant divided by $2n$ to
the largest upper nodal residual, and checks that the result is negative.
The included certificate and verifier prove these inequalities with
integer and rational arithmetic. Reflection gives the other digits.
Since $2769^{400}8^{247}<10000^{400}$ and $g_d$ is bounded above and away
from zero, iteration gives $\mathcal T_d^k1\ll X^{-247/400}$.
Thus the shifted-grid first norm is $O(X^{.3825})$. Integration in
the shift variable gives the continuous first-moment estimate.
\end{proof}


\subsection{Odd divisors and a quadratic Fourier approximation}
\label{sec:projection-typeI}

In this section $X=8^k$, $A=7^k$, $\rho=A/X$, and
$\delta=1-\log_8 7<129/2000$. Write $a=1_{\mathcal A_X}$ on
$\mathbb Z/X\mathbb Z$, and normalize the Fourier transform by
$\widehat a(b)=X^{-1}\sum_{n<X}a(n)e(-bn/X)$.


\begin{lemma}[A moment near Parseval]\label{octal:near-parseval}
Uniformly in the omitted digit and the real shift $v$,
\[
 \sum_{b<X}F_X((b+v)/X)^{127/64}
 \ll (11481/10000)^k\ll X^{\delta+141/64000}.
\]
\end{lemma}
\begin{proof}
Let $s=127/64$ and let $T_s g(u)=\sum_{j=0}^7 f_d((u+j)/8)^sg((u+j)/8)$.
The data in \path{supplement/base8/witnesses.py} specify a positive piecewise
linear periodic function $g_d$, for each $d=0,1,2,3$, with $256$ rational
nodes. The verifier proves $T_s g_d\le\lambda g_d$ everywhere, where
$\lambda=11481/10000$. Here are the finite enclosures and the guard used.
On the refined grid of $32768$ nodes, bound $f_d^2$ above using the
cosine expansion in Lemma~\ref{octal:uniformTypeI} and directed rational
enclosures of cosines. If $u$ is an integer upper bound for $Sf_d^2$,
with $S=2^{90}$, the least nonnegative integer $w$ with
$w^{128}\ge S u^{127}$ is an upper bound for $Sf_d^s$.
The verifier finds it by directed integer square roots.
The rational nodal residual uses this upper weight and the exact linear
interpolation of $g_d$.

For the extension from nodes, $\operatorname{Lip}(f_d)<15$ by centering
at $7/2$. On $0\le x\le1$, the function $x^s$ has Lipschitz constant
at most $2$, and $x^s-x^2$ has Lipschitz constant at most $2(2-s)$.
For the latter, differentiate $t x^{t-1}$ in $t\in[s,2]$ and use
$\sup_{0<x\le1}x^{s-1}|\log x|=1/(e(s-1))$. As $T_2 1=8/7$ by
orthogonality, $T_s 1$ has Lipschitz constant at most $30(2-s)$.
For $T_s(g_d-1)$ use the branch scale $1/8$ and sum over eight branches.
The full residual $T_sg_d-\lambda g_d$ has Lipschitz constant at most
\[
 30(2-s)+30\|g_d-1\|_\infty+(1+\lambda)\operatorname{Lip}(g_d).
\]
After adding the guard of this constant divided by $2\cdot32768$,
the largest upper nodal residual is negative for each digit.
Reflection gives $d=4,5,6,7$. Iteration gives the first bound since
$T_s^k1(v)$ is the shifted grid moment and each $g_d$ is bounded above
and away from zero. The second bound follows from the integer check
$(7\cdot11481)^{64000}<8^{141}(8\cdot10000)^{64000}$.
\end{proof}

\begin{lemma}[Quadratic approximation]\label{octal:quad-approx}
Put $f(n)=a(n)1_{n\ \mathrm{odd}}$. There is a periodic function $G$ such that
\[
 \sum_{1\le n<X}|f(n)-G(n)|\ll A X^{-1/128000},\qquad
 \sum_{b<X}|\widehat G(b)|\le \rho X^{697/2000}.
\]
\end{lemma}
\begin{proof}
Keep only the frequencies $b$ with $F_X(b/X)\ge X^{-142/1000}$ in
$a$ and denote the resulting Fourier projection by $P$.
The discarded squared mass of the normalized digital modulus is
\[
 \sum_{b\ \mathrm{discarded}} F_X(b/X)^2
 \ll X^{-142/64000}X^{\delta+141/64000}=X^{\delta-1/64000}.
\]
Parseval gives total squared mass $X^\delta$, so the number of kept
frequencies is at most $X^{\delta+284/1000}<X^{697/2000}$.
As $|\widehat a(b)|=\rho F_X(b/X)$, in probability normalization
$\|P-a\|_2\ll\rho^{1/2}X^{-1/128000}$ and $\|P\|_2\le\|a\|_2=\rho^{1/2}$.
Take $G=1_{\mathrm{odd}}|P|^2$ and use $|a|^2=a$ and Cauchy--Schwarz for
the first bound. For the second, $\|\widehat P\|_1^2\le X^{697/2000}\rho$
by Cauchy--Schwarz on its support. The Fourier transform of $|P|^2$
is a convolution, and multiplication by $1_{\mathrm{odd}}$ costs a
Fourier $\ell^1$ norm of exactly one.
\end{proof}

\begin{theorem}[Type I for octal words and their periodic projections]
\label{octal:lifted-odd-divisor-marks}\label{octal:strong-typeI}
Fix $\zeta>0$, $K,J>0$, and $\theta<.6515$.
Let $X^\zeta\le Q\mid X$ be dyadic and put $F_Q=E_{X,Q}f$.
Then, uniformly in $Q$,
\[
 \sum_{\substack{e\le X^\theta\\e\,\mathrm{odd}}}\tau(e)^K
    \sup_I\left|\sum_{n\in I\cap[1,X)} F_Q(n)
                       (1_{e\mid n}-1/e)\right|\ll A L^{-J}.
\]
Here $I$ ranges over intervals. The case $Q=X$ gives the estimate for $f$.
\end{theorem}

\begin{proof}
Use the Ramanujan expansion of the centered mark. For primitive
reduced denominator $s>1$, its total divisor weight is at most
$O(\tau(s)^K L^{O_K(1)}/s)$ by summing over multiples of $s$.
For $s\le Q^{.49}$ expand
$c_s(n)=\sum_{t\mid s}\mu(s/t)t(1_{t\mid n}-1/t)$.
A second harmonic summation leaves at most a logarithmic power
times the corresponding sum of centered marks at odd $t\le Q^{.49}$
with weights $\tau(t)^K$.
Let $Y$ be the least octal power divisible by $Q$. On each
$Y$-block $F_Q$ is a density-scaled average of at most four
translates of the word at length $Y$. Each translated block
splits into at most two intervals at its cyclic endpoint, and
each mark specifies one residue class on each interval.
Slice the terminal odd digit into its allowed classes and use
the maximum-residue estimate \eqref{arith:typeI-max} at
length $Y/8$, which applies since $Q^{.49}<(Y/8)^{1/2}$ for
large $X$. The sum of these errors over the blocks has scale $A$
and arbitrary logarithmic savings.

For larger $s$ use the quadratic approximation of the word
from Lemma~\ref{octal:quad-approx}, then average and lift
from length $Y$ as above. Its counting first-norm difference
from $F_Q$ is $O(AQ^{-\epsilon_0})$ for some fixed $\epsilon_0>0$.
Its probability-normalized Fourier first norm is at most
$O((A/X)Q^{.3485})$, since translating changes Fourier phases, averaging modulo $Q$ deletes
frequencies, and repeating a period inserts zeros in its normalized
Fourier array. The density of the digits above $Y$ multiplies both norms. Substitution in the large denominator
sum is justified before taking Fourier absolute values by
\[
 \sum_{s\le X^\theta}\frac{\tau(s)^K}{s}|c_s(n)|
    \ll \tau(n)^{K+1}L^{O_K(1)}\qquad (1\le n<X).
\]
This follows by expanding $c_s$ in divisors and summing over
$s$ divisible by each divisor of $n$. The pointwise divisor
bound absorbs the substitution error by a fixed power.

In a dyadic range $s\asymp S$, the distinct primitive fractions
have spacing $\gg S^{-2}$. Summing the geometric bound for the
interval test against any one Fourier frequency costs
$O(X+S^2 L)$ over these fractions. Including the divisor weight
and the Fourier first norm of the approximation, the cost
relative to $A$ is at most
\[
 L^{O(1)} Q^{.3485} (1/S+S/X) S^{\epsilon_1}
\]
for any fixed $\epsilon_1>0$. Both terms save a fixed power
since $S\ge Q^{.49}/2$, $Q\ge X^\zeta$, and
$\theta+.3485<1$. Summing the dyadic ranges proves the lemma.
\end{proof}


\paragraph{The upper estimate.}\label{sec:upper-all}
Taking $T=.640$ in Proposition~\ref{common:upper}, the number
$P_d(X)$ of prime words satisfies
\begin{equation}\label{octal:word-upper}
 P_d(X)\le \bigl((25/8)\gamma_d+o(1)\bigr) A/L .
\end{equation}
The same constant bounds the upper limit at ordinary endpoints,
now relative to $A_{8,d}(x)/\log x$.

For $q\mid X$ dyadic, $q\ge4$, set
$\mathcal D_q f=E_{X,q}f-E_{X,q/2}f$.
The flat reference is $f_0=E_{X,2}f=\gamma_d(A/X)1_{\rm odd}$.
These sources are periodic in Fourier calculations and restricted to
$[1,X)$ in the sieve. Put $c=\log_8 7$, $D>3/c$ fixed and
$I_X=[XL^{-D},X)$. Choose $0\le\omega\le1$ smooth, supported on
$[L^{-D},1]$, equal to one on $[2L^{-D},1-L^{-D}]$, with each fixed
derivative bounded by a power of $L$.

\section{A uniform octal prime count}
\label{sec:octal-prime-count}

\begin{lemma}[Marked octal prime shapes]\label{octal:balanced-shape-comparison}
Fix the rank, a positive roughness exponent, and a labelled proper group.
On the closure approached from $\sum t_i<1$ at $\sum t_i=1$, let the group's exponent lie in a compact
subinterval of $(.399,.601)$. For each affine family as in
\eqref{dt:sieve-flat-comparison}, including order restrictions and a
prescribed choice of marks, that comparison holds with $b=f-f_0$ on
$I_X$ and $\mathcal M=\gamma_d A/L$. It holds also after multiplication
by a fixed smooth function of the exponents.
\end{lemma}

\begin{theorem}\label{octal:prime-count}
For every omitted octal digit and sufficiently large real $x$,
\[
 .09967\frac{x^{\log_8 7}}{\log x}\le\pi_{8,d}(x)
 \le30\frac{x^{\log_8 7}}{\log x}.
\]
For $d\ne0$ the upper constant may be replaced by $26$.
\end{theorem}
\begin{proof}
Use Theorem~\ref{common:positive-sieve} with radix $8$, alphabet size
$7$, $u=|K_d|$, $T=.640$ and $[\alpha,\beta]=[2/5,1/2]$.
Indeed $1-\beta<T<.6515$ and $1-T/2=.680<c$.
Theorem~\ref{octal:strong-typeI} with $Q=X$ gives the interval
remainder with the actual interval mass, and odd multiples give the
flat remainder $O(A/X)$ per odd modulus. The marked-shape lemma supplies every compact interior group interval. The endpoint assertion of the common sieve and the triangle
calculation in Section~\ref{cost:octal} give the lower limit
$(\gamma_d/7)(1-\mathscr L_8)>.09967$, since $\gamma_d\ge6/7$.
For $8^k\le x<8^{k+1}$ we have $A_{8,d}(x)\le (7/6)7^{k+1}$
for all $d$, and $A_{8,d}(x)\le7^{k+1}$ when zero is allowed.
Thus the corresponding upper limits, in units of $x^c/\log x$,
are at most $(25/8)(8/7)(49/6)<30$ and $25$, respectively.
\end{proof}


\section{Fourier estimates for the octal prime comparison}
\label{sec:octal-fourier}

\begin{lemma}[Digital coefficients for dyadic characters]
\label{octal:semiprime-digital-char}
Let $q\mid X$ be a power of two. Write \(c=\log_8 7\), \(\sigma=.6175\), and \(\lambda=(c+\sigma-1)/2\). For
\(A_\chi=\sum_{n\bmod X}f(n)\overline{\chi(n)}\) one has
\begin{equation}\label{octal:semiprime-digital-moments}
 |A_\chi|\ll A q^{-\lambda}\quad(\chi\ {\rm primitive}\bmod q),
 \qquad
 \sum_{\chi\bmod q}|A_\chi|^2\ll A^2q^{1-c}.
\end{equation}
\end{lemma}

\begin{proof}
For the pointwise bound we may assume \(q\) exceeds an absolute
constant. Let \(Z\) be the least octal power divisible by \(q\),
let \(R\) be the least octal power at least \(\sqrt q\), and put
\(Q=q/R\) and \(V=Z/R\). Then \(R\mid q\), \(Q\mid R\),
\(Q\le V\le4Q\), and \(R,Q\asymp\sqrt q\). On \(1+R\mathbb Z\bmod q\)
the character is linear in the additive coordinate modulo \(Q\),
since \(R^2\equiv0\pmod q\). Primitivity gives
\(\chi(1+q/2)=-1\), so its frequency \(\ell\) is odd. With a
suitable sign for \(\ell\),
\[
 \overline{\chi(u+Rv)}=\overline{\chi(u)}e(\ell\bar u v/Q)
 \quad(u\ {\rm odd}).
\]
The lower block \(u\) of length \(R\) has its last digit in $K_d$.
The independent upper block \(v\) has length \(V\) and uses $D_d$
at every digit. Each residue of \(u\bmod Q\)
has \(O((R/Q)^c)\) allowed lifts, by fixing the digits below the
largest octal power dividing \(Q\). Inversion followed by
multiplication by \(\ell\) permutes the units modulo \(Q\).
The sum over \(v\) has absolute value \(V^cG_V(\ell\bar u/Q)\),
where \(G_V\) is the normalized absolute word transform.
Lemma~\ref{octal:uniformTypeI} and sampling give
\(\sum_{b\bmod Q}G_V(b/Q)\ll Q^{1-\sigma}\). Indeed \eqref{arith:sampling} gives the bound $(Q+V)\int G_V$
and $V\asymp Q$.
Digits above \(Z\) contribute the factor \((X/Z)^c\), giving
\[
 |A_\chi|\ll (X/Z)^c V^c(R/Q)^c Q^{1-\sigma}
       \ll A q^{(1-\sigma-c)/2}.
\]
For the second estimate of \eqref{octal:semiprime-digital-moments},
each residue of \(n\bmod q\) has \(O((X/q)^c)\) allowed lifts.
The sum of all residue counts is \(O(A)\). By orthogonality the
character square sum is \(\varphi(q)\) times the sum of the squares
of the unit residue counts, which proves the result.
\end{proof}


For odd \(n\) and \(q\ge4\), the residue class averages also coincide
with conditional expectations on the odd unit group. Their
difference retains precisely the characters of conductor \(q\).
Consequently
\begin{equation}\label{octal:dyadic-char-expansion}
 \mathcal D_q f(n)=\frac2X
      \sum_{\substack{\chi\bmod q\\\chi\ {\rm primitive}}}
                     A_\chi\chi(n).
\end{equation}
Indeed the multiplicative coefficient at \(\chi\) in \(E_{X,q}f\)
is \(q A_\chi/(X\varphi(q))=2A_\chi/X\).


\subsection{A principal model for the long group}
\label{sec:principal-grouped-models}

We will use the interval Type~I estimate for $f$ and
$F_Q=E_{X,Q}f$ with $Q\mid X$ dyadic and $Q\ge X^\zeta$ for a fixed
$\zeta>0$. These are odd-supported, nonnegative, periodic sources of
mass at most $A$ on one full period of length $X$. We use the estimate
also on subintervals of $[1,C_0X]$ for fixed $C_0$. For $F_Q$ the proof
of Theorem~\ref{octal:strong-typeI} gives this version
directly. In the small denominator estimate there are a bounded
multiple of the original number of octal blocks. In the large
denominator estimate the geometric interval bound is
$O_{C_0}(X+S^2L)$ and the pointwise Ramanujan bound applies to all
$1\le n\le C_0X$. The $\ell^1$ approximation error increases by only
a bounded factor from the additional periods. This also gives the
stated interval estimate for $f=F_X$. By linearity it applies to
$g=f-F_Q$ and to $g=\mathcal D_Yf=F_Y-F_{Y/2}$ if both periods in
the latter expression exceed $X^\zeta$.

Let a fixed number of prime factors of scales $P_i\gg X^\kappa$ be
partitioned into two nonempty groups $I_1,I_2$, where $\kappa>0$ is
fixed. Put $U_h=\prod_{i\in I_h}P_i$ and suppose that
$X L^{-C}\ll U_1U_2\ll X$ and $U_1\le U_2\ll X^{.601}$. Let each $w_i$ be
a smooth function supported in a fixed compact subset of $(0,\infty)$,
with every fixed derivative bounded by a power of $L$. The actual
coefficient at the grouped integer $v$ is
\[
 b_h(v)=\sum_{\substack{\prod_{i\in I_h}p_i=v\\ p_i\ \mathrm{prime}}}
                 \prod_{i\in I_h}(\log p_i)w_i(p_i/P_i).
\]
To count unweighted primes, include $1/\log(P_i t)$ in $w_i(t)$.
Choose one slot $i_h\in I_h$ and define
\[
 \Phi_h(v)=\int_{(0,\infty)^{I_h\setminus\{i_h\}}}
   w_{i_h}\left(\frac{v}{P_{i_h}\prod_{i\ne i_h}t_i}\right)
   \prod_{i\ne i_h}w_i(t_i/P_i)\,\frac{dt_i}{t_i},
\]
where the products in this formula are within group $h$. The empty
integral is the single-factor test. Write $c_d(v)$ for the Ramanujan
sum $\sum_{a\bmod d}^{*}e(av/d)$, with $c_1(v)=1$. For $D=L^{K_0}$,
where $K_0>0$ is fixed, put
\begin{equation}\label{eq:principal-grouped-coefficients}
 B_2(v)=2\,1_{v\ \mathrm{odd}}\,\Phi_2(v)
       \sum_{\substack{d\le D\\ d\ \mathrm{odd}}}
                   \frac{\mu(d)}{\varphi(d)}c_d(v),
 \qquad b_2^*(v)=b_2(v)-B_2(v).
\end{equation}
Here $\Phi_h(v)\,dv$ is the image under multiplication of
$\prod_{i\in I_h} w_i(t_i/P_i)\,dt_i$: at fixed companions the
Jacobian in the selected slot is $1/\prod_{i\ne i_h}t_i$.
Thus $\Phi_h$ does not depend on which slot is selected. All coefficients
remain supported on fixed enlargements of their group intervals.

\begin{proposition}\label{octal:principal-grouped-model}
Let $g$ be either of the signed sources above, and let $\Omega$ be a
smooth product test with fixed relative derivatives bounded by powers
of $L$. The constant test $\Omega=1$ is permitted. Then
\begin{align}
 \sum_{x,y} b_1(x)b_2(y)\Omega(xy/X)g(xy)
 &=\sum_{x,y} b_1(x)b_2^*(y)\Omega(xy/X)g(xy)\notag\\
 &\quad+2\sum_{n>0}g(n)\Omega(n/X)
                 \int_0^\infty \Phi_1(z)\Phi_2(n/z)\frac{dz}{z}
       +O_N(AL^{-N})                                           \label{eq:principal-grouped-representation}
\end{align}
for every fixed $N>0$. The squared counting norms of $b_1,b_2^*$ are
$O(U_1L^C),O(U_2L^C)$, respectively, and both coefficients are
$O(\tau(v)L^C)$. The exponent $C$ can be chosen independently of $K_0$.

For each fixed $K_1,N>0$, the coefficients satisfy the following
small-arc estimate uniformly in $(a,b)=1$, $1\le b\le D$, and
$U_2|\beta|\le L^{K_1}$:
\begin{equation}\label{eq:principal-grouped-low-arc}
 \left|\sum_v b_2^*(v)e\bigl((a/b+\beta)v\bigr)\right|
                    \ll_{K_0,K_1,N} U_2L^{-N}.
\end{equation}
The estimates include retained Mellin twists of polylogarithmic
height in the individual tests.
\end{proposition}

\begin{proof}
The bounds $|b_h(v)|\ll L^C$ follow from unique factorization
with boundedly many slots. The function $\Phi_h$ is supported at
scale $U_h$ and has derivative of order $j$ bounded by
$O_j(U_h^{-j}L^{C_j})$. For $B_2$ expand
$c_d(v)=\sum_{b\mid(d,v)}b\mu(d/b)$. Only squarefree $d$ occur,
so $\varphi(d)=\varphi(b)\varphi(d/b)$ for $b\mid d$. Hence
\[
 \left|\sum_{\substack{d\le D\\ d\ \mathrm{odd}}}
            \frac{\mu(d)c_d(v)}{\varphi(d)}\right|
 \le \sum_{b\mid v}\frac{b}{\varphi(b)}\sum_{e\le D}\frac1{\varphi(e)}
 \ll \tau(v)(1+\log v)(1+\log D).
\]
The last inequality follows from
$m/\varphi(m)=\sum_{s\mid m}\mu(s)^2/\varphi(s)$ and
$\sum_s\mu(s)^2/(s\varphi(s))<\infty$. Since
$1+\log D\ll_{K_0}L$, this proves the pointwise bound with a log
exponent independent of $K_0$. Also $\tau(v)^2\le\tau_4(v)$, by
checking prime powers, and $\sum_{v\le V}\tau_4(v)\ll V(1+\log V)^3$.
This gives the asserted squared norm.

The parity normalization can be read off before any estimates: $x$ and
the support of $g(n)$ are odd, so for $x\mid n$ the indicator of odd
$n/x$ is identically one, but the factor two in $B_2$ remains.
In the pairing of $b_1$ and $B_2$, hold $x$ fixed and expand the
Ramanujan sum at modulus $d$. The terms indexed by $b\mid d$ are
\[
 2b_1(x)\frac{\mu(d)b\mu(d/b)}{\varphi(d)}
     \sum_{n>0} g(n)\Omega(n/X)\Phi_2(n/x)1_{xb\mid n}.
\]
For each fixed $d,b$, replace the mark $1_{xb\mid n}$ by
$1/(xb)$ using interval Type~I and partial summation. The mark
is odd and $xb\ll U_1D$ has fixed reserve below the Type~I level.
The total variation and supremum of the real-variable test in $n$
are bounded by powers of $L$, uniformly in $x$. Indeed the relevant
product scale is between $X$ and $XL^{-C}$ up to constants, and
$x\asymp U_1$. The map $x\mapsto xb$ is injective for fixed $b$,
and $|b_1(x)|\ll\tau(xb)L^C$. Summing over the polylogarithmically
many choices of $d,b$ costs only powers of $L$, so the total
replacement error is $O_N(AL^{-N})$ for every $N$.
For $d>1$ the resulting main term is zero because
$\sum_{b\mid d}\mu(d/b)=0$.

For $d=1$ the pairing in which the other coefficient is genuine
becomes
\[
 2\sum_n g(n)\Omega(n/X)\sum_x \frac{b_1(x)}{x}\Phi_2(n/x).
\]
The ordinary prime number theorem and partial summation in the
prime slots give, uniformly in $n$, up to $O_N(L^{-N})$, the
continuous reciprocal sum
$\int_0^\infty\Phi_1(z)\Phi_2(n/z)\,dz/z$. To see the formula,
replace $(\log p_i)/p_i$ by $dt_i/t_i$ in each slot and change
variables from $t_{i_1}$ to the group product $z$, fixing the
other slots. The absolute mass of $g$ is $O(A)$ on the product
support. Inserting the continuous reciprocal sum into the pairing
with $B_2$ proves \eqref{eq:principal-grouped-representation}.

To prove \eqref{eq:principal-grouped-low-arc}, fix the companion
primes in group $2$, let their product be $c$, and abbreviate
$i_2=i$ and $P_i=P$. They are coprime to $b$ for large $X$.
The selected prime variable ranges through reduced classes modulo
$b$, since its scale is a fixed positive power of $X$. The
Siegel--Walfisz theorem and partial summation give
\begin{align*}
 \sum_p(\log p)w_i(p/P)e(acp/b+\beta cp)
 &=\frac1{\varphi(b)}\sum_{u\bmod b}^{*} e(acu/b)
               \int_0^\infty w_i(t/P)e(\beta ct)\,dt
       +O_N(PL^{-N})\\
 &=\frac{\mu(b)}{\varphi(b)}\int_0^\infty
                       w_i(t/P)e(\beta ct)\,dt+O_N(PL^{-N}).
\end{align*}
Here the modulus is polylogarithmic and the variation of the
selected-prime test is a power of $L$, since $|\beta|cP\ll
L^{K_1}$. After summing the companions the total error is
$O_N(U_2L^{-N})$. In the main term put $v=ct$ and use the
ordinary prime number theorem and partial summation in the
companion slots. The result for the genuine coefficient is
\begin{equation}\label{eq:principal-grouped-genuine-arc}
 \sum_v b_2(v)e((a/b+\beta)v)
     =\frac{\mu(b)}{\varphi(b)}\int_0^\infty
                         \Phi_2(v)e(\beta v)\,dv+O_N(U_2L^{-N}).
\end{equation}

For the model, use $2\,1_{v\ \mathrm{odd}}=1-e(v/2)$ and the
defining exponential sum for $c_d$. A rational phase can pair
with $e(av/b)$ to give an integer frequency only when its
reduced denominator matches. For odd $b$ this occurs in the
$1$ part exactly at $d=b$, with coefficient $\mu(b)/\varphi(b)$.
For $b=2d$ with $d$ odd it occurs only in the $-e(v/2)$ part,
with coefficient $-\mu(d)/\varphi(d)=\mu(b)/\varphi(b)$.
For $4\mid b$ there is no such phase and $\mu(b)=0$. In each
resonant case the required odd denominator is at most $D$.
Every nonresonant rational phase has distance $\gg D^{-2}$
from an integer, which is unchanged up to constants on adding
$\beta$. Apply Poisson summation to the smooth $\Phi_2$. Its
Fourier transform at frequency $\xi$ is
$O_j(U_2L^{C_j}(1+U_2|\xi|)^{-j})$ for every fixed $j$.
Since $U_2$ is a fixed positive power and all model denominators
are polylogarithmic, the nonresonant contributions and all
noncentral Poisson images of the resonant term are $O_N(U_2L^{-N})$.
This proves \eqref{eq:principal-grouped-genuine-arc} also for
$B_2(v)$ in place of $b_2(v)$. Their difference is
\eqref{eq:principal-grouped-low-arc}.
\end{proof}

\paragraph{Dependence on the shape and the product cutoff.}
For a fixed smooth shape depending on the normalized prime exponents,
choose a smooth partition over the eligible labelled groupings and
insert dyadic partitions in the prime slots. On each box one can
separate the resulting shape and the product cutoff by Fourier
inversion in the scaled logarithms of the prime variables. Smooth
compact extensions of these tests have fixed-order derivatives
bounded by powers of $L$. Their Fourier first norms are
polylogarithmic, and the integral can be truncated at a sufficiently
large fixed power of $L$ with arbitrary inverse-logarithmic error.
The error in the original pairing is bounded using the bounded
fixed-rank representation multiplicity and the $O(A)$ absolute
source mass. At each retained twist put the phases into the genuine
slot tests before defining \eqref{eq:principal-grouped-coefficients}.
Then \eqref{eq:principal-grouped-representation} applies with
$\Omega=1$. Its continuous term recombines by linearity and Fourier
inversion, with the same truncation estimate for the continuous
integral. This term is independent of the choice of group split
and of the prime slots selected to define $\Phi_h$.

More explicitly, for a fixed shape weight $W$ in $r$ labelled
prime slots and the product weight $\omega(n/X)$, let $D_W(n)$ denote
the recombined pushforward density. The term in
\eqref{eq:principal-grouped-representation} is $2\sum_n g(n)D_W(n)$, where
$D_W(n)$ is
\[
 \frac{\omega(n/X)}{\log n}
    \int_{u_1+\cdots+u_r=1} W(u_1,\ldots,u_r)
                        \prod_{i=1}^r\frac1{u_i}\,du_1\cdots du_{r-1}.
\]
Denote this simplex integral by $C_W$. Indeed use unweighted prime densities $dt_i/\log t_i$, fix the
first $r-1$ primes, and set $t_i=n^{u_i}$. Here the roughness
restriction is fixed, so the simplex integral is finite. Use a smooth partition of unity for the box boundaries. Sharp roughness and other affine cutoffs are imposed after recombination by the envelopes below.

For the whole source $g=f-E_{X,Q_*}f$ with
$Q_*=2^{\lfloor .741\log_2 X\rfloor}$, the original and projected
word have equal full mass. On discarded endpoint intervals the
original word contributes $o(A)$ by the digit interval bound.
The projected word has mass $O(AQ_*/X)$ per period, so an interval
of length $o(X)$ contributes $o(A)$ for that source as well.
Replacing $1/\log n$ by $1/L$ on the product support incurs
$o(A/L)$ there. Thus the entire continuous term for this shape
is $o(A/L)$ after restoring the product weight. The residual
pairing may be split into the exact conductors $Y>Q_*$, and the
small-arc estimate for their actual residual coefficients is
\eqref{eq:principal-grouped-low-arc}.


\subsection{Digital weights in the lower-period affine equation}
\label{sec:ordinary-group-gram-twists}

Let $q$ be the exact dyadic period of a lower word, $\rho_L$ its
density scale, and $Q\mid q$ an octal power. Write $a=.3825$.
The off-diagonal lower Gram, after completion in the long variable,
involves sums over $i,b\bmod q$ with $pi-p'b\equiv t\pmod q$.
Its Fourier indices and rows are odd for $q>2$.
Split each normalized digital amplitude at the upper octal block of
length $Q$, taking the least octal multiple of $q$ for the word length.
On residue classes $i_0,b_0\bmod Q$ the upper factors $U(i_0),V(b_0)$
are fixed. Each lift of $i_0$ determines a lift of $b_0$ by an odd
affine permutation. Shifted Cauchy--Schwarz in the lower blocks
therefore costs $O((q/Q)^\delta)$, with bounded aliases when $q$ is
non-octal. The remaining upper weights obey
\[
 \sum U,\ \sum V\ll Q^a,\qquad \sum U^2,\ \sum V^2\ll Q^\delta
\]
by the shifted moments. A bounded mask on the original amplitudes
can be ignored in these upper bounds.

Write $s=(i_0-b_0)\overline{b_0}\bmod Q$ and
$p'=p+j$. The affine equation is $ps-j\equiv t\overline{b_0}\pmod Q$.
The total Fourier weight on a slope $s$ is at most
$B(s)=\sum_{b_0\bmod Q}^{*}U((1+s)b_0)V(b_0)$, set to zero when
$1+s$ is even. For every octal divisor $Q_1\mid Q$, split the upper
blocks again to obtain
\begin{equation}\label{eq:ordinary-group-gram-overlap}
 B(s)\ll (Q/Q_1)^\delta B_{Q_1}(s\bmod Q_1),\qquad
       \sum_{\sigma\bmod Q_1}B_{Q_1}(\sigma)\ll Q_1^{2a}.
\end{equation}
Here $B_{Q_1}$ is the overlap of the two top subwords of length $Q_1$.
For each $b_0$ class modulo $Q_1$, multiplication by $1+s$ permutes
the lifts. Shifted Cauchy--Schwarz costs $(Q/Q_1)^\delta$ on the lower
blocks and the two first moments bound the sum of the top overlaps.


\paragraph{Sampling on separated cells.}\label{sec:zero-model-wide-source-bands}
For the amplitude $F_N$ of a normalized trigonometric polynomial of degree less than $N$,
let $I_j$ have length comparable to $N^{-1}$ and bounded overlap.
For $s\ge1$, the integral and variation on each cell give
\[
 \sum_j\sup_{t\in I_j}|F_N(t)|^s
 \ll N\int_0^1 |F_N|^s+\int_0^1\big|(|F_N|^s)'\big|
 \ll N\int_0^1 |F_N|^s .
\]
The last step is Bernstein's inequality for the underlying
trigonometric polynomial and H\"older. Thus a shifted counting
moment bound $O(N^{g_s})$ also bounds these sums of suprema by
$O(N^{g_s})$, on integrating the shifted moment. The argument covers
mixed words and cells with a bounded number of dyadic aliases.

\subsection{Upper digit factors and horizontal slopes}
\label{sec:short-row-weighted-offsets}

Consider an integer bilinear sum at group scales $P,M$, with $PM\asymp X$
up to powers of $L$, with $PM\ll X$. After separating smooth product weights, write its
source in the form $F(xy)U(\lfloor xy/R\rfloor)$, where $F$ is periodic
and the upper word $U$ satisfies
\[
 \sum_{|v|\le C X/R}|U(v)|^2\ll (X/R)\rho_U L^{O(1)}
\]
for every fixed $C$. The support intervals of the two coefficient
sequences are at fixed positive multiples of $P$ and $M$. In the
application, $U$ is the exact difference of two upper digital words
or its periodic lift, and $\rho_U$ is the original upper density.

\begin{lemma}\label{octal:short-row-upper-density}
Partition the row variable into intervals $I$ of integer length $H$,
where $H\le P$ and $HM\ll R$. If $\sum_y|v_y|^2\ll M L^{O(1)}$, then
\begin{align*}
 &\left|\sum_{x,y}u_xv_y F(xy)U(\lfloor xy/R\rfloor)\right|^2\\
 &\qquad\ll \frac{PM\rho_U}{H}L^{O(1)}
    \sum_{I,y}\max_{J\subseteq I}
                       \left|\sum_{x\in J}u_x F(xy)\right|^2,
\end{align*}
where the maximum is over subintervals. The same estimate holds if
the shorter scale is assigned to either coordinate.
\end{lemma}

\begin{proof}
For each $I,y$, the index of $U$ changes at most a fixed number of
times as $x$ traverses $I$. The modulus of the corresponding row
sum is therefore bounded by the indicated maximum times a fixed
multiple of the largest $|U|$ at one of those indices. Denote this
largest value by $A_I(y)$. For fixed $y$, the indices at the left
endpoints of $I$ increase by $Hy/R$ before taking the integer part.
Each integer index, including any of its bounded neighboring indices,
is used at most $O(1+R/(Hy))=O(R/(HM))$ times. They all lie within
$O(X/R)$ of the origin. Thus
\[
 \sum_I A_I(y)^2\ll \frac{P}{H}\rho_U L^{O(1)}.
\]
Apply Cauchy to the sum over $(I,y)$, using the factors
$|v_y|A_I(y)$ and the indicated maximum. This proves the inequality.
The argument counts indices over the row intervals for a fixed column,
so it does not require the individual row centres to be below $R$.
\end{proof}

To estimate the maximum on each $I$, decompose any $J\subseteq I$ into
at most two binary intervals at each length. Cauchy bounds the square
by a logarithmic factor times the sum of squares over all binary
intervals in $I$. At each length they are disjoint, and every row
pair they contain has difference of magnitude at most $H$. Thus the
usual ordinary completion reduces to the lower row Gram of
Section~\ref{sec:ordinary-group-gram-twists}, with this shorter
range of row differences. In particular its squared normalization
loses only one factor $\rho_U^{-1}$.

\begin{lemma}[Horizontal frequencies with their exact multiplicities]
\label{octal:horizontal-weighted-offset}
In a lower row Gram with a coarse octal modulus $Q$ comparable to
the lower period, let $g$ be the dyadic additive order of a horizontal
slope and $q_g=Q/g$. Take $g\ll P$, and keep row differences
$0<|j|\ll H$ and ordinary completion indices $|t|\ll T$. Let $B$
be a dyadic shell of reduced rational denominators associated to
$j b_0+t=0\pmod{q_g}$. If $B^2\max(1,T/B)\le c_*q_g$ for a sufficiently small absolute
constant $c_*$, then this shell costs at most
\begin{equation}\label{eq:horizontal-weighted-offset}
 g^{2a-1}B^{2\delta-1}\max(1,T/H)^\delta L^{O(1)}
\end{equation}
in the absolute lower Gram, after dividing by $PH$ and the square
of the lower density. Here $a=.3825$ and $\delta=1-\log_8 7$. Bounded common masks on the lower
additive coefficients are allowed. In the complementary range the
bound is
\begin{equation}\label{eq:horizontal-weighted-overflow}
 g^{2a-1}q_g^\delta(1/B+T/q_g)L^{O(1)}.
\end{equation}
\end{lemma}

\begin{proof}
We describe the horizontal reduction to record its counting factors.
Let $b$ be the second lower Fourier index, odd modulo $Q$, and let
$s$ be the slope, so that the first index is $(1+s)b$. In the affine
row equation $ps-j=t\bar b\pmod Q$ write $q_g=Q/g$, with $s/q_g$
a unit modulo $g$. Fixing $s,b,j,t$, there are $O(P/g)$ possible
rows since the equation fixes the row modulo $g$ whenever it is
soluble. Necessarily $jb_0+t=0\pmod{q_g}$ for $b_0=b\bmod q_g$.
For each $b_0$, sum the low subword amplitudes over both sets of
lifts modulo $Q$. For fixed $b$ the frequencies $(1+s)b$ run through
a subset of the lifts of $b_0$ once. The two shifted first norms
cost $g^{2a}L^{O(1)}$. The remaining normalized amplitudes, say
$V_1(b_0)V_2(b_0)$, are those of the upper subwords at modulus
$q_g$. Bounded aliases from the lower period divided by $Q$
are harmless by the shifted second moment. Dividing by $PH$
gives the bound
\[
 \frac{g^{2a-1}}{H}L^{O(1)}
 \sum_{b_0\bmod q_g}^{*} V_1(b_0)V_2(b_0)
       \sum_{\substack{0<|j|\ll H,\ |t|\ll T}}
                    1_{j b_0+t=0\ (q_g)}.
\]

Fix one reduced rational $a_0/B_0$, $B_0\asymp B$, obtained by
reducing $(j b_0+t)/(j q_g)$ as a real rational. Then $B_0\mid j$.
Writing $D_0=|b_0-q_g a_0/B_0|$, the exact identity
$|j|D_0=|t|$ shows that the number of possible $j$, divided by $H$,
is at most
\begin{equation}\label{eq:horizontal-count-per-offset}
 O\left(\frac1B \min\left(1,\frac{T}{H D_0}\right)\right).
\end{equation}
The statement includes $D_0=0$, by taking the minimum to be one.
For each such $j$ the value of $t$ is fixed by the real rational.
A nonempty band has $B\ll H$ and $D_0\ll T/B$.

If $D_0\le V$, the total $V_1V_2$ mass in the bands over the
denominator shell is at most
\[
                  O\bigl((B^2\max(1,V))^\delta L^{O(1)}\bigr).
\]
To prove this, split the word at a bottom block of length
comparable to $B^2$ and a terminal block of length comparable to
$\max(1,V)$. The rational centres are separated by $\gg B^{-2}$
on the circle, with only bounded repetitions from the real lifts
in this range. The sum of local suprema of the squared bottom
amplitude is bounded by its shifted square norm at this density,
using Bernstein's inequality as in
Section~\ref{sec:zero-model-wide-source-bands}. On the offset grid
the terminal block has square sum $O(\max(1,V)^\delta)$ by shifted
Parseval. Cauchy gives the claim for the two possibly different
amplitudes. Dyadic $q_g$ differing from an octal power by a factor
at most four gives only bounded aliases.

Use \eqref{eq:horizontal-count-per-offset} on $D_0\le1$ and each
dyadic shell $D_0\asymp V>1$. Since $0<\delta<1$, the weighted
sum of $V^\delta$ against $\min(1,T/(HV))$ over $V>1$ is
$O(\max(1,T/H)^\delta)$. This proves
\eqref{eq:horizontal-weighted-offset}. In the complementary range
one can count for each $b_0$ at most $O(1+BT/q_g)$ reduced real
rational centres within distance $O(T/(Bq_g))$ of $b_0/q_g$, by
their real separation. Each uses $O(H/B)$ row differences. The
complete square norms on $q_g$ then give
\eqref{eq:horizontal-weighted-overflow}. This counting also includes
the real lifts when an offset crosses the boundary of a period.
\end{proof}


\paragraph{A bilinear bound at a single rational approximation.}
\label{sec:dynamic-rational-source-bands}
For coefficients $u_x,v_y$ on intervals at scales $P\le M$, with squared
counting norms at most $PL^C$ and $ML^C$, Cauchy--Schwarz in the long
variable gives
\[
 \frac{\left|\sum_{x,y}u_xv_y e(\alpha xy)\right|^2}{(PM)^2}
 \ll \frac{L^{O(1)}}{PM}\sum_{|j|\ll P}\min(M,\|\alpha j\|^{-1}).
\]
The inner geometric sum is bounded by the minimum displayed, and the
coefficient autocorrelation at every $j$ is at most $PL^C$ by Cauchy.
For $\alpha=a/b+\beta$ and $S=X|\beta|$, suppose $b(1+S)\ll P$ with a
small fixed constant, allowing a logarithmic reserve when $PM\asymp X$
up to powers of $L$. For $j$ not divisible by $b$ the displacement stays
below a fixed fraction of $1/b$, and the geometric sums on these classes
total $O(P L^{O(1)})$. For $j=bk$ the arguments are $\beta bk$ without
wraparound. Summing their reciprocals with the cutoff at $M$ gives
$O((M+PM/[b(1+S)])L^{O(1)})$, by the trivial class bound when $S\le1$.
Thus the bilinear sum divided by $PM$ is
$O((b(1+S))^{-1/2}L^{O(1)})$ in this range.


\subsection{The full high-conductor comparison for grouped products}
\label{sec:near-balanced-high-conductors}

\begin{proposition}\label{octal:near-balanced-high-conductors}
Use the genuine short coefficient $u_x=b_1(x)$ and the long residual $v_y=b_2^*(y)$
of Proposition~\ref{octal:principal-grouped-model}, with any retained
logarithmic Mellin phases placed in the prime-slot weights. Let the shorter
group have reference scale $P=X^\nu$, with $.399\le\nu\le.500$, and let
$XL^{-C}\ll PM\ll X$. Take an exact dyadic conductor $Y=X^\eta\mid X$ with
$.740\le\eta\le1$. For every $J>0$, one may take the fixed power $K_0$ in
the principal model large enough that
\[
             \left|\sum_{x,y}u_xv_y\mathcal D_Y f(xy)\right|
                          \ll_J AL^{-J}.
\]
The choice is uniform in $Y$ and across all bounded sets of prime-slot
derivative bounds of the stated logarithmic size.
\end{proposition}

\begin{proof}
Use $a=.3825$, $h_0=.067$, and $\delta=1-\log_8 7<\delta_0=.0643$,
the last inequality following from $7^{10000}>8^{9357}$.
When the actual smaller scale exceeds $\sqrt X$ by a bounded factor,
take the reference scale to be $\sqrt X$ and absorb that factor in its
support constants. The coefficients are odd supported, have squared
counting norms $\ll PL^C,ML^C$, and are pointwise $\ll X^\epsilon$ for
every fixed $\epsilon>0$. The power of $L$ in the squared norms is
independent of $K_0$. The only further arithmetic information used in
this proof is that, on any specified polylogarithmic range for $b$ and
$S$, and for any fixed $N$, we may arrange
\begin{equation}\label{eq:near-balanced-small-model-arc}
    \left|\sum_y v_y e(\alpha xy)\right|\ll_N M L^{-N}
   \quad\bigl(\alpha=a_0/b+O(S/X),\ x\asymp P\bigr).
\end{equation}
Indeed multiplication by $x$ gives an angle with reduced denominator
at most $b$ and offset $O(S/M)$, up to powers of $L$, so this follows
from \eqref{eq:principal-grouped-low-arc}. The choice of the
polylogarithmic range is made after the logarithmic losses in the
elementary estimates below have been fixed.

Put $\mu=1-\nu$ and $\lambda=\max(.589,\mu+h_0)$, and choose
an octal cut $R\asymp X^\lambda$. The exact source
$(E_{X,Y}-E_{X,Y/2})f$ has physical factorization
\begin{equation}\label{eq:near-balanced-source-factorization}
   L_R(n\bmod R) U(\lfloor n/R\rfloor\bmod Y/R),
\end{equation}
times $(X/Z)^{-\delta}$, where $Z$ is the least octal multiple of $Y$.
Indeed $R\mid Y/2$. Put $N=Z/R$, $m=Y/R$, let $u$ be the
unprojected upper word of period $N$, and use probability-normalized
Fourier coefficients on each period. Its density is $\rho_U=N^{-\delta}$.
The upper factor is $U=(E_{N,m}-E_{N,m/2})u$; hence exactly
\[
 \widehat U(\ell)=1_{\ell\ {\rm odd}}\widehat u((N/m)\ell),
 \qquad |\widehat U(\ell)|/\rho_U
   =1_{\ell\ {\rm odd}}F_N(\ell/m).
\]
The bounded ratio $N/m=Z/Y\le4$ accounts for the nonoctal grid:
restriction to this subgrid never increases a positive moment sum.
Since averaging modulo $m$ or $m/2$ folds only a bounded number of
translates, the signed upper factor is bounded by one and supported
on $O(m\rho_U)$ cells, giving $O(m\rho_U)$ step boundaries.
The squared values on one period sum to $O((Y/R)\rho_U)$.
Decompose $L_R$ into exact dyadic periods $q\mid R$, combining
the periods at most two. Denote a lower component by $L_q$ and
its density scale by $\rho_L=(|K_d|/7)R^{-\delta}$, the density of
$L_R$ before projection. In units of this density the lower amplitude has the shifted word
moment bounds at length comparable to $q$. The digits above that
length contribute their density. In the least significant position,
the normalized $K_d$ factor has modulus at most one. Dropping it
in the first moment costs at most eight copies of a shortened
shifted grid. For the square moment, shifted Parseval costs the
factor $7/|K_d|$. These bounds are uniform since $|K_d|\ge3$.
In particular its squared mean is $O(\rho_L^2 q^\delta)$.
The comparison scale $X\rho_L\rho_U$ times the omitted density
$(X/Z)^{-\delta}$ is comparable to $A$.


\paragraph{Lower row completion.}
For a lower period $q$ put $\gamma=\log_X q$ and take
$H\asymp\max(X^{.065},q/M)$ up to powers of $L$, decreasing
it by a logarithmic factor if needed so that $HM\ll R$.
This is possible since $\lambda\ge\mu+.065$ and $q\le R$.
Here is the normalization of the completed correlation. Write
$w_i=q^{-1}\sum_{z\bmod q}L_q(z)e(-iz/q)$; a lower frequency mask
multiplies this array by its indicator. Apply
Lemma~\ref{octal:short-row-upper-density} and use its disjoint binary
intervals at each length. Majorize the column interval by a fixed
nonnegative $\psi(y/M)\in C_c^\infty$, equal to one there. For a row
pair $p,p'$ the Fourier expansion and Poisson summation give the kernel
\[
 M\sum_{i,b\bmod q}w_i\overline{w_b}
   \sum_{t\equiv pi-p'b\ (q)}\widehat\psi(-Mt/q),\qquad
 \widehat\psi(\xi)=\int\psi(v)e(-v\xi)\,dv .
\]
Decay permits restriction to $|t|\le T=(1+q/M)X^{\epsilon_0}$ with
error smaller than any prescribed negative power: all index ranges,
row weights and Fourier norms grow at most polynomially. At each
binary length every pair is counted at most once. After taking absolute
values off the diagonal, the weights $u_p\overline{u_{p'}}$ cost
$O_\epsilon(X^\epsilon)$ and we may sum over all $0<|p'-p|\ll H$.
Relative to $(PM\rho_L\rho_U)^2$ the prefactor of this index sum is
$O(X^\epsilon L^{O(1)}/(PH\rho_L^2\rho_U))$; in particular there is
only one inverse upper density. On the diagonal use physical space:
$\sum_{p,y}|L_q(py)|^2\ll X^\epsilon(PM+q)\rho_L^2q^\delta$
by the divisor bound and periodic Parseval, also after masking. Since
$q\le R\ll PM$, the relative diagonal is
$O(X^\epsilon L^{O(1)}q^\delta/(H\rho_U))$.
The exponents $\epsilon,\epsilon_0$ will be taken sufficiently
small with respect to all fixed margins below.

We now bound the contribution of nonhorizontal slopes to this Gram.
Choose an octal divisor $Q$ of $q$ with $Q\asymp\min(q,PH)$.
Use the least one exceeding a sufficiently large fixed multiple
of $PH$ if available, and otherwise the largest octal divisor
of $q$.
Slopes $s$ of additive order $g=Q/(s,Q)\ll P$ are called
horizontal, using a fixed comparison constant large enough
to include every slope with a horizontal difference of length
at most the diameter of the row box. For the other slopes the
absolute contribution at each fixed twist $t$, in units of
$\rho_L^2$, is at most
\begin{equation}\label{eq:near-balanced-capped-affine}
              (q/Q)^\delta Q^{2a}(1+PH/Q)L^C.
\end{equation}
To see this, rescale the difference lattice of the affine
equation by $P,H$. Its determinant is $Q/(PH)$. If $\ell$ is
the length of a shortest nonzero vector, the affine count in
the rescaled box is $O(1+PH/Q+1/\ell)$. Indeed parallel
lattice lines have spacing $Q/(PH\ell)$ and the points on
each have spacing $\ell$. If two points occur, $\ell\ll1$.
Use the full two digit first norms at multiplicities
$O(1+PH/Q)$. For multiplicity $J\gg1+PH/Q$, a shortest
vector gives $0<|v|\ll P/J$ and $0<|w|\ll H/J$. A vertical difference has unscaled height at least $Q\gg H$, and
a horizontal difference has width at least the additive order of $s$,
larger than the row box. Thus neither component vanishes. Put $c=(v,Q)$, which divides $w$, and
divide out the full greatest common divisor of $(v,w)$. Orient
the resulting primitive pair with first coordinate $v_0>0$.
Its first coordinate is odd and $s\equiv w_0/v_0\pmod{Q/c}$.
Choose an octal $Q_1\mid Q/c$ comparable to $PH/(Jc)^2$ with a
sufficiently large comparison constant. It exists on nonempty
vector ranges since $J\gg 1+PH/Q$. Equality of two ratios modulo
$Q_1$ would give a vanishing integer cross determinant, whose
absolute value is less than $Q_1$, hence the same primitive pair.
The ratio-overlap bound \eqref{eq:ordinary-group-gram-overlap} gives
$\ll c (Q/Q_1)^\delta Q_1^{2a}L^C$. Summing in the dyadic $c$,
with multiplicity $O(J)$, gives (the power of $c$ is
$1-2(2a-\delta)<0$, and $c$ runs over powers of two)
$\ll Q^\delta(PH)^{2a-\delta}J^{1-2(2a-\delta)}L^C$.
Since $2(2a-\delta)>1$, the tail sums to
$Q^{2a}(1+PH/Q)L^C$ on the coarse grid. In the original
congruence each lift of one index determines the other lift
by an odd affine permutation. Shifted Cauchy therefore gives
the factor $(q/Q)^\delta$ in \eqref{eq:near-balanced-capped-affine}.
Bounded masks on the
lower coefficients only decrease these positive majorants.

For $.119\le\gamma\le\mu-.0001$ take $H\asymp X^{.065}$ up to powers of $L$
and retain only $t=0$, by smooth completion. In this untwisted
case fix the odd greatest common divisor $c$ of a distinct row
pair and write $B_c=PH/c^2$. Here $c\ll H$ since it
divides the nonzero row difference. The reduced positive pair is
primitive. At a coarse modulus $Q_c\mid q$ chosen octal with
$Q_c\asymp\min(B_c,q)$, each ratio is attained by
$O(1+B_c/Q_c)$ primitive pairs. Rescale the primitive box
by $P/c,H/c$. If the shortest lattice vector is shorter than
a suitably small multiple of the determinant $Q_c/B_c$,
all points in the bounded box are collinear and at most one
is positive primitive. Otherwise counting points on parallel
lattice lines proves the asserted multiplicity. The digit
first norms and the lifted second norms give at fixed $c$
the bound $(q/Q_c)^\delta Q_c^{2a}(1+B_c/Q_c)L^C$.
Put $F=PH$ and $u=2a-\delta>1/2$. The bound at fixed $c$ is
$O(q^\delta B_c\min(q,B_c)^{u-1}L^C)$.
Since $\min(q,B_c)\ge c^{-2}\min(q,F)$ and $u<1$, summing $c^{-2u}$
gives $O(q^\delta F\min(q,F)^{u-1}L^C)$.
After dividing by $PH$ and including the upper density cost, the
relative exponent is bounded by $(2a-1)(.119)+.0643(1-\lambda)<0$
when $q\ll PH$, and by $.0643-.2993(.399+.065)<0$ when $q\gg PH$.
These bounds also apply at the constant-factor transition.
The diagonal exponent is at most $.0643-.065<0$.

For $\gamma\ge\mu-.0002$ put $h=\max(.065,\gamma-\mu)$
and retain the $T$ ordinary completions. If $q\gg PH$, the
relative nonhorizontal squared exponent from
\eqref{eq:near-balanced-capped-affine} is at most
\begin{equation}\label{eq:near-balanced-nonhorizontal-exponent}
 \delta(\gamma+\eta-\lambda)
        +(2a-\delta-1)(\nu+h)+(\gamma-\mu)_++O(\epsilon_0+\epsilon).
\end{equation}
We may use $\delta_0=.0643$ in this expression because the
net coefficient of $\delta$ is nonnegative when $q\ge PH$.
The expression is increasing in $\gamma$, so its largest
value is bounded by its value at $\lambda$, namely
\[
 \delta_0\eta-(1+\delta_0-2a)\nu+(2a-\delta_0)(\lambda-\mu).
\]
When $\lambda=\mu+h_0$ this is at most $-.000392$ uniformly
by evaluation at the left endpoint in $\nu$ and $\eta=1$. When $\lambda=.589$ its
maximum is at $\nu=.5$ and is smaller. If $q\ll PH$, the
corresponding exponent is
$(2a-1)\gamma+(\gamma-\mu)_++\delta(\eta-\lambda)$.
This is negative below $\mu$ and is bounded above $\mu$
by $2a\lambda-\mu+\delta(\eta-\lambda)<0$.
The transition with $q\asymp PH$ is covered by these
same bounds. The diagonal is paid since $H\gg X^{.065}L^{-C}$.

For a coarse horizontal slope, take the additive order $g$
of its true slope modulo $q$. It satisfies $g\ll(q/Q)P$.
Put $q_g=q/g$ and keep the exact affine equation rather than
projecting this contribution to $Q$. For $g\ll P$,
Lemma~\ref{octal:horizontal-weighted-offset} applies on the
full lower grid with at most bounded dyadic aliases.

For $g$ larger than a fixed multiple of $P$, write the two
indices as $b_0+k' q_g$ and $b_0+k q_g$. The exact order $g$
makes $k'-k$ odd. The equation for the rows $p,p+j$ becomes
\begin{equation}\label{eq:near-balanced-affine-true-order}
      p(k'-k)-j k=(t+j b_0)/q_g\pmod g,
               \qquad jb_0+t=0\pmod{q_g}.
\end{equation}
At $b_0,j,t$ fixed, split the digit transform on the lift
variable at a dyadic modulus $g_0\mid g$ comparable to $P$.
For each pair $k,k'\bmod g_0$, the congruence determines $p$
modulo $g_0$ because their difference is odd. Thus only
$O(1)$ integers $p$ in the actual row box occur. Fixing one,
each remaining lift of $k$ determines the lift of $k'$ by
an affine permutation, since both $p$ and $p+j$ are odd.
The shifted square norms on these lifts cost $(g/g_0)^\delta$.
The two upper first norms at the coarse lift scale cost
$g_0^{2a}L^C$. If $g$ is nonoctal, retain the lower block
through its greatest octal divisor $g_-$ and the upper block
starting from its least octal multiple. The retained lower
digits split at $g/g_0$, choosing $g_0/(g/g_-)$ octal.
The coarse first norm is over a bounded number of aliases
on the $g_0$ grid, and the remaining square sum is unchanged.
The top blocks at size $q_g$ are then treated as in the
horizontal lemma. Consequently both estimates in that lemma
hold in this range upon replacing $g^{2a-1}$ by
$P^{2a-1}(g/P)^\delta$.

Choose the constant in this large-order cutoff larger than four
times the coarse horizontal-order constant. If $Q$ were the largest
octal divisor of $q$, then $q/Q\le4$ and $g\ll(q/Q)P$ would place
the slope in the preceding range. Thus here $Q\asymp PH$, so
$g\ll q/H\ll M$ and $q_g\gg H$. Both the separated and
dense contributions save with room. In the separated case
use $(g/P)^\delta\ll(M/P)^\delta$. In the dense case
combine $(g/P)^\delta q_g^\delta=(q/P)^\delta$ and use
$T/q_g\ll X^{O(\epsilon_0)}$. After the upper density cost,
both exponents are at most
$\delta_0\eta-(1+\delta_0-2a)\nu+O(\epsilon_0+\epsilon)<0$.

For the remaining orders $g\ll P$, the separated part
of the horizontal bound with upper density included costs
\[
 X^{\delta(\eta-\lambda)+O(\epsilon_0+\epsilon)}
                   g^{2a-1} B^{2\delta-1}L^C,
\]
where $B$ is the denominator shell of the reduced centres
and we used $T/H\ll X^{O(\epsilon_0)}$. Define
\begin{equation}\label{eq:near-balanced-order-cutoff}
     \Gamma(\nu,\eta)=.002+\frac{\delta_0}{1-2a}(\eta-\lambda).
\end{equation}
The parts $g>X^{\Gamma(\nu,\eta)}$ or $B>X^{.037}$ save
powers, because
$(1-2a)\Gamma(\nu,\eta)-\delta_0(\eta-\lambda)\ge.00047$ and
$(1-2(.0643))(.037)=.0322418>.0264273$.
In the dense case of that lemma, $B\ll H$ and $T/H\ll
X^{O(\epsilon_0)}$ force $g\gg q/H^2$ up to small powers.
Thus $g>X^{.34}$ and $T/q_g\ll X^{.0002+O(\epsilon_0)}L^C$,
because $g\ll P$. The complete square bound in the lemma
pays the dense case with a fixed margin.

It remains to remove, before the row bound, the lower source
frequency mask
\begin{equation}\label{eq:near-balanced-lower-band}
 \left|\frac{i}{q}-\frac{u}{gd}\right|_{\mathbb R/\mathbb Z}
        \ll \frac{T}{dq},\qquad
               g\le C X^{\Gamma(\nu,\eta)}\text{ dyadic},\quad
               d\le C X^{.037}.
\end{equation}
Indeed both horizontal frequencies agree modulo $q_g$, and
the lower index condition in the horizontal lemma puts each
of them in this mask. Choose the fixed constants large
enough and include the half-period translates to preserve
odd physical support. Deleting this mask kills the indicated
terms in the remaining row Gram. We prove its comparison
directly with the unchanged upper word. More precisely, for this
mask $E$ split the lower Fourier polynomial as $L_q=L_{q,E}+L_{q,E^c}$
in the original linear pairing. Only the second pairing is squared
in the preceding Gram estimate: a horizontal pair with both indices
in $E$ is then absent identically. The first pairing retains the full
upper word and is estimated below by lifting its frequencies. Neither
the cardinality of $E$ nor an average over all lower indices bounds
the digital mass on $E$ by itself.

\paragraph{Fourier lifting of the lower mask.}
Group the reduced original centres in \eqref{eq:near-balanced-lower-band}
by the exact dyadic part $G$ of their denominator and the
dyadic size $D_0$ of its odd part. Also group $d\asymp B$.
Write $b=|\operatorname{den}(u/(gd))|\asymp GD_0$;
after reducing we may decrease $g,d$ and enlarge the bands.
Assign duplicates to one allowed representation of largest
width. Distinct centres are disjoint at scale $T/q$ by their
rational separation and the fixed upper bound $.1145+.037$
on their denominator exponent. For the squared mass keep an initial
block of length comparable to $G$, one of length comparable to $D_0^2$,
and a terminal block of length comparable to $1+T/B$. They fit below
$q$ by the fixed power bounds. Shifted Parseval applies to the initial
dyadic lifts and the terminal offset grid. After removing the dyadic
part, the odd fractions are separated by $\gg D_0^{-2}$. Sum the local
suprema of the squared middle amplitude using Bernstein's inequality,
as in Section~\ref{sec:zero-model-wide-source-bands}. The squared lower
mass in units of $\rho_L^2$ is consequently
\[
                  \ll G^\delta D_0^{2\delta}(1+T/B)^\delta L^C.
\]
Since $G/B\ll g$, $G\ll X^{.1145+.037}$ and
$T\le X^{.089+O(\epsilon_0)}$, its exponent is less than $.0180$
up to small completion losses. For $T\ge B$ the bound uses
$.0643(.1145+.089+2\cdot.037)<.0180$, and for $T<B$ it uses
$.0643(.1145+3\cdot.037)<.0180$.
For the first norm keep two middle blocks of length comparable
to $D_0^2$ and use Cauchy--Schwarz and the shifted square moment on
each. Octal dilation preserves the odd-denominator separation for
both blocks. The initial and terminal factors use shifted first norms. Thus the corresponding
first norm on a lower index radius $V\ll T/B$, enlarged to
$V\ge1$ when needed, is
\begin{equation}\label{eq:near-balanced-band-first-norm}
           G^a D_0^{2\delta} \max(1,V)^a L^C .
\end{equation}

Smooth the upper step word at scale $X^{-.020}R$ using
a compactly supported smooth approximate identity. In each
affected boundary interval, Cauchy and lower periodicity
bound the lower absolute mass by
\[
 (X^{-.020}R)^{1/2}(q+X^{-.020}R)^{1/2}
                       \rho_L X^{.0180/2+O(\epsilon_0)}L^C .
\]
There are $O((X/R)\rho_U)$ relevant boundaries in the
physical product range, including period repetitions.
Since $q\le R$, dividing the sum of these errors by $X\rho_L\rho_U$
costs at most $X^{(-.020+.0180)/2+O(\epsilon+\epsilon_0)}L^C$,
including the divisor bound for representations of each product and
both coefficient bounds. Thus smoothing is performed in the original
linear pairing, not inside a signed squared correlation.

Set $Q_R=Y/R$ and $Q_q=Y/q$. The upper step coefficient at
index $\ell$ is $\widehat U(\ell\bmod Q_R)$ times
\begin{equation}\label{eq:near-balanced-sinc-block}
 \frac1R\int_0^R e(-\ell u/Y)\,du
   = \left(\frac{q}{R}\sum_{j<R/q}e(-qj\ell/Y)\right)
                    \frac1q\int_0^q e(-\ell u/Y)\,du.
\end{equation}
The factor in parentheses is the normalized Fourier
transform of the unrestricted block from $q$ to $R$.
The smoothing multiplier cuts the tail past $Q_R X^{.025}$
to arbitrary power accuracy. The discrete coefficients on the integer
grid are the sums of the continuous coefficients in each class modulo
$Y$. The retained interval has length $o(Y)$ and the discarded aliases
have arbitrarily small total mass by the smoothing multiplier. Accordingly a full ordinary
frequency $\alpha=h/Y$ from one lower band remains within
effective offset $S=X\|\alpha-u/(gd)\|\ll X^{\nu+.0002+O(\epsilon_0)}/B$.
The padding exponent $1-\lambda+.025\le\nu-.042$ is
smaller than $\nu-.037$ throughout the range.

For the convolution, write $h=i_0 Q_q+r$ with $0\le r<Q_q$,
and in its lower coefficient put $i=i_0-j$. At
$\ell=r+jQ_q$ the last factor in
\eqref{eq:near-balanced-sinc-block} is $O((1+|j|)^{-1})$.
The other two upper factors are unchanged if $h$ is replaced
by $h-jQ_q$. The lower transform is evaluated at $(i_0-j)/q$,
within $1/q$ of that shifted argument divided by $Y$.
We can take local suprema on the retained lower digit blocks.
If $q$ is not octal, use its greatest octal divisor for the
last retained lower position and its least octal multiple
for the first retained unrestricted upper position, dropping
only bounded radix factors. The unrestricted geometric
factor and the upper word in
\eqref{eq:near-balanced-sinc-block} have complete first norm
$O(Q_q^a)L^C$ in their joint period and shifted local first
norm $O(V^a)L^C$ on any index interval of length $V\le Q_q$.
For the local norm retain the terminal digits of length comparable
to $V$ and bound all other normalized digit factors by one. On the
integer grid the exact upper projection is a frequency indicator, so
it can be discarded in this majorant; no derivative estimate for a
sharp mask is needed. The retained digits consist of at most one
unrestricted run and one digital run, up to bounded radix factors.
An unrestricted run of length $N$ has normalized kernel bounded by
$\min(1,(2N\|t\|)^{-1})$, hence shifted first sum $O(\log(2N))$.
For consecutive runs the shifted grid sum is at most the product of
the two uniform shifted bounds, by summing first over the lifts of
each residue. This proves the asserted estimate for the mixed block. For every shell of
effective offsets of size $S\ge1$ it follows that the
normalized first norm on the full grid is at most
\begin{equation}\label{eq:near-balanced-lift-first-norm}
                  G^a D_0^{2\delta}\max(1,YS/X)^a L^C .
\end{equation}
Put $V=\max(1,YS/X)$ and $W=1+V/Q_q$. Suppress the common
factor $G^aD_0^{2\delta}L^{O(1)}$. In the output interval there
are $O(W)$ upper period cells, each of first norm
$O(\min(V,Q_q)^a L^C)$. For lower indices within $O(W)$ of the
original centre, sum the sinc weight over those cells \emph{at fixed
lower index}: it costs $O(\log Y)$, not $O(W)$. Their cumulative
lower mass costs $W^a$ by \eqref{eq:near-balanced-band-first-norm}
(cap the radius at the width of the input mask). On a dyadic lower
offset $D_1\gg W$ the sinc factor is instead $O(D_1^{-1})$ in
every cell. The total lower and cell-count cost there is at most
$W\sum_{D_1\gg W}^{\rm dyadic}D_1^{a-1}\ll W^a$. Finally
$W^a\min(V,Q_q)^a\ll V^a$, as claimed.

\paragraph{Resolving secondary rational centres.}
For the inner range $b(1+S)\ll P$, the short-row bilinear
Cauchy estimate of Section~\ref{sec:dynamic-rational-source-bands}
gives a pointwise normalized bound $(b(1+S))^{-1/2}L^C$,
using only the counting second norms of $u_x,v_y$.
In conjunction with \eqref{eq:near-balanced-lift-first-norm}
this pays with arbitrary logarithmic savings except at
polylogarithmic $G,D_0,S$, since $a<1/2$ and $2\delta<1/2$.
There the coordinate model estimate
\eqref{eq:near-balanced-small-model-arc} applies to the
actual coefficient, uniformly in the opposite physical
coordinate.

Choose a reduced Dirichlet approximant $a'/b'$ to the
actual full angle, with $b'\le P$ and error
$\ll1/(Pb')$. If this is the original centre and
$b(1+S)\gg P$, it has $bS\ll M$. In this case the
short Cauchy estimate gives $O((P^{-1}+(b(1+S))^{-1})^{1/2}L^C)$.
Indeed the rational short-difference classes have displacement
$O(L^C/b)$. The classes at larger distance from the integers
cost $O(P L^C)$ in the sum of column geometric sums.
In the polylogarithmically many exceptional classes the phases have
gap $bS/X$ and cross at most polylogarithmically many integers. Their
cost is $O(M+X/(b(1+S)))L^C$, using the trivial estimate
when $S<1$. This proves the bound. For $bS\gg P$ its
$P^{-1/2}$ factor pays \eqref{eq:near-balanced-lift-first-norm}
by $a(\Gamma(\nu,\eta)+.0002)+(a-.5)\nu<0$. The maximum
for each parameter choice is at $\eta=1$ and the lower
endpoint of $\nu$.

Otherwise the cross determinant is nonzero and gives
\begin{equation}\label{eq:near-balanced-cross-determinant}
               1\ll b b' S/X + b/P.
\end{equation}
The $b/P$ term is power-small. In a secondary shell
$b'\asymp B'=X^\sigma$, write $S=X^\tau$ up to rounding,
$G=X^{\Gamma_G}$ and $D_0=X^\Delta$. Then
$\sigma+\tau+\Gamma_G+\Delta\ge1-O(L^{-1})$.
Moreover $\sigma\le\nu$ and $\log_X b\le .1515+o(1)$, so
\[
 \eta+\tau-1\ge\eta-\nu-\log_X b-o(1)>0.
\]
The last inequality uses $\eta\ge.740$ and $\nu\le.500$.
Thus the maximum in
\eqref{eq:near-balanced-lift-first-norm} equals $YS/X$ in this branch.
The same short Cauchy estimate at each such approximant
is $O((B')^{-1/2})L^C$, since $b'\le P$ and the
approximation error has size at most $1/(P b')$. Put
$\sigma_0=\min(\nu-.001,.490,\eta-\nu-.001)$.
For $\sigma\ge\sigma_0$ the direct bound
\eqref{eq:near-balanced-lift-first-norm} gives exponent
\[
  a(\Gamma(\nu,\eta)+\nu+.0002+\eta-1)-\sigma_0/2
     +O(\epsilon_0)\le-.00989+O(\epsilon_0).
\]
For the three choices in the minimum defining $\sigma_0$,
the maximizing endpoint pairs $(\nu,\eta)$ are respectively
$(\nu_0,1),(.5,1),(.5,.740)$, where $\nu_0$ is the applicable
lower endpoint. This follows by using the two affine pieces
of $\lambda$ and their respective monotonicities.

For $\sigma<\sigma_0$, fix one original centre, the
secondary denominator shell and the common integer $j$
in the sinc expansion. By
\eqref{eq:near-balanced-cross-determinant} the distinct
secondary centres all lie in an interval of length $O(S/X)$
about the original centre, and their mutual spacing is
$\gg (B')^{-2}$. Each full frequency window around them
has length $O(V')$ in index units, where $V'=Y/(PB')$.
Put an octal start position $B_0\asymp X/S$ and keep a
block of octal length $N_0\asymp (B')^2 S/X$ there, as
well as a terminal block of length comparable to $V'$.
Indeed $V'>X^{.001}$ by $\sigma<\eta-\nu-.001$, while the
determinant gives $N_0\gg B'/b$. Since $bS\ll
X^{\Gamma(\nu,\eta)+\nu+.0002+O(\epsilon_0)}$, it also gives
$B'\gg X^{1-\Gamma(\nu,\eta)-\nu-.0002-O(\epsilon_0)}$;
with $b\ll X^{.1515+o(1)}$ both retained lengths grow polynomially.
They fit disjointly since $B_0N_0V'\asymp YB'/P\ll Y$
with a fixed power margin. The original-centre bottom
blocks fit below $B_0$, since the exponent of $GD_0^4S$
is at most $.1145+.5002+3(.037)+O(\epsilon_0)<1$.

After dilation by $B_0$ the secondary centres lie in a
bounded real interval and have spacing $\gg N_0^{-1}$.
The common translate $-j/q$ preserves this property and
variation within their windows stays within the sampling
cells, since $B_0N_0V'/Y\ll B'/P$. A block entirely below
or above $q$ has sampled first norm $O(N_0^a)L^C$.
For a block below $q$, the displacement between the lower-grid
argument and $h/Y-j/q$ is at most $1/q$; its start times its
length is at most $q$, so this enlarges each sampling cell by only
a bounded factor. For a block above $q$ use the full angle with no
rounding (translation by $j/q$ is then immaterial). If the block crosses
$q$, split it into retained lengths $N_L,N_U$ on either
side, dropping a bounded radix factor if necessary.
Partition the dilated real angles into cells of width
$1/N_L$. In each cell the upper arguments lie in a
bounded interval and are $\gg N_U^{-1}$ separated. The
upper sampled first norm is $O(N_U^a)L^C$. Sum the
suprema of the lower factor on the enlarged cells at
cost $O(N_L^a)L^C$. Thus the joint norm is again
$O(N_0^a)L^C$. The upper part uses the unrestricted
middle digits in \eqref{eq:near-balanced-sinc-block};
they have the same shifted first norm bound. The determinant bound also gives $\nu+\sigma\ge1-\Gamma(\nu,\eta)-.0002-o(1)>\lambda$,
so the terminal block begins beyond the lower period. It contributes
$(V')^a L^C$ on each
secondary integer frequency window. Sum the bottom
suprema over original centres as in
\eqref{eq:near-balanced-band-first-norm}, and sum the
harmonic weights in $j$. The normalized joint first norm
is bounded by
\[
         G^a D_0^{2\delta}(N_0V')^a L^C
             \le G^a D_0^{2\delta}(B'S/P)^a L^C .
\]
Multiplying by the pointwise $(B')^{-1/2}$ and using
\eqref{eq:near-balanced-cross-determinant}, together with
$\tau\le\nu+.0002-\log_X d+O(\epsilon_0)$ and
$\Gamma_G\le\log_X g+\log_X d-\Delta+O(L^{-1})$,
gives the exponent
\[
  (a-1/2)(1-\nu)+.0001+\tfrac12\log_Xg
                       +(2\delta-a)\Delta+O(\epsilon_0)
          \le-.00142+O(\epsilon_0).
\]
This completes the comparison of the removed horizontal
mask with the original upper word.

\paragraph{The small lower periods.}
For $q\le X^{.120}$ use all the lower frequencies, without
a mask. Their squared mass in lower density units is at most
$q^\delta L^C$. Smooth the upper step at $X^{-.009}R$;
the boundary error again saves a power since $.120\delta<.009$.
Truncating at $Q_R X^{.012}$ now leaves disjoint source
frequency windows of effective width at most $X^{\tau_0+o(1)}$,
where $\tau_0=1-\lambda+.012\le\min(.423,\nu-.053)$,
around $i/q$. In each such window there is exactly one
lower coefficient, with total normalized first norm
$O(q^a)L^C$ over all centres. For $q>2$ the exact conductor forces $i$ to be odd,
so $i/q$ has reduced denominator $q$. Apply the same inner and
secondary estimates with $G=q$, $D_0=1$ and effective
offset exponent at most $\tau_0$; at the bounded periods
there are only boundedly many denominator families.
For any distinct secondary centre in this small-period case,
the determinant bound gives $\eta+\log_X S-1
\ge\eta-\nu-\log_X b-o(1)>0$, now with $\log_X b\le .120+o(1)$.
Thus the factor $\max(1,YS/X)^a$ equals $(YS/X)^a$.
The coincident-centre outer case uses
$a(.120+\tau_0)-\nu/2<0$. The direct large-secondary branch uses
$a(.120+\tau_0+\eta-1)-\sigma_0/2<0$, and the joint
small-secondary branch uses
$a(1-\nu)-1/2+(.120+\tau_0)/2<0$, uniformly with fixed
margins. For the direct estimate at $\sigma_0=\nu-.001$,
use $\tau_0\le\nu-.053$ and $\nu\ge.399$. For the other
two choices in $\sigma_0$, use $\tau_0\le .423$.
The latter sampling block starts past the lower
word and uses the unrestricted run to $R$ if needed.
This also includes the combined parity projection.

The small, untwisted, and twisted lower ranges overlap.
Summing their bounds over dyadic $q$ proves Proposition~\ref{octal:near-balanced-high-conductors}.
The exponent inequalities are checked in exact rational arithmetic by
\path{supplement/base8/verify_exponents.py}.
\end{proof}



\section{The grouped comparison for octal prime-factor regions}
\label{sec:octal-balanced-proof}

\begin{proof}[Proof of Lemma~\ref{octal:balanced-shape-comparison}]
First take a fixed smooth test $W$ on the rough simplex, extended
to a neighborhood on which the prescribed split has fixed reserve in
$(.399,.601)$. Use $u_i=\log_n p_i$, $n=\prod p_i$, and insert
$\omega(n/X)$. In the dyadic prime partitions the exponents of the
group scales differ from their simplex values by $O_D(\log L/L)$.
Use the smaller group as the short coordinate; its reference exponent
is in $[.399,.500]$, with the harmless balanced-scale convention of
Proposition~\ref{octal:near-balanced-high-conductors}. Separate the
tests in logarithmic coordinates as in Section~\ref{sec:principal-grouped-models}.

The distinctness condition can be discarded in such boxes. For
any of the nonnegative sources used below, use Type~I on the
divisors $p^2\le X^{.640}$ with $p\ge X^\kappa$, and use the
integer count $O(X^{1-.640/2}+X^{1/2})$ for larger squares.
The first main term is $O(A\sum_{p\ge X^\kappa}p^{-2})$, and
both bounds are $o(A/L)$. The fixed-rank representation count
and the interval cutoffs preserve that error.

Set $Q_*=2^{\lfloor .741\log_2 X\rfloor}$. Telescope
$f-f_0=(f-E_{X,Q_*}f)+\sum_{4\le q\le Q_*}^{\rm dyadic}\mathcal D_q f$.
We estimate the sum first; none of its characters is principal. For $q\le L^C$, fix all but one
prime and apply Siegel--Walfisz to the odd classes in that prime
column. The exact conductor contribution has zero mean on these
classes. The total absolute source mass over a period is
$O(Aq/X)$, and the reciprocal sum over fixed factor choices costs
only powers of $L$. For $L^C<q\le X^{c_0}$ with $c_0>0$ sufficiently
small in terms of $\kappa$, use the ordinary prime exponential
sum bound \cite[Theorem~23.8]{Koukoulopoulos2019} in the same column, removing the prime powers of exponent at least two at cost $O(N^{1/2}\log^2 N)$ and applying partial summation for the prime weight. In normalization
by the prime scale $N$, its three bounds are $q^{-1/2}$, $N^{-1/5}$, and
$(q/N)^{1/2}$, up to logarithms. The denominator remains $q$ since the
fixed factors are odd. The shifted digital first norm costs $q^{.3825}$.
Since $N\gg X^\kappa$, choosing $c_0$ sufficiently small gives arbitrary
inverse-logarithmic saving, on taking $C$ sufficiently large.

For $X^{c_0}\le q\le Q_*$ use \eqref{octal:dyadic-char-expansion}.
At $H=X^{3\epsilon}$ define $\beta_\chi$ to be the maximum of $1/2$
and the real parts of the zeros of $L(s,\chi)$ with $|\Im s|\le H$.
For a slot of scale $N_i$, the smooth explicit formula for the
nonprincipal primitive character bounds the log-weighted prime sum
by $N_i^{\beta_\chi}L^{O(1)}$. Indeed the Mellin transform of the
rescaled slot at $\sigma+it$ is $O_j(L^{C_j}(1+|t|)^{-j})$ uniformly
in $0\le\sigma\le1$, also after the retained polylogarithmic twists.
There are $O(\log(q(3+|t|)))$ zeros per unit height; use $j=2$ on
the retained zeros and then sufficiently large fixed $j$ beyond $H$.
The trivial-zero and prime-power terms are $O(N_i^{1/2}L^{O(1)})$.
The same $\beta_\chi$ controls every slot: their product costs only
$X^{\beta_\chi}L^{O(1)}$, not a product of independent zero counts.

Put $y=1-\beta$ for a layer with lower real-part endpoint $\beta$,
$\delta=1-\log_8 7$, $\sigma=.6175$ and $\lambda_0=(\sigma-\delta)/2$.
Ingham's fixed-modulus zero-density bound, in the form recorded in
\cite[(1.4)]{ChenGuptaLi2026}, gives at most
$X^{\gamma D_I(y)+O(\epsilon)}$ characters in the layer, where
$\gamma=\log_X q$ and $D_I(y)=3y/(1+y)$ for $0\le y\le1/2$.
Choose a finite real-part mesh of width at most $\epsilon$ and use its
upper endpoints in the prime bounds. The powers lost to the mesh and
to the height are $O(\epsilon)$. At the layer
with $\beta=1/2$ use the trivial character count, which agrees
with $D_I(1/2)=1$. If a layer above $1/2$ is nonempty, each of its characters accounts
for at least one zero counted at the lower endpoint of that layer.
For an arbitrary set $E$ of these primitive characters we have
\[
 A^{-1}\sum_{\chi\in E}|A_\chi|
 \ll\min\{|E|q^{-\lambda_0},\ |E|^{1/2}q^{\delta/2}\}.
\]
This is the pointwise bound and Cauchy--Schwarz on the actual set, not
an assumption that the digital mass on exceptional characters has its
average size. The pointwise and square norms in
Lemma~\ref{octal:semiprime-digital-char} give the relative exponent
\[
 -y+\gamma\min\{D_I(y)-\lambda_0,(D_I(y)+\delta)/2\}.
\]
At fixed $y$ this is affine in $\gamma$, so it suffices to check
$\gamma=c_0$ and $\gamma=.742$. When $c_0<1/3$, the first term
of the minimum bounds the exponent above by $-\lambda_0c_0$,
since $D_I(y)\le3y$. At $.742$ the two branches meet at
$y_*=\sigma/(3-\sigma)$. The derivative in $y$ of the pointwise
branch is positive up to $y_*$, and that of the square branch
is negative from $y_*$ onward. These inequalities follow from
$3(.742)/2<(1+y_*)^2<3(.742)$. The maximum is at most
\[
    -y_*+.742(\sigma+.0643)/2<0.
\]
The strict savings absorb the mesh and height losses. Summing
over the smooth boxes and conductors pays $E_{X,Q_*}f-E_{X,2}f$
with error $o(A/L)$, and $E_{X,2}f=f_0$.

For the remaining source $g=f-E_{X,Q_*}f$, apply the principal models
of Proposition~\ref{octal:principal-grouped-model} after separation
into smooth prime-slot weights. The reconstructed continuous term in
\eqref{eq:principal-grouped-representation} is $o(A/L)$ on the fixed
smooth shape, by equality of the total masses and the endpoint bounds.
In the residual pairing split $g$ into its exact conductors $Y>Q_*$.
The grouped coefficients are the same for all these conductors.
Proposition~\ref{octal:near-balanced-high-conductors} bounds each such
pairing with arbitrary logarithmic saving. The scales have a fixed
positive margin in the group bounds used there. Summation over all
retained twists, dyadic boxes, and conductors gives $o(A/L)$.

Only the recombined shape has a constant simplex integral; cancellation
of the continuous term has not been used on individual boxes.
For the flat source the ordinary prime number theorem and partial
summation in the same boxes give $(\gamma_d A/L)C_W+o(A/L)$.
Indeed the pushforward of $W(\boldsymbol u)\prod dy_i/\log y_i$
has density $C_W/\log n$; the product cutoff has integral $X(1+o(1))$.
Thus each nonnegative source has this same limit on smooth shapes.

For a sharp family first resolve forms constant on $\sum u_i=1$ by
their exact sign on the side $\sum t_i<1$, as in the common sieve.
The one-sided closure hypothesis leaves room to enlarge the support
while retaining both roughness and group reserve. On $I_X$ we have
$t_i=\log_X p_i=u_i+O_D(\log L/L)$. For each small fixed strip width
choose smooth envelopes for the remaining tests: on the ambiguous
strip the lower envelope is zero and the upper envelope is one, with
smooth transitions on a slightly larger neighborhood. They bound the
indicator pointwise for all sufficiently large $X$. The remaining
walls have simplex measure zero and the reciprocal density is bounded
on the enlarged support. Apply the envelopes to each nonnegative
source separately, then let the strip widths tend to zero to obtain the
comparison with the product cutoff. Fixed smooth weights follow by
linearity and uniform approximation; finite marked choices may be
partitioned by first eligible choice, or summed with their multiplicities.
Finally, on the two discarded endpoint intervals the digital mass is
$O(A L^{-Dc}+1)$ by the aligned word bound, and the flat mass is
$O(A L^{-D}+A/X)$. At fixed rank each product has at most $r!$
labelled representations. Both errors are $o(A/L)$, so the result is
for the full unweighted word on $I_X$.
\end{proof}

\part{Comparisons by cutting the word}\label{part:cuts}
The remaining three order estimates use one argument at two scales.
A cut separates the upper digits from the lower; averaging in a
long factor then measures overlaps of two translates and dilates
of the lower word. Small additive periods require an arithmetic
comparison rather than an absolute bound. We begin with two
decimal omissions, since their row estimate also supplies the
algebra for bases seven and six.

\section{Digital estimates for two omitted decimal digits}
\label{dt:digit-section}

Let \(\mathcal D\subset\{0,\ldots,9\}\) have cardinality eight. For a decimal
power \(Z=10^j\), let \(\mathcal W_Z\subset[0,Z)\) be the integers
represented by \(j\) digits from \(\mathcal D\), allowing initial zeros if
\(0\in\mathcal D\). We set \(\mathcal W_1=\{0\}\). Write \(c=\log_{10}8\),
\(\delta=1-c\), \(\e(t)=\exp(2\pi i t)\), and
\[
 p_{\mathcal D}(t)=\frac18\sum_{d\in\mathcal D}\e(dt),\qquad
 \Pi_Z(t)=Z^{-c}\sum_{n\in\mathcal W_Z}\e(nt),\qquad F_Z(t)=|\Pi_Z(t)|.
\]
In particular
\begin{equation}\label{dt:digit-product}
 \Pi_Z(t)=\prod_{0\le l<j}p_{\mathcal D}(10^lt),\qquad
 F_{UV}(t)=F_U(t)F_V(Ut)
\end{equation}
for decimal powers \(U,V\). In the estimates below, one may
prescribe a permitted terminal or leading digit, or restrict the
terminal digit to the units modulo \(10\). This changes the
corresponding digit polynomial to the normalized sum over the
restricted choices. A terminal restriction prescribes only one
position and a prescribed leading digit gives a phase. We detail
the effect of these modifications on each estimate.

\begin{proposition}[Shifted decimal moments]\label{dt:digit-moments}
Uniformly in \(\mathcal D\), in the decimal power \(Z\), and in \(u\in\mathbb R\),
\begin{equation}\label{dt:digit-decimal-moment}
 \sum_{j<Z}F_Z((j+u)/Z)^r\ll Z^{\gamma_r},
 \qquad
 \gamma_1=.428,\quad
 \gamma_{4/3}=.285,\quad
 \gamma_{7/5}=.258,\quad
 \gamma_{3/2}=.225,\quad
 \gamma_2=\delta.
\end{equation}
For \(r=2\) the sum is exactly \(Z^\delta\). The other bounds
also hold when every factor \(|p_{\mathcal D}(t)|\) is replaced by
\(\max\{|p_{\mathcal D}(t)|,\epsilon\}\), for all sufficiently small fixed
\(\epsilon>0\) independent of \(\mathcal D\).
\end{proposition}

\begin{proof}
At \(r=2\) this is orthogonality on the shifted \(Z\)-grid,
since the integers in \(\mathcal W_Z\) are distinct modulo \(Z\).
For the other cases consider the positive operator
\[
 (\mathcal T_r h)(u)
 =\sum_{d=0}^9 |p_{\mathcal D}((u+d)/10)|^r h((u+d)/10)
\]
on continuous periodic functions. The left side of
\eqref{dt:digit-decimal-moment} is \(\mathcal T_r^j1(u)\) at \(Z=10^j\).
For each of the 25 reflection classes, the supplementary certificate
\path{supplement/decimal_two/moments.py} for the shifted moments
verifies a positive piecewise linear periodic function \(h\)
with \(\mathcal T_r h<\ell_r h\). The caps are
\[
\begin{array}{c|cccc}
 r&1&4/3&7/5&3/2\\ \hline
 \ell_r&2.6789&1.924&1.8112&1.6785
\end{array}
\qquad \ell_r<10^{\gamma_r}.
\]
The reflection \(\mathcal D\mapsto9-\mathcal D\) preserves the single-digit magnitude.

Here are the details that make the certificate an inequality on the
whole circle. The nodes are \(i/N\), with \(N=6000\).
The branch points of the nodes have
denominator \(10N\). For such a point \(t\), set
\(C_j=\#\{d\in\mathcal D:d+j\in\mathcal D\}\) and use
\[
 64|p_{\mathcal D}(t)|^2=8+2\sum_{j=1}^{9}C_j\cos(2\pi jt).
\]
All \(C_j\) are nonnegative. The cosine bounds use an enclosure
of \(\pi\) from Machin's identity with alternating arctangent
sums, followed by interval evaluation of the degree \(40\)
Taylor polynomial on \([-\pi,\pi]\). Its error is at most
\((22/7)^{42}/42!\). Directed rounding and integer roots give
upper rational values for \(|p_{\mathcal D}(t)|^r\), including the fractional
powers. Iteration of the resulting positive node operator supplies
the nodal values for \(h\). The validity of the certificate depends
on their strict positivity and the checked inequalities.

On a node cell every inverse branch of \(\mathcal T_r\) lies in a
single linearity interval of \(h\), since \(10\mid N\). For the
centered branch polynomial
\[
 z_d(u)=8^{-1}\sum_{a\in\mathcal D}\e((a-9/2)(u+d)/10)
\]
we have \(|z_d|\le1\), \(|z_d'|<1.9\), and \(|z_d''|<8\).
Indeed the sum of the largest eight distances to \(9/2\) is
\(24\), and the sum of their squares is \(82\).
The convexity of \(z\mapsto |z|^r\) for \(r\ge1\) gives
\((|z_d|^r)''\ge-8r\) and \(|(|z_d|^r)'|\le1.9r\).
These bounds extend across zeros, for example by using
\((|z_d|^2+\varepsilon)^{r/2}\) and taking the distributional
limit. When \(h\) is normalized to have maximum at most one,
the product rule gives
\[
 (\mathcal T_r h-\ell_rh)''
 \ge -r(80+4\operatorname{Lip}h)
\]
inside each cell. A function with second derivative at least
\(-B\) lies below the chord of its endpoint values plus
\(B/(8N^2)\). The certificate checks that every nodal upper
residual plus this allowance is negative. The inequalities
\(\ell_r<10^{\gamma_r}\) are checked by raising to integer powers.

Comparison of \(h\) with positive constants and iteration prove
the estimates. The strict inequalities persist uniformly when
the digit magnitude is floored at a sufficiently small
\(\epsilon\), by compactness and the finiteness of the collection
of alphabets.
\end{proof}

\begin{proposition}[Shifted prime-power moments]\label{dt:digit-pure}
For \(p=2,5\), define
\[
 M_{p,r}(j)=\sup_{u\in\mathbb R}
            \sum_{i<p^j} F_{10^j}((i+u)/p^j)^r.
\]
Then \(M_{p,r}(j)\ll p^{j\beta_{p,r}}\), with the following
uniform exponents.
\[
\begin{array}{c|ccc}
  &r=1&r=7/5&r=2\\ \hline
 p=2& .028& .019& .004\\
 p=5& .2241& .137& .0053
\end{array}
\]
At \(r=7/5\) the same exponents apply with the sufficiently small
fixed floor in Proposition~\ref{dt:digit-moments}. For intermediate
\(r\) the prime-power moments obey the bounds obtained by
interpolating the first and second moments.
\end{proposition}

\begin{proof}
The moments are submultiplicative in \(j\). For \(j=s+t\)
write \(i=h+p^t b\) with \(0\le h<p^t\) and \(0\le b<p^s\).
The lower \(s\) factors sum over \(b\) at cost \(M_{p,r}(s)\).
The remaining factors depend only on
\((10/p)^s(h+u)/p^t\bmod1\), and the numerator multiplier
is a unit modulo \(p^t\). The sum over \(h\) therefore costs
\(M_{p,r}(t)\).

For the first moment, the certificate uses blocks of length \(12\)
or \(20\) at \(p=2\), except for the omitted pair \(\{2,6\}\)
where the block has length \(24\). The respective caps are
\(1.008,1.47,1.585\). At \(p=5\) the block has length \(8\)
with cap \(17.9\). In each case the cap is strictly smaller
than \(p^{j\beta_{p,1}}\). Here is a precise continuum
procedure for verifying the caps. Prepare upper bounds for
\(|p_{\mathcal D}(t)|\) on the uniform \(N=2^{17}\) grid, by the cosine
expansion and upward-rounded square roots in units of \(2^{-36}\).
On a dyadic shift cell \([b/2^d,(b+1)/2^d]\) of midpoint \(u_0\),
factor position \(l\) has argument \(v_l(i+u)/q_l\), where
\(q_l=p^{j-l}\) and \(v_l=(10/p)^l\). Its magnitude throughout the
cell is bounded by the value at the nearest prepared node to
\(v_l(i+u_0)/q_l\), plus
\[
 \frac{19}{2N}+\frac{19v_l}{2^{d+1}q_l}.
\]
This follows from \(\operatorname{Lip}(|p_{\mathcal D}|)\le6\pi<19\).
The bounds can be capped at one. Each multiplication by such a
bound is rounded upwards in units of \(2^{-36}\), then the
weights of residues with the same image modulo \(q_l/p\) are
added. All later factors depend only on that image. The last
scalar bounds the required moment throughout the cell. A cell
whose bound exceeds the cap is bisected. An upper bound
already obtained for a collection of initial factors on a
parent cell can be reused on all its descendants. Every leaf
in the verified subdivision satisfies the cap.

The square moment has an exact integer computation. The Fourier
coefficients of \(F_{10^j}^2\) are nonnegative. The sum on the
\(p^j\)-grid retains only harmonics with frequency divisible by
\(p^j\), so its maximum is attained at \(u=0\). With
\(C_h=\#\{d,d'\in\mathcal D:d-d'=h\}\) its value there is
\[
 \frac{p^j}{64^j}
 \sum_{\substack{h_0,\ldots,h_{j-1}\\
              p^j\mid \sum_{l<j}10^l h_l}}
 \prod_{l<j} C_{h_l}.
\]
After the first \(l\) factors retain the weighted counts of
\(z=p^{-l}\sum_{i<l}10^i h_i\bmod p^{j-l}\) on the part
where the numerator is divisible by \(p^l\). From state \(z\),
insert \(h_l\) only when \(p\mid z+(10/p)^l h_l\) and send
its count, with multiplicity \(C_{h_l}\), to the residue of
\((z+(10/p)^l h_l)/p\bmod p^{j-l-1}\). This gives the stated
sum starting at \(z=0\). The verifier gives \(M_{2,2}(24)<1.068\)
except at \(\{2,6\}\), where \(M_{2,2}(30)<1.084\), and
\(M_{5,2}(24)<1.223\). These caps are below the powers in
the proposition, as checked by integer exponentiation.

The prime-power certificate in the supplement checks eighteen representative
alphabets. Translations and reflections preserve \(|p_{\mathcal D}|\),
which also explains the agreement for omissions \(\{0,d\}\)
and \(\{0,10-d\}\) with \(1\le d\le9\). For the remaining
identifications, write \(B=\{0,1,2,3,4\}\setminus\{j\}\).
The allowed sets omitting \(\{j,j+5\}\) and \(\{2j,2j+1\}\)
are respectively \(B+\{0,5\}\) and \(2B+\{0,1\}\). If their
magnitudes carry subscripts \(D,D'\), the digit products satisfy
\[
 F_{D,Z}(2t)|\cos(\pi t)|=F_{D',Z}(t)|\cos(\pi Zt)|.
\]
This follows by multiplying the one-digit factorizations
and shifting the cosine factor by one decimal position.
The three applications use \((j,D')\) corresponding to
omissions \((0,\{0,1\})\), \((1,\{2,3\})\), and
\((2,\{4,5\})\), together with reflection. To transfer a
shifted grid, choose the representatives of its angles in
\([-1/2,1/2)\) and divide them by two. Splitting by the
parity of the numerator index gives two subsets of shifted
grids of the original size. The reciprocal cosine on the
divided angles is at most \(\sqrt2\). Thus the multiplicative
loss in the norm is at most \(2^{1+r/2}\), uniformly in the
entire word length. Apply the transfer after iterating
the block estimates.

Hölder's inequality supplies the \(7/5\)-moment because
\(7/5=(3/5)\cdot1+(2/5)\cdot2\), and
\((3/5)\beta_{p,1}+(2/5)\beta_{p,2}<\beta_{p,7/5}\).
Consequently a sufficiently long fixed block has this
moment strictly below \(p^{j\beta_{p,7/5}}\), also for the
transferred alphabets. Continuity in the floor on this
fixed block, uniform in the shifts, proves the floored
bound by submultiplicativity. A single floor works for both
propositions by taking the minimum over the finitely many
choices.
\end{proof}

\begin{lemma}[Mixed grids]\label{dt:digit-mixed}
Let \(Q=10^u p^v\), where \(p\in\{2,5\}\), and let \(G=10^{u+v}\)
be the least power of \(10\) divisible by \(Q\). For \(r\) as in
Proposition~\ref{dt:digit-moments},
\begin{equation}\label{dt:digit-mixed-moment}
 \mathcal M_r(Q)
 :=\sup_{x\in\mathbb R}\sum_{i<Q} F_G((i+x)/Q)^r
 \ll10^{u\gamma_r}p^{v\beta_{p,r}} .
\end{equation}
Here \(\beta_{p,r}\) is taken from Proposition~\ref{dt:digit-pure}
at \(r=1,7/5,2\) and by interpolation otherwise. In each case
\(\beta_{p,r}\le\gamma_r\), so the right side in
\eqref{dt:digit-mixed-moment} is \(O(Q^{\gamma_r})\).
The estimate at \(r=7/5\) holds with the floor as well.
For \(Z\ge G\) a decimal power the factors above \(G\)
are constant on the shifted \(Q\)-grid. If \(Z<G\) and \(G/Z\)
is bounded, the same right side also applies to the shifted
moment with \(F_Z\).
\end{lemma}

\begin{proof}
Write \(i=h+p^vb\) with \(h<p^v\) and \(b<10^u\). The first
\(u\) positions on summing over \(b\) give the shifted decimal
bound. The following \(v\) positions depend only on
\((h+x)/p^v\bmod1\), and give the pure prime-power bound.
At \(Z\ge G\), equation~\eqref{dt:digit-product} factors off
\(F_{Z/G}((G/Q)x)\).

These estimates also apply to any bounded number of prescribed
digit positions. Omit each restricted factor, which has modulus
at most one, and split the word at these positions. On the
decimal grid the inverse-branch summation estimates each intact
block by its shifted moment and costs at most ten at each
omission. On the mixed grid first separate the lower \(u\)
decimal positions from the upper \(v\) prime-power positions.
Use submultiplicativity to split the prime-power part and assign
the single-position bound \(p\) at each omission.
This also proves the assertion on dropping a bounded number
of high digits. A bounded summary of shifted envelopes can
be treated termwise using \(|\sum_{j=1}^K z_j|^r\le
K^{r-1}\sum_j|z_j|^r\).
\end{proof}

\begin{lemma}[Sampling of digital moments]\label{dt:digit-sampling}
Let \(r\) be as in Proposition~\ref{dt:digit-moments}. Then
\[
 \int_{\mathbb R/\mathbb Z} F_Z(t)^r\,dt\ll Z^{\gamma_r-1}.
\]
For a set \(\Theta\) of points on the circle with spacing at
least \(\Delta\), and for fixed \(C\), one has
\begin{equation}\label{dt:digit-separated}
 \sum_{\theta\in\Theta}\ \sup_{|t-\theta|\le C/Z}F_Z(t)^r
   \ll_{r,C} (Z+\Delta^{-1})Z^{\gamma_r-1}.
\end{equation}
In particular the complete shifted grid of arbitrary size \(Q\le Z\)
satisfies
\begin{equation}\label{dt:digit-arbitrary-grid}
 \sup_u\sum_{i<Q}F_Z((i+u)/Q)^r\ll_r Q^{\gamma_r}.
\end{equation}
If \(Q\mid Z\), the sum over any \(H\) consecutive indices
in this grid is \(O_r(H^{\gamma_r})\) for \(1\le H\le Q\).
All these estimates remain true with the terminal or leading
restrictions described above.
\end{lemma}

\begin{proof}
Integrate Proposition~\ref{dt:digit-moments} over the shift and
divide by $Z$. Lemma~\ref{common:sampling} now gives the separated
estimate and both grid estimates, in fact without $Q\mid Z$.
For a restricted word apply it to the restricted polynomial:
every retained block has the required integral bound by
Lemma~\ref{dt:digit-mixed}, splitting at its restricted positions.
\end{proof}

\begin{lemma}[Digital occupancy]\label{dt:digit-occupancy}
Every interval of length \(H\) in \([0,Z)\) contains
\(O((1+H)^c)\) members of \(\mathcal W_Z\). For
\(1\le q\le Z\), every residue class modulo \(q\) contains
\(O((Z/q)^c)\) members. The bounds hold for all subsets of
\(\mathcal W_Z\). In particular if \(f_Z\) is the indicator
of the allowed words restricted to units, and \(A_Z=\sum f_Z\),
then \(A_Z\asymp Z^c\) for \(Z\ge10\), with constants
independent of \(\mathcal D\).
\end{lemma}

\begin{proof}
Apply Lemma~\ref{common:occupancy}. At least two of the four unit
digits survive, so the terminal restriction retains at least one
quarter of the words.
\end{proof}

\begin{proposition}[Divisor distribution]\label{dt:digit-type-i}
Let \(X=10^k\), \(L=\log X\), and let \(f\) be the indicator of
the complete allowed word. Its terminal digit may be restricted
to a nonempty set of permitted digits, and a permitted leading digit
may be prescribed.
Write \(A_f=\sum_{n<X}f(n)\). For any fixed \(J,K>0\), uniformly
for \(1\le Q\le X^{.572}\),
\begin{equation}\label{dt:digit-type-i-equation}
 \sum_{\substack{q\le Q\\(q,10)=1}}\tau(q)^K
 \sup_{I,z}\left|\sum_{n\in I} f(n)
     \bigl(\ind_{n\equiv z\bmod q}-q^{-1}\bigr)\right|
 \ll_{J,K} A_f L^{-J}.
\end{equation}
The supremum runs over intervals \(I\subset[0,X)\) and
residues \(z\bmod q\). If \(Y\mid X\), the cyclic average of
\(f\) over translations by multiples of \(Y\) satisfies the
same estimate.
\end{proposition}

\begin{proof}
The transfer inequality in Proposition~\ref{dt:digit-moments} holds with
the cap $\ell_1=2.6789<10^{.428}$ before rounding the exponent.
Integration over the shift gives $\int F_Z\ll Z^{-\sigma}$ with
$\sigma=1-\log_{10}\ell_1$, so $1/2<.572<\sigma<c$. The eight
allowed digits include two consecutive integers, hence their differences
generate $\mathbb Z$. Proposition~\ref{arith:typeI} now applies with
radix ten and $T=.572$, including its end restrictions and its convex
averages of cyclic translates. Each nonempty end restriction costs
only a constant in the mass, so that $A_f\asymp X^c$.
\end{proof}

\begin{corollary}\label{dt:digit-prime-upper}
For two omitted decimal digits $S$, put $u=|\{1,3,7,9\}\setminus S|$.
Then, with $A_{10,S}(x)$ counting ordinary admissible integers,
\[
 \limsup_{x\to\infty}\frac{\pi_{10,S}(x)\log x}{A_{10,S}(x)}
 \le \frac{2}{.572}\,\frac{5u}{16}.
\]
In particular $\pi_{10,S}(x)\ll x^c/\log x$ uniformly for $x\ge3$.
\end{corollary}
\begin{proof}
Apply the common upper bound~\ref{common:upper} with
$\gamma=10u/(8\varphi(10))=5u/16$ and Proposition~\ref{dt:digit-type-i}.
Summing admissible words over lengths gives $A_{10,S}(x)\ll x^c$.
\end{proof}

\section{Fourier separation at two decimal cuts}\label{dt:section}

Retain the notation of Section~\ref{dt:digit-section}, and write
\[
 X=10^k,\qquad L=\log X.
\]
Let \(f_X\) be the indicator
of the length-\(k\) words whose last digit is a unit modulo \(10\).
We extend \(f_X\) periodically whenever we use conditional averages
on the integer circle of length \(X\). Its mass on one period is
\(u8^{k-1}\), where \(u=|\mathcal D\cap\{1,3,7,9\}|\ge2\).

\subsection{The comparison on prime boxes}\label{dt:comparison-setup}

The central analytic statement retains the actual digit weight in the
continuous reference. This distinction matters for boxes whose products
lie in a short interval. Only after aggregation over prime-factor shapes
will we replace the digit weight by its mean.

\begin{proposition}[Comparison of two groups of prime factors]\label{dt:grouped}
Fix a positive \(\kappa\), a positive \(D\), and a fixed number \(r\) of
prime variables at scales \(P_i\gg X^\kappa\). Equip the variables with
smooth weights of fixed-ratio support in their logarithmic coordinates,
and allow bounds by fixed powers of \(L\) for each prescribed derivative.
Allow also a smooth joint weight supported where the product lies
in \([XL^{-D},X)\), with the same derivative bounds in these coordinates.
Suppose that the variables split into two nonempty groups with product
scales
\[
 M=X^mL^{O(1)},\qquad P=X^{1-m}L^{O(1)},
 \qquad  .4195\le m\le .4801 .
\]
Let \(W\) denote the combined weight. Write \(\mathcal B(v)\) for the
pushforward density under \(\mathbf t\mapsto\prod_i t_i\) of the measure
\[
 W(\mathbf t)\prod_{i=1}^r \frac{\dd t_i}{\log t_i}.
\]
Then, for every fixed \(J>0\),
\begin{equation}\label{dt:box-equation}
 \sum_{p_1,\ldots,p_r}W(\mathbf p)f_X\Big(\prod_{i=1}^r p_i\Big)
   =\frac{10}{\varphi(10)}\sum_{v} f_X(v)\mathcal B(v)
     +O_J(X^cL^{-J}).
\end{equation}
The implied constants may depend on the fixed box parameters and the
weight bounds, and are uniform in the choice of the two omitted digits.
\end{proposition}

For clarity we describe the frequency splits before proving this
proposition. If \(Y\mid X\), let \(E_Y\) average an \(X\)-periodic
sequence over each residue class modulo \(Y\). Take
\(Y_0=10^{\lfloor .770k\rfloor}\). The initial term \(E_{Y_0} f_X\)
is treated by character estimates. For each successive decimal power
\(Y=X^\eta>Y_0\) we compare the increment \(E_Y f_X-E_{Y/10}f_X\).
On each such increment there are two choices of lower digit length.
One gives a smaller diagonal term in the row square, and the other
permits stronger estimates for the first moments of the coarse
frequency labels. The two estimates act on complementary multipliers
of the original increment.

Joint weights can be separated throughout by expansion into a Fourier
series on a fixed-ratio logarithmic box after a smooth extension by
zero. The coefficients decay faster than any fixed power of the
frequency, with factors bounded by powers of \(L\). For any desired
inverse-logarithmic accuracy only logarithmic powers of the
frequencies are needed. For the digit-weighted error, the number of
representations of an integer as an ordered product of \(r\) primes is
at most \(r!\), including repeated primes. The continuous error is
estimated in the same way by the logarithmic bounds on the pushforward
density of a box. We therefore first use product weights. Splitting a
single box into a bounded number of comparable boxes is harmless.

\subsection{Complementary filters at two decimal cuts}\label{dt:filters}

For the two cuts \(j=L,H\) use
\begin{equation}\label{dt:cut-parameters}
 h_L=.0918,\qquad h_H=.09693,\qquad
 \lambda_j=m+h_j+.00003,\qquad R_j=10^{\lfloor\lambda_jk\rfloor}.
\end{equation}
The rows have size \(P\) and are divided into intervals of length
\(\asymp X^{h_j}\). Thus the variation of their products with a fixed
column from the group of size \(M\) is \(o(R_j)\), with a fixed power
of \(X\) to spare. We have \(R_j\mid Y/10\) for large \(X\).

For \(R=R_j\) write \(n=w+Rz\) with \(0\le w<R\).
Let \(L_R(w)\) be the lower word indicator, including the unit last
digit, and let \(U_R\) be the periodic upper step function which on
the step \(Rz+[0,R)\) equals
\[
 \ind_{\mathcal W_{Y/R}}(z\bmod Y/R)
       -\frac{8}{10}\ind_{\mathcal W_{Y/(10R)}}(z\bmod Y/(10R)).
\]
Counting the free higher positions in the conditional average gives
the identity
\begin{equation}\label{dt:outer-factor}
 (E_Y-E_{Y/10})f_X(n)=(X/Y)^{-\delta} L_R(n\bmod R)U_R(n).
\end{equation}
The lower density is \(\rho_L\asymp R^{-\delta}\) and the upper
density scale is \(\rho_U=(Y/R)^{-\delta}\).
The upper steps have bounded height, and \(O((X/R)\rho_U)\)
of them are occupied on \([0,CX]\) for any fixed \(C\). The same
bound applies to the jumps. Fixed powers of \(L\) in product
intervals or cutoffs only introduce powers of \(L\) in subsequent
estimates. In particular,
\((X/Y)^{-\delta}X\rho_L\rho_U\asymp X^c\).

Set
\[
 v=\lambda_H/2+.00001,\qquad w_0=\lambda_L-v-.00015,\qquad
 V_*=10^{\lfloor vk\rfloor},\qquad B_*=10^{\lfloor w_0k\rfloor}.
\]
All uses of \(v,w_0\) as scale exponents allow the vanishing errors
from rounding. Write
\[
 \widetilde F_{B_*}(t)=
       \prod_{\ell<\log_{10}B_*}
         \max(|p_{\mathcal D}(10^\ell t)|,\varepsilon_*),
\]
with a sufficiently small fixed \(\varepsilon_*>0\).
This floor is used only to construct the filters and to majorize their
coarse moment weights. Propositions~\ref{dt:digit-moments} and
\ref{dt:digit-pure} give the same bounds for the required moments of
these floored factors. In the \(V_*\)-rescaled
frequency circle take the two parts of the threshold
\(\widetilde F_{B_*}(t)=B_*^{-.306}\).
Convolve their indicators with the same nonnegative trigonometric kernel
of integral one and degree \(O(B_*X^{.00005})\). For example, normalize
the square of the modulus of a Fourier sum with coefficients obtained
by sampling a nonzero smooth bump. Repeated summation by parts gives
the bound \(O_A(K(1+K\|t\|)^{-A})\) for each fixed \(A\), with
\(K\asymp B_*X^{.00005}\).

Write \(\omega_L(\theta),\omega_H(\theta)\) for the resulting
polynomials at \(t=V_*\theta\), with the low-amplitude region assigned
to \(L\). They lie in \([0,1]\) and sum to one. The log of the
floored polynomial is Lipschitz with constant \(O_{\varepsilon_*}(B_*)\).
Thus for any fixed \(C>0\),
\begin{align}
 \omega_L(\theta)&\ll
  \ind_{\widetilde F_{B_*}(V_*\theta)\le2 B_*^{-.306}}+X^{-C},\label{dt:low-filter}\\
 \omega_H(\theta)&\ll
  \ind_{\widetilde F_{B_*}(V_*\theta)\ge B_*^{-.306}/2}+X^{-C}.
 \label{dt:high-filter}
\end{align}
Indeed on a ball of radius \(X^{-.000025}/B_*\) the amplitude varies
by \(1+o(1)\), and integration outside the ball is \(O_C(X^{-C})\).
There is no change in degree when increasing the fixed order of
decay used in this reasoning. The physical convolution corresponding
to either multiplier involves only translations by multiples of \(V_*\)
of size \(O(V_*B_*X^{.00005})=o(R_L)\) by a fixed power.

To commute these convolutions with the upper step, split \(L_R\) into
at most eight pieces according to its leading digit \(d\).
For \(1\le d\le8\) the pieces occupy a subinterval of a step
separated from both endpoints by \(R/10\). For \(d=0\) translate
the upper step by \(-R/20\); for \(d=9\) translate it by \(R/20\).
In each case the product in \eqref{dt:outer-factor} for this piece is
unchanged, and the new step boundaries are at least \(R/20-O(1)\)
from its support. Therefore any convolution shift crossing a new
boundary has zero input lower weight. The filters commute exactly
with the translated upper step for each piece. Both factorizations
start from the same \(X\)-periodic increment. Applying the two
complementary multipliers in that ambient circle and adding therefore
recovers \eqref{dt:outer-factor} exactly, even though the factorizations
and boundary translations at the two cuts differ. All translations
preserve the unit last digit.

\subsection{Diagonal terms and small additive periods}\label{dt:diagonal}

For a lower piece denote by \(A_q(i)\) the normalized complex Fourier
coefficient at \(i/q\), for \(q\mid R_j\) and \((i,q)=1\), and by
\(G_q\) its exact period-\(q\) component after filtering. Its Fourier
coefficients on the \(q\)-circle are \(\rho_L A_q(i)\omega_j(i/q)\),
where the density now means that of the piece. Fixing the leading
digit replaces a single digit factor by a phase. A bound for a
shifted norm on a decimal or \(2,5\)-smooth grid persists with
constant loss when that factor occurs. One may split off the
at most ten residue cosets lost by cutting the highest position.
For a square sum, direct Parseval for the resulting shorter
polynomial gives the same \(O(q^\delta)\) bound. The unit last
digit is likewise a bounded modification to the norm estimates.

After short-row Cauchy--Schwarz and smooth column completion,
the diagonal of the row square, relative to
\((X\rho_L\rho_U)^2\), is bounded by
\begin{equation}\label{dt:diagonal-main}
 X^{\delta(\eta-\lambda_j)-h_j+o(1)}
       \sum_{(i,q)=1}|A_q(i)\omega_j(i/q)|^2 .
\end{equation}
Indeed the Cauchy factor is \(O((X/X^{h_j})\rho_U L^{O(1)})\).
In the diagonal, each integer product occurs at most a divisor
power number of times, which is \(X^{o(1)}\) for a bounded number
of factors. Summing over complete lower periods on an interval of
length \(O(X)\) then gives the usual normalized square mean by
Parseval. The same proof gives \eqref{dt:diagonal-main} for
subsets of the exact period. The unfiltered square sum
is \(O(q^\delta)\), so the \(H\)-cut diagonal has exponent
at most \(\delta\eta-h_H<0\).

For the \(L\)-cut diagonal split the digits at \(V_*\) and
\(V_*B_*\). Put \(D_1=(q,V_*)\), \(Q=q/D_1\), \(D_2=(Q,B_*)\),
and \(Q_2=Q/D_2\). At fixed residue modulo \(Q\), summing
the bottom squares on the shifted \(D_1\)-grid costs
\(O(D_1^\delta)\). The multiplier \(V_*/D_1\) of the
remaining grid is a unit modulo \(Q\).
At fixed residue modulo \(Q_2\), use \eqref{dt:low-filter} and
bound the square of the middle block by a constant times
\(B_*^{-.6(.306)}\widetilde F_{B_*}^{7/5}\).
Its shifted \(D_2\)-grid costs \(O(D_2^{.258})\).
The last block sees a shifted \(Q_2\)-grid, since \(B_*/D_2\)
is coprime to \(Q_2\). Also \(Q_2\mid R_j/(V_*B_*)\) by the
divisibility of these positional cuts. Its squares cost
\(O(Q_2^\delta)\). On the supported low shell this proves
\[
 \sum_{(i,q)=1}|A_q(i)\omega_L(i/q)|^2
   \ll D_1^\delta D_2^{.258} Q_2^\delta B_*^{-.6(.306)}
       +O_C(X^{-C}).
\]
The three grid sizes grow with the valuations of \(q\), so at
\(q\mid R_j\) their product with the indicated powers is at most
the value at \(q=R_j\). The bound is thus
\(O(R_j^\delta B_*^{.258-\delta-.6(.306)})\).
In \eqref{dt:diagonal-main} this gives the saving
\begin{equation}\label{dt:diagonal-saving}
 \delta+(.258-\delta-.6(.306))w_0-h_L<-.0005.
\end{equation}
The uniform power error from the complement of the low shell is
harmless throughout. The two endpoint checks in \(m\) suffice
because the expression is affine.

At either cut the low multiplier has power-small square sum
whenever \(q<X^{.192}\). In fact \(\log_X Q<.087+o(1)\):
at most one valuation of \(q\) can exceed that of \(V_*\),
and the excess exponent is bounded by
\(\max(\kappa_2(\lambda_H-v),.192-\kappa_5v)\), where
\(\kappa_p=\log_{10}p\).
The first bottom-grid square bound followed by the middle
\(7/5\) moment on its full shifted \(Q\)-grid gives exponent
\[
 .097(.192)+.258(.087)-.6(.306)w_0<-.005 .
\]
For this step retain the bottom block of decimal length \(Z_0\le Q<10Z_0\)
inside \(B_*\), dropping factors at most one. This estimate is for a
floored product, not a trigonometric polynomial. The log of the retained
product is Lipschitz with constant \(O_{\varepsilon_*}(Z_0)\), so its
\(7/5\)-power at each shifted grid point is at most a constant times its
average on the surrounding cell of length \(1/Q\). Integration of the
floored shifted decimal moment in \eqref{dt:digit-decimal-moment} gives
an integral \(O(Z_0^{.258-1})\); summing the disjoint cells therefore
costs \(O(Q^{.258})\), without any divisibility assumption. In a translated upper step the
lower \(q\)-period repeats exactly. Summing absolute values using
Parseval and Cauchy--Schwarz gives total
\(O(X\rho_L\rho_U X^{-.002+o(1)})\) for that component, both
with bounded-multiplicity product coefficients and against the
smooth density. Consequently small periods in the low branch
can be discarded. In the high branch one can replace its
multiplier by one at these periods: the omitted multiplier on
that cut is the same power-small low multiplier.

\section{The signed row Gram}\label{dt:gram}

Use the word polynomials \(\Pi_B,F_B\) and the unrestricted
sets \(\mathcal W_B\) of Section~\ref{dt:digit-section}.
Write \(z_+=\max(z,0)\) and \(\kappa_p=\log_{10}p\) for \(p=2,5\).
Use \(\Pi_B^\circ\) for the normalized polynomial after an
optional restriction on the ends of the word, and
\(\mathcal W_B^\circ\) for its word set and
\(F_B^\circ=|\Pi_B^\circ|\). In this section the lower
word at the cut has a unit terminal digit and may have a
prescribed leading digit. For the smallest block of a split
word only the terminal restriction is relevant.

Propositions~\ref{dt:digit-moments} and~\ref{dt:digit-pure} and
Lemmas~\ref{dt:digit-mixed} and~\ref{dt:digit-sampling} give the
shifted first exponents \(a=.428\) on decimal and mixed grids,
and \(a_2=.028,\ a_5=.2241\) on pure grids. The shifted square
norm on a grid of size \(D\) costs \(O(D^\delta)\) whenever it fits
in the word, including bounded end restrictions. Put
\begin{equation}\label{dt:gram-coarse-exponents}
 C=\frac{3c}{4}-a,\qquad L_p=(1-2a_p)\kappa_p,\qquad
                     \sigma=\min(C,L_2,L_5)>.249.
\end{equation}

\subsection{Completion and the signed lift}

Use row and column lengths \(P=X^{n+o(1)}\), \(M=X^{m+o(1)}\),
where \(n=1-m\) and \(.4195\le m\le .4801\). The \(o(1)\) here
permits fixed logarithmic distortions. At an outer decimal
period \(Y=X^\eta\), \(.770\le\eta\le1\), let \(R=X^{\lambda+o(1)}\)
be the lower decimal cut and \(H=X^{h+o(1)}\) the row interval
length. We will use
\begin{equation}
  h\in\{.0918,.09693\},\qquad \lambda=m+h+.00003,
  \qquad
  \rho_U=(Y/R)^{-\delta},\quad \rho_L\asymp R^{-\delta},\quad
  A_U=\delta(\eta-\lambda).
 \label{dt:gram-scales}
\end{equation}
The upper source is a bounded step function on intervals of
length \(R\), possibly translated by a fixed fraction of \(R\).
In a fixed-multiple product box it has
\(O((X/R)\rho_U X^{o(1)})\) occupied steps, and the same estimate
holds for the sum of their squared amplitudes. Row and column
coefficients are supported on units modulo \(10\) and are bounded
by \(X^{o(1)}\). For an exact additive period \(q\mid R\) the lower
source has the expansion
\begin{equation}
 G_q(t)=\rho_L\sum_{a\bmod q}^{*}
                    J(a)\Pi_R^\circ(-a/q)\e(at/q),
 \label{dt:gram-lower-source}
\end{equation}
where the star requires \((a,q)=1\), and \(|J|\le1\).
The factors \(J\) will include the taper and, where applicable,
the restriction to the complement of a union of frequency bands.

The following reduction explains our normalization. Partition the
rows into intervals \(I\) of length comparable to \(H\).
For a fixed column \(y\asymp M\), the upper source changes
at most \(O(1)\) times on each \(I\), because \(HM\ll R\).
Each occupied step meets \(O(1+R/(HM))\) such intervals.
Cauchy over the column and interval therefore bounds the
square of the original product sum by
\[
 X^{o(1)}\frac{PM\rho_U}{H}
   \sum_{y\asymp M}\sum_I
       \max_{J_0\subset I}
       \left|\sum_{p\in J_0}u_pG_q(py)\right|^2 ,
\]
where \(J_0\) ranges over subintervals. A dyadic decomposition of
each \(J_0\) costs a logarithmic factor and allows the maximum
to be replaced by the sum of squares over the dyadic intervals.
At every binary length these intervals are disjoint. The column
sum of these nonnegative squares is majorized by completing with
a nonnegative smooth cutoff at scale \(M\). Expand the squares
and apply Poisson summation to this cutoff on the residue classes
modulo \(q\). Apart from an arbitrarily power-small error, the
off-diagonal terms are supported on
\begin{equation}
 p a-(p+\ell)b\equiv t\pmod q,\qquad
      0<|\ell|\ll H,\qquad |t|\ll T,\qquad
      T=(1+q/M)X^\epsilon .
 \label{dt:gram-incidence}
\end{equation}
The coefficient of a twist has absolute value \(O(M)\), and is
independent of the labels and rows. The cutoff may be enlarged
to retain all integer twists in a symmetric interval of this
size for estimates by nonnegative weights. Here and below
\(\epsilon>0\) can be arbitrarily small. After dividing the
squared product sum by \((PM\rho_U\rho_L)^2\), an off-diagonal
incidence bound for the normalized word coefficients is divided
by \(PH\rho_U\), apart from \(X^{O(\epsilon)+o(1)}\).
For the diagonal there is no need to expand in twists.
Group by the integer product \(py\), use its divisor bound,
and apply Parseval on complete lower periods. The normalized
diagonal is at most
\begin{equation}
 X^{A_U-h+o(1)}
              \sum_{a\bmod q}^{*}|J(a)\Pi_R^\circ(a/q)|^2 .
 \label{dt:gram-diagonal}
\end{equation}
This form allows the use of the filtered square norm.

\paragraph{The correlation on one fibre.}
The part of the row estimate which uses the signs is independent of
the radix. Keep both labels on a complete fibre before taking absolute
values; otherwise the saved square root of the overlap is lost.

\begin{lemma}[Correlation on a complete lift]\label{dt:finite-signed-lift}
Let $S\ne\varnothing$ be a subset of $\{0,\ldots,V-1\}$,
$N=|S|$, $\Pi(t)=N^{-1}\sum_{x\in S}\e(tx)$, and let $r$ be a unit
modulo $V$. For arbitrary real shifts $\alpha,\beta$ put
\[
 C_r(w)=\sum_{y\bmod V}\Pi((ry+w+\alpha)/V)
                      \overline{\Pi((y+\beta)/V)}.
\]
Writing $[rx]_V$ for the representative in $[0,V)$, one has
\begin{equation}
 C_r(w)=\frac V{N^2}\sum_{x\in S,\,[rx]_V\in S}
       \e\big((wx+\alpha x-\beta[rx]_V)/V\big).
 \label{dt:gram-intercept-identity}
\end{equation}
If $w$ runs over $J\ge1$ successive terms of an arithmetic progression
of unit step modulo $V$, then, with $B_r=|S\cap r^{-1}S|$,
\[
 \sum_w |C_r(w)|\le J\,(V/N^2)\sqrt{B_r}\,(1+V/J)^{1/2}.
\]
The sum counts steps with multiplicity. For $V=1$ the singleton
identity applies without any unit-step condition.
\end{lemma}
\begin{proof}
Orthogonality in $y$ forces the second digit index to equal $[rx]_V$.
Parseval in $w$ gives the complete mean square $(V/N^2)^2B_r$.
Each intercept occurs at most $1+J/V$ times; sum the nonnegative
squares and apply Cauchy--Schwarz.
\end{proof}

\begin{lemma}[Signed lift and fold]\label{dt:gram-base}
For the incidence \eqref{dt:gram-incidence}, let \(g\) be the
exact order of \(a-b\) in the additive group modulo \(q\).
Write
\begin{equation}
 \begin{gathered}
 g=p_0^ip_1^j,\quad \{p_0,p_1\}=\{2,5\},\quad i\ge j,\qquad
 s=i/k,\quad d=j/k,\quad \xi=\log_X g,\\
 V=10^{\lfloor j/2\rfloor},\quad Q=q/V,\quad
 g_0=g/V,\quad Z=q/g,\quad G=10^i,\\
 K_0=p_1^{i-j},\quad d_2=(K_0,Z),\quad
 D_0=Z/d_2,\quad K=K_0/d_2.
 \end{gathered}
 \label{dt:gram-order}
\end{equation}
Suppose the two label multipliers are constant on the
\(V\)-lifts, that is, on increments of \(Q\) in the label.
The normalized contribution at fixed \(g\) is at most
\begin{equation}
 X^{A_U-Cd-L_{p_0}(s-d)+o(1)}
       \max(1,g/P)^{1/2}D_0^\delta(1+T/D_0).
 \label{dt:gram-base-bound}
\end{equation}
If the gaps are restricted to \((\ell,d_2)=e_0\), the
completion factor can be replaced by \(1+T/(e_0D_0)\).
The contribution can also be estimated by the finer folded
incidence over \(z_0\bmod D_0\) in the proof below.
In particular \eqref{dt:gram-base-bound} has exponent at most
\begin{equation}
 \delta(\eta-s)-Cd-L_{p_0}(s-d)+(s-m)_+
                                    +\tfrac12(\xi-n)_+,
 \label{dt:gram-base-exponent}
\end{equation}
with the usual arbitrarily small completion loss.
\end{lemma}

\begin{proof}
Fix the gap $\ell$, the completion twist $t$, and the coarse labels
$a_0,b_0\bmod Q$. Write $a=a_0+Qx$, $b=b_0+Qy$. Since $V^2\mid g$
and $Q=Zg_0$, the order condition on the entire fibre is
$a_0-b_0=Zu$, $(u,g_0)=1$: every lift adds a multiple of $g_0$
to $u$, and $g_0$ has the same prime divisors as $g$.
Primitivity of the labels is constant because every prime of $q$
divides $Q$. The coarse congruence first requires $Z\mid t+\ell b_0$;
with $t'=(t+\ell b_0)/Z$ it then reads $pu\equiv t'\pmod{g_0}$.
Thus the rows lie in one progression of step $g_0$, also when $g_0=1$.

For $V>1$ the supported rows $p,p+\ell$ are units modulo $V$.
On their progression $V\mid g_0$ fixes both residues. Solving the
original incidence modulo $q$ now gives $x=ry+w\pmod V$, where
\[
 r=(p+\ell)p^{-1},\qquad w=p^{-1}\big((t'-pu)/g_0\big)\pmod V.
\]
The quotient is taken in the integers. Incrementing $p$ by $g_0$
increments $w$ by the fixed unit $-up^{-1}$. In the bottom word take
$S=\mathcal W_V^\circ$, $\alpha=a_0/Q$, $\beta=b_0/Q$ in
Lemma~\ref{dt:finite-signed-lift}, conjugating if both Fourier signs
are reversed. All higher digit factors and both label multipliers are
constant on the respective lifts. Row weights and dyadic restrictions
are independent of $x,y$; only after performing the signed sum may
we bound these weights and enlarge to the whole row interval.

Indeed $\log_X g_0=\kappa_{p_0}(s-d)+d/2+o(1)\le .699\lambda+o(1)<n$.
The enlarged progression has $J\asymp P/g_0$ steps (a containing
interval of this length suffices). Since $N\asymp V^c$, $B_r\le N$
and $Vg_0=g$, the row sum is bounded, up to $X^{o(1)}$, by
\[
 (P/g_0)V^{\delta-c/2}\max(1,g/P)^{1/2}.
\]
If $V=1$ this is just row counting, including $g=1$. This is also
the fibre calculation for the other radices: the septenary argument
has no excess prime-power block, whereas radix six, like ten, must
retain that block and the common-residue congruence below.

Now sum the coarse variables on their common residue
\(z\bmod Z\). Above the \(V\)-block the coarse grid of size
\(g_0\) splits into a decimal grid of size
\(10^{j-\lfloor j/2\rfloor}\) and a pure \(p_0\)-grid of
size \(p_0^{i-j}\). In general for a decimal divisor
\(B\) of a grid of size \(B D\), write a grid index as \(u+Dv\).
The lower \(B\)-word is sampled on its shifted \(B\)-grid
at fixed \(u\), and all higher factors are independent of
\(v\). For a pure grid of size \(p^\ell\), the same splitting
at \(p^b\) for the first \(b\) positional digits leaves a
multiplication by the complementary unit \((10/p)^b\) on
the second grid. This proves the mixed first bound, in either
coarse variable, by
\begin{equation}
 O\left(
    (10^{j-\lfloor j/2\rfloor})^a
    (p_0^{i-j})^{a_{p_0}}
    F_{R/G}^\circ(Kz/D_0)\right).
 \label{dt:gram-first}
\end{equation}
A multiplier constant on each lift may remain inside these
first sums. In particular, a pointwise blockwise gain in a
coarse first norm is a gain in \eqref{dt:gram-base-bound}.
The product of the two ordinary first-norm factors, the row
density \(g_0^{-1}\), and the signed mean \(V^{\delta-c/2}\)
has exponent \(-Cd-L_{p_0}(s-d)+o(1)\).
Here the upper factor retains the prescribed leading digit
if there is one, and retains the terminal restriction when
\(G=1\). If an end restriction instead falls in a summed
block, the sampling losses are bounded by
Lemma~\ref{dt:digit-mixed}.

For the remaining square write \(z=z_0+D_0 b\) with
\(z_0\bmod D_0,\ b\bmod d_2\). The congruence is equivalent to
\begin{equation}
 D_0\mid t+\ell z_0,\qquad
 \ell b\equiv -(t+\ell z_0)/D_0\pmod{d_2}.
 \label{dt:gram-fold}
\end{equation}
At a fixed \(e_0=(\ell,d_2)\) it requires \(e_0\mid t\)
and there are \(e_0\) possible \(b\)'s when it holds.
There are \(O(H/e_0)\) nonzero gaps. The \(e_0\) lifts
cancel this multiplicity after division by \(H\).
For fixed \(z_0,\ell\) there are \(O(1+T/(e_0D_0))\)
quotients \((t+\ell z_0)/(e_0D_0)\).
The number of choices of \(e_0\mid d_2\) is \(O(\log X)\).
The high square depends only on \(z_0\) and is sampled
on the \(D_0\)-grid. The valuation identities
\begin{equation}
 GD_0=qK,\qquad (K,D_0)=1,\qquad D_0\mid R/G
 \label{dt:gram-fold-identities}
\end{equation}
follow directly from \eqref{dt:gram-order}. They give the
factor \(O(D_0^\delta)\) and prove \eqref{dt:gram-base-bound}.
Since \(D_0\le R/G\) and \(q/D_0\le G\), we may bound the
square and the completion factor by their separate worst
values at \(q=R\). This gives \eqref{dt:gram-base-exponent}.
\end{proof}

\begin{lemma}[The zero twist]\label{dt:gram-middle}
Suppose \(q<MX^{-c_0}\) with fixed \(c_0>0\). The
off-diagonal bound, for multipliers constant on the signed
lifts, has exponent at most
\begin{equation}
                A_U-\sigma\log_X q+o(1).
 \label{dt:gram-zero-exponent}
\end{equation}
\end{lemma}

\begin{proof}
The smooth column completion retains only \(t=0\) up to
an arbitrarily small power error. Since \(z\bmod Z\) is
a unit, \eqref{dt:gram-incidence} forces \(Z\mid\ell\).
There are \(O(H/Z)\) such nonzero gaps. There is no incomplete
row factor, as \(g\le q<M\le P\) up to bounded scale factors.
Use the unrestricted high square trace in \(z\), of size
\(O(d_2D_0^\delta)\) by \eqref{dt:gram-fold-identities}.
The remaining high factor after dividing by \(H\) is
\(D_0^{\delta-1}\). Since \(\log_XD_0\ge(\log_Xq-s)_+\),
the exponent is at most
\(A_U-\sigma s-c(\log_Xq-s)_+\le A_U-\sigma\log_Xq\).
\end{proof}

\subsection{Joint affine averaging}

\begin{lemma}[Joint affine overlap and intercept mean]\label{dt:gram-joint}
In the setting of Lemma~\ref{dt:gram-base}, suppose
\(\xi\ge m+.020\) and
\begin{equation}
                 h+\xi-m<d/2-\epsilon_1
 \label{dt:gram-joint-fit}
\end{equation}
for a fixed \(\epsilon_1>0\). Then
\eqref{dt:gram-base-exponent} is improved by subtracting
\begin{equation}
         \delta(\xi-n)_+
             +\frac{\delta}{2}\min\{d/2,\,2(h+\xi-m)\}.
 \label{dt:gram-joint-gain}
\end{equation}
The improvement is uniform in the coarse label pair and
can be combined with a gain inside the two coarse first sums.
\end{lemma}

\begin{proof}
Use complete unrestricted words at auxiliary scales. For
decimal lengths \(J,U\mid W\) put
\[
 M_r(U)=\max_{b\bmod U}\#\{x\in\mathcal W_U:
                                    rx+b\bmod U\in\mathcal W_U\}.
\]
We first prove, with uniform average over units \(r\bmod W\),
\begin{equation}
        \mathbb E_r[B_r(J)M_r(U)]\ll J^{2c-1}U^{2c-1}.
 \label{dt:gram-joint-overlap}
\end{equation}
For \(q_1\mid W\), condition on any unit class \(r\bmod q_1\).
The induced distribution modulo \(J\) is uniform on the unit
class modulo \(g_1=(J,q_1)\). For fixed \(x\) with
\((x,J)=m_0\), the images of \(x\) lie in one residue class
modulo \(d_0=(g_1m_0,J)\), and each image has probability at
most \(O(d_0/J)\). Indeed the unrestricted progression of
ratios gives multiplicity \(d_0/g_1\), and its unit part has
cardinality comparable to \(J/g_1\). A residue class modulo a
divisor \(b\) of \(J\) contains \(O((J/b)^c)\) words: fix the
high block after a decimal power comparable to and below \(b\).
Using \(d_0\le g_1m_0\), and summing the resulting costs over
\(m_0\mid J\), yields
\begin{equation}
 \mathbb E[B_r(J)\mid r\bmod q_1]\ll
  J^{2c-1}g_1^\delta\sum_{m_0\mid J}m_0^{1-2c}
                       \ll J^{2c-1}g_1^\delta .
 \label{dt:gram-joint-conditional}
\end{equation}
To bound the affine overlap, Fourier inversion and grouping
the additive labels by exact order give
\[
 M_r(U)\le U^{2c-1}\sum_{q_1\mid U}
              \sum_{b\bmod q_1}^{*} F_U(b/q_1)F_U(rb/q_1).
\]
The inner nonnegative sum depends only on \(r\bmod q_1\).
Its mean is \(O(q_1^{2a-1})\) by transitivity of the unit
group and the mixed first norm. Combining with
\eqref{dt:gram-joint-conditional} proves
\eqref{dt:gram-joint-overlap}, since
\(\sum_{q_1=2^u5^v}q_1^{2a-c}<\infty\).
Moreover, if \(W\mid V\) and \(U\mid V/W\), fixing the
lowest \(W\)-word leaves an affine action on the next block.
Consequently
\begin{equation}
          B_r(V)\ll (V/W)^c B_r(W)M_r(U)/U^c .
 \label{dt:gram-affine-lift}
\end{equation}
Restrictions at the lowest digit of the largest word can be
discarded in all these upper bounds.

Fix the coarse labels and abbreviate
\(\mathcal A=HT/Z,\ \mathcal L=\max(1,g/P)\).
The argument up to \eqref{dt:gram-joint-factor} needs two scale
conditions: $1\le\mathcal L\le\mathcal A$ and
$\mathcal A\ll VX^{-\epsilon_1/2}$. Here they follow from
$\xi\ge m+.020$, $P\ge M$ up to scale constants, and
\eqref{dt:gram-joint-fit}, on taking the completion loss small enough.
The arithmetic inputs to the overlap mean are residue occupancy,
$2c>1$ and $2a<c$ for convergence of the two smooth-divisor sums.
For fixed coarse labels set $t'=(t+\ell b_0)/Z$. The supported rows
give both $t'\equiv pu$ and $t'+\ell u\equiv(p+\ell)u$ units modulo
$V$, so $r=1+\ell u/t'$ is a unit there. In particular a common
divisor of $\ell,t'$ is invertible modulo every $W\mid V$, as needed
when reducing the ratio.
Group gap--twist pairs by the positive gcd \(e_0=(\ell,t')\).
For fixed \(e_0\), choose a decimal \(W\mid V\) comparable
to \(\mathcal A/e_0^2\), sufficiently large to exceed the
cross-determinants, when this scale exceeds a fixed constant.
The reduced primitive pairs \((\ell/e_0,t'/e_0)\) inject
up to sign into the resulting ratios \(r\bmod W\).
Indeed equality of ratios forces their determinant to be
a multiple of \(W\), whereas this determinant equals
\((\ell_1t_2-\ell_2t_1)/(Ze_0^2)\) and is smaller than \(W\).
At scales below a sufficiently small constant the determinant
vanishes automatically and there are at most two pairs.
Bounded transition scales may be included with \(W\asymp1\).

For the mean of \eqref{dt:gram-intercept-identity} along the
row progression, apply Cauchy and majorize the intercept
interval of length comparable to \(\min(P/g_0,V)\) by a
translated triangular weight of comparable length. Its
normalized Fourier kernel has decay \(O(2^{-2b})\) on a
physical circular interval of radius
\(L_b=2^b\mathcal L\) in the unit-step-dilated differences.
After expanding the squared signed identity, for each \(x\)
the possible \(x'\) in this interval project into the
overlap at any decimal divisor \(W_b\mid V\). A circular
interval of length \(L_b\) of the dilated residues contains
at most \(O((L_b/W_b)B_r(W_b))\) of these points if \(W_b\le L_b\),
since the intercept step is a unit. Thus the squared mean
is at most
\begin{equation}
 \ll (V/|\mathcal W_V^\circ|^2)^2 B_r(V)
        \sum_{b\ge0}2^{-2b}(L_b/W_b)B_r(W_b).
 \label{dt:gram-fejer}
\end{equation}
The bounds also follow by a complete period if \(P/g_0\ge V\).
In the fixed-gcd sum choose
\(U=\min(W,V/W)\) and \(W_b\asymp\min(L_b,W)\) dividing \(W\).
By \eqref{dt:gram-affine-lift} and Cauchy over \(r\bmod W\),
use the mean bound \(W^{2c-1}\) for \(B_r(W)\) and
\eqref{dt:gram-joint-overlap} for \(M_r(U)B_r(W_b)\).
Summing the square roots in \eqref{dt:gram-fejer} using its
dyadic decay yields, for this gcd, the bound
\begin{equation}
 \ll V^{\delta-c/2}\,
           W^{1-\delta/2}U^{-\delta/2}
                 \mathcal L^{1/2}\min(\mathcal L,W)^{-\delta}
 \label{dt:gram-joint-per-gcd}
\end{equation}
on the sum of the signed intercept means.

The \(W\ge\mathcal L\) terms in
\eqref{dt:gram-joint-per-gcd} sum, apart from
\(V^{\delta-c/2}\mathcal L^{1/2-\delta}\), to
\(O(\mathcal A\min(V,\mathcal A^2)^{-\delta/2})\).
On \(W\ge\sqrt V\) the summand is \(W V^{-\delta/2}\).
On \(W\le\sqrt V\) it is \(W^{1-\delta}\). If
\(\mathcal A\ge\sqrt V\) the latter tail starts at
\(e_0\gg(\mathcal A/\sqrt V)^{1/2}\) and costs
\(O(\mathcal A^{1/2}V^{1/4-\delta/2})\).
If \(\mathcal A\le\sqrt V\) the full sum is
\(O(\mathcal A^{1-\delta})\).

For \(W<\mathcal L\), the same summation starting at
\(e_0\gg(\mathcal A/\mathcal L)^{1/2}\) costs
\(O(\mathcal A^{1/2}\mathcal L^{1-2\delta})\) when
\(W\le\sqrt V\), and
\(O(V^{-\delta/2}\mathcal A^{1/2}\mathcal L^{1-\delta})\)
otherwise, apart from \(V^{\delta-c/2}\).
They are both absorbed by the preceding bound.

Finally there are \(O(H)\) below-threshold gcds and
at most two gap--twist pairs at each. In
\eqref{dt:gram-fejer} use \(W_b\asymp\min(L_b,V)\)
and \(B_r(W_b)\le W_b^c\) for these, giving
\(O(HV^{\delta-c/2}\mathcal L^{c/2})\).
Dividing by the ordinary bound
\(\mathcal A V^{\delta-c/2}\mathcal L^{1/2}\),
the relative gain in the intercept and incidence factor is
\begin{equation}
   \min(V,\mathcal A^2)^{-\delta/2}\mathcal L^{-\delta}
                       +(H/\mathcal A)\mathcal L^{-\delta/2}.
 \label{dt:gram-joint-factor}
\end{equation}

On this sector, for \(x=\xi-m\ge .020\),
\(\log_X\mathcal A=h+x+O(\epsilon)+o(1)\) and
\(\log_X\mathcal L=(\xi-n)_+\le x\). The second term of
\eqref{dt:gram-joint-factor} is absorbed by the first,
since $-x+\delta(h+1.5x)<0$ for these decimal cuts.
Without this last inequality the two terms in \eqref{dt:gram-joint-factor}
must be kept separately, as in the six-adic argument.
The signed estimate is uniform in the fixed coarse labels
before taking their first sums. The high square factor
over \(z\bmod Z\) is \(d_2D_0^\delta\). Its multiplicity
\(d_2\) combines with the main normalized incidence \(T/Z\)
to give \(T/D_0\). Thus the usual folded high bound applies
and the gain in \eqref{dt:gram-joint-gain} follows.
\end{proof}

\subsection{A terminal average for orders with excess powers of two}

\begin{lemma}[Averaged pure tail]\label{dt:gram-tail}
Use the notation of Lemma~\ref{dt:gram-base} with \(p_0=2\).
Suppose \(s-d\ge.05\), and write
\(\omega=(i-v_5(q))_+/k\) and \(u=(\log_X(T/D_0))_+\).
If \(\omega\ge.024\), the completion exponent \((s-m)_+\)
in \eqref{dt:gram-base-exponent} can be replaced by
\((s-m-.024\kappa_5)_+\), up to the completion loss.
If \(u\le.075\), it can be capped at \(.075\).
If \(\omega<.024\) and \(u>.075\), the bound
\eqref{dt:gram-base-exponent} gains a power of at least
\(.0158+o(1)\).
\end{lemma}

\begin{proof}
The first two claims follow from \(q/D_0=G5^{-k\omega}\)
and the actual twist count in \eqref{dt:gram-base-bound}.
For the last claim, \(u>.075\) implies \(q>M\) by a power,
and the true completion exponent is
\(s-m-\kappa_5\omega+O(\epsilon)+o(1)\).
Group the gaps by \(e_0=(\ell,d_2)\) and put
\(\vartheta=\log_X e_0\). If \(\vartheta>.0158\), the
factor \(1+T/(e_0D_0)\) pays the required gain by itself.
When this factor grows, its savings over \((s-m)_+\) is
\(\vartheta+\kappa_5\omega\) up to completion losses.
When it is bounded, the saving is \(s-m>.075-o(1)\).

For \(\vartheta\le.0158\), reserve the last \(L\) pure \(2\)
positions in the two coarse norms \eqref{dt:gram-first},
where \(L=\lfloor (.05-\vartheta/\kappa_5)k\rfloor\).
The first decimal positions and the first \(i-j-L\) pure
\(2\) positions are summed at their usual shifted norms.
The remaining factor in either coarse first norm is
\[
 S(z)=\sum_{b\bmod 2^L}
       F_{10^L}^\circ\big((b+5^{i-j-L}z/Z)/2^L\big).
\]
The common high factor beyond \(10^i\) is unaffected.
On writing \(z=z_0+D_0 b_0\), the shift in this sum
changes by \(1/J_5\) per unit increase of \(b_0\), where
\(J_5=5^{L-k\omega}\). In fact
\(5^{i-j-L}/d_2=1/J_5\).
Here \(L/k>.027\), so \(L>k\omega\), and
\(J_5\mid d_2/e_0\), since \(i-j\ge.05k\).
In the congruence \eqref{dt:gram-fold}, with \(z_0,\ell\)
fixed, \(b_0\bmod d_2/e_0\) traverses a unit progression
as the integer quotient \((t+\ell z_0)/(e_0D_0)\) varies.
The \(e_0\) free lifts have the same shift.
The available quotient interval contains many periods
on the \(J_5\)-grid, because
\[
 \log_X(T/(e_0D_0))>.075-\vartheta
    >\kappa_5(.05-\omega)-\vartheta+o(1)=\log_X J_5.
\]
Completing the quotient interval is therefore valid with a
bounded factor. All weights at this point are nonnegative,
so any further restrictions on the quotient may be dropped.

Cauchy in the pure \(2\) grid and the shifted square norm
on the combined grid give
\[
 \frac1{J_5}\sum_{b_0\bmod J_5}S(z_0+D_0 b_0)^2
       \ll \frac{2^L}{J_5}(2^L J_5)^\delta .
\]
The combined grid has size at most \(10^L\).
Compared with \(2^{2a_2 L}\), the saving in this bound is
\(F L/k-(1-\delta)\kappa_5\omega\), where
\[
 F=(2a_2-1)\kappa_2+\kappa_5-\delta.
\]
Adding the twist saving \(\vartheta+\kappa_5\omega\),
the total is
\[
 .05 F+(1-F/\kappa_5)\vartheta+\delta\kappa_5\omega+o(1)
                         >.0158+o(1).
\]
The strict inequality uses \(.05F>.0158\) and \(F<\kappa_5\).
The choices of \(e_0\) cost only a logarithmic factor.
Dropping the bounded label multipliers after the signed lift
preserves this argument.
\end{proof}

\subsection{Rational frequency bands at small difference orders}

\begin{lemma}[Small-order rational bands]\label{dt:gram-mask}
Use the scales \eqref{dt:gram-scales}, with \(q\ge MX^{-.0001}\).
There is a union \(\mathfrak B_q\) of rational frequency
bands with the following properties.

\begin{enumerate}
\item Apart from an arbitrarily small fixed completion loss,
each family has parameters \(s,b_s,b_c,v_0\ge0\), with
\begin{equation}
 .249s+.806(b_s+b_c)+.903v_0\le .0477,\qquad
 E=\max(10^i,10^{\lfloor .117 k\rfloor}),\quad s=i/k.
 \label{dt:gram-mask-budget}
\end{equation}
Its frequencies lie on complete orbits under addition of
\(1/E\). A representative has the form
\(\alpha=u/(E d_s d_c)+\beta\), with \(d_s\) supported on
\(2,5\), \((d_c,10)=1\), \((u,d_s d_c)=1\),
\(d_s=X^{b_s}\), \(d_c\asymp X^{b_c}\) and
\(|\beta|\asymp V_0/q\), where \(v_0=(\log_X V_0)_+\).
Here \(d_s\mid E^2\) and \(E\mid q\). The offsets for a
fixed center are spaced by at least \(1/q\).
Writing \(\tau=(\log_Xq-m)_+\), the offsets also satisfy
\begin{equation}
 (h+\log_X V_0-\tau)_+
              \le .0477-.249s-.806(b_s+b_c)+.097v_0.
 \label{dt:gram-mask-offset}
\end{equation}

\item For small difference orders \(s\le.29\) whose lifts
divide \(10^{\lfloor.117k\rfloor}\), the off-mask signed
Gram saves a fixed power. For \(s\le.29\) and
\(\xi\ge.234-o(1)\), the small-order estimate saves without
any mask.

\item The normalized full first mass of \(\mathfrak B_q\)
is at most \(X^{a(.192)+o(1)}\). For any order \(g\) with
\(\xi\ge.234-o(1)\), after the unmasked estimate the absolute
correction for excluding one or both masked labels also
saves a fixed power.
\end{enumerate}
\end{lemma}

\begin{proof}
At \(s\le.29\) there is no incomplete row cost. By
\(\gamma=\log_Xq\ge m-.0001+o(1)\) and
\(v_2(q),v_5(q)\le\lambda k+O(1)\), the two valuations
of \(q\) exceed \(.193 k-o(k)\) and \(.370 k-o(k)\),
respectively. For example
\(v_2(q)/k\ge(m-.0001-\kappa_5\lambda)/\kappa_2+o(1)\),
with the analogous inequality for \(5\).
In particular the residual multiplier \(K\) equals one
for \(2\)-heavy orders, and otherwise is a power of \(2\)
with \(\log_XK\le.302s\).
Also
\(\log_XD_0-\log_X T\ge m-.0001-s-O(\epsilon)+o(1)\),
by \eqref{dt:gram-fold-identities}. Thus \(D_0\gg T\)
by a power.

Retain the folded incidence \eqref{dt:gram-fold} rather
than summing the high square on its full grid. For fixed
\(e_0=(\ell,d_2)\), write \(\ell=e_0\ell'\),
\(t=e_0t'\), \(H'=H/e_0,\ T'=T/e_0\).
The \(e_0\) lifts have canceled the normalization in the
gap count. For \(z_0\bmod D_0\), the reduced rational number
\[
             \frac{\ell' z_0+t'}{\ell'D_0}
\]
has denominator \(d_*\mid\ell'\) since
\(D_0\mid\ell'z_0+t'\). Let \(V_0\) be a dyadic size
for the absolute displacement \(|t'/\ell'|\), with
zero displacement treated separately. A nonzero shell
requires \(|\ell'|\ll T'/V_0\).
Project the centers by multiplication by \(K\).
For fixed \(C_*=(K,d_*)\) the projected denominator is
\(d'=d_*/C_*\), and each projected center has at most \(C_*\)
preimages at that fixed scale. Group \(d'\asymp B\).
For each original center and each \(z_0\), the twist is fixed
by the gap. The remaining gap multiplicity after division
by \(H'\) is
\[
     O\bigl((C_*B)^{-1}\min(1,T/(HV_0))\bigr).
\]
For zero displacement use \(O((C_*B)^{-1})\). The factors
\(C_*\) cancel on projection. The divisor groupings
over \(C_*,e_0\) have only logarithmically many choices,
and positive displacement shells also have this property,
since a nonzero displacement is at least \(1/|\ell'|\).

Suppose first that
\(B^2\max(1,KV_0)\) is smaller than a sufficiently small
constant times \(D_0\). In the high squared word keep
an initial block of length comparable to \(B^2\).
The projected centers are \( \gg B^{-2}\)-separated and
the variation of the projected angle is \(O(KV_0/D_0)\).
The sum of their local initial square suprema is
\(O(B^{2\delta})\) by separated sampling.
At each original center also keep a terminal block
beginning at a decimal position \(10^j\) with
\(D_0/(D_0,10^j)\asymp 1+V_0\). Such a position exists
within bounded factors because each multiplication by
ten reduces the grid size by a factor of at most ten.
It is at least a constant times \(D_0/(1+V_0)\),
and the remaining word has room for an initial block
comparable to the resulting arithmetic grid.
The two blocks are disjoint by the assumed separation.
At each center the integer \(z_0\)'s in the shell run over
an interval of \(O(1+V_0)\) points, so the terminal square
sum costs \(O((1+V_0)^\delta)\) by sampling on that grid.
The multiplier \(K\) permutes it because \((K,D_0)=1\).
The normalized folded incidence at this shell therefore
costs at most
\begin{equation}
        O\bigl(B^{2\delta-1}(1+V_0)^\delta
                              \min(1,T/(HV_0))\bigr).
 \label{dt:gram-separated}
\end{equation}
Use one for the final two factors at zero displacement.

In the other case, at fixed \(z_0\) the projected centers
are in an interval of width \(O(KV_0/D_0)\) about
\(Kz_0/D_0\), with \(O(1+B^2KV_0/D_0)\) possible
centers. Combine the preceding gap multiplicity, the
complete squared norm \(O(D_0^\delta)\), and \(BV_0\ll T\).
The resulting bound is at most
\begin{equation}
 X^{.00003+O(\epsilon)+o(1)}D_0^\delta
          \left(\frac{1}{B\max(1,V_0)}+\frac{KT}{D_0}\right).
 \label{dt:gram-overflow}
\end{equation}
The factor \(X^{.00003+O(\epsilon)+o(1)}\) allows
\(T/H\) in the first term and can be omitted in the second.
At zero displacement only the first term is needed.
Since separation fails, \(B\max(1,V_0)\gg(D_0/K)^{1/2}\).
Multiplying by the coarse saving and \(\rho_U^{-1}\),
the first exponent in \eqref{dt:gram-overflow} is at most
\(.0474-.403\gamma+.305s+.00003+O(\epsilon)+o(1)<0\).
For the second term use \(KT/D_0=GT/q\) and
\(A_U+\delta\log_XD_0\le\delta(\eta-s)\).
Its exponent is at most
\(.097+.655s-m+.0001+O(\epsilon)+o(1)<0\).
Both inequalities hold on \(.4195\le m\le.4801\)
and \(s\le.29\).

Return to \eqref{dt:gram-separated} and write
\(b=\log_X B\), \(v=\log_X V_0\).
We only need to delete the families satisfying
\begin{equation}
 .249s+(1-2\delta)b+(h+v-\tau)_+-\delta v_+
                                       \le A_U+.00015.
 \label{dt:gram-active-mask-budget}
\end{equation}
At zero displacement omit the \(v\)-terms. Each label
in such an incidence lies within \(O(V_0/q)\) of a
rational center with denominator dividing \(Gd'\),
since \(G\alpha=Kz_0/D_0\bmod1\) and the scaled
projected displacement gives \(KV_0/(GD_0)=V_0/q\).
Exact zero displacement cannot occur in a deleted
family, since \eqref{dt:gram-active-mask-budget} would
make its center denominator smaller than \(q\) by a
power, contrary to primitivity of \(\alpha\).
Conversely for all remaining families the inequality
opposite to \eqref{dt:gram-active-mask-budget} saves a
fixed power. When \(s\ge.234-o(1)\) none need be
deleted, since already \(.249s>A_U+.00015+.097(\tau-h)_+\).

Enlarge each family by all translates of \(1/E\) with
\(E=\max(G,10^{\lfloor.117k\rfloor})\).
The projected denominator \(d_s d_c\) is the reduced
denominator of \(E\) times the rational center, so it
can only decrease from \(d'\). Write
\(b_s=\log_X d_s\) and group \(b_c=\log_X D_c\) dyadically
with \(d_c\asymp D_c\).
Since \(A_U<.0474\) and \(\tau-h\le.00003+o(1)\),
the bounds \eqref{dt:gram-mask-budget} and
\eqref{dt:gram-mask-offset} follow from
\eqref{dt:gram-active-mask-budget}. For the first one
use \((h+v-\tau)_+\ge v_+-.00003-o(1)\).
The budget gives \(s<.192\) and \(b_s+b_c<.060\).
Thus \(10E\mid q\) for large \(X\), so translation by \(1/E\)
preserves primitive support. Also \(d_s\mid E^2\) since
\(.060<2\kappa_2(.117)\).
The frequencies can be partitioned into entire orbits
by taking the families in decreasing order of \(E\)
and rejecting the points already assigned. Such rejection
is compatible with \(1/E\)-translation, since the
previous orbit lengths are divisible by \(E\).
On each selected orbit choose a covering center
consistently under translation. This proves part 1.
The mask is invariant under the increments of any signed
lift of size at most \(10^{\lfloor.117k\rfloor}\).
This gives part 2.

For the first mass of an enlarged family put \(e=\log_XE\).
At fixed projected center and offset, the lowest \(E\)-block
has the shifted first cost \(E^a\) on the orbit. The higher
coefficient is the word above \(E\) at
\(u/(d_s d_c)+E\beta\). Keep an initial block of length
comparable to \(d_s\), paying \(d_s^a\) on the smooth
grid at each coprime center. Skip to a decimal position
divisible by \(d_s\) and use two successive blocks of length
comparable to \(D_c^2\) on the coprime rational centers.
Digit dilation permutes these centers at each coprime
denominator, preserving separation between distinct reduced
fractions on this scale. The two squared local-supremum sums
are \(O(D_c^{2\delta})\) individually. Taking square roots
and using Cauchy between the two blocks charges
\(D_c^{2\delta}\) to the first mass. At each original rational
center use a terminal block on the offset grid of size
comparable to \(1+V_0\), costing \((1+V_0)^a\).
For the fit of these blocks, the initial blocks end at a
position at most \(E d_s^{1/\kappa_2}D_c^4\ll
X^{e+4(b_s+b_c)+o(1)}\) in the original coordinates.
The terminal offset block begins at a position at least
a constant times \(q/(1+V_0)\), by the same grid argument
used for \eqref{dt:gram-separated}. The budget gives
\[
 e+4(b_s+b_c)+v_0
             \le .117+4(.0477)/.806+o(1)<.354<\gamma .
\]
This also guarantees that the earlier block arguments vary
only within their permitted local sampling neighborhoods
on the corresponding offset shells. The first mass is
therefore at most
\begin{equation}
  X^{a(e+b_s+v_0)+2\delta b_c+o(1)}
                                     \le X^{a(.192)+o(1)} .
 \label{dt:gram-mask-first}
\end{equation}
For the last bound use \(2\delta<a\). For \(s\le.117\)
one has \(e+b_s+b_c+v_0\le .117+.0477/.806+o(1)\).
If \(s\ge.117\), maximize the affine bound
\(s+(.0477-.249s)/.806\) over \(s\le.0477/.249\).

If \(V\) does not divide \(10^{\lfloor.117k\rfloor}\),
then \(\xi\ge.234-o(1)\). Estimate the unmasked Gram,
then subtract incidences with one or both labels in
\(\mathfrak B_q\) using absolute values. For a fixed
masked label the partner at order \(g\) lies on a
shifted \(g\)-grid with first norm \(O(g^a)\).
At fixed labels and gap there are \(O(1+T/Z)\) twists,
and each permits \(O(1+P/g)\) rows by the exact order.
After division by \(PH\rho_U\), using
\eqref{dt:gram-mask-first}, the exponent is at most
\[
 .0474+a(.192)-(1-a)\xi+(\xi-n)_++(\xi-m)_+
                                      +O(\epsilon)+o(1).
\]
By convexity in \(\xi\), test the endpoints \(.234\) and
\(\lambda\), and at the second endpoint the two values
of \(h\) and the two values of \(m\) in
\eqref{dt:gram-scales}. Both tests give strict negative
margins. This proves part 3.
\end{proof}

\subsection{High-period bounds at the two cuts}

\begin{proposition}[The complementary higher periods]\label{dt:gram-high}
At either cut \eqref{dt:cut-parameters}, with its associated
multiplier, the normalized row Gram saves a fixed power on
the complement of \(\mathfrak B_q\) for all
\(q\ge MX^{-.0001}\).
For \(X^{.192}\le q<MX^{-.0001}\) it saves without any mask.
\end{proposition}

\begin{proof}
The diagonal is covered at both cuts by
\eqref{dt:diagonal-main}--\eqref{dt:diagonal-saving}.
The medium range follows from Lemma~\ref{dt:gram-middle}
and \(A_U<.0474<.249(.192)\).
For the high range below the difference-order cutoff
\(s=.29\), apply Lemma~\ref{dt:gram-mask}.
For all larger \(s\) first bound the unmasked Gram
and then make the absolute mask correction from that
lemma when necessary. The taper itself is constant on
every signed lift, since \(V\mid V_*\) by
\eqref{dt:cut-parameters} and the definition of \(V_*\).

In the higher-amplitude branch the bound
\eqref{dt:high-filter} gives the majorant
\[
         \omega_H(\theta)
     \ll B_*^{.4(.306)}\widetilde F_{B_*}(V_*\theta)^{.4}
                      +O_C(X^{-C}).
\]
Retain its digit factors where they overlap each of the
coarse first sums in \eqref{dt:gram-first}. The decimal part
of the coarse grid ends at position \(10^j\) and is followed
by its pure \(p_0\)-grid. Each splits at the taper endpoints
as in the proof of Lemma~\ref{dt:gram-base}, with powers
of the complementary radix permuting the pure grid.
At a retained
position the actual digit factor times the floored taper
factor is at most its floored \(7/5\)-power. Outside these
positions discard the taper factors and retain the ordinary
first norm. The top square starting at \(10^i\) is unchanged.
Write \(f_*=.258,\ f_2=.019,\ f_5=.137\) for the
\(7/5\) exponents and use the taper position and length exponents \(v,w_0\)
from Section~\ref{dt:filters}.
The saving in each coarse first norm is at least
\begin{equation}
 \begin{aligned}
 \Delta(s,d)=\Big[&
 (a-f_*)\bigl(\min(v+w_0,d)-v\bigr)_+\\
 &+(a_{p_0}-f_{p_0})\kappa_{p_0}
      \bigl(\min(v+w_0,s)-\max(v,d)\bigr)_+
                  -.4(.306)w_0\Big]_+ .
 \end{aligned}
 \label{dt:gram-filter-gain}
\end{equation}
The outer positive part uses the alternative
\(|\omega_H|\le1\). The error term in the filter bound is
negligible after all finite sums. Since both coarse sums
are taken after the signed lift, the total gain is
\(2\Delta(s,d)\), compatible with
Lemma~\ref{dt:gram-joint}.

Use \(\eta\le1\) in \eqref{dt:gram-base-exponent}
and denote the resulting exponent by \(E_0\).
For \(p_0=5\) and for \(p_0=2,\ s-d\le.05\), use the
joint affine gain on \(\xi\ge m+.020\), and use the
ordinary bound on \(\xi\le m+.020\).
On the active sector the hypothesis
\(h+\xi-m<d/2-.010\) holds throughout.
On the \(H\) cut subtract \(2\Delta(s,d)\) in both
sectors. For \(p_0=2,\ s-d\ge.05\), use respectively
the two twist substitutions and the terminal gain
of Lemma~\ref{dt:gram-tail}. The upper bounds, common
to the two cuts, are
\[
\begin{array}{l|r}
 \text{orders and estimate} & \text{exponent upper bound}\\ \hline
 p_0=5\text{ or }s-d\le .05,\ \xi\le m+.020 & -.00349\\
 p_0=5\text{ or }s-d\le .05,\ \xi\ge m+.020 & -.00068\\
 p_0=2,\ s-d\ge .05,\ \omega\ge.024 & -.00345\\
 p_0=2,\ s-d\ge .05,\ u\le.075 & -.00104\\
 p_0=2,\ s-d\ge .05,\text{ before the terminal gain} & .0134
\end{array}
\]
Here the last bound is reduced by more than \(.0158\).
The finitely many tests use exact rational arithmetic
as detailed below. Each negative entry permits a
uniform choice of the small completion loss.
This gives the desired bounds at all orders.
There are only logarithmically many scales and at most
\(O(k^2)\) possible difference orders for each \(q\).
\end{proof}

\paragraph{Verification of the large-order inequalities.}
The entries in the preceding table can be verified
without numerical approximation of the maxima.
Use \(.096910<\delta<.096911\) and
\(.301029<\kappa_2<.301030\), with the complementary
interval for \(\kappa_5\).
These enclosures follow from the convergent identity
\[
 \log b=2\sum_{\ell\ge0}
                 \frac{((b-1)/(b+1))^{2\ell+1}}{2\ell+1}
\]
and its geometric remainder bound, applied to \(b=2,5/4\).
Thus \(\log10=3\log2+\log(5/4)\) and \(\delta=1-3\kappa_2\).
At each cut and choice of \(p_0\) the variables are
\((m,s,d)\) in the rational polytope
\[
 .4195\le m\le.4801,\qquad .29\le s\le m+h+.00003,
                                              \quad 0\le d\le s.
\]
When cutting by \(\xi\ge m+.020\) use the rational upper
endpoint \(\xi_+\) for \(\xi\), and for the inactive
constraint \(\xi\le m+.020\) use the lower endpoint
\(\xi_-\). In \(E_0\) replace \(\delta\) by its upper
endpoint, each directional \(\kappa_p\) by its lower
endpoint, and \(\xi\) by \(\xi_+\) in the incomplete
progression term. For the active correction subtract
\[
 .096910(\xi_+-n)_+
 +\frac{.096910}{2}\min\{d/2,\,2(h+\xi_--m)\}.
\]
This is a valid upper bound after combining the two
incompletion terms, since \(1/2-\delta>0\).
Use \(\xi_+\) to check \(h+\xi-m<d/2-.010\),
and use the lower directional endpoint in the
nonnegative subtracted taper gain. For a wide \(2\)-heavy
order the first modified twist uses the lower endpoint
for \(\kappa_5\), and the second uses
\(\min((s-m)_+,.075)\).

The upper bounds are affine on cells separated by the
zero planes of the displayed positive parts, by
the equality of the two arguments of the minimum in
\eqref{dt:gram-joint-gain}, and, if the taper is in use,
by the clipping points \(d,s=v,v+w_0\) and the zero
planes of the possible affine arguments of
\eqref{dt:gram-filter-gain}. In the two wide-order tests
also insert the planes for the modified twist.
On each cell it is enough to test its vertices.
Equivalently, intersect every triple of independent
planes with the polytope, retaining only points satisfying
its defining inequalities. Evaluating the actual
piecewise affine majorant at all such points entails no
choice of branch in the enumeration and proves the same
upper bound. The accompanying rational verifier
\path{supplement/decimal_two/exponents.py} performs
these tests and the elementary endpoint checks
in Lemma~\ref{dt:gram-mask}.

\section{Decimal projections and Mellin comparison}
\label{dt:analytic-section}

Put \(L=\log X\), \(c=\log_{10}8\), and \(\delta=1-c\).
All prime boxes below have the smooth weights and the fixed positive
constituent scale bound in \eqref{dt:box-equation}. We use \(e(u)=\exp(2\pi i u)\)
and the Mellin transform in the coordinate \(\log(x/X)\), with phase
\((x/X)^{-it}\). Factors that are powers of \(L\) are allowed in the box
scales and in the derivative estimates. The digit density of a word of
length \(B\) is \(\asymp B^{-\delta}\); a terminal unit digit or a prescribed
admitted leading digit changes this by a bounded factor.

\subsection{Zero counts and prime polynomials}

For \(\sigma\geq 1/2\), let \(N_q(\sigma,T)\) count the zeros with
\(\Re\rho\geq\sigma\) and \(|\Im\rho|\leq T\), over the primitive characters
of conductor \(q\), with multiplicity. For \(0\leq y\leq\tfrac12\)
and \(T\geq2\), we use the fixed-modulus bounds
\begin{align}\label{dt:analytic-density}
 N_q(1-y,T)&\ll (qT)^{(12/5)y}(\log(2qT))^C,\notag\\
 N_q(1-y,T)&\ll_\epsilon(qT)^{3y/(1+y)+\epsilon}.
\end{align}
The first is Huxley's estimate \cite{Huxley} in the logarithmic-loss
fixed-modulus form recorded in \cite[(3.3)]{BanksShparlinski2019}. For the second
use the Ingham exponent of Forti and Viola \cite[(2.5), p.~386]{FortiViola1973}:
for characters of one modulus their difference conductor divides $q$, so
their parameter $D$ is at most $q(2T+1)$. The resulting bound is also
recorded in \cite[(1.4)]{ChenGuptaLi2026}. At \(y=1/2\) the bound
\(O(qT\log(2qT))\) suffices. Below \(y=1/4\) we always use the first
bound, including in partial summation arbitrarily close to one.
Iwaniec's zero-free region \cite[Theorem~2]{Iwaniec1974}, in the
form \cite[(3.1)]{ThornerZaman}, bounds a nonexceptional zero by
\begin{equation}\label{dt:analytic-zero-free}
 1-\Re\rho \gg
 \big(\log \operatorname{rad}(q)
       +(\log(q(|\Im\rho|+3)))^{3/4}
        (1+\log\log(q(|\Im\rho|+3)))^{3/4}\big)^{-1}.
\end{equation}
For \(q=1\) use the Vinogradov--Korobov region for the zeta function
\cite[Theorem~1]{Ford2002ZeroFree}.
For \(q\mid 10^k d\) with \(d\le L^C\), the conductor of an exceptional
real primitive character is \(\ll d\). This follows by factoring the
unit group at each prime: a primitive character of order two has odd
prime exponents at most one and 2-exponent at most three. Siegel's
bound \cite[Theorem~5.28]{IwaniecKowalski2004}, with its exponent
chosen in terms of \(C\), gives \(1-\Re\rho\gg_C L^{-1/4}\) for such
an exceptional zero. In particular \(X^{-b(1-\Re\rho)}\ll_{b,J}L^{-J}\)
for every fixed \(b>0,J>0\) in these families and at polynomial heights.

We record the explicit-formula reduction for any fixed number of
prime variables. Let \(P_i\) be their scales and \(X'=\prod_i P_i\).
For a weight \(w_i\), the polynomial
\(S_{i,\chi}(t)=\sum_p w_i(p/P_i)\chi(p)p^{it}\) has a pole term only
for the principal character, and that term is
\(\int w_i(x/P_i)x^{it}\,dx/\log x\). After removal of that term,
the nonprincipal or principal remainder satisfies, on \(|t|\leq 2T\)
with \(T\ge X^\epsilon\),
\[
 |S_{i,\chi}^{\,\mathrm{rem}}(t)|\ll L^{C}
 \left(P_i^{1/2}+
   \sum_{\substack{\rho=\beta+i\gamma\text{ of }L(s,\chi^*)\\
                    \beta\ge1/2,\ |\gamma|\le4T}}
           P_i^\beta J(t+\gamma)\right),\qquad J(u)=(1+|u|)^{-3}.
\]
Here \(\chi^*\) is the primitive inducing character. Use the
von Mangoldt explicit formula with \(w_i(x/P_i)x^{it}/\log x\)
and integrate by parts in \(\log x\); derivatives give arbitrary
inverse powers of \(1+|t+\gamma|\). The prime powers cost
\(P_i^{1/2}L^C\), trivial zeros cost \(L^C\), and the number of
zeros of one character in a unit ordinate interval is \(O(L)\)
at polynomial heights. These observations also bound omitted
zeros with \(\beta<1/2\) by \(P_i^{1/2}L^C\) and permit the
truncation. The same bound holds if finitely many prime divisors
of any polynomial-size modulus are excluded from any slot:
their direct contribution is \(O(L^C)\). On the central interval
\(|t|\le X^\epsilon\), one can instead truncate at \(X^{2\epsilon}\).

In the product of the remainders, apply the weighted arithmetic--geometric
mean inequality to the factors \(P_i^{\beta_i}\), with weights
\(\log P_i/\log X'\). Terms with at least one zero are bounded by a fixed
sum of \(X'^\beta\) times one retained zero kernel; summing the other
kernels at fixed \(t\) costs \(L^C\). Terms with no zeros cost \(X'^{1/2}L^C\).
Thus a zero layer of deficit \(y=1-\beta\) contributes \(X'^{1-y}\)
times the remaining kernel, whose supremum is \(O(L)\), whose integral
is \(O(N_\chi)\), and whose squared integral is \(O(LN_\chi)\).
Here \(N_\chi\) is the number of zeros of \(\chi^*\) on the layer.
This proves the fixed-rank reduction used below. Mixed pole and remainder
terms on noncentral shells have arbitrary power decay from the smooth
pole factor. On the central interval they may be estimated in the same
way, with at least \(X^{-\kappa y}\) saving from the slots with remainders,
where each \(P_i\ge X^{\kappa+o(1)}\). Choose \(2\epsilon\) so that
\((12/5)2\epsilon<\kappa/2\). The central mixed terms use only zeta
zeros, and \eqref{dt:analytic-zero-free} and partial summation give
arbitrary logarithmic gains. A term with only poles is the continuous
product integral by Mellin inversion. The transform of a smooth
product cutoff alone decays to all fixed orders and may instead be
truncated at \(X^\epsilon\), with power-small error.

\subsection{Digital character coefficients and the initial projection}

For a decimal word of length \(B\) and mass \(A_B\), let \(b_\chi\)
be its sum against \(\overline\chi\), divided by \(A_B\). The word
may be twisted by a fixed additive phase. For a primitive character of
conductor \(q=2^i5^j\mid B\), put
\(R=10^{\lceil\max(i,j)/2\rceil}\), \(D=(q,R)\), and \(Q=q/D\).
Use \(m=w+Ru\), \(0\le w<R\), on the units. The kernel of reduction
from units modulo \(q\) to units modulo \(D\) has the additive
parameter \(1+Rz\), \(z\bmod Q\), since \(q\mid R^2\). Consequently
\(\chi(w+Ru)=\chi(w)e(\ell \overline w u/Q)\), where \(\ell\) is a
unit modulo \(Q\) by primitivity. This is the linear character identity on the kernel; on each prime-power component it is the linear case of Postnikov's formula \cite{Postnikov}. In any residue class modulo
\(d\le R\) there are \(O((R/d)^c)\) permitted lower words, by
splitting off the largest decimal power at most \(d\). The same
bound holds after a common translation. Inversion permutes the
unit classes modulo \(Q\). The shifted first moment for the
upper word on that grid therefore bounds \(|b_\chi|\) by \(Q^{-c}\)
times the first moment on the grid. In the block above the cut the
least decimal multiple of \(Q\) divides \(B/R\). If its leading digit
is fixed, discard that digit's phase; the possible loss in the number
of grids is at most ten. Multiplicative Parseval and the same residue
bound give \(\sum_\chi |b_\chi|^2\ll q^\delta\).

The mixed grid exponents of Proposition~\ref{dt:digit-pure} and
Lemma~\ref{dt:digit-mixed} give the useful stronger pointwise estimate
\begin{equation}\label{dt:analytic-linear}
 \sup_{\chi\ {\rm primitive}\bmod q}|b_\chi|\ll q^{-.201}.
\end{equation}
Indeed, if the larger valuation is \(i\) at \(p\), with
\(\kappa_p=\log_{10}p\), the saving in base-ten logarithms is
\((c-a_p)\kappa_p i/2+O(1)\) for \(j\le i/2\), and
\((c-a)(j-i/2)+(c-a_p)\kappa_p(i-j)+O(1)\) for \(j\ge i/2\).
Compare to \(.201(\kappa_p i+\kappa_{p'}j)\) at \(j=i/2,i\), and
use linearity on the two intervals. The endpoint inequalities are
strict for \(a=.428,\ a_2=.028,\ a_5=.2241\).

For the cubic norm group extensions with a common slope \(\ell\),
which are indexed by the characters of \((\mathbb Z/D\mathbb Z)^*\).
If \(M_\ell\) is the sum of the \(3/2\)-powers of the upper Fourier
amplitudes over allowed lower words \(w\), counting Hausdorff--Young
at exponent \(3/2\), followed by the lower lift occupancy, gives for
that slope a cubed norm \(\ll D^\delta R^{-2c}M_\ell^2\).
In detail, the counting Fourier bound is
\(\sum_\psi|\widehat f(\psi)|^3\le \varphi(D)(\sum_r |f(r)|^{3/2})^2\)
when the transform is a sum; Hölder in each residue class of \(w\)
costs one power of the maximum lift multiplicity after squaring.
The upper \(3/2\)-moment with exponent \(.225\) gives
\(\max_\ell M_\ell\ll (R/Q)^c Q^{.225}\) and
\(\sum_\ell M_\ell\ll R^c Q^{.225}\), by interchanging \(w,\ell\).
It follows that
\begin{equation}\label{dt:analytic-cubic}
 \sum_{\chi\ {\rm primitive}\bmod q}|b_\chi|^3
   \ll q^\delta Q^{-.55}.
\end{equation}

We include the quartic bound since its applicability to unequal
valuations is needed. Now take
\(R=10^{\max(2,\lceil\max(i,j)/3\rceil)}\),
\(D=(q,R)\), \(Q=q/D\), \(K=(R,Q)\), \(W=Q/K\).
Then \(W\mid K\mid R\), \(Q\mid R^2\) and \(R^2W\mid B\) for large \(B\)
in the application below. The linearizing coordinate on \(1+Rz\)
is \(z-(R/2)z^2\bmod Q\). It is additive for the multiplication
law \(z+z'+Rzz'\) and is invertible since the quotient of a
difference by \(z-z'\) is a unit. Thus write a word as \(z+Rm+R^2h\).
Its upper factor in a slope fibre has Fourier amplitude \(U(a)\)
on a shifted \(W\)-grid, with \(a=\ell\bar z\bmod W\); multiplication
by \(R/K\) only permutes this grid. The middle normalized polynomial is
\[
 B_\ell(z)=R^{-c}\sum_{m\ {\rm admitted}}e\!\left(
     \frac{\ell\bar z m}{Q}-\frac{R\ell\bar z^2m^2}{2Q}+\alpha Rm\right)
\]
for any fixed additive twist \(\alpha\).
For fixed unit \(a\bmod W\), put \(\ell=z(a+Wb)\) and sum
\(|B_\ell(z)|^4\) over \(b\bmod K\) and all units \(z\bmod R\).
On expansion, with \(S=m_1+m_2-m_3-m_4\) and
\(V=m_1^2+m_2^2-m_3^2-m_4^2\), \(b\)-orthogonality imposes \(K\mid S\)
since \(V\equiv S\bmod2\). The sum over inverse units modulo \(W\)
is a Ramanujan sum, forcing \(W/\operatorname{rad}(W)\mid V\)
up to a bounded factor at 2. As \(W\mid K\),
\(V\equiv-2(m_3-m_1)(m_3-m_2)\bmod W\).
Split a divisor \(W'\mid W\) of bounded index between the two
differences; residue occupancy gives \(O(R^{3c}(W')^{-c}L^C)\) choices
of the first three variables. The fourth has \(O(R^c K^{-c})\) choices.
After normalization the resulting fourth sum is
\(\ll R K^\delta W^{-c}L^C\). Fixed phase factors in the expansion
do not change this bound. For \(W=1\) the same proof uses only
orthogonality in \(b\).

The fibre Hausdorff--Young estimate at exponent \(4/3\) is
\[
 \sum_\psi|b_{\chi_\ell\psi}|^4\ll
 D^\delta R^{-3c}
 \left(\sum_{z\ {\rm admitted}}|U(\ell\bar z)B_\ell(z)|^{4/3}\right)^3 .
\]
By Hölder the cube is at most
\((\sum_z |U|^{4/3})^2 \sum_z |U|^{4/3}|B_\ell(z)|^4\).
Use occupancy and the shifted \(4/3\)-moment to bound the first sum
by \((R/W)^c W^{.285}\). To sum the other term over slopes, enlarge
the \(z\)-set to all units, group \(\ell\bar z\) by \(a\bmod W\), and
use the preceding fourth sum. The remaining sum over \(a\) again
costs \(W^{.285}\). Since \(DK=q/W\), the result is
\begin{equation}\label{dt:analytic-quartic}
 \sum_{\chi\ {\rm primitive}\bmod q}|b_\chi|^4
   \ll (qR)^\delta W^{3(.285-c)-\delta}L^C.
\end{equation}
For the moments on comparable grids one can also retain the first
decimal block of length comparable to the grid size and apply
separated sampling with an \(L^1\)-bounded reproducing kernel.
This costs only a constant because the discarded digit factors
have modulus at most one.

\begin{proposition}[The initial decimal projection]\label{dt:analytic-projection}
The comparison in \eqref{dt:box-equation} holds with the digit source
projected onto its lowest \(\lfloor .770 k\rfloor\) digits.
\end{proposition}

\begin{proof}
Expand the projected source on the units modulo \(10^{\lfloor .770k\rfloor}\).
For each primitive \(q\), a smooth prime polynomial with only low Mellin
twists is bounded on the zero layer \(y\) by \(P_i^{1-y}L^C\).
Group characters by their rightmost zero at height \(X^{2\epsilon}\);
if there is none above \(1/2\), use \(y=1/2\) and count all characters.
The total product has relative size \(X^{-y}\), and the cardinality
of a zero layer is bounded by \eqref{dt:analytic-density}.
All fixed powers in the height and the real-part mesh in the strict
power-saving ranges below can be taken arbitrarily small.

Write \(\gamma=\log_X q\). For \(.01\le\gamma\le.565\), use
\eqref{dt:analytic-linear} up to \(y=.300\), with relative exponent
\(-y+\gamma(2.4y-.201)\). It is negative at both endpoints. Thereafter
use the square norm with exponent \(-y+\gamma(1.2y+.0485)\),
which decreases and is negative at \(.300\).
For \(0\le\gamma\le.01\) the linear density exponent even with the
height included is \(<.03 y\) for a sufficiently small height parameter.
Sum in the real part by partial summation, keeping the factor
\(X^{-y}\) and using \eqref{dt:analytic-zero-free}, with no fixed
mesh loss near one. This costs only logarithms and gives arbitrary
logarithmic savings. The principal character is treated separately.

For \(.565\le\gamma\le.665\), use \eqref{dt:analytic-cubic} and
\(Q\gg qX^{-.770/2}\). Hölder and the Ingham branch of
\eqref{dt:analytic-density} give relative exponent at most
\[
 -y+\frac{2\gamma y}{1+y}-.151\gamma+\frac{.55(.770)}6
 \le 1+1.849\gamma-2\sqrt{2\gamma}+\frac{.55(.770)}6.
\]
The last bound is convex in \(\gamma\), and is below \(-.006\)
at each endpoint. For \(.665\le\gamma\le.770\), use
\eqref{dt:analytic-quartic}, \(R\ll X^{.770/3}\), and
\(W\ge q/R^2\). Its cubically split radix satisfies \(R^2W\mid X\).
The fourth norm with Hölder gives relative exponent at most
\[
 1+1.7865\gamma-3\sqrt\gamma+
 \left(\frac{.097}{12}+\frac{1.951}{6}\right).770 .
\]
This convex bound is below \(-.00025\) at both endpoints. The rational
comparisons are in the \texttt{projection} certificate. Sum the
error over the \(O(L^2)\) conductors. The principal coefficient is the
unit mean, and the all-pole continuous integral is the desired
integer-tested continuous model up to a power-small relative error:
average over complete periods \(10^{\lfloor .770k\rfloor}\) and
use partial summation against the smooth product density. The normalized
absolute mass of this period is bounded.
\end{proof}

\subsection{Mellin estimates for the upper step}

Let \(R\) be either cut in \eqref{dt:cut-parameters},
\(Z=X/R\), \(B=Y/R\), \(u=1-\eta\), \(z=1-\lambda\), and
\(\rho_U=B^{-\delta}\). Either shift of the upper step is permitted,
and the signed upper increment can be treated as a difference of
two words of lengths \(B\) and \(B/10\). Let \(\Psi(x/X)\) be a
smooth product test. In logarithmic coordinates the upper test is
\(\Psi(x/X)U_R(x)e(-\beta x)\), with a possible smoothing of \(U_R\).
Write \(\tau=\log_X T\) on a shell \(|t|\asymp T\), and
\(r_T=(\tau-u)_+\). The elementary upper-word estimates are
\begin{align}
 \|\widehat V_0\|_{L^2(|t|\le 2T)}
    &\ll \rho_U X^{\delta r_T/2}L^C, \label{dt:analytic-upper-square}\\
 \sup_{|t|\asymp T}|\widehat V_\beta(t)|
    &\ll \rho_U L^C X^{-B_T},
 \qquad B_T=
 \begin{cases}
  .572r_T,&\tau\le 2u,\\
  .2375\tau+.0969u,&\tau>2u .
 \end{cases} \label{dt:analytic-upper-point}
\end{align}
The second bound is uniform in the modulation and holds also when
\(T>Z\) in the retained shells of the mask. In addition, if
\(T\le c_0 X|\beta| L^{-C}\) with \(C\) large enough on the
product support, the normalized bound is
\(O((1+R|\beta|)^{-1}L^C)\).

For \eqref{dt:analytic-upper-square}, a logarithmic interval of
length \(1/T\) meets physical length \(O(X/T)\) up to powers of \(L\).
The absolute step mass on such an interval, in logarithmic coordinates,
is \(O(\rho_U T^{-1}(1+TY/X)^\delta L^C)\). When the physical interval
exceeds \(Y\) this follows by counting periods; between \(R\) and \(Y\)
it follows from the bound \(O(H^c)\) for admitted words in an upper
interval of length \(H\). Below \(R\) the bound follows from bounded
height. The total log mass is \(O(\rho_U L^C)\). A smooth frequency
majorant of \([-2T,2T]\) with kernel \(O(T(1+T|v|)^{-2})\) and
summation over its translates of length \(1/T\) give
\eqref{dt:analytic-upper-square}.

For the pointwise estimate, write a physical step as \(x=R(j+\theta)\)
with \(0\le\theta\le1\), absorbing its fixed offset into \(\theta\).
After partitioning the compact product range if needed, use aligned
decimal blocks of the upper index \(j\) of length
\(B_0\asymp\min(B,Z/\sqrt T)\), shortening by at most a power of \(L\)
to absorb box distortions. At the left endpoint \(j_0\) remove
the linear logarithmic and modulating phase. The residual derivatives
of fixed order in \((j-j_0)/B_0\) and in \(\theta\) are bounded by
powers of \(L\), since \(TB_0^2/Z^2\ll L^C\).
Integrate in \(\theta\) first. For a large step frequency
\(\vartheta_j=t/j_0-2\pi\beta R\), integration by parts gives
an additional \(O((1+|\vartheta_j|)^{-1}L^C)\) in each fixed-order
derivative of the residual block weight. Indeed the change in this
phase derivative within the block is \(O(TB_0/Z^2)\); mixed derivatives
after scaling are bounded. Fourier expansion of the smooth block
weight, extended by a cutoff to a fixed larger period, then gives
summable coefficients with the same bound. Fourier indices multiply
the final estimate by only a logarithmic power.

The range of the linear frequencies \(B_0\vartheta_j\) has length
\(W\asymp TB_0/Z\), up to powers of \(L\), and a unit frequency cell
has preimage of length \(O(Z/W)\). The sum of normalized upper amplitude suprema over these cells,
weighted by the step factor \((1+|\vartheta_j|)^{-1}\) (taking its
supremum on each cell), is \(O((1+W)^a L^C)\), \(a=.428\).
If \(W\le B_0\), retain the terminal block of length comparable to
\(1+W\) and use separated shifted first-moment sampling. Otherwise
each cycle of \(B_0\) cells has first norm \(O(B_0^aL^C)\).
The reciprocal step factor \((1+|\vartheta_j|)^{-1}\) sums over
the cycles at a logarithmic cost, uniformly in the translation.
This also proves the cell bound when \(T>Z\).

If \(T\le (X/Y)^2\), then \(B_0\asymp B\) up to logarithms and
there are \(O(Z/B)\) blocks. Each unit cell contains at most
\(O(1+Z/(BW))\) starts, with \(W\ll Z/B\) up to logarithms, giving
relative decay \((1+TY/X)^{a-1}\).
In the other range \(B_0\asymp Z/\sqrt T\) and each unit cell
contains \(O(L^C)\) starts. Normalize by the number of occupied
blocks, of density scale \((Z/B)(B/B_0)^c\).
This gives \((Y/X)^\delta T^{-(c-a)/2}\), as in
\eqref{dt:analytic-upper-point} since \((c-a)/2>.2375\) and
\(\delta>.0969\). Absolute constants in the transition change
only the logarithmic factor. The separate nonstationary bound
comes from integration inside each step when the additive phase
dominates the logarithmic phase. Smoothing by averaging common
logarithmic translations preserves all these estimates: after
changing variables the modulation is dilated by the translation,
and the product test has the same derivative bounds.

\subsection{Small exact lower periods}

\begin{proposition}[Small lower periods]\label{dt:analytic-small-periods}
In the high cut, the unfiltered exact additive lower components
of periods \(q<X^{.192}\), multiplied by their shifted upper steps,
satisfy the comparison in \eqref{dt:box-equation}. This includes the
components of each lower word with a forced leading digit.
\end{proposition}

\begin{proof}
Use the expansion on units modulo \(R\) of
\(G_q=\sum_{d\mid q}\mu(q/d)E_d L_R\), where \(E_d\) averages on
the full integer circle. A primitive coefficient \(b_\chi(q_0)\)
of the original word gets multiplied, in units of \(\rho_L\), by
\[
             \sum_{\substack{d\mid q\\q_0\mid d}}\mu(q/d)d/\varphi(d).
\]
Indeed each unit residue modulo \(d\) has
\(\varphi(R)/\varphi(d)\) unit lifts to \(R\). Averaging
\(\overline\chi\) over those lifts gives zero unless its conductor
divides \(d\). The multiplier is bounded (only \(2,5\) divide \(q\)),
and it vanishes unless \(q_0\mid q\). This handles also additive
periods lacking one of the two prime divisors of the base. The
forced leading digit retains the bound \eqref{dt:analytic-linear}
and the square bound \(q_0^\delta\).

Smooth the step in logarithmic coordinates at width
\(Z^{-1}X^{-\epsilon}\), with \(\epsilon>0\) arbitrarily small and fixed.
On a physical jump interval \(I\) the absolute mass of \(G_q\)
is \(O(\rho_L(|I|+q))\), since the conditional projections in its
definition each have this bound on their own periods.
Each interval has length \(O(RX^{-\epsilon}L^C)\). Sum these
bounds over \(O((X/R)\rho_U L^C)\) jumps and use the bounded
multiplicity of a product of a fixed number of primes. As
\(q/R\) is power-small, this proves that smoothing costs a fixed
power, also in the continuous reference. The resulting Fourier
cutoff is \(T\ll ZX^{2\epsilon}\).

Put \(\gamma=\log_X q_0\) and \(D(y)=\min(2.4y,3y/(1+y))\).
The two relative zero bounds after the fixed-rank explicit formula
are
\[
 P_0(y)=-y+(\gamma+\tau)D(y)-.201\gamma-B_T,\qquad
 Q_0(y)=-y+\tfrac12(\gamma+\tau)D(y)+.0485(\gamma+r_T).
\]
The first uses the pointwise coefficient; the second uses the
upper square norm and the squared integral of the remaining zero
kernel, followed by character Cauchy--Schwarz.
With \(\gamma\le.192,\ z\le .4836\), at \(y=.234\) the first bound
is less than \(-.009\) in the first curvature range and less than
\(-.008\) in the second. In the first range maximize
\(2.4y\tau-.572(\tau-u)_+\) at \(\tau=u\) and use \(u\le.230\).
In the second use \(\tau\le.4836+2\epsilon+o(1)\) and discard the
negative term in \(u\). The square bound at this split is less
than \(-.011\) and decreases thereafter on both density regimes.
The no-zero term is bounded by
\(-1/2+.5485\gamma+\tau/2+.0485r_T<-.10\).
These numerical bounds allow the fixed positive height and mesh
reserves used away from \(y=0\).

For \(0\le y\le .234\), the pointwise bound with linear density is
affine in \(y\), nonpositive at zero, and uniformly negative at the
upper endpoint. Hence it is at most \(-c_1y\) for a fixed \(c_1>0\).
Use real-part partial summation and \eqref{dt:analytic-zero-free},
with the logarithmic-loss Huxley estimate throughout. In particular
the smoothing reserve is included in the height multiplying \(y\),
without an additive fixed loss in the exponent near zero.
On the central interval, count zeros up to \(X^{2\epsilon}\)
and use the trivial upper bound; the density slope is
\(2.4(.192+2\epsilon)<1\). The same argument bounds these zeros.
For principal mixed terms use the central and noncentral arguments
following \eqref{dt:analytic-zero-free}. The all-pole term tests the
unit mean. The shifted steps still have length \(R\), divisible
by \({\rm lcm}(q,10)\); average on these periods and use partial
summation against the smooth product density. This recovers the
integer-tested continuous model, with relative error
\(O((q/R+q/X)L^C)\). Finally restore the unsmoothed step using its
already bounded jump error.
\end{proof}

\subsection{Comparison on the rational masks}

\begin{proposition}[Arithmetic of the mask orbits]\label{dt:analytic-mask}
At either cut in \eqref{dt:cut-parameters}, consider the assigned
rational-mask orbits at exact lower periods \(q\ge MX^{-.0001}\).
These orbits satisfy the comparison in \eqref{dt:box-equation},
including for a forced leading digit and either taper.
In the normalization used below their error is
\(O_J(X\rho_U L^{-J})\), before restoring the lower density and
the scalar \((X/Y)^{-\delta}\) of \eqref{dt:outer-factor}.
\end{proposition}

\begin{proof}
We give the arithmetic argument with its precise orbit hypotheses.
A family has a decimal orbit modulus \(E\mid q\), with
\(e_*=\log_X E=\max(s,.117)+o(1)\). The projected denominator of
a covering rational center is \(d=d_s d_c\), where \(d_s\) is
supported on \(2,5\) and \((d_c,10)=1\). Fix \(d_s\), a dyadic
size \(D_c\asymp d_c\), and a dyadic offset size \(V_0\) on the
\(q\)-grid. Write \(b_s=\log_Xd_s,\ b_c=\log_XD_c\), and
\(v_0=(\log_X V_0)_+\). The rational masks supply
\begin{equation}\label{dt:analytic-mask-budget}
 .249s+.806(b_s+b_c)+.903v_0\le .0477,\qquad
 \ell_*=.0477-.249s-.806(b_s+b_c)+.097v_0 .
\end{equation}
In particular \(e_*+b_s+b_c\le .192+o(1)\),
\(d_s\mid E^2\), and \(b_c<.060\).
The displacement \(\beta\) from a covering rational center has
\(|\beta|\asymp V_0/q\); zero displacement is absent by primitivity
at the high additive period. For a fixed numerator the offsets
are spaced by \(1/q\). The assigned families retain full \(E\)-orbits.
The two tapers are constant on these orbits, since \(E\mid V_*\).
The forced leading lower digit has constant Fourier phase on an
orbit since \(E\mid R/10\); it changes only the common scalar.

Write a representative as \(\alpha_0=u/(Ed)+\beta\), with
\((u,d)=1\). Inversion over the orbit \(\alpha_0+j/E\) gives its
normalized physical source exactly as
\[
 A_{u,\beta}\frac{E}{N_E}\ind_{\mathcal D_E}(w)
       e(-u\lfloor x/E\rfloor/d)e(-\beta x)e(\beta w),
 \qquad w=x\bmod E .
\]
Here \(\mathcal D_E\) is the lower \(E\)-word with unit last digit,
\(N_E=|\mathcal D_E|\), and, before including either bounded taper
and the forced digit phase, \(A_{u,\beta}\) is the upper positional
factor at \(E\alpha_0\). Finite additive orthogonality forces the
bottom \(E\)-word to be \(w\). One may omit \(e(\beta w)\) at
power-small error: \(|\beta|E\ll X^{-.16+o(1)}\), and the total
first mass of the orbit scalars is at most \(X^{.026+o(1)}\) as
proved below. On each occupied upper step the normalized absolute
periodic profile has mass \(O(R+E)\). The same argument applies on
the continuous reference side, using the smooth density.

Split \(u/d=u_s/d_s+u_c/d_c\) into primitive rational terms via
the Chinese remainder theorem. For each fixed \(u_c,d_c\), expand
the periodic profile in characters \(\psi\bmod Ed_s d_c\). Its
normalized coefficients \(b_{\psi,u_s}\) satisfy
\begin{align}
 \sum_\psi\left|\sum_{u_s} b_{\psi,u_s}f_{u_s}\right|^2
     &\ll E^\delta\sum_{u_s}|f_{u_s}|^2, \label{dt:analytic-mask-fibre-square}\\
 |b_{\psi,u_s}|
     &\ll d_c^{-1/2}L^C
       \min\{X^{-.201\gamma_s+b_s},\
                    d_s^{-1/2}(E,d_s)^{-(c-1/2)}\} .
             \label{dt:analytic-mask-fibre-point}
\end{align}
Here \(X^{\gamma_s}\mid Ed_s\) is the smooth part of the inducing
conductor. For \eqref{dt:analytic-mask-fibre-square}, Parseval on
the units reduces the left side to a mean square of the periodic
profiles. At fixed unit \(w\), the quotient \(\lfloor x/E\rfloor\)
runs through all residues modulo \(d_s\), independently of the
admissible classes on the coprime side. Additive orthogonality
therefore gives \(\sum|f_{u_s}|^2\). The squared bottom profile
has mean \(O(E^\delta)\).

On the coprime side, write
\(e(-u_c\lfloor x/E\rfloor/d_c)=e(-u_c\overline E x/d_c)
 e(u_c\overline E w/d_c)\). The first factor supplies a normalized
Gauss coefficient \(O(d_c^{-1/2}L^C)\). Indeed for a primitive
additive numerator the imprimitive Gauss sum is zero unless the
quotient over the inducing conductor is squarefree and coprime
to it, in which case its absolute value is the square root of the
inducing conductor. This follows from inclusion--exclusion of the
unit restriction and the primitive Gauss sum formula. The second
factor is a harmless additive twist in \(w\).

For the first smooth option in \eqref{dt:analytic-mask-fibre-point},
fix the quotient modulo \(d_s\) and apply the linear fibre to the
translated word in \(w\). If \(d_s\mid E\), the half cut \(R_0\) for
the smooth conductor fits in \(E\). The slope modulus \(Q\) may
exceed the available upper grid, in which case split it into
\(Q/(Q,E/R_0)\) shifted grids; this number divides \(d_s\).
The smaller shifted moment bounds do not increase the first norm,
as each mixed-grid exponent is nonnegative in both valuations.
This loses at most \(d_s\) relative to
\eqref{dt:analytic-linear}. Translation changes the inverse residue
labels but leaves the same occupancy bound. If \(d_s\nmid E\), this
option is trivial because \(b_s\ge(\log_{10}2)e_*+o(1)\) and
\(.201\gamma_s-b_s\le .201(e_*+b_s)-b_s<0\).

For the second smooth option, the quotient is a complete quadratic
sum since \(d_s\mid E^2\). More explicitly at fixed \(w\) the
character on \(1+Ez\bmod Ed_s\) has a slope \(\ell\) in the
coordinate \(z-(E/2)z^2\bmod d_s\). After inserting the primitive
quotient numerator \(u_s\) the quadratic coefficient is
\(\pm\ell E/(2w^2)\) and the linear coefficient is
\(\pm\ell/w-u_s\), up to conjugation. If the slope is not a unit
at a prime of \(d_s\), the primitive linear coefficient makes the
sum zero there. For a unit slope, the elementary quadratic Gauss
bound is \(O(\sqrt{(E,d_s)/d_s})\) after normalization. Its support
requires the linear coefficient to vanish modulo \((E,d_s)\), up
to a bounded loss at 2. This follows by translating the summation
variable by multiples of \(d_s/(E,d_s)\), or by squaring and using
additive orthogonality. Thus \(w\) is confined to \(O(1)\)
residue classes of that modulus, up to that bounded loss.
Digit occupancy costs another \(O((E,d_s)^{-c})\). This proves
\eqref{dt:analytic-mask-fibre-point}. Since
\(\log_X(E,d_s)\ge\min(b_s,.301 e_*)+o(1)\), the smooth saving
is at least
\[
 \max\{.201\gamma_s-b_s,\ .5b_s+.403\min(b_s,.301 e_*)\}
\]
up to rounding.

The scalar masses needed above and below are
\begin{align}
 \sum_{d_c\asymp D_c}\sum_{u_c,u_s,\beta}|A_{u,\beta}|
        &\ll d_s^a D_c^{2\delta}(1+V_0)^a L^C, \label{dt:analytic-mask-lone}\\
 \sum_{d_c\asymp D_c}\sum_{u_c}
      \left(\sum_{u_s,\beta}|A_{u,\beta}|^2\right)^{1/2}
        &\ll d_s^{\delta/2}D_c^{2\delta}(1+V_0)^{\delta/2}L^C .
                \label{dt:analytic-mask-ltwo}
\end{align}
Here is the digit-block calculation. First sample the smooth
centers, of spacing \(1/d_s\), with an initial block of length
comparable to \(d_s\), using the shifted first or square bound.
Then skip to a decimal position divisible by \(d_s\). Use two
successive blocks of lengths comparable to \(D_c^2\) on the
coprime reduced centers. Dilation by either radix power
preserves their separation \(\gg D_c^{-2}\). The square mass on
each of the two grids is \(O(D_c^{2\delta})\); Cauchy gives the
same bound for the product of the two factors. In
\eqref{dt:analytic-mask-ltwo} these factors are placed in local
suprema before summation over smooth centers and offsets.
At a given center the offsets form a grid of mesh \(1/q\)
over width \(O(V_0/q)\); a terminal block starting at
\(\gg q/(1+V_0)\) in original positions has arithmetic period
comparable to \(1+V_0\), by taking a decimal cut of the mixed
denominator. Use its first or square bound.
The separated initial sampling ends at exponent at most
\(e_*+4(b_s+b_c)\), and
\(e_*+4(b_s+b_c)+v_0<.354\) by
\eqref{dt:analytic-mask-budget}, strictly below \(\log_X q\).
Thus sample suprema on the corresponding short perturbation
intervals using the bounded reproducing kernel; the terminal
factor is disjoint. Omitting or reassigning orbits only decreases
these absolute mass bounds.

Let \(W_1=.428(b_s+v_0)+.194b_c\). Smooth only the upper step in
logarithmic coordinates at width \(Z^{-1}X^{-W_1-.00004}\).
At a jump the absolute mass of one normalized profile in the
affected interval is \(O(RX^{-W_1-.00004}L^C+E)\).
Use \eqref{dt:analytic-mask-lone} and sum over all jumps. This
is power-small since \(e_*+W_1\ll\lambda\) with fixed reserve.
Multiplication by the modulation is performed after smoothing.
Differentiation in the log coordinate bounds all Mellin tails
beyond \(X^{.00004}\max(ZX^{W_1+.00004},O(XV_0/q))\) to arbitrary
power accuracy. In the orbit budget the retained height satisfies
\begin{equation}\label{dt:analytic-mask-height}
 \tau\le z+\max(\ell_*,W_1)+.00032,\qquad z=1-\lambda.
\end{equation}
Indeed the incidence budget bounds
\((h_j+\log_X V_0-(\log_Xq-m)_+)_+\) by \(\ell_*\).
Rounding and fixed log factors fit within the height reserve.

At a fixed rational numerator and denominator, Plancherel and
the spacing \(1/q\) of the offsets give
\[
 \left\|\sum_\beta A_{u,\beta}\widehat V_\beta\right\|_2^2
 \ll \rho_U L^C\sum_\beta|A_{u,\beta}|^2.
\]
In fact on each physical upper step of length \(R\) the separated
exponential vectors have squared integral at most
\(O(R\sum_\beta|A_{u,\beta}|^2)\), by a smooth majorant at scale
\(R\) and the decay of its transform. The same estimate holds on
commonly translated log steps, since the physical scale and
frequency spacing change by bounded factors. Sum over occupied
steps and use the log Jacobian and smoothing by averaging.
Combine this bound with \eqref{dt:analytic-mask-fibre-square}
before applying character Cauchy--Schwarz, and then use
\eqref{dt:analytic-mask-ltwo}. For the pointwise option use
\eqref{dt:analytic-mask-fibre-point},
\eqref{dt:analytic-mask-lone}, and \eqref{dt:analytic-upper-point}.
The resulting relative bounds on a layer of deficit \(y\) are
\begin{align}
 P_1(y)&=-y+D(y)(\gamma_s+b_c+\tau)-B_T+W_1-.5b_c\notag\\
       &\quad-\max\{.201\gamma_s-b_s,\ .5b_s+.403\min(b_s,.301e_*)\},
            \label{dt:analytic-mask-point-exp}\\
 Q_1(y)&=-y+\tfrac12D(y)(e_*+b_s+b_c+\tau)
                  +.0485(e_*+b_s+v_0+\eta-\lambda)+.194b_c .
            \label{dt:analytic-mask-square-exp}
\end{align}
For each \(d_c\) use the fixed-modulus density estimate at the
inducing conductor. The summed norm
\eqref{dt:analytic-mask-ltwo} accounts for all coprime numerators
and denominators without counting them a second time.

Here are the exponent checks needed for these bounds. In
\([.180,.233]\) use \(D(y)\le2.4y\) and
\(\gamma_s\le e_*+b_s\); the pointwise bound increases with
\(\gamma_s\) in this interval. In the first curvature regime,
\(D(y)<.572\); maximize its height contribution at \(\tau=u\)
and use \(u\le.230\). In the second \(D(y)>.2375\);
maximize at \eqref{dt:analytic-mask-height}. For the second
regime and for the square bound it suffices to set
\(\eta=1\), because the coefficients of \(u=1-\eta\)
with \(z=1-m-h_j-.00003\) held fixed are nonpositive:
with our constant cuts, the pointwise coefficient is
\(-.0969\) and the square coefficient is \(-.0485\).
Then use \(z\le.48867\). The \texttt{exponents} certificate
checks the first-curvature pointwise bound at both ends;
it checks the high-curvature pointwise bound at \(.180\),
the square bound at \(.233\), and a convex combination of
these last two exponents on each deficit subinterval of
width \(10^{-4}\). It checks the combination using the more
expensive endpoint for each exponent and tests all 110
vertices from the regions of affinity of the maxima and minima
in the mask budget. This certifies strictly negative margins
on the full interval. The square term decreases thereafter,
since \(e_*+b_s+b_c+\tau<.750+.097(.230)<5/6\);
after \(y=.25\) its derivative using the Ingham branch is
smaller still. The term without zeros is paid by the same
square argument with trivial count \(O(Ed_s d_c T)\) and
\(y=1/2\), which is checked separately.

At \(y=0\) in \eqref{dt:analytic-mask-point-exp} the quadratic
option pays at least \(.5b_s\) and the Gauss coefficient pays
\(.5b_c\), leaving strict savings \(.072b_s\) and \(.306b_c\).
The term \(.428v_0\) is paid by the nonstationary estimate if
\(T\ll X|\beta|L^{-C}\), since \(R|\beta|\gg V_0\).
Otherwise, for \(V_0\ge1\), \(\tau\ge z-o(1)\), up to a
logarithmic factor only. In the second regime
\(.2375z>.100\); in the first use
\(.572(z-.230)>.100\). These bounds absorb \(.428v_0\).
For \(V_0<1\) that cost is zero. Interpolate the linear density
bound between this endpoint and the strict bound at \(.180\)
for the same shell and family. The result is \(-c_1y\) with
an additional negative multiple of \(b_c\) when \(y\le.090\).
Here the height reserve occurs in \(D(y)\tau\), not as an
additive loss at \(y=0\). Use partial summation with the
logarithmic-loss Huxley bound, splitting the \(d_c\)-ranges
at a sufficiently large power of \(L\).
Above that threshold the factor \(d_c^{-c_2}\) pays arbitrary
logarithms and the \(d_c^{o(1)}\) number of inducing divisors.
Below it the logarithmic divisor bound and
\eqref{dt:analytic-zero-free} suffice. On the central interval
take the height \(X^{2\epsilon}\) for zeros and use the same
argument, since \(2.4(.192+2\epsilon)<1\).

We detail the contributions of slots not coprime to \(d_c\).
For each \(d_c\), select precisely those slots whose prime
divides \(d_c\), and fix their values. There are
\(O_{\kappa,r}(1)\) choices, since each selected prime is at
least \(X^{\kappa+o(1)}\). Write \(a_0>1\) for the product of
these fixed slots with multiplicity, and \(x=a_0 x'\).
Put \(d_0=(a_0,d_c)\), \(d_c'=d_c/d_0\),
\(\alpha_0=\log_X a_0\), and \(\zeta_0=\log_Xd_0\le\alpha_0\).
All residual prime polynomials retain the restriction \(p\nmid d_c\),
even if a prime dividing \(d_c\) was erased in \(d_c'\).
This differs from the primitive inducing prime polynomial by
only \(O(L^C)\). The coprime quotient phase in \(x'\) has
primitive numerator at modulus \(d_c'\); expand it there.
The smooth part of the profile remains character-expanded
modulo \(Ed_s\), with its coefficients multiplied by the unit
phase \(\chi(a_0)\). The square bound
\eqref{dt:analytic-mask-fibre-square} still holds by orthogonality
in the smooth coordinate, since the coprime exponential has
modulus one. All upper-step and scalar estimates are unchanged
because they still test \(x\) at the original physical scale.
The pointwise coprime Gauss coefficient now costs
\((d_c')^{-1/2}L^C\). The remaining product scale is \(X/a_0\)
up to logarithms, so for terms with no poles the change in the
pointwise exponent is at most
\[
          -\alpha_0(1-y)+\zeta_0/2
                    \le -\alpha_0(1/2-y)\le0 .
\]
In the square branch there is no Gauss loss. The estimates
already proved therefore bound these terms as well.
In the low-deficit part they even save a fixed power since
\(\alpha_0\ge\kappa+o(1)\) and \(y\le .180\).
If all slots are fixed, the absolute first mass
\eqref{dt:analytic-mask-lone} and the bound \(O(E^\delta)\)
for one profile give a power saving directly against \(X\rho_U\).

Next consider the terms with a principal pole. Its profile
coefficient vanishes when \(d_s>1\), by the average in the
smooth quotient coordinate. For \(d_s=1\) the coprime principal
coefficient is \(O(1/\varphi(d_c))\), or
\(O(1/\varphi(d_c'))\) for fixed nonunit slots. In the latter
case the normalization of the residual product supplies
\[
             a_0^{-1}/\varphi(d_c')
                      \ll d_c^{-1}L^C
\]
because \(d_0\le a_0\). This applies to the all-pole and mixed
principal terms. On central shells the pointwise upper bound
is \(O(\rho_U(1+V_0)^{-1}L^C)\); use the nonstationary
estimate for \(V_0\ge1\) and the trivial bound otherwise.
Thus \eqref{dt:analytic-mask-lone} sums the principal coefficients
on each family at cost
\[
             O(D_c^{2\delta-1}(1+V_0)^{a-1}L^C).
\]
Mixed terms are negligible by the fixed-rank argument at the start
of the section; the coprime and offset sums contribute only
logarithms. In the case of a fixed nonunit slot, \(D_c\gg X^\kappa\),
so the all-pole contribution also saves a fixed power.

For the remaining, all-unit tuples the all-pole term is the
continuous box density tested against the mean on the units
of \(Ed_s d_c\) of the corresponding normalized profile.
We compare this with the integer-tested continuous model on
units of 10. On each shifted step average over periods of
length \(Ed_s d_c\); either normalization has absolute mass
\(O(|I|+Ed_s d_c)\) on an interval \(I\).
Partial summation against the smooth density and the modulation
has relative error \(O(Ed_s d_c(R^{-1}+|\beta|)L^C)\) at an
orbit. This saves a fixed power after summing the orbit masses
by \eqref{dt:analytic-mask-lone}; the relevant bounds are
\(Ed_s d_c\ll X^{.192+o(1)}\),
\(|\beta|\ll X^{-.360+o(1)}\), and \(R\gg X^{.511}\).
The period means at both unit normalizations are zero if
\(d_s>1\). When \(d_s=d_c=1\) the normalizations coincide.
Otherwise \(d_s=1,d_c>1\): the complete coprime residue mean
in the integer-tested model is zero by additive orthogonality,
whereas the unit mean has size \(O(1/\varphi(d_c))\).
At height zero the upper integral is
\(O(\rho_U(1+V_0)^{-1}L^C)\).
For \(d_c\ge L^C\) the bound just proved for principal coefficients
pays arbitrary logarithms by choosing \(C\) large.

For \(1<d_c<L^C\), the orbit scalar has further decay
\(\exp(-c_2\log X/\log(2d_c))\). To prove this use a fixed
positive fraction of the first digits in the block above \(E\).
Its argument is the reduced fraction \(u_c/d_c\) perturbed by
\(E\beta\); the perturbation remains power-small throughout
that fraction of the block because \(E|\beta|\ll X^{-.16+o(1)}\).
Under multiplication by ten a nonzero residue modulo \(d_c\)
remains nonzero. Within \(O(\log(2d_c))\) multiplications its
distance to the nearest integer must exceed \(1/20\):
until that event the distance grows by a factor of ten and
starts at no less than \(1/d_c\). At such a digit position the
one-digit amplitude is uniformly bounded away from one.
Indeed the alphabet of eight digits necessarily contains two
consecutive digits. Thus each such run gives a uniform
contraction. A forced leading digit and the bounded taper do
not affect the contraction on this initial portion.
There are \(O(d_c^2(1+V_0))\) orbits at a fixed small
coprime size, which costs only logarithms after the upper
factor \((1+V_0)^{-1}\). This proves the model comparison in
all denominator ranges. Restore the unsmoothed step using
the earlier jump estimate and restore the suffix phase using
its power-small bound. This completes the comparison on the
assigned orbits and proves the proposition.
\end{proof}

\subsection{Reconstruction on prime boxes}\label{dt:comparison-proof}

\begin{proof}[Proof of Proposition~\ref{dt:grouped}]
Apply Proposition~\ref{dt:analytic-projection} to the initial
projection \(E_{Y_0}f_X\), and telescope in the outer period
\(Y=10Y_0,\ldots,X\). For each increment use the two cuts and
leading-digit pieces of Section~\ref{dt:filters}. Put each piece
on the additive period grid of its own cut. All digital comparisons
for an increment are in units of \(X\rho_L\rho_U\), before the
factor \((X/Y)^{-\delta}\) in \eqref{dt:outer-factor}.

For \(q<X^{.192}\) discard the low-filtered piece and remove the
high multiplier from the high-filtered piece, as justified in
Section~\ref{dt:diagonal}. Proposition~\ref{dt:analytic-small-periods}
then gives the reference for the unfiltered high-cut piece.
Throughout this decomposition the reference acts on a source \(g\) as
\((10/\varphi(10))\sum_{(v,10)=1}g(v)\mathcal B(v)\). Individual
additive period components need not vanish on nonunits. After summation
the restriction is redundant because \(f_X\) is supported on units.
For \(X^{.192}\le q<MX^{-.0001}\), the zero-twist part of
Proposition~\ref{dt:gram-high} gives a power saving. The remaining
range \(q\ge MX^{-.0001}\) is split into the rational bands and
their complement. For that range, the same proposition makes
the complementary terms power-small.
Proposition~\ref{dt:analytic-mask} supplies the reference
for the bands. In both cases use
the true prime products, with the fixed-rank interpretation of
the constituent weights. Congruence completions require only
divisor bounds on the grouped coefficients, and the Mellin
comparisons separate the actual prime variables.

For the medium periods and the complementary higher periods the
continuous reference is itself power-small. This uses no comparison
theorem on the deleted frequency sets. On the units, a periodic function
with exact additive frequency denominator \(q>X^{.192}\) has zero
mean over \(\operatorname{lcm}(q,10)\) consecutive integers.
Indeed the Fourier support of the unit indicator has denominators
dividing \(10\); its product with a primitive period-\(q\) Fourier
mode has zero constant coefficient. Also
\(\operatorname{lcm}(q,10)\mid R_j\). Each translated upper step
of length \(R_j\) is therefore a union of complete periods of mean
zero. The normalized absolute average of the lower component on
any such step is \(O(q^{\delta/2})\) by Parseval, even with an
arbitrary complementary mask. A smooth box density varies by at most
\(O((R_j/X)L^{O(1)})\) over a step, by differentiating its
pushforward in logarithmic coordinates. Summing over the occupied
steps gives a relative error \(O((R_j/X)q^{\delta/2}L^{O(1)})\),
which saves a power because \(q\le R_j\) and
\(\lambda_j\le.57706\).

The weights and the reference in \eqref{dt:box-equation} are linear
in the digital source. Adding the two cut contributions and all
leading-digit pieces restores the original increment by the exact
convolution identity. There are only \(O(k^2)\) possible smooth
lower periods at each level and \(O(k)\) outer levels, so arbitrary
inverse powers of \(L\) survive these sums. The initial projection
and the increments sum exactly to \(f_X\), which proves the proposition.
\end{proof}

\subsection{Application to two decimal omissions}
For decimal words omitting two digits take $B=10$, $s=8$,
$\gamma=5u/16$ and $T=.560$. Propositions~\ref{dt:digit-type-i}
and~\ref{dt:grouped} give the inputs of Theorem~\ref{common:positive-sieve}:
$[.420,.480]\subset(.4195,.4801)$, $1-.480<T$, and $1-T/2<c$.
The loss in Section~\ref{dt:sieve-certificate} is less than $.934$.
This gives the lower bound of Theorem~\ref{main:two-decimal} at
ordinary endpoints; Corollary~\ref{dt:digit-prime-upper} gives its upper bound.

\clearpage

\section{Primes omitting one septenary digit}\label{sec:seven}
\begin{theorem}\label{seven:main}
For $0\le d\le6$, let $\pi_{7,d}(x)$ count primes $p<x$ whose
ordinary base-seven expansion omits $d$. There exist absolute
$C_2>C_1>0$ and $x_0$ such that, for $x\ge x_0$ and every $d$,
\[
 C_1\frac{x^c}{\log x}\le \pi_{7,d}(x)\le C_2\frac{x^c}{\log x},
 \qquad c=\frac{\log6}{\log7}.
\]
In the lower bound one can take $C_1=1/100$.
\end{theorem}

Unlike the decimal-pair comparison, the septenary comparison uses a single
cut depending on the shorter group scale. Its arithmetic input concerns
products of genuine primes, not unrestricted Type II coefficients.

Throughout, $X=7^k$, $L=\log X$, $\mathcal D=\{0,\ldots,6\}\setminus
\{d\}$, and $\delta=1-c$. A word at radix scale $B=7^j$ consists
of $j$ positions, including the leading positions even if their
digits are zero. The integer parameter $B$ is its period. Write
$f_X$ for the $X$-periodic indicator of permitted $k$-position
words whose units digit is nonzero. Its mass on a period is
$u6^{k-1}$, where $u=6$ for $d=0$ and $u=5$ otherwise. Unless
indicated, implicit constants can depend on fixed test functions,
their derivative bounds, and the indicated error exponent. All
seven alphabets share the constants by finiteness. Decimal
numbers specifying exponents in the proof are exact rationals.

The first Fourier moment already gives divisor distribution and hence the
upper bound. The lower bound requires more: Proposition~\ref{seven:grouped}
compares products of primes while retaining the digital weight in the
reference density. We first compare a projection onto a shorter word,
then restore one digit at a time. At each stage the signed row estimate
handles the complement of a set of rational frequencies; the comparison
on that set uses the prime factors individually.

\subsection{Shifted Fourier moments of the septenary word}
\label{b7:fourier-section}

For \(D=\{0,\ldots,6\}\setminus\{d\}\), \(Z=7^k\), and
\(\delta=1-\log_7 6\), write
\[
 f_d(t)=\tfrac16\left|\sum_{a\in D}e^{2\pi i at}\right|,
 \qquad F_{d,Z}(t)=\prod_{j=0}^{k-1} f_d(7^jt).
\]

\begin{proposition}\label{b7:word-moments}
For each entry of the tables, uniformly in \(d\in\{0,\ldots,6\}\),
\(k\in\mathbb Z_{\ge0}\) and \(u\in\mathbb R\),
\[
 \sum_{l=0}^{Z-1} F_{d,Z}((l+u)/Z)^s\ll\lambda_s^k\le Z^{a_s}.
\]
\[
\begin{array}{c|rrrr}
 s&1&4/3&3/2&2\\ \hline
 \lambda_s&2.17457&1.64337&1.46669&7/6\\
 a_s&.39922&.2553&.19684&\delta
\end{array}
\]
\[
\begin{array}{c|rrr}
 s&7/5&127/64&3\\ \hline
 \lambda_s&1.56675&1.17229&1.03\\
 a_s&.23075&\delta+.159/64&.016
\end{array}
\]
For each nonsquare row the bound persists when each digit amplitude
is floored at a sufficiently small fixed positive constant: the strict
transfer inequality below is stable under uniform perturbations.
For \(s=2\) the unrestricted sum is exactly \(Z^\delta\).
Prescribing nonempty subsets of allowed digits at a bounded number
of positions preserves the upper bounds, with digit polynomials normalized
by their numbers of allowed choices. The constant may depend on the number
of prescribed positions.
\end{proposition}

\begin{proof}
For \(s=2\), orthogonality on the shifted grid gives \(Z6^{-k}\),
because the words give distinct residues modulo \(Z\). Write
\[
 T_s h(u)=\sum_{r=0}^6 f_d((r+u)/7)^s h((r+u)/7).
\]
Iteration identifies \(T_s^k1(u)\) with the shifted sum. Reflection
\(a\mapsto6-a\) preserves \(f_d\). The self-contained verifier
\path{verify_moments.py} in the septenary supplement constructs one strictly positive periodic
piecewise linear \(h\le1\) for each \(d=0,1,2,3\) and each nonsquare row (with constant test function at \(s=3\)) and checks \(T_s h<\lambda_s h\). The estimate follows from
iteration and comparison of \(h\) with constants. For prescribed
positions, bound the corresponding amplitudes by one. In inverse-branch
summation each altered position then costs at most seven in supremum
norm, and each intervening run of length \(m\) costs
\(\|T_s^m1\|_\infty\ll\lambda_s^m\). This also works at \(s=2\).

We detail the reduction of the continuous inequality to rational
calculations. The \(N=2100\) nodes are \(i/N\) on the circle. On each
node cell all \(u\mapsto h((r+u)/7)\) are affine. Center the single-digit
polynomial by writing \(z(u)=\frac16\sum_{a\in D}
e^{2\pi i(a-3)(r+u)/7}\). Since \(\sum_{a\in D}|a-3|\le12\) and
\(\sum_{a\in D}(a-3)^2\le28\), its first and second derivatives have
magnitudes less than 2 and 4. A supporting functional for the convex
map \(w\mapsto |w|^s\) at \(w=z(u)\) can be chosen of norm at most
\(s\). Applying it at \(z(u\pm t)\) gives
\[
 |z(u+t)|^s+|z(u-t)|^s-2|z(u)|^s\ge -4s t^2.
\]
Thus \((|z|^s)''\ge-4s\) distributionally, and \(|z|^s\) has Lipschitz
constant at most \(2s\). On each open node cell, the product rule gives
\((T_s h-\lambda_s h)''\ge -s(28+4\operatorname{Lip}h)\).
Consequently the residual is below the interpolation of its two nodal
values plus \(s(28+4\operatorname{Lip}h)/(8N^2)\). The two endpoints
of the full mesh agree by periodicity of \(T_s h\) and of \(h\).

For branch nodes \(t=j/(7N)\) the verifier uses
\[
 36 f_d(t)^2
 =6+2\sum_{v=1}^6 \#\{a\in D:a+v\in D\}\cos(2\pi vt).
\]
All coefficients are nonnegative. Directed integer cosine bounds use
Machin's formula with alternating arctangent series and the degree-41
cosine Taylor polynomial on \([0,\pi]\), with error at most
\((22/7)^{42}/42!\). The verifier rounds the square weight up in units
of \(2^{-32}\), caps it at one, then uses a ceiling integer root to
bound its \(s/2\)-power. Finite matrix iteration selects nodal values
of \(h\) which are integers in units of \(2^{-40}\).
The proof uses only the positivity and final residual tests on that
vector. Direct evaluation of the seven branches with the upper weights,
linear interpolation and the stated curvature allowance certifies a
strictly positive residual in every case. Comparison of integer powers gives \(\lambda_s<7^{a_s}\) in
every nonsquare row (using \(7^\delta=7/6\)). This proves the proposition.
\end{proof}

\begin{lemma}[Sampling on separated sets]\label{b7:digit-sampling}
Write \(F_Y\) for any of the normalized word amplitudes in
Proposition~\ref{b7:word-moments} at scale \(Y\). For each displayed \(s\),
\[
 \int_{\mathbb R/\mathbb Z}F_Y(t)^s\,dt\ll Y^{a_s-1}.
\]
Let \(\Theta\subset\mathbb R/\mathbb Z\) have mutual circular distances
at least \(\Delta>0\). For every fixed \(C\ge0\),
\[
 \sum_{\theta\in\Theta}\sup_{\|t-\theta\|\le C/Y}F_Y(t)^s
 \ll_{s,C}(Y+\Delta^{-1})Y^{a_s-1}.
\]
For any integer \(1\le Q\le Y\) one consequently has
\[
 \sup_{u\in\mathbb R}\ 
 \sum_{j=0}^{Q-1}\ \sup_{|v|\le C/Q} F_Y((j+u)/Q+v)^s
 \ll_{s,C}Q^{a_s}.
\]
If \(Q\mid Y\), the last sum restricted to any \(H\) consecutive
indices is \(O_{s,C}(H^{a_s})\), for \(1\le H\le Q\).
\end{lemma}

\begin{proof}
Integrating Proposition~\ref{b7:word-moments} over the shift and
dividing by $Y$ gives the integral bound at every block length.
Both sampling assertions follow from Lemma~\ref{common:sampling}.
\end{proof}

\begin{lemma}[Physical occupancy]\label{b7:digit-occupancy}
Let \(\mathcal W_Y\) be the set of complete words at scale \(Y=7^k\)
and \(c=\log_7 6\). Any interval of length \(H\) in \([0,Y)\)
contains \(O((1+H)^c)\) words. Any residue class modulo an integer
\(1\le q\le Y\) contains \(O((Y/q)^c)\) words.
Both inequalities hold for every subset of \(\mathcal W_Y\).
At \(k\ge1\) the number of words prime to seven is exactly
\(\#(D\cap\{1,\ldots,6\})\,6^{k-1}\).
\end{lemma}

\begin{proof}
Lemma~\ref{common:occupancy} gives both inequalities. Each permitted
unit digit is followed to its left by $k-1$ free positions.
\end{proof}


\subsection{Divisor distribution on arbitrary intervals}
\label{b7:type-i-section}

Put $c=\log_7 6$, $X=7^k$, and $D=\{0,\ldots,6\}\setminus\{d\}$.
Let $f$ be the indicator of the words $n<X$ with all $k$ digits in $D$.
In addition, the lowest and highest digits may each be restricted to a
nonempty subset of $D$. Put $A_f=\sum_{0\le n<X}f(n)$.
In particular the lowest digit may be required to be a unit modulo seven.
For $k\ge2$ we have $6^{-2}X^c\le A_f\le X^c$.

\paragraph{Interval Type I estimate.}
For any fixed $\theta<1-\log_7(2.17457)$, $J,K>0$, and $Q\le X^\theta$,
\begin{equation}
 \sum_{\substack{q\le Q\\(q,7)=1}}\tau(q)^K
 \sup_{I,a}\left|\sum_{n\in I} f(n)
       \left({\bf1}_{n\equiv a\bmod q}-\frac1q\right)\right|
 \ll_{\theta,J,K} A_f(\log(2X))^{-J}.
 \label{b7:type-i}
\end{equation}
Here $I$ ranges over intervals in $[0,X)$ and $a$ over residues modulo $q$.
This includes $\theta=.60078$ and $\theta=.560$ with fixed reserve. It holds as well with
$f$ replaced by any convex average of its translates on $\mathbb Z/X\mathbb Z$.
All constants are uniform in the omitted digit and in the two restrictions.

Indeed, integrating the shifted first moment of Proposition~\ref{b7:word-moments}
over one shift and dividing by the block length $Z$ gives
$\int_0^1 F_{d,Z}(t)\,dt\ll Z^{-\sigma}$ with
$\sigma=1-\log_7(2.17457)$. We have $1/2<\sigma<c$, and after any one
omission the six allowed digits still contain a consecutive pair.
Proposition~\ref{arith:typeI} therefore applies, with the actual
interval mass of the restricted word. For fixed $\theta<\sigma$ choose
$\epsilon>0$ with $(1+\epsilon)\theta<\sigma$ in \eqref{arith:typeI-max}.
The two prescribed end digits cost only a constant in the mass; averaging
cyclic translates uses linearity of the interval remainder.

\paragraph{Prime upper bound.}
Let $A_{7,d}(x)$ count positive integers below $x$ whose ordinary
expansion omits $d$, without a restriction on the units digit. Use
\eqref{b7:type-i} on the word ending in a unit, with $K=2$ and any
$0<T<\sigma$. Since $T<\min(1,2c)$, Proposition~\ref{common:upper}
applies with $B=7$, $s=6$, $u=6-\mathbf1_{d\ne0}$ and
$\gamma=7u/36$. Letting $T\uparrow\sigma$ gives
\[
 \limsup_{x\to\infty}\frac{\pi_{7,d}(x)\log x}{A_{7,d}(x)}
 \le \frac{2\gamma}{1-\log_7(2.17457)}.
\]
This limit uses only fixed levels strictly below $\sigma$. Since
$A_{7,d}(x)\ll x^c$ by summing over ordinary lengths, it implies the
upper bound in Theorem~\ref{seven:main}, also when zero is forbidden.

\subsection{Comparison on products of primes}\label{seven:setup}
For fixed $r\ge2$, take $r$ primes at scales $P_i\ge X^\kappa$
with $\kappa>0$ fixed, and allow any fixed powers of $L$ in the
scale inequalities. Their weights are smooth with compact support in
fixed-multiple intervals about the $P_i$. Allow a joint smooth weight
with product support in $[X L^{-D},X)$, $D$ fixed. The mixed
derivatives of each prescribed order in the coordinates
$\log(t_i/P_i)$ are $O(L^{C})$, with $C$ depending on the order.
Let $W$ be the combined weight and let $\mathcal B(x)\,dx$ be the
pushforward of $W(\boldsymbol t)\prod_i dt_i/\log t_i$ by multiplication.
It has dyadic support and
$|\mathcal B(x)|+x|\mathcal B'(x)|\ll L^C$ there, by solving for the
last variable in the product. For a digital component $G$, use
the linear reference functional
\[
  \mathcal R(G)=\frac76\sum_{n\ge1,\,7\nmid n} G(n)\mathcal B(n).
\]

\begin{proposition}\label{seven:grouped}
Suppose that the prime variables split into two groups, of scale
products $M=X^m L^{O(1)}$ and $P=X^{1-m} L^{O(1)}$, with
$.4195\le m\le .4801$. For every fixed $J>0$ one has
\[
 \sum_{p_1,\ldots,p_r}W(\boldsymbol p) f_X(p_1\cdots p_r)
   =\mathcal R(f_X)+O_J(6^k L^{-J}).
\]
\end{proposition}

A smooth partition on the logarithmic box and a Fourier series on a
slightly larger box separate joint weights. The smooth Fourier tail,
truncated to $L^{C(J)}$ in each variable, is uniformly $O(L^{-J'})$
for any prescribed $J'$, by integration by parts. Restore fixed smooth
one-variable cutoffs equal to one on the joint support before making
this expansion. For $f_X$ the cost of this approximation is
$O(6^k L^{-J'})$ times at most a fixed logarithmic power, since
each integer has at most $r!$ representations as an ordered product
of $r$ primes. The continuous density of the box satisfies the
bound above. Thus it suffices to treat separated smooth boxes. On
those boxes retain a smooth physical cutoff supported at the product
scale and equal to one on their support, including when part of a box
is outside the original product interval. There $x\asymp\prod P_i
=X L^{O(1)}$. All estimates below permit these fixed logarithmic
distortions and the periodic extension of the digit indicator.

For a divisor $Y$ of $X$, let $E_Y$ average an $X$-periodic
function on each residue class modulo $Y$. Put
$Y_0=7^{\lfloor .770k\rfloor}$ and telescope from $E_{Y_0}f_X$
through successive $Y=7Y_0,\ldots,X$. At each level take
\[
 h=.09693,\quad H=X^h,\quad \lambda=m+.09696,\quad
 R=7^{\lfloor\lambda k\rfloor},\quad \eta=\log_XY .
\]
For large $k$ we have $7R\mid Y$. With $F$ the lower $R$-word
indicator including its units condition and $U(n)$ the upper
word at $\lfloor n/R\rfloor$ of period $Y/R$ minus $6/7$ times the
upper word of period $Y/(7R)$, one has identically
\[
 (E_Y-E_{Y/7}) f_X(n)=(X/Y)^{-\delta}F(n\bmod R)U(n).
\]
Indeed the free upper digits have conditional density $6/7$ each.
The density scales are $\rho_L=R^{-1}\sum_{w<R}F(w)\asymp
R^{-\delta}$ and $\rho_U=(Y/R)^{-\delta}$. There are
$O((X/R)\rho_U L^C)$ occupied steps on the product range, counting
both signs, by counting periods of each upper word. In particular
$(X/Y)^{-\delta}X\rho_L\rho_U\asymp X^c$.
On the lower word use the exact additive period decomposition
$F=\sum_{q\mid R}G_q$, where
\[
 G_q=\sum_{e\mid q}\mu(q/e)E_e F
 =\rho_L\!\sum_{a\bmod q}^{*}\Pi_R^{\circ}(-a/q)e(an/q)
\]
in the last expression evaluated at $n$. Here $e(t)=\exp(2\pi i t)$
and $\Pi_R^{\circ}$ is the normalized exponential sum of the
permitted lower word. We establish the comparison for these sources
before reconstructing $f_X$.


\subsection{The initial projection on prime products}
\label{b7:initial-section}

Put \(c=\log_7 6\), \(X=7^k\), \(L=\log X\), and
\(Y=7^{\lfloor .770k\rfloor}\). Write \(w_B\) for the admissible
word on a radix period \(B\) with its units digit restricted to
\(\{1,\ldots,6\}\), extended periodically. Put
\(A_B=\sum_{0\le n<B}w_B(n)\). Thus, for \(B\ge7\),
\(A_B=uB^c/6\), where \(u\) is the number of allowed nonzero digits.
The projection \(E_Y w_X\), the average on each class modulo \(Y\),
equals \((X/Y)^{c-1}w_Y\).

\subsubsection{A twisted cubic character moment}
\label{b7:refined-cube-section}

Let \(q=7^v\), let \(B\ge q\) be a seven power, and take a complete word
in \([0,B)\) restricted to units, with any bounded number of prescribed
digit positions. Denote its mass by \(A_B\). Allow an arbitrary real
additive twist \(\alpha\). Write
\[
 b_\chi=A_B^{-1}\sum_{\text{allowed }n<B}
             e^{2\pi i\alpha n}\chi(n).
\]
Take \(a=.19684\), \(b=.016\), \(\delta=1-\log_7 6\), and set
\[
 h=(1-2a)/2-\delta,\qquad j=.007,\qquad h'=h+j/2>.22743.
\]
We show
\begin{equation}\label{b7:refined-cube}
 \sum_{\chi\bmod q}^{*}|b_\chi|^3\ll q^{-h'}.
\end{equation}
The star restricts the sum to primitive characters.

The two additive inputs are the $3/2$ and cubic rows of
Proposition~\ref{b7:word-moments}; in particular the cubic exponent
$b$ refers to the sum of cubes, not its cube root. Both rows allow the
shifts and prescribed digits in the present statement.

We derive first the preliminary bound with \(h\) in place of \(h'\).
It suffices to consider \(v\ge2\). Put \(R=7^{\lceil v/2\rceil}\),
\(Q=q/R\), and \(U=(\mathbb Z/Q\mathbb Z)^*\).
For \(n=w+Ru\), a primitive character satisfies
\(\chi(n)=\chi(w)\exp(2\pi i\ell\overline w u/Q)\) for some
\(\ell\in U\). Indeed the kernel \(1+Rz\bmod q\) is additive in
\(z\bmod Q\), and a slope divisible by seven would make the character
trivial on the kernel of reduction modulo \(q/7\). For each slope the
characters are either absent or a coset of the lifts of characters
modulo \(R\). In fact every character on this subgroup extends, though
only the upper bound on the fibre sum is needed.

Write \(W\) for the lower allowed unit words in \([0,R)\),
\(K=|W|\asymp R^c,\ c=1-\delta\), and write \(F\) for the normalized
upper amplitude of the word at scale \(B/R\). Define on \(U\)
\[
 A(x)=F(x/Q+\alpha R)^{3/2},\qquad
 M(\ell)=\sum_{w\in W} A(\ell\overline w).
\]
Since \(Q\mid B/R\), we may bound all factors above \(Q\) by one
and apply the shifted digital moments. This gives
\(\sum_{x\in U} A(x)\ll Q^a\), \(\sum_{x\in U} A(x)^2\ll Q^b\).
Counting Hausdorff--Young on the units modulo \(R\) in each slope
fibre, for the unnormalized sum transform, gives
\[
 \sum_{\chi\text{ of slope }\ell}|b_\chi|^3
       \le \varphi(R)K^{-3} M(\ell)^2 .
\]
Each lower residue has at most \(R/Q\le7\) lifts. Consequently
\(\max_\ell M(\ell)\ll Q^a\) and
\(\sum_{\ell\in U}M(\ell)\ll KQ^a\).
It follows that the total cubed moment is
\(\ll R K^{-2}Q^{2a}\ll q^{\delta-(1-2a)/2}=q^{-h}\).
This preliminary result holds for all complete words in the statement,
including the untwisted lower word used next.

Let \(\mu\) be the pushforward to \(U\) of uniform probability on
\(W\), so \(M=K(\mu*A)\) for unnormalized convolution on \(U\).
Write \(\widehat{\ }\) for the sum transform on \(U\). Parseval gives
\[
 \sum_{\ell\in U}M(\ell)^2=
 \frac{K^2}{|U|}\sum_{\psi\in\widehat U}
                  |\widehat\mu(\psi)|^2|\widehat A(\psi)|^2 .
\]
Split at conductor \(Q^{.91}\). The characters with smaller
conductors (including equality) all factor through the largest seven
power \(Q_0\le Q^{.91}\). Each residue modulo \(Q_0\) has mass
\(O(Q_0^{-c})\) under \(\mu\), by its low digital block or
Lemma~\ref{b7:digit-occupancy}. Parseval gives
\(\sum_{\operatorname{cond}\psi\le Q^{.91}}|\widehat\mu(\psi)|^2
\ll Q^{.91\delta}\). With \(|\widehat A|\ll Q^a\), this part of
\(\sum M^2\) is \(O(Q^{c+2a-.09\delta})\).

For the larger conductors, apply the preliminary cubed moment to the
lower word at every exact primitive conductor \(q_0\mid Q\).
There is one character in the family per primitive character of that
conductor. Summing a geometric series yields
\(\sum_{\operatorname{cond}\psi>Q^{.91}}|\widehat\mu(\psi)|^3
\ll Q^{-.91h}\).
Also
\(\sum_\psi |\widehat A(\psi)|^6
\le\|\widehat A\|_\infty^4 |U|\sum_{x\in U}A(x)^2
\ll Q^{4a+1+b}\).
H\"older with exponents \(3/2,3\) shows that the second part of
\(\sum M^2\) is
\[
 \ll Q^{2c-1-(2/3)(.91)h+(4a+1+b)/3}
 = Q^{c+2a-((\delta-b)/3-.06h)}.
\]
Both \(.09\delta\) and \((\delta-b)/3-.06h\) exceed \(j\), by
\(.07921<\delta<.07923\). Thus \(\sum M^2\ll Q^{c+2a-j}\).
Reusing the slope-fibre estimate saves \(Q^{-j}\asymp q^{-j/2}\)
over the preliminary result and proves \eqref{b7:refined-cube}.


In particular, for the untwisted word at $B=q$ the primitive
coefficient vector has $\ell^3$ norm $O(q^{-h'/3})$.

\begin{proposition}\label{b7:initial-projection}
Fix a number \(r\ge1\) of prime variables on scales
\(P_i\ge X^\kappa\), where \(\kappa>0\), and assume
\(\prod_i P_i=X L^{O(1)}\). Take smooth compactly supported weights
\(W(\mathbf t)\) on a product of fixed-ratio intervals \(t_i\asymp P_i\),
with each fixed logarithmic derivative bounded by a fixed power of \(L\).
Write \(\mathcal B\) for the pushforward density of
\(W(\mathbf t)\prod_i dt_i/\log t_i\) by multiplication.
For each fixed \(J>0\),
\[
 \sum_{p_1,\ldots,p_r} W(\mathbf p)(E_Yw_X)(\textstyle\prod_i p_i)
 =\frac{7}{6}\sum_{n\ge1}(E_Y w_X)(n)\mathcal B(n)
   +O_J(X^c L^{-J}).
\]
\end{proposition}

\begin{proof}
First take a product of the one-variable weights. Expand \(w_Y\)
in multiplicative characters modulo \(Y\). After taking out
\(A_Y/\varphi(Y)\), its constant coefficient is one. For a character
induced by a primitive \(\chi\) of conductor \(q\ge7\) its normalized
coefficient is the corresponding conjugated character coefficient of \eqref{b7:refined-cube}, since the positions above \(q\)
are free admissible digits. All prime variables are units for large
\(X\). It suffices to prove an error \(O_J(XL^{-J})\) in the resulting
normalized sum.

Here are the density and truncation details for the nonprincipal
coefficients. Write \(N_q(1-y,T)\) for the number of zeros with
real part at least \(1-y\) and ordinate of absolute value at most
\(T\), over primitive characters of conductor \(q\), with multiplicity.
Use the fixed-modulus Ingham bound \eqref{dt:analytic-density} at
fixed deficits $0<y<1/2$, with $\nu(y)=3y/(1+y)$. Thus the count is
$O_\varepsilon((qT)^{\nu(y)+\varepsilon})$; at the final mesh endpoint
we bound the number of characters, not zeros, by $q$.

For \(q>L^B\) truncate the explicit formula separately at
\(T=q^{\xi}\), with \(\xi=10^{-5}\). Let \(1-y_\chi\) be the
maximum of \(1/2\) and the real parts of zeros of \(L(s,\chi)\) up
to this height. For a smooth weight \(w_i\) as above, the explicit
formula gives, for each fixed \(H\ge3\),
\[
 \left|\sum_p w_i(p/P_i)\chi(p)\right|
 \ll_H L^{C_H}\big(P_i^{1-y_\chi}+P_i T^{1-H}\big).
\]
To see this, use the von Mangoldt explicit formula with the weight
\(w_i(t/P_i)/\log t\), integrating by parts \(H\) times in
\(\log(t/P_i)\). Its zero term at \(\beta+i\gamma\) is
\(O(P_i^\beta L^{C_H}(1+|\gamma|)^{-H})\). The zero count in a
unit interval is \(O(\log(q(2+|\gamma|)))\). Summing above the
cutoff produces the indicated tail. The prime powers, trivial
zeros and the explicit formula's conductor term are covered by the
first bound, since \(1-y_\chi\ge1/2\) and \(P_i\ge X^\kappa\).
Take \(H\) so large that \(T^{1-H}\le q^{-4}\).
Terms in the product using any tail are \(O(XL^{C}q^{-4})\) per
character. All constants and derivative orders here are fixed before
choosing \(B\).

The remaining product is \(O(XL^C X^{-y_\chi})\). Since
\(q\le X^{.770}\), \(X^{-y_\chi}\le q^{-y_\chi/.770}\).
By H\"older and \eqref{b7:refined-cube}, the exponent in
base \(q\) on a zero layer at deficit \(y\), apart from the height,
density loss and mesh width, is bounded above by
\[
 E(y)=-y/.770+2y/(1+y)-h'/3.
\]
By differentiation its maximum on $[0,1/2]$, multiplied by $.770$,
is $1+2(.770)-2\sqrt{2(.770)}-.770h'/3<-.00030$.
Partition \([0,1/2]\) in intervals of length \(10^{-4}\), using
the right endpoint for the zero count and the left for the weight.
In the last interval bound the number of characters by \(q\),
which agrees with \(\nu(1/2)=1\). In the first
interval use the density bound at \(y=10^{-4}\). As \(\nu\le1\),
the extra exponent from the width, \(T=q^\xi\), and the density
bound with \(\varepsilon=10^{-5}\) is at most
\[
 10^{-4}/.770+\tfrac23(\xi+(1+\xi)10^{-5})<.00015 .
\]
Thus all layers at this conductor give \(O(XL^{C}q^{-.00020})\).
One can recover the stronger $q^{-.002}$ saving without changing the
cut. For $T=q^\xi\ge7$, $q$ is $T$-smooth. By
\cite[Theorem~1.2 and (1.4)]{ChenGuptaLi2026} take instead
$\nu(y)=\min(30y/13,3y/(1+y))$. Even with $h'/3$ weakened to
$.0746$, the maximum of $-y/.770+2\nu(y)/3-.0746$ occurs at
$y=.3$: the first piece increases and the derivative of the second
is at most $-1/.770+2/(1.3)^2<0$. The maximum is less than
$-.0026$, before the same mesh and height costs. This smooth-modulus
improvement is not needed for the prime-count theorem.
There are \(O(L)\) conductors. Choosing \(B\) large gives the
required logarithmic saving and also ensures \(T\ge7\).

For \(q\le L^B\), Siegel--Walfisz and partial summation give
arbitrary logarithmic savings for every nonprincipal constituent
prime sum, uniformly since \(\log P_i\asymp L\). The number of
characters in these conductors is \(O(L^{B+1})\). The principal
character contribution is the continuous product integral with the
same accuracy by the prime number theorem with its classical error.

To pass back to the integer-tested continuous model, the partial
sums of \(w_Y-A_Y/Y\) have absolute value at most \(2A_Y\).
The density \(\mathcal B\) has supremum and total variation
\(O(L^C)\), by integration in the logarithmic coordinates of the
fixed number of variables. Summation by parts and comparison of sums
with integrals therefore give an error \(O(A_YL^C)\) before
multiplication by \((X/Y)^{c-1}\). This is
\(O(X^c(Y/X)L^C)\). Since \(A_Y/\varphi(Y)=(7/6)A_Y/Y\),
the main terms have the stated normalization.

For a joint weight, extend it smoothly by zero on a fixed larger
logarithmic box and expand in a Fourier series with separated modes.
Integration by parts to any fixed order \(h>r+2\) bounds each
coefficient by \(O(L^{C_h}(1+\|\mathbf n\|_\infty)^{-h})\).
For any prescribed \(A\), a cutoff at \(L^{K}\) in each coordinate
therefore gives supremum error \(O(L^{-A})\) when
\(K(h-r)>C_h+A\). The remaining
logarithmic weights satisfy the preceding derivative estimates. The
number and the total absolute coefficient mass cost only powers of
\(L\). For the truncation error each integer has at most \(r!\)
ordered representations as a product of \(r\) primes. Its total
projected mass in the box is \(O(X^cL^{C})\), by counting complete
periods. For the continuous model the same bound follows from the
pushforward of a smooth product majorant. This proves the proposition.
\end{proof}


\subsection{The row square at the lower digit cut}

\subsubsection{Averaging the signed digital lift in base seven}
\label{sec:seven-signed-lift}

Put $c=\log_7 6$ and $\delta=1-c$. Write $e(t)=\exp(2\pi i t)$.
For a power $D$ of seven let $\mathcal W_D$ be the word set on
$\log_7D$ digits, and let
\[
 \Pi_D(t)=|\mathcal W_D|^{-1}\sum_{w\in\mathcal W_D}e(wt).
\]
A superscript $\circ$ permits restriction of the terminal digit to units,
or of the leading digit to one prescribed allowed digit. A block inherits
only the restrictions on its own ends. We use the shifted first moment
$\sum_{b\bmod D}|\Pi_D((b+v)/D)|\ll D^a$ uniformly in real $v$, with
$a<1/2$. Bounded end restrictions preserve this bound because dropping a terminal
factor leaves seven copies of a shifted grid for the shortened word. Dropping
a leading factor partitions the original grid into seven such shifted grids.
For the square moment, discrete orthogonality gives
$\sum_b|\Pi_D^\circ((b+v)/D)|^2=D/|\mathcal W_D^\circ|\ll D^\delta$.

Take $M=X^{m+o(1)}$, $P=X^{n+o(1)}$, $m+n=1$, $m\le n$, where
the scale distortions are at most powers of $\log X$. Let $R\mid Y\mid X$
be powers of seven, with $R=X^{\lambda+o(1)}$, $Y=X^{\eta+o(1)}$.
The upper source is bounded and supported on
$O((X/R)\rho_U X^{o(1)})$ blocks of length $R$ on the product range,
where $\rho_U=(Y/R)^{-\delta}$. This support estimate follows for the
projected difference by counting the two upper words separately.
Take $H=X^{h+o(1)}$ with $HM\ll R$,
and write $A_U=\delta(\eta-\lambda)$. In the application one has
$\lambda=m+h+.00003$ and $\lambda/2<n$. For a lower exact period $q\mid R$
write the lower source as
\[
 G_q(t)=\rho_L\sum_{a\bmod q}^{*} J(a)\Pi_R^\circ(-a/q)e(at/q),
             \qquad |J(a)|\le1.
\]
Here $\rho_L$ is the density of the lower word. The modulus $q=1$ has
one character by convention.

Let the rows $p\asymp P$ and columns $y\asymp M$ have coefficients bounded
by $X^{o(1)}$ and supported on units modulo seven. Apply the column completion leading to
\eqref{dt:gram-incidence}: since $HM\ll R$, there are
$O((PM/H)\rho_U X^{o(1)})$ occupied column--row-interval pairs.
Cauchy--Schwarz, dyadic subdivision of the row intervals and Poisson
summation of the resulting nonnegative column squares give
\begin{equation}
 p a-(p+\ell)b\equiv t\pmod q,\qquad
       0<|\ell|\ll H,\qquad |t|\le T=(1+q/M)X^\epsilon .
 \label{eq:seven-signed-incidence}
\end{equation}
The coefficient of each twist is $O(M)$ and depends on $t$ alone.
Discarded twists have arbitrarily large power saving by rapid decay,
since all remaining sums have at most polynomial size. The row weights
and dyadic interval conditions depend on $(p,\ell)$ alone. In the
squared estimate normalized by $(PM\rho_L\rho_U)^2$, each normalized
off-diagonal incidence sum is divided by $PH\rho_U$, up to $X^{o(1)}$.
The diagonal can be grouped by $py$, whose divisor multiplicity
is $X^{o(1)}$ and $\sum_{u\ll X}|G_q(u)|^2\ll
X\rho_L^2\sum_{a\bmod q}^{*}|J(a)\Pi_R^\circ(a/q)|^2$. Thus the
relative squared diagonal is
\begin{equation}
 \ll X^{A_U-h+o(1)} \sum_{a\bmod q}^{*}|J(a)\Pi_R^\circ(a/q)|^2
 \ll X^{\delta\eta-h+o(1)} .
 \label{eq:seven-signed-diagonal}
\end{equation}

\begin{lemma}[Exact difference order]\label{lem:seven-signed-lift}
Let $g=7^i$ be the exact additive order of $a-b$ in
\eqref{eq:seven-signed-incidence}. Put
\[
 V=7^{\lfloor i/2\rfloor},\quad g_0=g/V,\quad Z=q/g,\quad Q=q/V,
 \quad C=3c/4-a,\quad s=\log_Xg,\quad \gamma=\log_Xq .
\]
If the label multipliers are constant on increments of $Q$ in each
label, the relative squared off-diagonal contribution of this order is
\begin{equation}
 \ll X^{A_U+O(\epsilon)+o(1)}g^{-C}\max(1,g/P)^{1/2}
                    Z^\delta(1+T/Z).
 \label{eq:seven-lift-basic}
\end{equation}
One can retain the following more precise factor in place of
$Z^\delta(1+T/Z)$:
\begin{equation}
 \frac1H\sum_{z\bmod Z}^{*}|\Pi_{R/g}^\circ(z/Z)|^2
      \sum_{0<|\ell|\ll H,\ |t|\le T}1_{Z\mid t+\ell z}.
 \label{eq:seven-lift-precise}
\end{equation}
If $q<MX^{-b}$ for fixed $b>0$, choose $\epsilon<b/2$. Then the relative
squared off-diagonal contribution of any such order is
$O(X^{A_U-C\gamma+o(1)})$, assuming $0<C\le c$.
\end{lemma}

\begin{proof}
Say $g>1$ first. Write $a=A+Qx$, $b=B+Qy'$, with coarse labels
$A,B\bmod Q$ and $x,y'\bmod V$. Exact order amounts to $A-B=Zu$
with $u$ a unit modulo $g_0$. Indeed $Q=Zg_0$, $V\mid g_0$,
and a unit modulo $g_0$ is a unit modulo $g$. Primitivity of $a,b$
is constant on each lift. The necessary congruences on the coarse
labels and row are $Z\mid t+\ell B$ and
$p u\equiv(t+\ell B)/Z\pmod{g_0}$, respectively. The rows lie in
a progression modulo $g_0$. On it, $(p+\ell)/p\equiv r\pmod V$
is constant. The ratios refer to the supported unit rows. When
$V>1$ their unit conditions are constant on the progression.
The remaining condition on $(x,y')$ is
$x\equiv r y'+w\pmod V$, with
\[
 w\equiv p^{-1}\big((t+\ell B)/Z-pu\big)/g_0\pmod V ,
\]
where the parenthesized integer is first divided by $g_0$ and only then
multiplied by $p^{-1}$. Thus consecutive rows $p,p+g_0$ have intercepts
differing by $-up^{-1}\pmod V$, a fixed unit. Coarse labels and $t,\ell$
are held fixed at this point.

Apply Lemma~\ref{dt:finite-signed-lift} with the bottom word
$S=\mathcal W_V^\circ$, $N=|S|\asymp V^c$ and shifts $A/Q,B/Q$.
Indeed $x=ry'+w\bmod V$ is the fibre in \eqref{dt:gram-intercept-identity};
the higher digit factors and both label multipliers are constant on it.
Only after summing this fibre with its signs do we bound the row weights
and drop the dyadic row restrictions. The enlarged row interval is
contained in $J\asymp P/g_0$ successive progression steps since
$g_0\ll R^{1/2}\ll P$. With $|S\cap r^{-1}S|\le N$ and $Vg_0=g$,
the lemma bounds the row sum of the absolute correlations by
\[
 \ll (P/g_0)V^{1-3c/2}\max(1,g/P)^{1/2}.
\]
For $V=1$ the same estimate follows by counting rows. No overlap
estimate specific to a mixed radix enters here.

Both coarse labels have common residue $z$ modulo $Z$. The next
$g_0$ digits in positional \emph{length} (that is, a block of
$\log_7g_0$ digits above $V$) contribute two independent shifted
first sums, each $O(g_0^a)$, after dropping the congruence restrictions
on the coarse labels above $z$. More explicitly, $A=z+Zb_1$
gives the argument $(b_1+z/Z)/g_0$ in this block. Above position $g$
each frequency leaves the same absolute factor
$|\Pi_{R/g}^\circ(z/Z)|$. The constant multipliers on the lifts
can now be bounded by one. From $g_0,V\asymp\sqrt g$ we obtain
$g_0^{2a-1}V^{1-3c/2}\ll g^{a-3c/4}$. Dividing the row bound by
$PH\rho_U$ gives \eqref{eq:seven-lift-precise}.
For $g=1$, both labels are equal and $Z=q$. The same conclusion
follows by bounding the rows by $O(P)$ directly.

At fixed $z,\ell$ there are $O(1+T/Z)$ twists.
Since $Z\mid R/g$, the transform of the digits at and above length
$Z$ is constant on $z/Z$, and the shifted square sum below $Z$ costs
$O(Z^\delta)$. This proves \eqref{eq:seven-lift-basic}. If
$q<MX^{-b}$ the Poisson cutoff leaves only $t=0$ with an arbitrarily
power-small remainder, since each nonzero twist has Fourier argument
of magnitude at least $M/q>X^b$. Now $z$ is a unit and $Z\mid\ell$.
There are $O(H/Z)$ nonzero such gaps, including when $Z\gg H$
and there are none. As $g\le q\ll P$, the bound becomes
$X^{A_U+o(1)}g^{-C}Z^{-c}\le X^{A_U+o(1)}q^{-C}$.
\end{proof}

Since $Z\le R/g$ and $q/Z=g$, the exponent of
\eqref{eq:seven-lift-basic} is at most
\[
   \delta(\eta-s)-Cs+(s-m)_++\tfrac12(s-n)_+
                 +O(\epsilon)+o(1).
\]
This bound and the $t=0$ conclusion require only a first digital
moment and Parseval.


\subsection{Complement of the rational bands in the seven-adic row square}
\label{sev:high}

Write $X=7^k$, $c=\log_7 6$, $\delta=1-c<.080$,
$a=.401$, and $C=3c/4-a>.289$. All decimals in this
section are rational. We use a shifted first-moment bound with
exponent $a$ for the normalized amplitude $F_B$ of a word
of length $\log_7 B$, and the shifted square bound with exponent
$\delta$. Either bound permits restriction to a unit last digit
and to a specified leading digit at constant cost. The same bounds
hold on one local supremum per separated point at spacing
$\gg1/B$, with perturbations $O(1/B)$, by convolving the trigonometric
polynomial with a smooth reproducing kernel bounded by
$O(B(1+B\|x\|)^{-2})$ and summing the translates, applying
Cauchy--Schwarz with this positive majorant in the square case.
The needed continuous integral is bounded by averaging the
corresponding shifted grid bound and dividing by $B$.

Use $M=X^{m+o(1)},P=X^{n+o(1)}$, $n=1-m$,
$.4195\le m\le .4801$. The $o(1)$ in these scales allows
logarithmic factors. At outer period $Y=X^\eta$, take
\[
 h=.09693,\quad H=X^h,\quad \lambda=m+.09696,\quad
 R=7^{\lfloor\lambda k\rfloor},\quad
 A=\delta(\eta-\log_X R)<.0387+o(1).
\]
The upper step has density scale $\rho_U=(Y/R)^{-\delta}$.
Let $q\mid R$ be the exact additive period, $g\mid q$ the additive
order of a difference of frequency labels, $s=\log_X g$,
$D=q/g$, $V=7^{\lfloor (\log_7 g)/2\rfloor}$, and
$T=(1+q/M)X^\epsilon$. Only row gaps $0<|\ell|\ll H$ and integer
twists $|t|\ll T$ enter. Fixed multiples and factors $X^{O(\epsilon)+o(1)}$
are understood in power-saving estimates and $\epsilon$ may be fixed
arbitrarily small. For normalization divide the row square by
$(PM\rho_L\rho_U)^2$, with $\rho_L$ the density of its lower word.

We apply the signed lift of Lemma~\ref{lem:seven-signed-lift}. In the bound for fixed $g$
the coarse factor is
$X^{A-Cs}\max(1,g/P)^{1/2}$, and the remaining factor is bounded by
\begin{equation}
 \frac1H\sum_{0<|\ell|\ll H}\sum_{|t|\ll T}
 \sum_{\substack{z\bmod D,\ (z,D)=1\\D\mid\ell z+t}}
 F_{R/g}(z/D)^2 .
 \label{sev:fold}
\end{equation}
If a frequency multiplier is present this estimate applies when it is
constant on all $V$ lifts, that is, under translation of the frequency
by $1/V$. The coarse summands supported on the removed lifts are absent.
The variable $z$ is the common residue of the two labels modulo $D$.
Bounding the $t$ sum at fixed $z,\ell$ and applying the square bound gives
$D^\delta(1+T/D)$. The diagonal is
$O(X^{A-h}q^\delta)=O(X^{-.016+o(1)})$.
In the zero-twist range $q<MX^{-.0001}$ only $t=0$ remains,
up to a negligible error. The residue $z$ is a unit whenever
$D>1$, so $D\mid\ell$ and the factor in \eqref{sev:fold} is
$O(D^{\delta-1})$. There is no incomplete row loss in this range.
Hence for $q\ge X^{.192}$ its exponent is at most
$A-C\log_X q<-.016+o(1)$.

For $q\ge MX^{-.0001}$ and $s\ge .220$ the unmasked exponent is
at most
\[
 .080(1-s)-.289s+(s-m)_++\tfrac12(s-n)_+.
\]
Indeed $D\le R/g$ and $T/D\ll X^{\epsilon}(1+g/M)$.
This convex function on the trapezoid
$.4195\le m\le .4801$, $.220\le s\le m+.09696$
is bounded by its maximum at a vertex. At $s=.220$ it equals
$-.00118$, while at $s=m+.09696$ it is less than $-.013$
and $-.007$ at the smaller and larger $m$, respectively.

For $q\ge MX^{-.0001}$, $s\le .220$, we describe explicitly the
exceptional bands which require an arithmetic comparison.
There is no incomplete-row factor. Write $\gamma=\log_X q$,
$\tau=(\gamma-m)_+$. Then $\gamma\ge .4194-o(1)$ and
$D/T\ge X^{m-.0001-s-O(\epsilon)+o(1)}$.
For a summand of \eqref{sev:fold}, reduce the rational
$(\ell z+t)/(\ell D)$, taken modulo one, to a center $a_0/d$
with $d\mid\ell$. Group $d\asymp B$ dyadically and
$|t/\ell|\asymp U$ dyadically with zero displacement separate.
At fixed center and $z$, $t$ is determined by $\ell$, so the
gap multiplicity divided by $H$ is
\[
 O(B^{-1}\min(1,T/(HU))).
\]
For zero displacement omit the minimum. The estimate counts only
nonzero multiples of $d$ and retains no additive one for a short
interval starting at the origin.

Since $d\mid\ell$, we have $B\ll H$ and $BU\ll T$.
In the present range
\[
 B^2\max(1,U)\ll H\max(H,T)\le
 X^{.19389+O(\epsilon)+o(1)},\qquad
 D\ge X^{.1994-o(1)} .
\]
This also holds for zero displacement with $U=0$. Thus all
relevant centers are separated at the required scale.
In every shell,
use an initial block of radix length comparable to $B^2$ in the
polynomial $F_{R/g}(z/D)$. Here and below radix lengths may be
rounded to a power of seven. The centers are $\gg B^{-2}$ separated
and $z/D$ differs from its center by $O(U/D)$, so the sum of the
local square suprema is $O(B^{2\delta})$. Choose a power of seven
$J\mid D$ for which $D/J\asymp1+U$. The block beginning at $J$ of
radix length $D/J$ is disjoint from the initial one and fits since
$D\mid R/g$. At fixed center the possible $z$ form at most a bounded
number of consecutive intervals of length $O(1+U)$. The terminal
block on this set has square mass $O((1+U)^\delta)$ by the grid bound.
Consequently the factor in \eqref{sev:fold} on this shell is
\begin{equation}
 O(B^{2\delta-1}(1+U)^\delta\min(1,T/(HU))).
 \label{sev:separated}
\end{equation}
At zero displacement both $U$ factors are replaced by one.
Bounded lengths at either endpoint may be absorbed into the constant.

Now write $b=\log_X B$, $v=\log_X U$,
$v_0=\max(v,0)$. Mark a family when
\begin{equation}
 .249s+(1-2\delta)b+(h+v-\tau)_+-\delta v_+
                 \le A+.00015 .
 \label{sev:active}
\end{equation}
At zero displacement omit the $v$ terms. Every unmarked family has
a fixed power saving by \eqref{sev:separated}, since $C>.249$.
Each label lies within $O(U/q)$ of a center of denominator
dividing $gd$, since $g(a/q)=z/D\bmod 1$. For zero displacement
this contradicts the primitivity of the label, as $gd\ll
X^{.220+h+o(1)}=o(q)$.

Enlarge these marked sets of labels by translations by $1/E$, where
$E=\max(g,7^{\lfloor.117k\rfloor})$ and $e=\log_X E$. Each
enlarged family has a covering rational center whose denominator
after multiplication by $E$ is $d_s d_c$, with $d_s$ a power of
seven and $(d_c,7)=1$. Write $d_s=X^{b_s}$ and
$d_c\asymp X^{b_c}$ with $d_c$ grouped dyadically.
In each $E$-orbit the frequencies have the form
$u/(Ed_s d_c)+j/E+\beta$, $(u,d_s d_c)=1$,
$|\beta|\asymp U/q$. For fixed rational center the offsets are
spaced by at least $1/q$. Since $E$ is a multiple of $g$,
$b_s+b_c\le b+o(1)$. From $\tau-h\le .00003+o(1)$ and
\eqref{sev:active} we obtain, with room to absorb the rounding errors,
\begin{align}
 .249s+.806(b_s+b_c)+.903v_0&\le .0477,\label{sev:budget}\\
 (h+v-\tau)_+&\le .0477-.249s-.806(b_s+b_c)+.097v_0 .
 \label{sev:height}
\end{align}
In particular $s<.192$, $b_s+b_c<.060$, $7E\mid q$, and $d_s\mid E$.
The orbit translations preserve primitive support. Disjoint assignment
of the families can be made in decreasing order of $E$. An earlier
family is invariant under the smaller translation group of each later
family. Choose one covering center per orbit and translate that choice
on the orbit. There are only logarithmically many choices of $g$,
$B$, $d_s$, the dyadic scale of $d_c$, and the dyadic offset.
The union of all assigned orbits is the mask.

The normalized first mass of the mask is
\begin{equation}
 \sum_{a/q\,\mathrm{masked}}F_R(a/q)
        \ll X^{.192a+o(1)} .
 \label{sev:firstmask}
\end{equation}
Here are details, also giving the scalar bounds needed in the arithmetic
comparison. Factor off the lowest $E$-word and put
$A_{u,\beta}=\Pi_{R/E}(u/(d_s d_c)+E\beta)$ for the remaining
normalized complex amplitude, with an allowed leading-digit
restriction. The first mass of the lowest block on an orbit is
$O(E^a)$. Before this factor the two useful bounds are, with $D_c=X^{b_c}$,
\begin{align}
 \sum_{d_c\asymp D_c,u_c,u_s,\beta}|A_{u,\beta}|
    &\ll X^{o(1)} d_s^a D_c^{2\delta}(1+U)^a,\label{sev:lone}\\
 \sum_{d_c\asymp D_c,u_c}
        \left(\sum_{u_s,\beta}|A_{u,\beta}|^2\right)^{1/2}
    &\ll X^{o(1)} d_s^{\delta/2}D_c^{2\delta}(1+U)^{\delta/2}.
 \label{sev:ltwo}
\end{align}
In these two inequalities the $X^{o(1)}$ can be taken to be a fixed
power of $\log X$. Use Chinese remainder coordinates
$u/(d_s d_c)=u_s/d_s+u_c/d_c\pmod1$. A first block of length
$d_s$ handles the $u_s$ grid by the shifted first or square estimate,
using local suprema for the perturbation. After dilation by $d_s$,
the following block of length $B'\asymp D_c^2$ samples a
$\gg D_c^{-2}$ separated set since
$u_c\mapsto d_su_c\pmod {d_c}$ permutes primitive numerators.
Use a second consecutive block of the same length. The new numerator
is again multiplied by a unit. The two squared local-supremum sums
are $O(D_c^{2\delta})$ each, and Cauchy--Schwarz bounds the sum of
their products by $O(D_c^{2\delta})$. The low smooth-grid bound is
uniform for every coprime center before this summation.
For each center put a final block of length $\asymp1+U$ starting at
a position $\asymp q/(1+U)$ in the original coordinates and ending
at $q$. Each retained offset interval consists of $O(1+U)$ consecutive
mesh points with step $1/q$ (omissions only reduce the mass), so this
block contributes $O((1+U)^a)$ or $O((1+U)^\delta)$ for the
corresponding first or square sum. The earlier blocks, including
their perturbations, fit strictly before this one since
\[
 e+4(b_s+b_c)+v_0
    \le .117+4(.0477)/.806+o(1)<.354<\gamma.
\]
For $s\le.117$ this follows directly from \eqref{sev:budget}.
For $s>.117$ the upper bound $s+4(.0477-.249s)/.806+o(1)$
is decreasing, since $4(.249)>.806$.
Also $e+b_s+b_c+v_0\le .192+o(1)$, by
\eqref{sev:budget}, maximizing at one endpoint of each of the
intervals $s\le.117$ or $.117\le s\le .0477/.249$.
Now $2\delta<a$, so \eqref{sev:lone} proves \eqref{sev:firstmask}.

Finally the mask is invariant under the lift whenever
$V\mid 7^{\lfloor.117k\rfloor}$. For such $V$, removing the mask
removes every marked incidence, and the foregoing bounds apply to
the remainder. If the divisibility fails, $s\ge .234-o(1)$, so
\eqref{sev:active} is impossible by \eqref{sev:budget}. Bound these
orders with the mask absent and then estimate the absolute correction.
For a fixed masked label the partner lies on a shifted $g$ grid,
with first cost $O(g^a)$ by retaining a lowest block of length $g$.
For each gap and fixed pair there are $O(1+T/D)$ possible twists,
each giving $O(1+P/g)$ rows since the difference has exact order $g$.
Dividing by $PH\rho_U$ and using \eqref{sev:firstmask}
bounds the correction exponent by
\[
 .0387+.192a-(1-a)s+(s-n)_++(s-m)_+ .
\]
It is convex in $(m,s)$. On
$.4195\le m\le .4801$, $.234\le s\le m+.09696$
its values on the lower edge are less than $-.024$, and on the
two upper vertices are less than $-.096$ and $-.074$, respectively.
This proves a fixed power saving at every medium period and on
the complement of the rational mask at every large period.


\subsection{Character estimates for the remaining frequencies}

\subsubsection{The fixed-conductor zero estimates}
\label{bseven:external-density}
Let $N^*_q(1-y,T)$ count, with multiplicity, the zeros $\rho$ with
$\Re\rho\geq 1-y$ and $|\Im\rho|\leq T$, summed over the primitive
Dirichlet characters of conductor $q$. For $0\leq y\leq 1/2$, $T\geq2$
and any fixed $\xi>0$, the following two inequalities hold:
\[
 N_q^*(1-y,T)\ll_\xi (2qT)^{(12/5+\xi)y},
 \qquad
 N_q^*(1-y,T)\ll_\xi (2qT)^{3y/(1+y)+\xi}.
\]
For the first bound we use the all-character, single-modulus
form of \cite[Theorem~2.1]{ThornerZaman}; restriction to primitive
characters decreases the count. Conductor one is covered by the principal
character modulo two. For the second see \eqref{dt:analytic-density}.

Here is how the density inputs enter the real-part exponents.
Write $K=\log(2qT)/\log X$, so $0\leq K\leq K_0$ for a fixed
$K_0$. In an upper estimate that uses $(12/5)Ky$, the first theorem
adds at most $K_0\xi y$ to the exponent. If the bound before this
change is $-c_0 y$, choose $K_0\xi<c_0/2$. Bounds using the second
inequality have an additional exponent at most $K_0\xi$ and apply
on ranges with a fixed power saving. The estimates on compact
parameter ranges with a strict margin accommodate either change by
taking $\xi$ below that margin divided by $2K_0$.
Thus the linear density branch at $y\downarrow0$ and the Ingham
branch for $y\geq1/4$ have the required forms. For completeness, put $L=\log X$, and suppose $q\mid7^k d$ with
$d\leq L^A$ for fixed $A\geq0$, $7^k\leq X^{O(1)}$ and
$T\leq X^{O(1)}$. Iwaniec's region \cite[Theorem~2]{Iwaniec1974}, in the form
\cite[(3.1)]{ThornerZaman}, gives, apart from a possible simple real
zero of a real character,
\[
 1-\Re\rho\gg
 \bigl(\log(2\operatorname{rad}q)+
  [\log(3qT)\log\log(9qT)]^{3/4}\bigr)^{-1}.
\]
The zeta zeros satisfy this estimate by the Vinogradov--Korobov
zero-free region. For any exceptional primitive character of
conductor $f\mid 7^k d$, the $7$-part of $f$ is at most $7$,
since the kernel of reduction modulo $7$ has odd order on
$(\mathbb Z/7^j\mathbb Z)^\times$. Thus $f\leq 7d$. Siegel's
theorem gives $1-\beta\gg_{\theta}f^{-\theta}$ for every
$\theta>0$ in this case (see \cite[Theorem~5.28]{IwaniecKowalski2004}).
Taking $\theta=1/(4(A+1))$ gives $1-\beta\gg_A L^{-1/4}$.
As $\operatorname{rad}q\leq7d$, every counted deficit satisfies
\[
     y_\rho:=1-\Re\rho\gg_A
                    L^{-3/4}(\log(2L))^{-3/4}.
\]
In a bound with exponent $-c_0y$, partition the deficits into the
intervals $[j/L,(j+1)/L)$, using the larger endpoint for the
zero count. Every exponent linear or Lipschitz in $y$ changes by
$O(1/L)$, hence by a bounded multiplicative factor. There are
$O(L)$ intervals and $L y_\rho/\log L\to\infty$ uniformly, giving
$O_{A,J}(L^{-J})$ for any fixed $J$. This argument is used for
the polylogarithmic coprime parts of the conductor. For the larger
coprime parts the separate negative power of $d$ in the orbit
estimate pays the logarithmic losses on the small-deficit layers.


\subsubsection{Small additive periods at the single cut}\label{seven:small}

Let \(X=7^k,\ c=\log_7 6,\ \delta=1-c,\ L=\log X\). Here
\(.9207<c<.9209,\ .0791<\delta<.0793\). Assume that at every radix scale
\(B\) the shifted first moment of a normalized word is \(O(B^{.41})\).
Let \(E_e\) denote averaging on residue classes modulo \(e\) on the
integer circle. Consider a fixed smooth prime-product box with
\(P_i\gg X^\kappa\), \(.4195\le m\le.4801\), and constituent lengths
\(P_i\), such that \(X'=\prod_iP_i=XL^{O(1)}\). Powers of \(L\)
in estimates may depend on the fixed smoothness parameters.

Put \(h=.09693,\ R=7^{\lfloor(m+h+.00003)k\rfloor},\ Z=X/R,\ 
z=\log_X Z\le.4836+o(1)\). Write \(Y=7^j\) with
\(.770+O(1/k)\le\eta=\log_X Y\le1\), \(u=1-\eta\).
Let \(F\) be the lower word indicator with unit last digit, of mean
\(\rho_L\asymp R^{-\delta}\). Let \(U\) be the upper step equal to the 
upper word of period \(Y/R\) minus \(6/7\) times the upper word
of period \(Y/(7R)\), both evaluated at \(\lfloor n/R\rfloor\).
It has density scale \(\rho_U=(Y/R)^{-\delta}\), bounded height and
\(O((X/R)\rho_U L^C)\) jumps on the product support.
Set \(G_q=\sum_{e\mid q}\mu(q/e)E_e F\).

\begin{proposition}\label{seven:small-compare}
For \(q\mid R,\ q<X^{.192}\), the weighted prime sum for the
source \(G_q(n)U(n)\) agrees, up to \(O_J(X\rho_L\rho_U L^{-J})\), with
\[
 \frac{7}{6}\sum_{n:\,7\nmid n}G_q(n)U(n)\,\mathcal B(n).
\]
Here \(\mathcal B\) is the product pushforward of the smooth box
against \(\prod_i dt_i/\log t_i\).
\end{proposition}

\begin{proof}
The unit restriction in the reference is used on each source
separately. It acts as the identity on the reconstructed original digit word.

\par\smallskip\noindent\textit{Character expansion and principal coefficient.}  On \((\mathbb
Z/R)^*\), for each inducing conductor \(q_0\), write \(b_\chi=(R\rho_L)^{-1}\sum_{0\le
w<R:\,F(w)=1}\bar\chi(w)\). The normalized coefficient of \(G_q/\rho_L\) is
\[
 b_\chi\sum_{\substack{e\mid q\\q_0\mid e}}\mu(q/e)e/\phi(e).
\]
Indeed \(E_e F\) is constant on the lifts of each class modulo \(e\). Summing \(\bar\chi\) over
its unit lifts vanishes unless \(q_0\mid e\), then each class of a permitted word (automatically a
unit) has \(\varphi(R)/\varphi(e)\) lifts. This multiplier is bounded since at most two terms
occur. At \(q=1,7\) its principal values are respectively \(1\) and \(1/6\). For \(q\ge49\) the
principal value vanishes. The \(q=7\) principal coefficient is required because this component
need not vanish off the units. For characters of primitive conductor \(q_0\) the coefficients
satisfy
\[
 \sup |b_\chi|\ll q_0^{-.201},\qquad \sum_\chi |b_\chi|^2\ll q_0^\delta .
\]
For the first split the digital word at \(D=7^{\lceil\log_7(q_0)/2\rceil}\). Since \(q_0\mid
D^2\), on a unit suffix \(w\) the upper linear phase has slope \(\ell\bar w\bmod Q\), \(Q=q_0/D\)
with \(\ell\) a unit, by primitivity. Each suffix class mod \(Q\) has \(O((D/Q)^c)\) admissible
lifts. Inversion permutes the unit slopes. Upon summation the shifted first bound gives \(\ll
Q^{.41-c}\ll q_0^{-(c-.41)/2}\). For bounded conductors enlarge the constant. The square bound
follows by grouping the word by residues mod \(q_0\), bounding its maximum normalized residue mass
by \(O(q_0^{-c})\), then applying character Parseval.

\par\smallskip\noindent\textit{Mellin bounds for the upper step.}  Write \(V(v)=\Psi(e^v)U(Xe^v)\)
for the compactly supported test in \(v=\log(x/X)\). Powers of \(L\) in the support and
derivatives are allowed as above. If desired partition into fixed-ratio boxes. Write
\(V_\beta(v)=V(v)\exp(-2\pi i\beta Xe^v)\) for any additive modulation. On shells \(|t|\asymp
T=X^\tau\) with \(2\le T\le Z^{1.9}\), write \(r=(\tau-u)_+\). One has
\[
 \|\widehat V\|_{L^2(|t|\le2T)}\ll\rho_U L^C X^{\delta r/2},\qquad
 |\widehat V_\beta(t)|\ll \rho_U L^C X^{-B(\tau)},\quad\]
\[
 B(\tau)=\begin{cases}.572r,&\tau\le2u,\\.2375\tau+.0969u,&\tau>2u.\end{cases}
\]
The square bound uses the absolute step mass. Each log interval of length \(1/T\) contains
\(O(\rho_U L^C T^{-1}(1+TY/X)^\delta)\) in absolute mass by counting complete periods, then first
differing digits on shorter intervals (an interval of length \(K\ge1\) meets \(O(K^c)\) word
points). Use a smooth Fourier majorant of \([-2T,2T]\) with kernel \(O(T(1+T|v|)^{-2})\) and the
global absolute mass \(O(\rho_U L^C)\). For the pointwise estimate consider each of the two unsigned words.
Write $X'$ for the product scale of the box, $Z'=X'/R$, and $B=Y/R$.
In this paragraph replacing $X'$ by $X$ in power bounds costs only
powers of $L$. For bounded or polylogarithmic $T$ the total absolute
mass proves the stated upper bound, since $B(\tau)\ll\tau$.
Otherwise choose a power of seven $B_0\asymp\min(B,Z'/\sqrt T)$,
rounding down and, if necessary, shortening by a fixed factor.
This is at least one for the displayed heights. Write $x=R(j+\theta)$,
$0\le\theta<1$, and take aligned intervals of $B_0$ consecutive
indices $j$. On each such interval meeting the compact support,
its start $j_0\asymp Z'$ is positive. The upper digit constraint is
a fixed choice (zero or one) on the positions above $B_0$ times the
full word on the remaining positions.

Remove the linear term in the expansion of
$\pm t\log (j+\theta)-2\pi\beta R(j+\theta)$ at $j_0$.
Its slope in radians is $\omega=\pm t/j_0-2\pi\beta R$.
The remaining smooth amplitude, including the log Jacobian
with its size $O(1/Z')$ taken out, has derivatives of every fixed
order in $(j-j_0)/B_0$ and $\theta$ bounded by $L^{O(1)}$.
Indeed $T B_0^2/(Z')^2\ll1$. The same bound applies on a fixed
larger coordinate interval. Integrating first in $\theta$ gives
an additional factor $(1+|\omega|)^{-1}$ in every one of the
bounds for the derivatives in the first coordinate, by one
integration by parts against $\exp(i\omega\theta)$.
Multiply by a cutoff on the enlarged interval and use a Fourier
series on an interval of length four in the scaled first coordinate.
For any fixed $A$ the contribution of a block is at most
\[
 \frac{B_0^c L^{O_A(1)}}{Z'(1+|\omega|)}
  \sum_{v\in\mathbb Z}(1+|v|)^{-A}
 F_{d,B_0}\!\left(\frac{\omega}{2\pi}+\frac{v}{4B_0}\right).
\]
The shifts are the same for every block. The following bounds for
a fixed shift are uniform, as the reciprocal weight uses the
unshifted $\omega$.

The frequencies $\omega/(2\pi)$ range through a real interval of
length $O(W/B_0)$, $W=TB_0/Z'$. A cell of length $1/B_0$ contains
at most $O(1+Z'/(B_0W))$ starts, by monotonicity and the derivative
size $\asymp T/(Z')^2$. For $1\le W\le B_0$, retain the highest
digit block of radix length $D\asymp W$ in $F_{d,B_0}$.
Dilation by $B_0/D$ sends these cells to consecutive cells of
size $1/D$, with bounded multiplicity modulo one. Their local
suprema sum to $O(W^{.41})$ by Lemma~\ref{b7:digit-sampling}.
For $W<1$ use the total mass. For $W>B_0$, count by unit intervals
of the unshifted real frequency. The supremum sum in each is
$O(B_0^{.41})$ for any common shift, and the reciprocal step
factor sums over these intervals at cost $O(\log(2+W))$, uniformly
in their location. Thus the weighted sum of cell suprema is
$O((1+W)^{.41}L^C)$.

When $T\le (Z'/B)^2$, we have $B_0\asymp B$ and $W\ll Z'/B_0$.
Multiplying the block bound by the cell occupancy and summing yields
$O(\rho_U L^C(1+TB/Z')^{.41-1})$.
When $T>(Z'/B)^2$, each cell contains $O(1)$ starts.
In this case the bound, divided by $\rho_U$, is at most
\[
 L^C B^\delta\frac{(Z'/\sqrt T)^c}{Z'}(\sqrt T)^{.41}
 =L^C(B/Z')^\delta T^{-(c-.41)/2}.
\]
The two power bounds agree on the transition up to logarithms,
including when the threshold is written $(X/Y)^2$.
This proves the displayed piecewise estimate since, for $\tau>2u$,
$[(c-.41)/2-.2375]\tau+ (\delta-.0969)u\ge0$.
It also gives the stronger raw exponent
$(c-.41)\tau/2+\delta u$ used on the rational orbits.
Replacing $Y$ by $Y/7$ changes only fixed factors. When
$T\ll X|\beta|L^{-C}$ with $C$ sufficiently large for the support,
the first derivative test within each step gives the separate bound
$O(\rho_U L^C/(1+R|\beta|))$. Smoothing the step by averaging
common logarithmic translations preserves all these bounds. For each translate change
the log integration variable in each translate, dilate the modulation
and translate the smooth cutoff. The constants are uniform for the
small translations used here.

The absolute mass of \(G_q\) on a physical jump interval \(I\) is \(O(\rho_L(|I|+q))\), since each
of at most two conditional projections has this bound over its complete periods. Smooth at log
width \(Z^{-1}X^{-\epsilon}\). Summing over the jumps costs \(O(X\rho_U\rho_L
L^C(X^{-\epsilon}+q/R))\). The discrete prime product has bounded multiplicity at every integer
(at most \(r_0!\) for \(r_0\) factors), and the continuous model's smooth pushforward is
\(O(L^C)\). High-frequency tails above \(ZX^{2\epsilon}\) cost arbitrary powers by repeated
differentiation.

\par\smallskip\noindent\textit{Zero estimates.}  Use the fixed-modulus density estimate
\(N_{q_0}(1-y,T)\ll_\xi (2q_0T)^{(2.4+\xi)y}\) of subsection~\ref{bseven:external-density}, with
arbitrarily small fixed \(\xi>0\). Below we write \(2.4y\) in numerical tests. Replacing \(2.4\)
by \(2.4+\xi\) changes any exponent by \(O(\xi y)\), absorbed in the strict margins even for small
deficits. The nonexceptional zeros of these primitive characters at polynomial heights have \(y\gg
L^{-3/4}(1+\log L)^{-3/4}\) by the Iwaniec smooth-modulus region. The only real primitive
characters here have bounded conductor, hence their exceptional real zeros, if any, cause no loss.
For the principal character use the usual zeta zero-free region.

Apply the fixed-rank explicit-formula reduction of
Section~\ref{dt:analytic-section} to the product of pole-subtracted
prime polynomials. The weights here have the same logarithmic derivative
bounds. On a shell $T\ge X^\epsilon$ retain zeros up to $4T$; in the
central range truncate at $X^{2\epsilon}$. At a zero of maximal real
part $1-y$ the product of scale factors is at most $(X')^{1-y}$;
the other kernels sum pointwise to $L^{O(1)}$. The retained kernel for
a character with $N_\chi$ zeros on the layer has integral $O(N_\chi)$
and squared integral $O(LN_\chi)$. The term with no zeros is
$O((X')^{1/2}L^C)$. Terms involving a principal pole are treated separately.

Let \(\gamma=\log_X q_0\le.192\). Taking the sup norm of the transform and the digital coefficient, then summing the kernel in \(L^1\), gives the relative exponent
\[
 P=-y+2.4y(\gamma+\tau)-.201\gamma-B(\tau).
\]
Taking the \(L^2\) Mellin norm for each character and then Cauchy over characters gives
\[
 Q=-y+1.2y(\gamma+\tau)+.0485(\gamma+r).
\]
Here \(.0485>\delta/2\). With \(\tau\le z+2\epsilon\), at \(y=.234\) the upper bounds for \(P\) on
the two curvature ranges are, before arbitrarily small height losses, \(-.03559\) and \(-.00803\).
The upper bound for \(Q\) is \(-.01152\). These follow by setting \(\gamma=.192\), bounding
\(2.4y\tau-.572(\tau-u)_+\) by \(2.4yu\) in the first range, and setting \(\tau=.4836\) in the
second. The slope of \(Q\) is negative for all larger \(y\), since \(1.2(.192+.4836)<1\). For the
no-zero term, the number of characters is at most \(q_0\), the length of the Mellin integral
\(O(T)\). Two Cauchy inequalities give exponent
\(-1/2+.5485\gamma+\tau/2+.0485r<-.1294+O(\epsilon)\). All strict gaps accommodate smoothing and
fixed positive discretization losses.

For \(0\le y\le .234\), \(P(0)\le0\) and its affine dependence yields \(P(y)\le-c_1y\) for
sufficiently small fixed \(\epsilon\), uniformly over all parameters. Integrate in real part by
partial summation (or mesh \(1/L\)) to avoid a constant power loss near one. Iwaniec's region then
supplies arbitrary powers of \(L^{-1}\). At central heights, take the trivial upper transform,
count zeros to \(X^{2\epsilon}\), and use \(2.4(.192+2\epsilon)<1\) in the same calculation. There
are \(O(L)\) inducing conductors/shells separately.

Only the principal character has a pole. On noncentral shells \(|t|\ge X^\epsilon\), every pole
factor is \(O_A(P_i L^C(1+|t|)^{-A})\) by smoothness, so all mixed terms are negligible. On
central heights their remaining slots have combined scale \(X^\xi\) with \(\xi\ge\kappa+o(1)\),
and only ordinary zeta zeros occur. Retaining the rightmost zero gives saving \(X^{-\xi y}\), with
kernel integral bounded as above and count at most \(X^{(2.4)(2\epsilon)y}L^C\). Take
\(4.8\epsilon<\kappa/2\). The term without zeros in the remainder slots costs at most
\(X^{-\xi/2+\epsilon}L^C\). Thus all mixed terms have arbitrary log savings.

The all-pole term is the product pushforward integral against \(U\) times the unit mean of
\(G_q\). On each occupied step of length \(R\), compare it with the stated integer-tested
reference by summation by parts. The mean-zero fluctuation of \((7/6)1_{7\nmid n}G_q(n)\) has
cumulative sums \(O(\rho_L\operatorname{lcm}(7,q))\), by the periodic absolute-mass bound, and the
smooth density and its first derivative are \(O(L^C)\) and \(O(X^{-1}L^C)\). Therefore the
relative error is \(O((q+7)R^{-1}L^C)\). The smooth endpoint replacement has already been paid on
both sides. This proves the low-period comparison.

\end{proof}


\subsubsection{Comparison on the rational frequency orbits}\label{seven-orbits}

We give the argument for a single cut. Write \(L=\log X\),
\(R=X^\lambda\), \(Y=X^\eta\), \(z=1-\lambda\), \(u=1-\eta\),
\(c=\log_7 6\), \(\delta=1-c<.080\), \(a=.428\).
Here and below, \(L^C\) denotes a fixed logarithmic power, which may
change from line to line. The required shifted first-moment exponent
is \(a\). The shifted square exponent is \(\delta\).
We use \(0\le u\le .230+O(1/L)\), \(.422<z\le .48847+O(1/L)\).
All restrictions that follow from rounding powers of seven allow
bounded multiplicative factors. The scalar representing the density
of the lower word has been taken outside the comparison.
The upper step \(U\) has bounded height and \(O((X/R)\rho_U L^C)\)
occupied steps and jumps, where \(\rho_U=(Y/R)^{-\delta}\).

Here are the precise inputs from the frequency decomposition. At the
exact additive period \(q\mid R\), \(q\ge MX^{-.0001}\),
an assigned family consists of complete translates by \(1/E\) with
\(E\mid q\) a power of seven, \(E\asymp X^e\), \(e=\max(s,.117)\).
Each orbit has a representative
\[
 \alpha_0=\frac{u_0}{E d_s d_c}+\beta,\qquad
 (u_0,d_s d_c)=1,\quad d_s=7^j,\quad (d_c,7)=1 .
\]
Fix \(d_s=X^{b_s}\), a dyadic size \(D_c\asymp d_c\), \(b_c=\log_XD_c\),
and a dyadic displacement size \(V\asymp q|\beta|\), \(v=(\log_X V)_+\).
At each center the offsets are \(1/q\)-spaced, and the assignment
uses each orbit at most once. With \(q=X^\gamma\), we use
\begin{align*}
 .249s+.806(b_s+b_c)+.903v&\le .0477,\\
 (h+\log_X V-(\gamma-m)_+)_+&\le
 \ell=.0477-.249s-.806(b_s+b_c)+.097v.
\end{align*}
In particular \(d_s\mid E\), since \(b_s<.060<.117\).
The finite checks below give \(e+b_s+b_c<.192\),
\(e+4(b_s+b_c)+v<.354\). Denominators \(E d_s d_c\) are smaller than
\(q\) by a fixed power, so \(\beta\ne0\). There are \(O(L)\) dyadic
offset magnitudes, including those with \(V<1\), by rational separation.

Write \(w=x\bmod E\) in \([0,E)\), and \(D_E\) for the set of
admitted lower words with unit terminal digit, of size \(N_E\asymp E^c\).
Finite Fourier inversion over the \(E\) translates gives the orbit's
physical source (with either simultaneous choice of conjugations)
\[
 A_{u_0,\beta}\,\frac{E}{N_E}{\bf1}_{D_E}(w)
 e\!\left(-\frac{u_0\lfloor x/E\rfloor}{d_sd_c}\right)
 e(-\beta x)e(\beta w).
\]
Indeed the translate sum selects the actual bottom word modulo \(E\).
The scalar is the normalized digit polynomial above \(E\) at \(E\alpha_0\),
including any forced leading digit phase. Write
\(u_0/(d_sd_c)=u_s/d_s+u_c/d_c\bmod 1\). The unit numerator at
modulus one denotes the trivial additive character.

\paragraph{Two scalar sums.}
Uniformly for the assigned subsets,
\begin{align}
 \sum_{d_c,u_c,u_s,\beta}|A_{u_0,\beta}|
 &\ll d_s^aD_c^{2\delta}(1+V)^a L^C,\label{seven-orb-lone}\\
 \sum_{d_c,u_c}\left(\sum_{u_s,\beta}|A_{u_0,\beta}|^2\right)^{1/2}
 &\ll d_s^{\delta/2}D_c^{2\delta}(1+V)^{\delta/2}L^C .
 \label{seven-orb-ltwo}
\end{align}
These are \eqref{sev:lone} and \eqref{sev:ltwo}: $V$ here denotes the
offset scale $U$ there, and the scalar excludes the lowest $E$-word.
For clarity, sum the terminal block over offsets first, taking local
suprema for the other blocks by Lemma~\ref{b7:digit-sampling}.
A forced top digit costs only a constant by Proposition~\ref{b7:word-moments}.
Next sum the bottom $d_s$-block over $u_s$ at fixed coprime center.
For the square estimate take the square root at this point, before
applying Cauchy--Schwarz to the two separated $D_c^2$-blocks in the
outer sum over $(d_c,u_c)$. This explains why the coprime-center cost
is $D_c^{2\delta}$ in both estimates.

Put \(W=.428(b_s+v)+.16b_c<.026\).
We may omit \(e(\beta w)\), since \(E|\beta|\ll X^{-.16}\) and the
normalized absolute mass of the bottom profile on an interval \(I\)
is \(O(|I|+E)\). Sum over the occupied steps and then apply
\eqref{seven-orb-lone}. Each integer occurs as a product of the fixed
number \(r\) of primes at most \(r!\) times.
Smooth the upper step, before modulation, by convolution in logarithmic
coordinates at width \((R/X)X^{-W-.00004}\). A jump changes an
interval of length \(O(RX^{-W-.00004}L^C)\). The same mass estimate
and \eqref{seven-orb-lone} make the smoothing error power-small
relative to \(X\rho_U\), also in the reference. With
\(V_\beta(x)=\Psi(x/X)U_{\rm sm}(x)e(-\beta x)\), write its Mellin
transform in the coordinate \(\log(x/X)\). Integration by parts
retains shells of exponent
\[
 \tau=\log_XT \le z+\max(\ell,W)+.00032 .
\]
In fact \(T\ll X^{.00004}\max((X/R)X^{W+.00004},X|\beta|L^C)\).
For the second term use \(1-\gamma+\log_XV-z=\lambda-\gamma+\log_XV\),
\(\lambda=m+h+.00003+O(1/L)\) and \(\gamma\ge m-.0001+O(1/L)\).
The normalized pointwise bound for the upper transform on this shell is
\(L^CX^{-B}\), where the normalization is \(\rho_U\) and
\[
 B=.572(\tau-u)_+\quad(\tau\le2u),\qquad
 B=.245\tau\quad(\tau>2u).
\]
These are the upper-step linearization bounds
\((1-a)(\tau-u)_+\) and \((c-a)\tau/2+\delta u\);
\((c-a)/2>.245\).
We also use \(O((1+R|\beta|)^{-1}L^C)\) when
\(T\ll X|\beta|L^{-C}\), by integration within a step.

\paragraph{Character coefficients and their order of summation.}
For the moment let every prime slot be coprime to \(d_c\).
At fixed \(u_c,d_c\), expand the normalized periodic profile, excluding
the modulation, on the units of \(Ed_sd_c\) with coefficients
\(b_{\chi,j}\), \(j=u_s\), using probability counting measure.
Write \(q_s\mid Ed_s\) for the conductor of the seven-part of \(\chi\).
Then
\begin{align}
 \sum_\chi\left|\sum_j b_{\chi,j}f_j\right|^2
    &\ll E^\delta\sum_j|f_j|^2,\label{seven-orb-unitary}\\
 |b_{\chi,j}|&\ll d_c^{-1/2}L^C
       \min\{q_s^{-.245}d_s,\ d_s^{-.92}\}.\label{seven-orb-point}
\end{align}
For the first assertion put \(x=w+Et\). At fixed admissible \(w\),
\(t\bmod d_s\) is unrestricted and independent of the coprime
coordinate. Parseval first in \(\chi\), then in \(t\bmod d_s\),
gives the bound, since the squared normalized bottom profile has
mean \(O(E^\delta)\). For the second, split the coprime phase as
\(e(-u_c\bar E x/d_c)e(u_c\bar E w/d_c)\).
The first factor gives a normalized Gauss sum
\(\ll d_c^{-1/2}L^C\), since an imprimitive Gauss sum at a primitive
additive numerator vanishes unless the quotient by the conductor is
squarefree and coprime to the conductor, and otherwise has magnitude
the square root of the conductor. This follows by imposing
coprimality by inclusion--exclusion. The second factor is a fixed
additive twist of \(w\).

Because \(d_s\mid E\), the kernel \(1+E t\bmod Ed_s\) is
additively parametrized by \(t\bmod d_s\). A character restricts
there to \(e(l t/d_s)\). Averaging in the quotient at fixed \(w\)
selects \(l\bar w=\pm j\bmod d_s\). An admissible \(w\) lies in a
single residue class modulo \(d_s\). Its relative occupancy is
\(O(d_s^{-c})\). This gives the second smooth estimate, including
\(d_s=1\). For the first, fix the quotient \(t\bmod d_s\) and
split \(w\) at \(R_0=7^{\lceil\log_7(q_s)/2\rceil}\le E\).
Put \(Q=q_s/R_0\asymp q_s^{1/2}\). On writing \(w=y+R_0n\), the
primitive smooth character is linear in \(n\), of slope
\(l'\bar y/Q\) with \(l'\) a unit mod \(Q\). This follows from
\(q_s\mid R_0^2\), and \(Q\mid R_0\mid E\).
The translated bottom word has at most \(O((R_0/Q)^c)\) admissible
\(y\)'s in each class modulo \(Q\). The upper word has length
\(E/R_0\). If this is smaller than \(Q\), its first-moment bound
on the \(Q\)-grid costs at most \(Q/(E/R_0)\le d_s\) extra by
splitting into shifted grids. Otherwise retain its first \(Q\)-block.
Thus the average is \(\ll d_s Q^{a-c}\ll d_s q_s^{-.245}\),
uniformly in the fixed additive twist and quotient. Bounded
conductors cause only a constant change.

Fix a shell and a zero deficit \(y=1-\Re\rho\).
Use the same pole-subtracted product estimate as in
Subsection~\ref{seven:small}, with the character now induced from
the two conductor coordinates. Omitting the primes dividing $7d_c$
changes each prime polynomial by $O(L^C)$ and is absorbed in its
$P_i^{1/2}L^C$ baseline. The kernel and the no-zero term are unchanged.

For the pointwise estimate use absolute values in the scalar sum
\eqref{seven-orb-lone}. At fixed \(d_c\) and seven-part conductor
\(q_s\) the characters in the expansion inject into all characters
modulo \(q_sd_c\), preserving the nontrivial zeros, so the all-character, single-modulus estimate used in
Subsection~\ref{bseven:external-density} bounds their total zero count
by $O_\xi((2q_s d_c T)^{(2.4+\xi)y})$. We use $2.4y$ in the numerical
inequalities: the extra exponent is $O(\xi y)$, absorbable even as
$y\downarrow0$, unlike a fixed $O(\xi)$ loss. Only
\(O(L)\) seven-conductor choices are summed. No coprime conductor
divisor multiplicity is required.

For the square estimate first keep \(d_c,u_c\) fixed. The modulation
frequencies at each \(j\) are \(1/q\)-spaced with \(q\mid R\).
On any interval of length \(O(R)\),
\(\int|\sum_\beta A_{j,\beta}e(-\beta x)|^2dx
 \ll R\sum_\beta|A_{j,\beta}|^2\) by the separated large sieve
(or integration against a smooth majorant of width \(R\)).
Mellin Plancherel, summing over occupied steps, consequently gives
\(\int|\sum_\beta A_{j,\beta}\widehat V_\beta(t)|^2dt
 \ll \rho_U L^C\sum_\beta|A_{j,\beta}|^2\).
Logarithmic smoothing preserves this inequality by Jensen and the
same calculation on commonly dilated steps. Apply
\eqref{seven-orb-unitary} at each \(t\), integrate, and use
Cauchy in \((\chi,t)\) against \(K_\chi\).
The second factor costs
\((E d_s d_c T)^{1.2y}L^C\), by the preceding square-integral
bound and the same density estimate, now on all characters modulo
\(Ed_s d_c\). Finally sum the resulting scalar norms over \(u_c,d_c\)
using \eqref{seven-orb-ltwo}. This order retains orthogonality both
in the offsets and in the smooth numerator.

Writing \(g_s=\log_X q_s\), the resulting exponents relative to
\(X\rho_U\), apart from logarithmic factors, are bounded by
\begin{align*}
 P(y)={}&-y+2.4y(g_s+b_c+\tau)-B+W-\tfrac12 b_c
             -\max\{.245g_s-b_s,.92b_s\},\\
 Q(y)={}&-y+1.2y(e+b_s+b_c+\tau)
                    +.04(e+b_s+v+\eta-\lambda)+.16b_c .
\end{align*}
At \(y=1/5\), \(P\) increases with \(g_s\le e+b_s\).
For the low-curvature case maximize
\(.48\tau-.572(\tau-u)_+\) at \(\tau=u\le.230\).
For high curvature maximize \((.48-.245)\tau\) at the shell cutoff.
For \(Q\), use \(\eta-\lambda\le z\) and the same cutoff.
The septenary certificate \path{verify_orbits.py} bounds
all three exponents by \(-.003\). It enumerates vertices of the
nonnegative budget simplex cut by \(s=.117\) and by
\(.245(e+b_s)-b_s=.92b_s\) on each piece of \(e\). On each resulting
polytope the functions tested are convex, since the shell cutoff is a
maximum of affine functions and its coefficients in these tests are
positive. For \(y\ge1/5\), \(Q\) decreases because the certified bound
\(e+b_s+b_c+\tau<.80\) gives slope less than \(-.04\).
For the term with no zeros, Cauchy costs
\((Ed_s d_cT)^{1/2}\) from counting characters and integrating the
constant kernel. Its exponent, also tested, is less than \(-.06\).

At \(y=0\) use the second smooth option and the coprime Gauss sum.
Their savings absorb \(W\) with at least \(.34b_c\) to spare,
apart from the term \(.428v\). If \(V\ge1\) and
\(T\ll X|\beta|L^{-C}\), the nonstationary step bound supplies
\((1+R|\beta|)^{-1}\ll(1+V)^{-1}\). Otherwise \(\tau\ge z-O(\log L/L)\).
Both curvature bounds exceed \(.10-O(\log L/L)\), whereas
\(.428v<.023\). For \(V<1\) there is no offset cost. The best
upper-factor saving is independent of \(y\), hence interpolation
in \(0\le y\le 1/5\) yields \(-c_1y\), with an additional
\(-.17b_c\) when \(y\le .1\), up to logarithmic factors.
On the central interval \(|t|\le X^\epsilon\) use the unattenuated
upper bound unless the nonstationary bound applies. At \(V\ge1\)
it always applies for large \(X\), and \(2.4(.192+2\epsilon)<1\).
Thus the same negative multiples follow directly for central zeros
by choosing \(\epsilon\) sufficiently small.

To implement these bounds, layers of width \(1/L\) cost \(O(L)\)
and change the exponentiated zero weights by bounded factors. For
\(d_c\ge L^A\) the spare \(d_c^{-.17}\) pays any prescribed
logarithmic losses close to one, upon choosing \(A\) accordingly.
For $d_c<L^A$ the radical of the modulus is at most $7L^A$.
The zero-free estimate of Subsection~\ref{bseven:external-density},
including its treatment of a possible real exceptional zero, makes
$X^{-c_1y}$ smaller than every fixed negative power of $L$.

\paragraph{Prime slots meeting the coprime denominator and pole terms.}
For each \(d_c\), choose exactly the slots with prime dividing \(d_c\)
and fix their primes. There are \(O_{\kappa,r}(1)\) choices because
all slot scales exceed \(X^{\kappa+o(1)}\). Let \(a_0\) be the product
of these primes with multiplicity, \(d_0=(a_0,d_c)\),
\(d'_c=d_c/d_0\). The case \(a_0=1\) is allowed. In the residual variable
\(x'=x/a_0\), the coprime additive phase has primitive numerator
modulo \(d'_c\). Exclude primes dividing the original \(d_c\)
from each remaining slot, an \(O(L^C)\) modification to the prime
polynomials. The smooth character expansion is transported by the
permutation \(x'\mapsto a_0x'\) of smooth units, so its square norm
bound persists. In Mellin inversion keep the test in $x=a_0x'$:
the change of log variable translates it by $\log a_0$, so the common
Mellin phase has modulus one. In particular the separated modulation
estimate still uses the original spacing $1/q$ on an $x$-interval
of length $O(R)$. The
smaller product scale costs the relative factor \(a_0^{-(1-y)}L^C\)
in all zero terms, whereas the loss in the pointwise Gauss bound is
at most \(d_0^{1/2}\). Since \(d_0\le a_0\) and \(y\le1/2\),
this does not increase either exponent. Density counts can use
\(q_s d'_c\) and \(Ed_s d'_c\) in the two branches respectively,
with no extra multiplicity at fixed \(d_c,a_0\).
If every slot is fixed the contribution per center is
\(O(E^\delta L^C)\), so \eqref{seven-orb-lone} gives a power
saving relative to \(X\rho_U\).

Only the principal character gives poles. Its smooth coefficient
vanishes if \(d_s>1\), by the complete quotient sum. For \(d_s=1\)
its absolute value is \(O(1/\varphi(d'_c))\), since the coprime
mean of a primitive exponential is \(\mu(d'_c)/\varphi(d'_c)\).
The remaining normalized bottom mean is bounded. Including the
product scale gives
\(a_0^{-1}/\varphi(d'_c)\ll d_c^{-1}L^C\).
Each term with at least one pole has arbitrary power decay outside
the central range from that smooth pole factor. Centrally the
upper transform costs \(O(\rho_U (1+V)^{-1}L^C)\).
After \eqref{seven-orb-lone} the common principal cost is bounded by
\(D_c^{2\delta-1}(1+V)^{a-1}L^C\).
For mixed pole and remainder terms the remainder slots have product
scale at least \(X^{\kappa+o(1)}\). Use the explicit formula on these
slots, the zeta zero-free region and \(2.4(2\epsilon)<\kappa/2\).
This pays arbitrary logarithmic losses even when \(D_c\) is bounded.
When some slot is fixed \(D_c\gg X^{\kappa+o(1)}\), so the same
principal cost makes the all-pole term power-small as well.

It remains to identify the all-pole reference for \(a_0=1\).
It tests the continuous product density against the periodic mean
on units of \(Ed_sd_c\). The desired model instead uses \(7/6\)
times the counting mean on all integers, since the bottom profile already
vanishes at multiples of seven. On each occupied step, partial
summation against the smooth density times \(e(-\beta x)\) replaces
either periodic profile by its corresponding mean at relative cost
\(O(Ed_sd_c(R^{-1}+|\beta|)L^C)\) per orbit. This follows by
subtracting its mean and bounding the primitive of the resulting
periodic sequence by \(O(Ed_sd_c)\) using absolute mass per period.
Sum using \eqref{seven-orb-lone}. Since \(Ed_sd_c\ll X^{.192}\),
\(|\beta|\ll X^{-.36}\), \(R\gg X^{.511}\), this costs a fixed power.
Both means vanish for \(d_s>1\). When \(d_s=d_c=1\), the unit mean
and \(7/6\) times the counting mean coincide. For \(d_s=1,d_c>1\),
the latter mean vanishes by complete coprime-coordinate orthogonality,
and the former is \(O(1/\varphi(d_c))\). For \(d_c\ge L^A\) the
principal cost just proved handles this correction.

For \(1<d_c<L^A\) the orbit scalar has a further factor
\(\exp(-c_2\log X/\log(2d_c))\) as an absolute bound. Indeed the
first fixed positive fraction of digit positions above \(E\) sees
the reduced fraction \(u_c/d_c\) with perturbation initially
\(O(E|\beta|)\) and still power-small throughout these positions.
Among every \(O(\log(2d_c))\) successive multiplications by seven
an unperturbed residue has distance at least \(1/14\) from an integer:
until this occurs the distance grows sevenfold from at least \(1/d_c\).
Our alphabet contains two consecutive digits, so on the resulting
closed subarc away from zero its one-digit amplitude has supremum
strictly less than one, uniformly in the forbidden digit.
The number of pairs \((d_c,u_c)\) in the dyadic range is \(O(D_c^2)\).
Each admits \(O(1+V)\) offsets. The central upper bound cancels the
latter factor and the contraction pays the remaining logarithms.
This proves comparison on all assigned orbits with relative error
\(O_J(L^{-J})\) in units of \(X\rho_U\), including restoration of
the omitted suffix phase and of the original unsmoothed upper step.


\subsubsection{Reconstruction on each prime box}
\begin{proof}[Conclusion of the proof of Proposition~\ref{seven:grouped}]
Proposition~\ref{b7:initial-projection} treats $E_{Y_0}f_X$.
For each subsequent increment use the product $G_q U$ at the single
cut of subsection~\ref{seven:setup}. Proposition~\ref{seven:small-compare}
gives its reference for $q<X^{.192}$.
For $X^{.192}\le q<MX^{-.0001}$, Subsection~\ref{sev:high} gives a
power saving on the prime side. Here the coefficients of the grouped
variables are sums over factorizations into a fixed number of primes,
so they obey the required divisor bound. On the larger periods,
Subsection~\ref{sev:high} gives a power saving after the specified
rational bands have been removed. Subsection~\ref{seven-orbits} supplies
the comparison on all removed orbits.

For the medium periods and the complementary higher frequencies the
reference itself has a power saving. In fact multiplication of an
exact period-$q$ additive mode by $\mathbf 1_{7\nmid n}$ has mean
zero when $q>7$ because the latter function has frequencies of denominator
dividing seven. Every upper step consists of whole $q$-periods,
and the normalized lower component on such a step has mean absolute
value $O(q^{\delta/2})$ by Parseval, including after deletion of any
mask. Replace $\mathcal B(n)$ by its value at the start of each
step at absolute cost $O((R/X)L^C)$ per unit of weight. Summing
over occupied steps, the relative cost in units of
$X\rho_U\rho_L$ is $O(q^{\delta/2}(R/X)L^C)$, which is power-small
since $q\le R\ll X^{.578}$.

The validation of these component comparisons on separated boxes
allows arbitrary logarithmic savings. The prime sum and its reference
are both linear in the digit source. Add the disjoint orbit
contributions and the complement at each $q$, sum over the $O(k)$
periods at each level, restore $(X/Y)^{-\delta}$ and sum the
$O(k)$ levels. The result is the asserted estimate for $f_X$.
\end{proof}

\subsection{Application of the common positive sieve}\label{sieve:seven}
In Theorem~\ref{common:positive-sieve} take $B=7$, $s=6$,
$Q=6^k=X^c$, $u=6-\ind_{d\ne0}$ and $\gamma=7u/36$.
Use $f=f_X$ and $f_0(n)=(\gamma Q/X)\ind_{7\nmid n}$ on $[1,X)$,
both zero outside, and $\mathcal M=\gamma Q/L$. Their full masses
are both $uQ/6$. With fixed $D>3/c$ and $I_X=[XL^{-D},X)$,
Lemma~\ref{b7:digit-occupancy} gives the endpoint mass difference
$O(QL^{-Dc}+Q/X)$; the same occupancy bound controls smoothing
at the right endpoint.

Take $T=.560$, $\alpha=.420$, $\beta=.480$. For every fixed $J,K>0$,
\eqref{b7:type-i} on $I_X$ gives the required $\tau(d')^K$-weighted
sum of errors $O_{J,K}(Q L^{-J})$ for $d'\le X^T$ prime to seven
at residue zero. The discrepancy for $f_0$ is $O(Q/X)$ per modulus
by periodic unit counting. We have $1-\beta=.520<T$ and
$1-T/2=.720<c$, as required for the small-prime cutoff and for
discarding square factors on rough products, respectively.
Proposition~\ref{seven:grouped} supplies the other hypothesis:
its fixed positive roughness and fixed rank can be chosen afresh for
each application; its smooth window $[.4195,.4801]$ contains a
neighborhood of $[\alpha,\beta]$. Its model is exactly
$(7/6)\sum_n f(n)\mathcal B(n)$, since $f$ vanishes on nonunits.
It is this identity, rather than a direct replacement of $f$ by its
mean inside each prime box, that permits the common sieve to apply.

The simplex measure and stopping tests are unchanged: the roughness
exponent is $.060$ and the reflected interval is $[.520,.580]$.
Thus the very same enumeration in Section~\ref{dt:sieve-certificate}
gives $\mathscr L<\ell_*:=933363153278560/10^{15}$. In particular
\[
 \liminf_{k\to\infty}\frac{L}{6^k}\sum_{p\in I_X}f(p)
 \ge \gamma(1-\mathscr L).
\]
For $7^k\le x<7^{k+1}$ these primes are strictly below $x$.
Removing padding zeros is harmless if zero is allowed; if zero is
forbidden the word already has ordinary length $k$. Since $x^c<6Q$
and $\log x\ge L$, we obtain the digit-dependent lower limit
\[
 \liminf_{x\to\infty}\frac{\pi_{7,d}(x)\log x}{x^c}
 \ge \frac{7u}{216}(1-\mathscr L)
 >\frac{35}{216}(1-\ell_*)>\frac1{100}.
\]
Taking the maximum of the seven thresholds and using the upper sieve
following \eqref{b7:type-i} proves Theorem~\ref{seven:main}.


\clearpage

\section{One omitted digit in base six}
\begin{theorem}\label{six:main}
Let $\pi_{6,d}(x)$ count the primes $p<x$ whose ordinary base-six
expansion omits $d$. With $c=\log_6 5$, there are absolute $C,x_0$
such that for every $d\in\{0,1,2,3,4,5\}$ and every real $x\ge x_0$,
\[
 \frac1{600}\frac{x^c}{\log x}\ \le\ \pi_{6,d}(x)
 \ \le\ C\frac{x^c}{\log x}.
\]
\end{theorem}
All lengths in the proof are based on $X=6^k$, $L=\log X$.
We use $e(t)=\exp(2\pi i t)$ and $\delta=1-c$. For integers greater
than three the possible units digits of a prime are 1 and 5. If $u$ of
these survive the exclusion, the complete $k$-position word with a
unit final digit has mass $u5^{k-1}$, and its relative density on the
units modulo six is three times its density among all integers.
The two possible values of $u$ account for the local factor in the sieve.
A complete word treats a leading zero as a digit. We return to ordinary
expansions at the end.

We use the prime-product comparison format of Proposition~\ref{dt:grouped}
and the common positive sieve, Theorem~\ref{common:positive-sieve}.
For base six, the full-grid moments give the upper sieve and the initial projection.
The joint mean on prime-power grids then allows a different cut for
each range of Fourier amplitudes;
Proposition~\ref{six:box} is the comparison needed for the positive
lower sieve. Certificate paths in this part are relative to the base-six
certificate directory. Decimal specifications of parameters in inequalities
and computer checks denote exact rational numbers. The constants
in analytic estimates may be ineffective.

\section{Fourier estimates for six-adic words}
\label{six:sec-six-word-moments}

Put $c=\log_6 5$, $\delta=1-c$, and $D=\{0,\ldots,5\}\setminus\{d\}$.
For $Y=6^k$, write
\[
 f(t)=\frac15\left|\sum_{a\in D}\exp(2\pi iat)\right|,\qquad
 F_Y(t)=\prod_{0\le j<k} f(6^jt).
\]
The normalized Fourier amplitude of a word is the modulus of its
exponential sum divided by its number of words. A restriction at a
position replaces $D$ there by a specified nonempty subset of $D$.
The constants in the following bounds are uniform in $d$ and in the
positions of any fixed number of restrictions.

\paragraph{Shifted moments.}
In the order $s=1,4/3,3/2,7/4,2$, the following exponents give
uniform bounds for all $k\ge0$ and $u\in\mathbb R$:
\begin{align}
 \sum_{0\le r<6^k} F_{6^k}((r+u)/6^k)^s
     &\ll 6^{k a_s},\label{six:eq-six-full-moments}\\
 \sum_{0\le r<p^k} F_{6^k}((r+u)/p^k)^s
     &\ll p^{k b_{p,s}} \qquad(p=2,3).\label{six:eq-six-pure-moments}
\end{align}
\[
\begin{array}{c|rrrrr}
 &1&4/3&3/2&7/4&2\\ \hline
 a_s\ (d=0,5)&.38197&.25813&.20833&.14755&\delta\\
 a_s\ (d=1,4)&.42608&.28468&.22604&.15440&\delta\\
 a_s\ (d=2,3)&.43389&.28965&.22947&.15580&\delta\\ \hline
 b_{2,s}\ (\text{all }d)&.108&.052&.038&.025&.00581\\
 b_{3,s}\ (\text{all }d)&.207&.096&.067&.040&.01105
\end{array}
\]
The unrestricted second moment in \eqref{six:eq-six-full-moments} is
exactly $6^{k\delta}$ by orthogonality, since distinct words give
distinct residues modulo $6^k$. The shift gives only phases of modulus
one, which cancel on the diagonal. Interpolation with H\"older's
inequality allows linear interpolation between successive exponent
entries as a function of $s$.

\paragraph{Finite certificate for the full grid.}
Define the positive operator
\[
 (T_sh)(u)=\sum_{r=0}^5 f((u+r)/6)^s h((u+r)/6).
\]
The full shifted sum equals $T_s^k1(u)$ by induction. The program
\texttt{cert/verify\_moments.py} constructs, for each $d=0,1,2$
and each of the first four moments, a continuous periodic positive
piecewise affine function $h\le1$ and checks
$T_sh<\lambda h$ for a rational $\lambda<6^{a_s}$.
Positive iteration and comparison with $h/\min h$ prove the bound.
The reflection $a\mapsto5-a$ preserves all amplitudes.

For the $N=1800$ grid apply Lemma~\ref{prelim:operator-cert} to the six
branches $z_r(u)=5^{-1}\sum_{a\in D}e((a-5/2)(u+r)/6)$, where
$e(t)=\exp(2\pi it)$. The sums of $|a-5/2|$ and $(a-5/2)^2$ are at
most $17/2$ and $69/4$, so $|z_r'|<2$ and $|z_r''|<4$.
The branch magnitudes cycle under $u\mapsto u+1$ and $h$ is periodic.
The required node margin is therefore $s(24+4\operatorname{Lip}h)/(8N^2)$.

At each node the single-digit evaluations use
\[
 25 f(t)^2
 =5+2\sum_{j=1}^5\#\{a\in D:a+j\in D\}\cos(2\pi jt).
\]
All cosine coefficients are nonnegative. Machin's formula for $\pi$
and alternating arctangent series give rational endpoint bounds for
$\pi$. After reflection to $[0,\pi]$, the degree-$41$ Taylor
polynomial for cosine has absolute remainder at most
$(22/7)^{42}/42!$. Interval calculations round outward. Upper square
weights are rounded to units of $2^{-32}$ and raised to $s/2$ by
integer ceiling roots. Iterations on the finite grid select integer
heights in units of $2^{-40}$. Only the positivity of this vector and
its final residual enter the proof. After allowing for the continuous
interpolation error, every minimum residual exceeds $10^{-5}$.
The comparisons $\lambda<6^{a_s}$ are checked by integer powers.

\paragraph{Pure-prime and mixed grids.}
Let $B_p(k)$ be the supremum over $u$ of the left side of
\eqref{six:eq-six-pure-moments}, with $s$ fixed. Then
$B_p(k+l)\le B_p(k)B_p(l)$. Indeed write $r=h+p^l b$ with
$0\le h<p^l$, $0\le b<p^k$. The first $k$ digits contribute a
shifted $p^k$ sum over $b$. The other digits begin at angle
$(6/p)^k(h+u)/p^l$ modulo one. Multiplication by $(6/p)^k$
permutes the $h$ grid since it is a unit modulo $p$. This proves
submultiplicativity. Thus any certified inequality
$B_p(K)<p^{K b_{p,s}}$ at a fixed $K$ gives the stated bound upon
partitioning into blocks of length $K$ and a bounded remainder.

For the first four pure moments, use $K=7$ for $p=2$ and $K=6$ for
$p=3$. The program \texttt{cert/verify\_pure\_refined.py} computes an
upper table for $f(i/L)^s$, $0\le i<L=65536$, by the preceding
rational trigonometric bounds. The centered single-digit polynomial
has derivative of modulus at most $2\pi(17/2)/5<11$, so $f^s$ has
Lipschitz constant at most $11s$. Set $m=p^K$. On a shift interval
$[(h-1)/n,(h+1)/n]$, for each branch $a\bmod m$ and digit $j<K$,
round $6^j(a+h/n)/m$ on the circle to the nearest node $i/L$.
Throughout that interval the digit contributes at most
\[
 \min\left(1,f(i/L)^s+11s\left(\frac{6^j}{mn}+\frac1{2L}\right)\right).
\]
The program multiplies these upper weights along each branch and sums
over $a$, rounding every operation upward. The result bounds the
whole interval. Shifts have period one, and reflection makes the sum
even, so the initial interval is $[0,1/2]$. An unresolved interval
is bisected. Every final cell is below the fixed-point lower
approximation to $p^{K b_{p,s}}$ in units of $2^{-32}$, with equality
against that lower approximation harmless. The resulting bound is
strict against $p^{K b_{p,s}}$, which is irrational for every one of
these eight entries.

For the second pure moment use $K=14$ for $p=2$ and $K=11$ for
$p=3$. Let $n(r)$ count the $K$-digit words in the residue class
$r\bmod p^K$. Addition of one digit corresponds exactly to adding
the five translates of the preceding integer residue counts.
Orthogonality and the triangle inequality show that
\[
 B_p(K)=p^K\sum_{r\bmod p^K} n(r)^2/25^K,
\]
where the maximum is attained at shift zero. The integer program
\texttt{cert/verify\_mixed.py} compares this rational number with
$p^{K b_{p,2}}$ by integer powers. The same program includes a
four-digit calculation supplying the weaker first four pure moment
bounds $(.151,.081,.056,.036)$ for $p=2$ and
$(.231,.119,.085,.051)$ for $p=3$. For this additional check the
continuum guard on an $N$-interval mesh is
$137s\sum_{j=0}^3 36^j/(8p^4 N^2)$. In fact, centering four
independent digits about their means gives a sum of variance at
most $(69/20)\sum_{j=0}^3 36^j$. At denominator $p^4$ its centered
Fourier polynomial has second derivative bounded by
$137\sum_{j=0}^3 36^j/p^8$ in the shift coordinate.
Summing the lower second-derivative bound after taking $s$-powers
gives the displayed interpolation guard.

For a general $q\mid 6^k$, write $q=6^v p^j$ with $v+j\le k$ and
$p=2$ or $3$. Drop all factors above the first $v+j$ digits.
Writing the residue index as $h+p^j b$, first sum over $b\bmod 6^v$
by \eqref{six:eq-six-full-moments}; the last $j$ digits see only
$(h+u)/p^j$. Hence
\begin{equation}
 \sup_u\sum_{i=0}^{q-1} F_{6^k}((i+u)/q)^s
    \ll 6^{v a_s}p^{j b_{p,s}}.\label{six:eq-six-mixed-moments}
\end{equation}
In each block decomposition the following change permits finitely
many restricted positions. Bound the normalized amplitude at such a
position by one, split into runs of unchanged positions, and use the
uniform shifted bound on each run. Each removed digit has at most
six branches and there are at most one more unchanged runs than
removed digits. The implied constant may therefore depend on the
number of restrictions.

\paragraph{Positive floors.}
For every positive $\epsilon$, in the full-grid square column replace
$\delta$ by $\delta+\epsilon$. All the other exponents in the first
table may be retained without change after replacing $f$ at every
unrestricted position by $f_\zeta=\max(f,\zeta)$ for all sufficiently
small fixed $\zeta>0$. Indeed the finitely many strict transfer
inequalities on the circle persist by uniform convergence of the
one-digit amplitudes. At $s=2$ one has $T_21=6^\delta$, so the
floored operator satisfies
$T_{2,\zeta}1\le 6^\delta+6\zeta^2<6^{\delta+\epsilon}$.
Each fixed-length strict pure-prime block inequality likewise
persists, since its sum converges uniformly in the shift. The same
submultiplicative, mixed-grid and finite-restriction arguments apply.
These floor bounds concern absolute amplitudes; square orthogonality
itself is used only for the original digital polynomial.

\paragraph{Separated sampling.}
For $Y=6^k$ and every $s$ in the table, one has
$\int_0^1 F_Y(t)^s\,dt\ll Y^{a_s-1}$ by integrating
\eqref{six:eq-six-full-moments} in $u$. If the circular distances in
$\Theta\subset\mathbb R/\mathbb Z$ are at least $\Delta>0$, then
for each fixed $C\ge0$,
\begin{equation}
 \sum_{\theta\in\Theta}\sup_{\|v\|\le C/Y} F_Y(\theta+v)^s
 \ll_{s,C}(Y+\Delta^{-1})Y^{a_s-1}.\label{six:eq-six-separated-sampling}
\end{equation}
For actual word polynomials, including restricted ones with the
corresponding integral bounds, this is \eqref{arith:sampling}.

For floored products, $\log f_\zeta$ is Lipschitz with constant at
most $11/\zeta$. The logarithm of the product varies by $O_{C,\zeta}(1)$
over any interval of length $(1+C)/Y$ by the geometric sum of the
frequencies. Bound the local supremum by $O_{C,s,\zeta}(Y)$ times
the integral of the $s$-power over an interval of radius $1/Y$
about each center. Such intervals overlap $O(1+(Y\Delta)^{-1})$
times. This proves the analogous floored estimate with its stated
exponents. Restricted digit factors can be bounded by one in this
argument. For the mixed bound \eqref{six:eq-six-mixed-moments} with
positive floors, the same logarithmic variation permits a separate
perturbation $O(6^{-(v+j)})$ at every grid point: discard all factors
above the leading $v+j$-digit block before applying the bound.
Unfloored factors are dominated by these floored factors. This version
retains the original full-grid entries except for the arbitrarily
small exponent allowance in the square column.

In particular, for integers $1\le H\le Q\le Y$ and any $H$
consecutive indices $j$, the sum of the individual local suprema of the $s$th powers
at $(j+u)/Q$, with displacement at most $C/Q$, is $O(H^{a_s})$.
This is the consecutive-block argument of Lemma~\ref{common:sampling}:
use \eqref{six:eq-six-separated-sampling} on the retained block,
with the preceding local bound in place of polynomial reproduction
for the floored products.

\paragraph{Physical occupancy.}
The interval bound $O((1+H)^c)$ and the single-class bound
$O((Y/Q)^c)$ follow from Lemma~\ref{common:occupancy}. For $q=6^v p^j\mid Y$
there is also the bound
\[
 5^{-k}\max_{r\bmod q}\#\{w\text{ an admissible word}:w\equiv r\bmod q\}
 \ll 5^{-v}p^{(-1+e_p)j}.
\]
Here $e_2=.0348$ and $e_3=.0518$.
The full $6^v$ residue fixes the lowest $v$ digits. On the remaining
$j$-digit pure portion the maximum normalized residue occupancy
$p^j\max_r\Pr(w\equiv r\bmod p^j)$ is submultiplicative by
conditioning on a low block and using that $6/p$ is a unit modulo $p$.
At lengths $K=14,11$ as above, \texttt{cert/verify\_mixed.py}
checks $p^K\max_r n(r)/5^K<p^{Ke_p}$. Additional high positions
give a convex combination of translates of these residue laws.
Finite digit restrictions cost at most $5$ per position by conditioning.

\section{Distribution of divisors and the prime upper bound}
Let $\mathcal A_{k,d}$ consist of the ordinary $k$-digit allowed numbers
coprime to six, $k\ge2$. Its size $M=h_0 u5^{k-2}$ has
$h_0=5-\mathbf1_{d\ne0}$ allowed nonzero first digits and
$u=2-\mathbf1_{d\in\{1,5\}}$ allowed unit last digits. Take
$\sigma=1-a_1(d)\ge1-.43389>.56$, where $a_1(d)$ is the full-grid
first-moment exponent in the table for the omitted digit. The integrated first moment from
Section~\ref{six:sec-six-word-moments} is $O(Y^{-\sigma})$. We prove
for every $J,K>0$ and each fixed $0<\theta<\sigma$ that
\begin{equation}\label{six:interval-type-i}
 \sum_{\substack{q\le X^\theta\\(q,6)=1}}\tau(q)^K
 \sup_{I,z}\left|\sum_{n\in \mathcal A_{k,d}\cap I}
       (\mathbf1_{n\equiv z\bmod q}-1/q)\right|\ll_{J,K,\theta}M L^{-J}.
\end{equation}
Here $I\subset[0,X)$ is any interval and $z$ any residue. Apply
Proposition~\ref{arith:typeI} with $b=6$, $s=5$. Deleting one digit
leaves two consecutive digits, and the shifted first moment gives
$\int_0^1 F_Y\ll Y^{-\sigma}$ at every radix length, with
$.56611\le\sigma\le.61803<c=\log_6 5$. Both endpoint sets are
nonempty and $M\asymp5^k$. In \eqref{arith:typeI-max} choose
$\epsilon>0$ so that $\theta(1+\epsilon)<\sigma$; this proves
\eqref{six:interval-type-i}, also with the full allowed leading set
for padded words. Convex averages of cyclic translates have the same
bound with their own interval mass as mean.

For the upper bound take the complete word with unit last digit in
Proposition~\ref{common:upper}. Its divisor hypothesis is
\eqref{six:interval-type-i} with $z=0$, $K=2$ and any fixed
$0<T<\sigma<\min(1,2c)$. The unit density relative to the integers
is $(u/5)/(2/6)=3u/5$. Letting $T\uparrow\sigma$ gives the upper
constant $6u/(5\sigma)$ for both ordinary and complete-word counts.

\section{A quantitative two-almost-prime consequence of divisor distribution}
Let $X=6^k$, $A=5^k$, $L=\log X$, $c=\log_6 5$, and let
$f$ indicate the permitted $k$-place words with a unit last digit.
Write $j=2-\mathbf1_{d=1\ {\rm or}\ d=5}$ and $\gamma=3j/5$.
Then $\sum f=j A/5=\gamma A/3$. All integers in the support are
positive and prime to six. Set
\[
  (\theta,u,v)=(14/25,7/50,49/100),\qquad
  (D,z,y)=(X^\theta,X^u,X^v).
\]
The interval estimate \eqref{six:interval-type-i} applies at $\theta=.56$.
\begin{theorem}[Rough primes and semiprimes]\label{six:rough}
With $P^-(1)=+\infty$ and $\lambda=(3-1/v)^{-1}=49/47$,
put $w(p)=1-\lambda(1-\log_X(p)/v)_+$.
For primes $p,q$ one has
\[
 \sum_{z\le p<X} f(p)+\sum_{z\le p<q,\ pq<X}w(p) f(pq)
 \ge (C_2+o(1))\,\gamma A/L,
\]
where
\[
 C_2=\frac{25}7\left(\log3-\frac{49}{47}
        \left(\log(7/2)-\tfrac17\log6\right)\right)>.212 .
\]
In particular this many words are squarefree with $P^-(n)\ge X^{7/50}$
and at most two prime factors, counted with multiplicity.
\end{theorem}
\begin{proof}
For $p\ge z$, $0<1-\lambda(1-u/v)\le w(p)\le1$.
Put $S_e(z)=\sum_{e\mid n,\,P^-(n)\ge z} f(n)$
and $W=S_1(z)-\lambda\sum_{z\le p<y}(1-\log_X(p)/v)S_p(z)$.
For a squarefree rough integer with $r\ge3$ factors the subtracted
sum of weights before multiplication by $\lambda$ is at least
$r-\log_X(n)/v\ge3-1/v$. Its coefficient in $W$ is therefore
nonpositive. On primes the coefficient is at most one; on semiprimes
it is at most $w(p)$ for the smaller factor $p$ since the contribution of the
larger factor is nonnegative. Rough integers with a repeated prime
factor have total count at most
$X\sum_{h\ge z}h^{-2}\ll X^{1-u}=o(A/L)$; their individual weights
are bounded by a constant. The contribution of $1$ is at most one.

For the lower linear sieve on $S_1(z)$ take level $D$, with
density $1/p$ at primes $p\ge5$ and density zero at $2,3$. The
Mertens product at $z$ equals $(3e^{-\gamma_E}+o(1))/(uL)$,
where $\gamma_E$ denotes Euler's constant. The dimension-one lower
sieve function satisfies $s f_1(s)=2e^{\gamma_E}\log(s-1)$ for
$2\le s\le4$, with $s=\theta/u=4$ here. For
$S_p(z)$ take the upper linear sieve on the cofactor with level
$D/p$ and sieving limit $z_p=(D/p)^{1/3}\le z$. Its upper function
at $s=3$ has value $2e^{\gamma_E}/3$. The inequality obtained is
in the required direction because every original rough cofactor is
also $z_p$-rough. Both $z_p$ and $D/p$ tend to infinity at a fixed
power of $X$ uniformly. Since $p\ge z\ge z_p$, the local densities
in this upper sieve are unchanged and its main mass is
$\gamma A/(3p)$. The error terms in both sieves are negligible:
the remainder modulus in the upper sieve is $pe\le D$ and the
number of possible markings of each such modulus is at most its
divisor count. The dimension-one sieve functions used here are the
classical Rosser--Iwaniec lower and upper functions.

Partial summation over the primes gives
\[
 \frac{LW}{\gamma A}\ge \frac2\theta\log(\theta/u-1)
 -2\lambda\int_u^v\frac{1-t/v}{t(\theta-t)}\,dt+o(1)
 =\frac{25}{7}
 \left(\log3-\frac{49}{47}
    \left(\log(7/2)-\tfrac17\log6\right)\right)+o(1).
\]
The explicit constant exceeds $.212$. For a rational check, for
each of $a=3,7/2,6$ set $t=(a-1)/(a+1)$ and bound $\log a$ between
$S=2\sum_{i=0}^{29}t^{2i+1}/(2i+1)$ and
$S+2t^{61}/(61(1-t^2))$. Substitute the two lower endpoints for
the positive logarithms and the upper endpoint for $\log(7/2)$.
This calculation is also verified by \path{check_weights.py}.
\end{proof}

The same proof applies to ordinary words of exactly $k$ digits on replacing
$\gamma A$ by $3M$: only the mass and its divisor estimate enter the two
sieves. For $6^k\le x<6^{k+1}$ the complete-word conclusion supplies
$\gg x^c/\log x$ ordinary admissible squarefree integers below $x$ with
at most two prime factors and least prime factor at least $(x/6)^{7/50}$.
This does not give a lower bound for semiprimes alone.

\section{Prime products with the lowest 19/25 of the positions tested}
\label{six:sec-six-projection}
Set \(X=6^k,\ L=\log X,\ c=\log_6 5,\ \delta=1-c,\ \eta=19/25\)
and \(Y=6^{\lfloor\eta k\rfloor}\). For a radix length \(B=6^j\ge6\)
let \(f_B\) be the periodic indicator of the words in \([0,B)\)
with all \(j\) digits different from \(d\) and last digit in \(\{1,5\}\).
Digits at the left end are read as part of the word, including zeros.
Write \(A_B=\sum_{0\le n<B}f_B(n)=uB^c/5\), where
\(u=\#(\{1,5\}\setminus\{d\})\). For an \(X\)-periodic weight \(G\)
let \(E_YG\) denote its average on each residue class modulo \(Y\).

\begin{theorem}\label{six:thm-six-low-projection}
Fix \(r\ge1,\ \kappa>0\), and scales \(P_i\ge X^\kappa\)
with \(\prod_{i=1}^rP_i=XL^{O(1)}\). Let \(W(\boldsymbol t)\) be
smooth and compactly supported on a box \(t_i/P_i\in I_i\), with
fixed compact intervals \(I_i\subset(0,\infty)\). Each fixed mixed
derivative in the variables \(\log(t_i/P_i)\), including order zero,
is \(O(L^{C})\), where \(C\) may depend on the order. Then for every
fixed \(J>0\),
\begin{equation}\label{six:eq-six-low-principal}
 \sum_{p_1,\ldots,p_r} W(\boldsymbol p)f_Y(\textstyle\prod_i p_i)
 =\frac{A_Y}{\varphi(Y)}
  \left(\int W(\boldsymbol t)\prod_i\frac{dt_i}{\log t_i}
             +O_J(XL^{-J})\right).
\end{equation}
Let \(\mathcal B(x)\,dx\) be the pushforward of
\(W(\boldsymbol t)\prod_i dt_i/\log t_i\) by multiplication.
For \(G=f_X\), and also for the periodic indicator \(G\) of
\(\mathcal A_{k,d}\), one has
\begin{equation}\label{six:eq-six-low-model}
 \sum_{p_1,\ldots,p_r} W(\boldsymbol p)(E_YG)(\textstyle\prod_i p_i)
 =3\sum_{\substack{n\ge1\\(n,6)=1}}(E_YG)(n)\mathcal B(n)
       +O_J(X^cL^{-J}).
\end{equation}
The bounds are uniform in \(d\).
\end{theorem}

\subsection{Fractional moments and character coefficients}
In the notation of the Fourier section, put
\(s_3=3/2,\ s_4=4/3,\ a_3=.22947,\ a_4=.28965\).
For every \(B=6^j\) and \(h=3,4\) the certificate gives
\begin{equation}\label{six:eq-six-low-sampling}
 \sup_{x\in\mathbb R}\sum_{0\le b<B}F_B((b+x)/B)^{s_h}
 \ll B^{a_h}.
\end{equation}
For $1\le v\le B$ let $Z$ be the largest power of six not exceeding
$v$. Each residue modulo $v$ contains $O((B/v)^c)$ words: once the
digits above $Z$ are fixed, the low block has at most one lift.
Dropping the factors above $Z$ in $F_B$ and applying
\eqref{six:eq-six-separated-sampling} at spacing $1/v$ likewise gives
$\sup_x\sum_{b<v,\ b\ge0} F_B((b+x)/v)^{s_h}\ll v^{a_h}$, since $Z\asymp v$.

For \(q\mid Y\) and primitive \(\chi\bmod q\) define
\[
 b_\chi=A_Y^{-1}\sum_{0\le n<Y}f_Y(n)\overline{\chi(n)} .
\]
These are the coefficients of the induced characters modulo \(Y\)
with the same normalization, since \(f_Y\) is supported on units.
Conjugating all characters will have no effect on the estimates.
Write \(q=2^i3^j,\ t=\max(i,j)\), and absorb bounded \(q\) in
implicit constants.

\paragraph{Cubic bound.}
Put $R=6^{\lceil t/2\rceil}$, $D=(q,R)$ and $Q=q/D$.
Prime by prime, $Q\mid D$, $(R/D,Q)=1$ and $RQ\mid Y$.
Thus the half-cut argument for \eqref{dt:analytic-cubic} applies:
$v\mapsto1+Rv$ identifies the reduction kernel with the additive
group modulo $Q$, and primitivity forces its slope $\ell$ to be a
unit (test the subgroup of order $p$ for each $p\mid Q$).
For a unit lower word $z$, the upper amplitude is
$U_\ell(z)=F_{Y/R}(\ell\bar z/Q)$. With $M_\ell=\sum_z U_\ell(z)^{3/2}$,
Hausdorff--Young on the units modulo $D$ and the lift multiplicity
$O((R/D)^c)$ bound the cube sum in this slope by
$D^{1-c}R^{-2c}M_\ell^2$. Sampling modulo $Q$ gives
$\max_\ell M_\ell\ll(R/Q)^c Q^{a_3}$; sampling first in the unit
slope gives $\sum_\ell^*M_\ell\ll R^c Q^{a_3}$. Their product yields

\begin{equation}\label{six:eq-six-low-cube}
 \sum_{\chi\bmod q}^{*}|b_\chi|^3\ll q^\delta Q^{2a_3-1},
 \qquad Q\gg\max(q^{.278},qY^{-1/2}).
\end{equation}
In fact, if \(q=p^t(p')^v,\ v\le t\), before rounding the cut the
ratio \(\log Q/\log q\) is minimized at \(v=t/2\), as is seen by
dividing its numerator and denominator by \(t\) in each of the two
ranges \(v\le t/2,\ v\ge t/2\). Its minimum is at least
\(\log2/\log12>.278\). Rounding costs a bounded factor.
The other bound follows from \(D\le R\ll Y^{1/2}\).

\paragraph{Quartic bound.}
Now take \(R=6^{\max(2,\lceil t/3\rceil)},\ D=(q,R),\
Q=q/D,\ K=(R,Q),\ W_0=Q/K\).
Comparison of valuations gives \(Q\mid R^2,\ W_0\mid K\mid R,\
(R/K,W_0)=1\), and \(R^2W_0\mid Y\) for large \(k\).
In the same reduction kernel the parameter has law
\(v\circ v'=v+v'+Rvv'\). The map
\(\mathcal L(v)=v-(R/2)v^2\bmod Q\) identifies it with the additive
group modulo \(Q\). Expansion proves additivity using \(Q\mid R^2\);
the factor \(1-(R/2)(v+v')\) in a difference is a unit modulo \(Q\).
The slopes of primitive characters are again units.
Writing \(n=z+Rm+R^2v\), use the character on
\(1+R\bar z(m+Rv)\). Its upper factor has magnitude
\[
 U_\ell(z)=F_{Y/R^2}\big((R/K)\ell\bar z/W_0\big).
\]
Up to conjugation the normalized middle sum is
\[
 B_\ell(z)=R^{-c}\sum_{m\ {\rm allowed},\,0\le m<R}
 e\bigl((\ell\bar z m-(R/2)\ell\bar z^2 m^2)/Q\bigr).
\]
All cross terms with \(v\) in the quadratic part vanish modulo \(Q\).

For fixed \(a\bmod W_0\) coprime to \(W_0\),
\begin{equation}\label{six:eq-six-low-four-fibre}
 \sum_{z\bmod R}^{*}\sum_{b\bmod K}|B_{z(a+W_0 b)}(z)|^4
       \ll RK^{1-c}W_0^{-c}(1+\log Y)^2.
\end{equation}
Use the fourth-power expansion in the proof of \eqref{dt:analytic-quartic},
with $W=W_0$. Write $S=m_1+m_2-m_3-m_4$ and
$T=m_1^2+m_2^2-m_3^2-m_4^2$. Orthogonality in $b$ imposes
$K\mid S-(R/2)\bar zT$, equivalently $K\mid S$: if $K\nmid R/2$,
$K$ and $R/2$ are even, so either congruence forces $S$ and hence $T$
even. The inverse-unit sum modulo $W_0$ then has Ramanujan argument
$-a(R/K)T/2$, an integer. Its support requires
$V=W_0/\operatorname{rad}(W_0)\mid T$, since $(aR/K,W_0)=1$.
As $W_0\mid S$, one has $T\equiv-2(m_3-m_1)(m_3-m_2)\pmod{W_0}$.
Thus $W_1=V/(V,2)$ divides the product of differences and
$W_0/W_1\le12$. Splitting $W_1$ between them, residue occupancy
gives $O(R^{3c}W_0^{-c}\tau(W_1))$ triples and $O((R/K)^c)$
choices of $m_4$. Multiplication by $RK/R^{4c}$ proves
\eqref{six:eq-six-low-four-fibre}, since $W_1$ is supported on $2,3$.

The slope summation for \eqref{dt:analytic-quartic} now uses $s=4/3$,
with $a_4=.28965$ in place of $.285$. In its H\"older bound the two
upper-word sums each cost $(R/W_0)^c W_0^{a_4}$; grouping the remaining
fourth powers by $a=\ell\bar z\bmod W_0$ costs the fibre bound times
$W_0^{a_4}$. Multiplying by the Hausdorff--Young factor
$D^{1-c}R^{-3c}$ and using $DK=q/W_0$ gives

\begin{equation}\label{six:eq-six-low-quart}
 \sum_{\chi\bmod q}^{*}|b_\chi|^4
 \ll(qR)^\delta W_0^{3(a_4-c)-\delta}(1+\log Y)^2,\qquad
 R\ll Y^{1/3},\quad W_0\gg qY^{-2/3}.
\end{equation}

\subsection{Summation over zero layers}
First take \(W=\prod_i w_i(t_i/P_i)\). Expand \(f_Y\) on the unit
group modulo \(Y\), factoring out \(A_Y/\varphi(Y)\). The principal
coefficient is one. Each induced nonprincipal coefficient is
\(b_\chi\) for a unique primitive character modulo some \(q\mid Y\).
For fixed \(0<y<1/2\) the number of zeros of real part at least
\(1-y\), height at most \(T\), summed over these primitive
characters at conductor \(q\), is
\[
 \ll_\epsilon(qT)^{\nu(y)+\epsilon},\qquad
 \nu(y)=3y/(1+y).
\]
This is the classical fixed-modulus density bound, as stated in
\cite[(1.4)]{ChenGuptaLi2026}; see also \cite[(2.5)]{FortiViola1973}.

For \(q>L^B\) choose \(T=q^\xi\), with \(\xi>0\) fixed and small.
Let \(1-y_\chi\) be the maximum of \(1/2\) and the real parts of the
zeros of \(\chi\) up to height \(T\), with value \(1/2\) if there is
no such zero. The smooth explicit formula gives, for every fixed \(H>2\),
\[
 \left|\sum_p w_i(p/P_i)\chi(p)\right|
 \ll_H L^{C_H}(P_i^{1-y_\chi}+P_i T^{1-H}).
\]
Indeed, apply the von Mangoldt explicit formula with the test divided
by \(\log t\). Mellin integration by parts gives each zero
\(\beta+iv\) weight \(O(P_i^\beta L^{C_H}(1+|v|)^{-H})\); the
number of zeros in each unit height interval is \(O(\log(q(2+|v|)))\).
The sum above the cutoff is \(O(P_i L^{C_H}T^{1-H})\) after enlarging
\(C_H\). The omitted prime powers cost \(O(P_i^{1/2}L^{C_H})\),
and the trivial zeros and conductor term fit this bound.
Take \(H>1+4/\xi\). Products with a tail term cost
\(O(XL^Cq^{-4})\) per character. The other product is
\(O(XL^C X^{-y_\chi})\).

Put \(\gamma=(\log q)/L\). Using H\"older and
\eqref{six:eq-six-low-cube}, an ideal layer at \(y\) has exponent in
powers of \(X\), relative to \(X\), at most
\begin{equation}\label{six:eq-six-low-expcube}
 -y+\tfrac23\gamma\nu(y)+
 \tfrac13\bigl(.102\gamma-(1-2a_3)
       \max(.278\gamma,\gamma-\eta/2)\bigr).
\end{equation}
For \(\gamma>2\eta/3\), \eqref{six:eq-six-low-quart} alternatively gives
\begin{equation}\label{six:eq-six-low-expquart}
 -y+\tfrac34\gamma\nu(y)+
 \tfrac14\bigl(.102(\gamma+\eta/3)-1.925(\gamma-2\eta/3)\bigr).
\end{equation}
These expressions omit mesh, density and height losses, which we
now control. The constants use \(c>.898\) and
\(3(c-a_4)+\delta=1+2c-3a_4>1.925\).
For \(\gamma\le.5\) use \(X^{-y}\le q^{-2y}\) before grouping;
\(-2y+(2/3)\nu(y)\le0\), and the cubic coefficient exponent in
powers of \(q\) is less than \(-.016\). For the remaining ranges,
the \(y\)-dependent term \(-y+Ay/(1+y)\) is bounded by
\((\sqrt A-1)^2\). Combining this with the affine terms yields
convex functions of \(\gamma\). The following bounds hold at both
ends of their indicated intervals, so also throughout:
\[
\begin{array}{c|c|c}
 \gamma &\text{bound used}&\text{upper bound}\\ \hline
 {}[.5,.52]&(\sqrt{2\gamma}-1)^2-.016\gamma&-.0075\\
 {}[.52,.67]&(\sqrt{2\gamma}-1)^2+
   (.102\gamma-(1-2a_3)(\gamma-\eta/2))/3&-.0045\\
 {}[.67,.76]&(\sqrt{9\gamma/4}-1)^2+
   (.102(\gamma+\eta/3)-1.925(\gamma-2\eta/3))/4&-.0014
\end{array}
\]
The integer comparisons and rational squarings are given in
\texttt{bilinear/verify\_projection.py}. For example,
\((\sqrt A-1)^2+z<C\) follows by verifying \(1+A+z-C>0\)
and \((1+A+z-C)^2<4A\).
Use a fixed mesh on \([0,1/2]\), the upper endpoint for the density
and the lower for the weight. On the last cell use the trivial
character count \(q\), corresponding to \(\nu(1/2)=1\), and on
the first cell use its positive upper endpoint. Since \(\nu\) has
Lipschitz constant three, the change in exponent is
\(O(\text{mesh width}+\xi+\epsilon)\) in powers of \(q\) for
\(\gamma\le.5\), and in powers of \(X\) for \(\gamma\ge.5\).
Choose the mesh width, \(\xi\) and \(\epsilon\) sufficiently small
in this order. This gives \(O(XL^Cq^{-\epsilon_0})\) at every
conductor \(q>L^B\) for a fixed \(\epsilon_0>0\).
There are \(O(L^2)\) conductors. Taking \(B\) large gives the
required logarithmic saving. For \(q\le L^B\), Siegel--Walfisz and
partial summation give arbitrary powers of \(L^{-1}\) for each
normalized nonprincipal prime sum. There are \(O(L^{B+2})\) such
characters. For large \(X\) all prime variables are units modulo
six. The prime number theorem and partial summation for the principal
character finish \eqref{six:eq-six-low-principal} for products of tests.

For a joint test expand in a Fourier series in the logarithms on a
fixed larger box, then reinsert separate smooth cutoffs equal to one on
its support. Integration by parts gives coefficients
\(O_H(L^{C_H}(1+\|\boldsymbol m\|)^{-H})\) at every fixed \(H\).
Truncation at \(L^{C'}\) in each coordinate, with \(C'\) sufficiently
large, gives a supremum error \(O(L^{-A})\) for any chosen \(A\).
Each integer has at most \(r!\) ordered representations as a product
of \(r\) primes. Complete periods therefore bound the counting cost
of the discarded test by \(O(L^{-A}(1+XL^C/Y)A_Y)\); the principal
integral after normalization has the same bound. The retained modes
have logarithmic derivative growth and their absolute coefficient
sum costs at most a fixed power of \(L\). This proves
\eqref{six:eq-six-low-principal} for the stated boxes.

Finally \(E_Y f_X=(X/Y)^{-\delta} f_Y\). For the periodic ordinary
word \(G=\mathbf1_{\mathcal A_{k,d}}\) there is the extra factor
\(h/5\), where \(h\) counts its admissible nonzero first digits.
Write \(E_YG=\alpha f_Y\) in either case. Since
\(\varphi(Y)=Y/3\), the principal term after multiplication by
\(\alpha\) is \(3\alpha(A_Y/Y)\int\mathcal B(x)\,dx\).
Solving for one variable in the product shows that \(\mathcal B\)
has supremum and total variation \(O(L^C)\) on its support. For
instance, its expression in the other variables has measure
\(\prod_{i<r}dt_i/t_i\) on fixed logarithmic intervals; the first
logarithmic derivative of the resulting integrand is \(O(L^C)\).
The function
\(\sum_{0\le n\le x}f_Y(n)-(A_Y/Y)x\) is \(O(A_Y)\), by
periodicity. Partial summation thus replaces the integral mean by
\(\sum_{n\ge1} f_Y(n)\mathcal B(n)\) with error \(O(A_Y L^C)\)
before multiplication by \(3\alpha\). Since
\(\alpha A_Y\ll X^c(Y/X)\) and \(Y\le X^{.76}\), this error is
\(O_J(X^c L^{-J})\). This proves \eqref{six:eq-six-low-model}.
In particular, by smooth majorants and minorants with transition
width \(XL^{-H}\), \eqref{six:eq-six-low-principal} gives the same
formula for the sharp prime interval \([X/2,X]\) when \(r=1\):
the two integrals differ by \(O(XL^{-H-1})\), and \(H\) can be
chosen after \(J\).

\section{Comparison on products of primes}\label{six:setup}
Let $r\ge2$, $\kappa>0$ be fixed and $W$ smooth with compact support in a product of
fixed-ratio positive intervals $t_i\asymp P_i$, $P_i\ge X^\kappa$,
$\prod_i P_i=XL^{O(1)}$. Each fixed mixed derivative of $W$ in
$\log(t_i/P_i)$ is bounded by a power of $L$. Write $\mathcal B(x)\,dx$
for the pushforward under multiplication of $W(\boldsymbol t)
\prod_i dt_i/\log t_i$. Solving for one variable in log coordinates gives
$|\mathcal B(x)|+x|\mathcal B'(x)|\ll L^{O(1)}$ and the corresponding
bounds on its compact support. For a component $G$ put
\[
 {\cal R}(G)=3\sum_{x\ge1,\,(x,6)=1}G(x)\mathcal B(x).
\]
Here $3=6/\varphi(6)$. Let $f_X$ indicate the complete allowed word
with unit final digit, extended periodically.
\begin{proposition}\label{six:box}
If the indices split into two groups of product scales
$M=X^m L^{O(1)}$ and $P=X^{1-m}L^{O(1)}$, $.4229\le m\le.4801$,
then for every fixed $J>0$,
\[
 \sum_{p_1,\ldots,p_r} W(\boldsymbol p) f_X(p_1\cdots p_r)
       ={\cal R}(f_X)+O_J(X^cL^{-J}).
\]
\end{proposition}
Extend the joint weight smoothly by zero on a larger fixed log box
and expand in Fourier series, reinserting separate smooth cutoffs which
are one on its support. Integrating by parts to a fixed order $>r+1$
and truncating to $L^{K(J)}$ in every variable gives any prescribed
inverse-log uniform error. Its cost on the prime side is
$O(X^cL^{O(1)})$ times this error, by the bound of $r!$ ordered
prime factorizations per integer and by counting full word periods.
The pushforward of a smooth box majorant has density $O(L^{O(1)})$,
giving the same cost for the reference. Thus it suffices to use
separated smooth boxes. In each group the coefficients are supported
on units and bounded by $r! L^{O(1)}$. When taking Mellin transforms
retain a smooth product-scale cutoff equal to one on each such box.

\subsection{Separation of upper and lower positions}
For $Y\mid X$ let $E_Y$ be averaging over the residue classes modulo
$Y$ on the integer circle of length $X$. Start at
$Y_0=6^{\lfloor .76k\rfloor}$ and telescope using the increments
$(E_Y-E_{Y/6})f_X$, $Y=6Y_0,\ldots,X$.
For the remainder of the comparison put $\eta=\log_XY$ at each
such level. In particular $.76<\eta\le1$ even when $.76k$
is not an integer.
The initial comparison is Theorem~\ref{six:thm-six-low-projection}.
For every radix power $R\mid Y/6$, write $F$ for the lower $R$-word
restricted to units. On the steps indexed by $\lfloor n/R\rfloor$ let
$U$ be the upper word of period $Y/R$ minus $5/6$ times the upper
word of period $Y/(6R)$, extended periodically. Then
\[
 (E_Y-E_{Y/6})f_X(n)=(X/Y)^{-\delta}F(n\bmod R)U(n).
\]
A block of radix scale $R$ has $\log_6 R$ positions. The lower mean is
$\rho_L\asymp R^{-\delta}$ and the upper density scale is
$\rho_U=(Y/R)^{-\delta}$. Counting upper words over covering periods
gives $O((X/R)\rho_U L^{O(1)})$ occupied steps and jumps.
Prescribing one allowed end digit changes only the absolute constants.
Throughout, $(X/Y)^{-\delta}X\rho_L\rho_U\asymp X^c$.

Section~\ref{six:adaptive} specifies a partition of the frequency
circle into a fixed finite number of amplitude ranges. Each has a
different cut $R$, but they use a common block, of scale
$B=6^{\lfloor wk\rfloor}$ above position $V_*=6^{\lfloor vk\rfloor}$.
Here $v+w=m+.05901$ and $R\asymp X^{m+h+.00003}$,
$.077\le h\le .101856$. Floor each one-digit magnitude at a small
fixed positive constant before forming the filter block amplitude at
$t=V_*\alpha$. Its logarithm has Lipschitz constant $O(B)$.
Convolve all range indicators on the rescaled circle with the same
even nonnegative kernel of mass one and degree $O(K)$, $K\asymp
B X^{.00005}$. Take the normalized squared modulus of an exponential
sum with coefficients sampling a fixed real even smooth bump.
Repeated summation by parts gives $O_b(K(1+K\|t\|)^{-b})$ for every
fixed $b$. Outside a range enlarged by a factor of two in amplitude
its multiplier is $O_A(X^{-A})$ for every fixed $A$, by using the
neighborhood of radius $X^{-.000025}/B$. The resulting even multipliers
are between zero and one and sum to one.

Apply each multiplier to the same full increment on the ambient
$X$-circle. Physical convolution uses shifts by multiples of $V_*$
of size $O(X^{m+.05906})$. Split the lower word by its highest digit.
For the pieces with digit zero and five shift the upper step boundaries
respectively left and right by $\lfloor R/12\rfloor$. For the interior
digits keep the boundaries in place. The product for each piece is
unchanged, and the support of its lower factor now avoids every step
boundary by $\gg R$. If a convolution shift crosses such a boundary,
the lower factor at the input vanishes. Thus the filtered product is
exactly the shifted upper step times the filtered lower piece. This
gives exact reconstruction across different cuts and shifts. Write the
filtered lower word as the sum of its exact additive periods $q\mid R$.
If $J$ is the multiplier and $\Pi_R^\circ$ the normalized complex word
sum, its $q$ term is $\rho_L\sum_{a\bmod q}^*J(a/q)\Pi_R^\circ(-a/q)
e(an/q)$. The star indicates reduced numerators (one term for $q=1$).
At $J=1$ this term is also $\sum_{e\mid q}\mu(q/e)E_e F$.
Our error unit for comparison of such an individual component times
its upper step is $X\rho_L\rho_U$.

\section{Completion over columns and the signed row sum}\label{six:dt-gram}
Let $\kappa_p=\log_6p$, $n=1-m$, and $A_U=\delta(\eta-\lambda)$.
The $R$-word has normalized complex exponential sum $\Pi_R^\circ$,
where the end positions may be restricted, and absolute value $F_R^\circ$.
Write $F_R$ and $\mathcal W_R$ for the absolute sum and set of the unmodified
word (with the common excluded digit). Put
\[
 \begin{gathered}
 a=.43389,\quad a_2=.108,\quad a_3=.207,\quad C=3c/4-a,\\
 L_p=(1-2a_p)\kappa_p,\quad \sigma=\min(C,L_2,L_3)>.2395 .
 \end{gathered}
\]
These are the full and pure first-moment bounds from the digital
certificates. For a general divisor grid use either the arbitrary-grid
bound or the sharper pure bound on a block of single-prime excess.
We use $m\in[.4229,.4801]$, $.76-o(1)\le\eta\le 1$,
$.077\le h\le .101856$, $\lambda=m+h+.00003$,
$R=6^{\lfloor\lambda k\rfloor}$ and $H=X^h$.
Denote the row and column lengths by $P=X^{n+o(1)}, M=X^{m+o(1)}$;
in their sizes the distortions are at most logarithmic powers.
The positive part of $z$ is denoted by $z_+$.
The upper source is a bounded step function on intervals of
length \(R\), possibly translated by a fixed fraction of \(R\).
In a fixed-multiple product box it has
\(O((X/R)\rho_U X^{o(1)})\) occupied steps, and the same estimate
holds for the sum of their squared amplitudes. Row and column
coefficients are supported on units modulo \(6\) and are bounded
by \(X^{o(1)}\). For an exact additive period \(q\mid R\) the lower
source has the expansion
\begin{equation}
 G_q(t)=\rho_L\sum_{a\bmod q}^{*}
                    J(a)\Pi_R^\circ(-a/q)\e(at/q),
 \label{six:dt-gram-lower-source}
\end{equation}
where the star requires \((a,q)=1\), and \(|J|\le1\).
The factors \(J\) will include the taper and, where applicable,
the restriction to the complement of a union of frequency bands.

Apply the column completion preceding \eqref{dt:gram-incidence} to
\eqref{six:dt-gram-lower-source}. Indeed $HM/R=X^{-.00003+o(1)}$,
so the translated upper steps have $O(1)$ jumps per row interval.
The dyadic subdivision and smooth completion are independent of the
radix and of the label weights. They give
\begin{equation}
 p a-(p+\ell)b\equiv t\pmod q,\qquad
      0<|\ell|\ll H,\qquad |t|\ll T,\qquad
      T=(1+q/M)X^\epsilon .
 \label{six:dt-gram-incidence}
\end{equation}
Each twist coefficient is $O(M)$ and independent of rows and labels.
As in \eqref{dt:gram-diagonal}, normalization by $(PM\rho_U\rho_L)^2$
divides the off-diagonal incidence sum by $PH\rho_U$ up to
$X^{O(\epsilon)+o(1)}$; the normalized diagonal is bounded by
\begin{equation}
 X^{A_U-h+o(1)}
              \sum_{a\bmod q}^{*}|J(a)\Pi_R^\circ(a/q)|^2 .
 \label{six:dt-gram-diagonal}
\end{equation}
This form allows the use of the filtered square norm.

\begin{lemma}[Signed lift and fold]\label{six:dt-gram-base}
For the incidence \eqref{six:dt-gram-incidence}, let \(g\) be the
exact order of \(a-b\) in the additive group modulo \(q\).
Write
\begin{equation}
 \begin{gathered}
 g=p_0^ip_1^j,\quad \{p_0,p_1\}=\{2,3\},\quad i\ge j,\qquad
 s=i/k,\quad d=j/k,\quad \xi=\log_X g,\\
 V=6^{\lfloor j/2\rfloor},\quad Q=q/V,\quad
 g_0=g/V,\quad Z=q/g,\quad G=6^i,\\
 K_0=p_1^{i-j},\quad d_2=(K_0,Z),\quad
 D_0=Z/d_2,\quad K=K_0/d_2.
 \end{gathered}
 \label{six:dt-gram-order}
\end{equation}
Suppose the two label multipliers are constant on the
\(V\)-lifts, that is, on increments of \(Q\) in the label.
The normalized contribution at fixed \(g\) is at most
\begin{equation}
 X^{A_U-Cd-L_{p_0}(s-d)+o(1)}
       \max(1,g/P)^{1/2}D_0^\delta(1+T/D_0).
 \label{six:dt-gram-base-bound}
\end{equation}
If the gaps are restricted to \((\ell,d_2)=e_0\), the
completion factor can be replaced by \(1+T/(e_0D_0)\).
The contribution can also be estimated by the finer folded
incidence over \(z_0\bmod D_0\) in the proof below.
In particular \eqref{six:dt-gram-base-bound} has exponent at most
\begin{equation}
 \delta(\eta-s)-Cd-L_{p_0}(s-d)+(s-m)_+
                                    +\tfrac12(\xi-n)_+,
 \label{six:dt-gram-base-exponent}
\end{equation}
with the usual arbitrarily small completion loss.
\end{lemma}

\begin{proof}
Use the coarse-label parametrization in the proof of Lemma~\ref{dt:gram-base},
replacing $10$ by $6$ and $5$ by $3$ in its radix factors. Namely,
$a=a_0+Qx$, $b=b_0+Qy$, $a_0-b_0=Zu$, $(u,g_0)=1$ and
$pu\equiv(t+\ell b_0)/Z\pmod{g_0}$ (the quotient must be integral).
Since $V\mid g_0$, primitivity is constant on a lift and for $V>1$
the unit ratio $r=(p+\ell)/p\bmod V$ is fixed. With
$\alpha=a_0/Q$, $\beta=b_0/Q$ and $w=p^{-1}((t-pa_0+(p+\ell)b_0)/Q)$
modulo $V$, the signed lift sum is exactly \eqref{dt:gram-intercept-identity}
for the restricted six-adic word:
\begin{equation}
 C_r(w)=\frac{V}{|\mathcal W_V^\circ|^2}
  \sum_{\substack{x\in\mathcal W_V^\circ\\rx\bmod V\in\mathcal W_V^\circ}}
              \e(wx/V)\e((\alpha x-\beta(rx\bmod V))/V).
 \label{six:dt-gram-intercept-identity}
\end{equation}
Here $C_r(w)$ is the correlation of the two low-block polynomials in
that identity. The intercept advances by the unit $-u/p$ and
$g_0\ll X^{.614\lambda+o(1)}\ll P$: indeed $\log_6 p_0\le\log_6 3<.614$
and $\log_X g_0=\kappa_{p_0}(s-d)+d/2+o(1)$. Thus the same Parseval
bound costs $(P/g_0)V^{\delta-c/2}\max(1,g/P)^{1/2}$; at $V=1$ the lift
has one term, including $g=1$, when the labels coincide. The two label
multipliers factor out before this signed sum. Above the $V$-block the
coarse grid splits as in \eqref{dt:gram-first}: the complementary
multiplier $(6/p_0)^b$ after $b$ positions is a unit modulo $p_0$.
The shifted full and pure first moments therefore give, per coarse variable,
\begin{equation}
 O\left(
    (6^{j-\lfloor j/2\rfloor})^a
    (p_0^{i-j})^{a_{p_0}}
    F_{R/G}^\circ(Kz/D_0)\right).
 \label{six:dt-gram-first}
\end{equation}
A multiplier constant on each lift may remain inside these
first sums. In particular, a pointwise blockwise gain in a
coarse first norm is a gain in \eqref{six:dt-gram-base-bound}.
The product of the two ordinary first-norm factors, the row
density \(g_0^{-1}\), and the signed mean \(V^{\delta-c/2}\)
has exponent \(-Cd-L_{p_0}(s-d)+o(1)\).
Here the upper factor retains the prescribed leading digit
if there is one, and retains the terminal restriction when
\(G=1\). If an end restriction instead falls in a summed
block, the sampling losses are bounded by
the digital sampling and product bounds.

Decompose $z=z_0+D_0 b$ as in \eqref{dt:gram-fold}. The resulting fold is
\begin{equation}
 D_0\mid t+\ell z_0,\qquad
 \ell b\equiv -(t+\ell z_0)/D_0\pmod{d_2}.
 \label{six:dt-gram-fold}
\end{equation}
For $e_0=(\ell,d_2)$ the $e_0$ lifts cancel the $1/e_0$ in the
gap count, leaving $O(1+T/(e_0D_0))$ admissible twists per gap and
$z_0$. No decimal valuation is needed for \eqref{dt:gram-fold-identities}:
$G=gK_0$ and $\operatorname{lcm}(G,q)=GD_0\mid R$ give
\begin{equation}
 GD_0=qK,\qquad (K,D_0)=1,\qquad D_0\mid R/G
 \label{six:dt-gram-fold-identities}
\end{equation}
so the shifted square bound on the $D_0$-grid costs $O(D_0^\delta)$,
proving \eqref{six:dt-gram-base-bound}.
Since \(D_0\le R/G\) and \(q/D_0\le G\), we may bound the
square and the completion factor by their separate worst
values at \(q=R\). This gives \eqref{six:dt-gram-base-exponent}.
\end{proof}

\begin{lemma}[The zero twist]\label{six:dt-gram-middle}
Suppose \(q<MX^{-c_0}\) with fixed \(c_0>0\). The
off-diagonal bound, for multipliers constant on the signed
lifts, has exponent at most
\begin{equation}
                A_U-\sigma\log_X q+o(1).
 \label{six:dt-gram-zero-exponent}
\end{equation}
\end{lemma}

\begin{proof}
Apply the $t=0$ calculation of Lemma~\ref{dt:gram-middle}: since $z$
is a unit, $Z\mid\ell$; also $g\le q<M\ll P$. Thus $O(H/Z)$ gaps and
the high trace $O(d_2D_0^\delta)$ leave $D_0^{-c}$ after normalization.
Here $C,L_p\ge\sigma$ and $\log_XD_0\ge(\log_Xq-s)_+$, which gives
\eqref{six:dt-gram-zero-exponent}.
\end{proof}

\subsection{Joint affine averaging}

\begin{lemma}[Joint affine overlap and intercept mean]\label{six:dt-gram-joint}
In the setting of Lemma~\ref{six:dt-gram-base}, suppose $\xi\ge m+.00001$ and
$h+\xi-m<d/2-.00001$. Writing $x=\xi-m$ and $l=(\xi-n)_+$,
the exponent in \eqref{six:dt-gram-base-exponent} may be reduced by
\begin{equation}\label{six:dt-gram-joint-gain}
 \min\{\delta l+(\delta/2)\min(d/2,2(h+x)),\ x+(\delta/2)l\}.
\end{equation}
This improvement is uniform in the coarse label pair. It can be
combined with a saving inside the two coarse first sums.
\end{lemma}

\begin{proof}
Use complete unrestricted words at auxiliary scales. For
base-six lengths \(J,U\mid W\) put
\[
 M_r(U)=\max_{b\bmod U}\#\{x\in\mathcal W_U:
                                    rx+b\bmod U\in\mathcal W_U\}.
\]
Write $B_r(J)=\#\{x\in\mathcal W_J:rx\bmod J\in\mathcal W_J\}$.
The overlap calculation \eqref{dt:gram-joint-conditional}--\eqref{dt:gram-affine-lift}
gives the same bounds for these sets. In particular, averaging over units modulo $W$,
\begin{equation}
        \mathbb E_r[B_r(J)M_r(U)]\ll J^{2c-1}U^{2c-1}.
 \label{six:dt-gram-joint-overlap}
\end{equation}
Indeed for $g_1=(J,q_1)$, the proportion of units in any nonempty
unit fibre is bounded below by $1/3$, and a residue class modulo
$b\mid J$ contains $O((J/b)^c)$ words. The conditional count is therefore
\begin{equation}
 \mathbb E[B_r(J)\mid r\bmod q_1]\ll
  J^{2c-1}g_1^\delta\sum_{m_0\mid J}m_0^{1-2c}
                       \ll J^{2c-1}g_1^\delta .
 \label{six:dt-gram-joint-conditional}
\end{equation}
Both divisor sums used in the decimal proof still converge: the moduli
are $2^i3^j$, $1-2c<0$ and $2a-c<0$ ($a=.43389$, $c>.898$).
Thus averaging the Fourier expression for $M_r(U)$ gives
\eqref{six:dt-gram-joint-overlap}. Conditioning on the lowest $W$-word
when $W\mid V$ and $U\mid V/W$ gives
\begin{equation}
          B_r(V)\ll (V/W)^c B_r(W)M_r(U)/U^c .
 \label{six:dt-gram-affine-lift}
\end{equation}
An end restriction may be dropped for an upper bound on these counts.
Put $\mathcal A=HT/Z$, $\mathcal L=\max(1,g/P)$ as in Lemma~\ref{dt:gram-joint}.
Since $q\ge g\ge MX^{.00001+o(1)}$, their exponents are respectively
$h+x+O(\epsilon)+o(1)$ and $l+o(1)$; $l\le x$ since $n\ge m$.
Taking $\epsilon$ small in the strict fit gives $1\le\mathcal L\le\mathcal A\ll V X^{-.000005}$.
These, not the decimal cutoff $x\ge .020$, are the hypotheses of the
calculation through \eqref{dt:gram-joint-factor}. In the fixed-gcd
injection take $e_0=(\ell,(t+\ell b_0)/Z)$ and round $W\asymp\mathcal A/e_0^2$
to a power of six above the determinant bound when this scale is nontrivial.
The determinant is still $(\ell_1t_2-\ell_2t_1)/(Ze_0^2)$; bounded rounding
costs only a constant. With $L_b=2^b\mathcal L$, the unit intercept step
therefore gives the same squared mean as \eqref{dt:gram-fejer}:
\begin{equation}
 \ll (V/|\mathcal W_V^\circ|^2)^2 B_r(V)
        \sum_{b\ge0}2^{-2b}(L_b/W_b)B_r(W_b).
 \label{six:dt-gram-fejer}
\end{equation}
Choose $U=\min(W,V/W)$ and $W_b\mid W$ a radix power comparable to
$\min(L_b,W)$. The Cauchy estimate and dyadic sum in the derivation of
\eqref{dt:gram-joint-per-gcd}, with \eqref{six:dt-gram-affine-lift}
and \eqref{six:dt-gram-joint-overlap}, bound the sum of intercept means at this gcd by
\begin{equation}
 \ll V^{\delta-c/2}\,
           W^{1-\delta/2}U^{-\delta/2}
                 \mathcal L^{1/2}\min(\mathcal L,W)^{-\delta}
 \label{six:dt-gram-joint-per-gcd}
\end{equation}
after summing the triangular-kernel tails. The gcd summation from
\eqref{dt:gram-joint-per-gcd} to \eqref{dt:gram-joint-factor} splits at
$W=\sqrt V$ and $W=\mathcal L$; it uses $W\asymp\mathcal A/e_0^2$
and $0<\delta<1/2$, not the decimal cut exponents. Including the
$O(H)$ below-threshold gcds, the factor relative to
$\mathcal A V^{\delta-c/2}\mathcal L^{1/2}$ is bounded by
\begin{equation}
   \min(V,\mathcal A^2)^{-\delta/2}\mathcal L^{-\delta}
                       +(H/\mathcal A)\mathcal L^{-\delta/2}.
 \label{six:dt-gram-joint-factor}
\end{equation}

Here one must keep \emph{both} terms: the decimal absorption after
\eqref{dt:gram-joint-factor} does not follow for $x\ge .00001$.
Their respective savings are $\delta l+(\delta/2)\min(d/2,2(h+x))$
and $x+\delta l/2$, giving \eqref{six:dt-gram-joint-gain} after folding
(the trace multiplicity $d_2$ changes $T/Z$ to $T/D_0$). This is uniform
before the coarse first sums. At our cuts $V\mid V_*$ by $\lambda/2<v$,
so the amplitude filters, functions of $V_*\alpha$, are constant on the
lifts. Orbit masks are used here only when their lengths are at least
$V$; shorter ones are treated separately in Section~\ref{six:sixmask-geometry}.
\end{proof}

\section{A joint digital mean for powers of two and three}
\label{six:joint-digital}

For \(p\in\{2,3\}\) put \(p'=6/p\), \(\kappa_p=\log_6p\).
Let \(F_{6^l}(t)=\prod_{0\le j<l} f(6^jt)\), using the
normalized one-digit magnitude \(f\) of the Fourier estimates.
For \(v>0\) define
\[
 C_{p,l}(v)=\sup_{x\in\mathbb R}\frac1{p'^l}
 \sum_{u\bmod p'^l}
 \left(\sum_{\substack{0\le n<6^l\\n\equiv u\bmod p'^l}}
 F_{6^l}\left(\frac{n+x}{6^l}\right)^v\right)^2.
\]
\begin{theorem}\label{six:joint-tail-moment}
For \(L_1,L_2\ge0\),
\(C_{p,L_1+L_2}(v)\le C_{p,L_1}(v)C_{p,L_2}(v)\).
Uniformly in the forbidden digit and \(l\ge0\),
\[
 C_{2,l}(1)\ll6^{-.300l},\qquad C_{3,l}(1)\ll6^{.135l}.
\]
\end{theorem}
\begin{proof}
Expand the square into pairs \((n_1,n_2)\) having a common
residue modulo \(p'^l\). The same upper bound holds for a
prescribed residue difference \(t\) and separate shifts
\(x_1,x_2\). Indeed, Cauchy's inequality in the residue class
\(u\) bounds their sum by the geometric mean of the two
unshifted-residue-difference square sums.

Set \(L=L_1+L_2\) and write
\(n_i=h_i+6^{L_2}m_i,\ 0\le h_i<6^{L_2}\).
The condition \(n_1-n_2\equiv t\bmod p'^L\) first requires
\(h_1-h_2\equiv t\bmod p'^{L_2}\).
For each such pair it leaves a fixed difference for
\(m_1-m_2\bmod p'^{L_1}\), because \(p^{L_2}\) is a unit
there. Put \(y_i=(h_i+x_i)/6^{L_2}\). Each word separates as
\[
 F_{6^L}\left(\frac{n_i+x_i}{6^L}\right)
 =F_{6^{L_2}}(y_i)F_{6^{L_1}}\left(\frac{m_i+y_i}{6^{L_1}}\right).
\]
Apply the preceding translated-class inequality first to the
sum on the \(m_i\), then to the sum on the \(h_i\).
The two normalizations multiply to \(p'^{-L}\), proving the
product inequality. The same argument gives a product of
bounds \(C_{p,L_j}(v_j)\) for adjacent blocks with different
positive powers \(v_j\) on corresponding positions.
Different powers on the two words can first be separated by
Cauchy's inequality. The estimates allow a prescribed highest
digit or a unit-digit restriction at the bottom. Split at the
at most two affected positions. Each of their normalized factors
has magnitude at most one and its one-position coset mean is at
most \(p^2\), even with a translated class and separate shifts.
The bounded number of intervening ordinary blocks then gives the
same exponent, at a constant cost. The same splitting applies
to first and square norms of the coarse word.

For \(L=6\), the script
\texttt{bilinear/verify\_pure\_tail.py} certifies
\(C_{2,6}(1)<6^{-1.800}\) and \(C_{3,6}(1)<6^{.810}\).
Here are the details of the certificate. Shifts modulo one
suffice by permutation of the residue classes. Evenness of
the Fourier magnitude reduces the interval to \([0,1/2]\).
Partition it into intervals of length \(1/H\), \(H=256\).
At all \(N=2^{15}\) uniform phase nodes bound \(f\) above by
the square root of the upper endpoint of
\[
 \frac1{25}\left(5+2\sum_{j=1}^5
  \#\{a:a,a+j\text{ allowed}\}\cos(2\pi j t)\right).
\]
The cosine bounds use Machin's identity, the alternating
arctangent series, and directed interval arithmetic in the
cosine polynomial through degree 40. The error of that
polynomial on \([-\pi,\pi]\) is at most \(\pi^{42}/42!\),
by Taylor's formula through the identically vanishing
degree-41 term. All square roots and products are rounded
up in units of \(2^{-34}\).
For a factor having denominator \(D=6^j,\ 1\le j\le6\),
use the nearest phase node at the shift midpoint, adding
\(6/N+6/(HD)\). Indeed \(\operatorname{Lip}(f)<12\) by
centering the five phases at \(5/2\); the two phase errors
are at most \(1/(2N)\) and \(1/(2HD)\).
Summing the resulting upper products on each residue class
and squaring gives the certified respective bounds.
It suffices to check three digits by reflection. Iterate
length-six blocks and absorb the five possible nonempty
remainders into the implied constant.
\end{proof}

\subsection{A consequence for congruence sums}
Here is the arithmetic form of the estimate used in
completing a congruence. It records explicitly which
other factors can remain inside the same bound. Set
\(X=6^k\), and let \(p^i p'^j\mid q\mid X\) with \(i\ge j\).
Write
\[
 Z=q/(p^i p'^j),\quad D_2=(p'^{i-j},Z),\quad D_0=Z/D_2,\quad
 s=i/k,\ d_*=j/k,\ \omega=(i-v_{p'}(q))_+/k .
\]
Then \(q/D_0=6^i p'^{-k\omega}\).
For \(M=X^{m+o(1)}\) and completion
\(T=(1+q/M)X^\varepsilon\), fix a nonzero integer gap
\(h\) and \(z_0\bmod D_0\). Put \(e_0=(h,D_2)\) and
\(\theta=\log_X e_0\).
Writing \(z=z_0+D_0b\) in the congruence
\(t+hz\equiv0\bmod Z\) for an integer \(|t|\ll T\) gives
\[
 t+hz_0\equiv0\bmod e_0D_0,\qquad
 b\equiv-\overline{h/e_0}\,\frac{t+hz_0}{e_0D_0}
                 \pmod{D_2/e_0}.
\]
For a terminal block of \(l\le i-j\) pure positions
after summation of the preceding positions, its first
norm at this label is
\[
 S(b)=\sum_{a\bmod p^l}
 F_{6^l}\left(\frac{a+p'^{i-j-l}z/Z}{p^l}\right).
\]
This formula follows by writing an index on the pure
\(p^{i-j}\)-grid as \(a+p^lb'\): the lower \(i-j-l\)
positions are summed first in \(b'\), and the factor
\(p'^{i-j-l}\) multiplying \(a\) is permuted away modulo
\(p^l\). First norms of all preceding blocks can be bounded
uniformly in this remaining label.

If \(l>k\omega\) put \(J=p'^{l-k\omega}\).
Since \(D_2=p'^{i-j-k\omega}\), the shift in \(S(b)\) is
\((b+z_0/D_0)/J\), and \(J\mid D_2\). Hence (with arbitrary
shift \(\beta\) in place of \(z_0/D_0\)) one has
\[
 J^{-1}\sum_{b\bmod J}
 \left(\sum_{a\bmod p^l} F_{6^l}
    \left(\frac{a+(b+\beta)/J}{p^l}\right)\right)^2
 \le p'^{k\omega} C_{p,l}(1).
\]
To see this, multiply numerator and denominator of the nested
grid by \(p'^{k\omega}\), then retain the corresponding
\(J\) residue classes in the positive sum defining \(C_{p,l}\).

When \(T/(e_0D_0)\gg J\), intervals of quotients in the
congruence above can be enlarged to full \(J\)-periods
at bounded cost. Taking every solution \(b\bmod D_2\),
their per-lift mean is the displayed \(J\)-grid mean.
If \(J\mid D_2/e_0\), the quotient traverses a unit
progression modulo \(J\). If \(D_2/e_0\mid J\), the free
lifts fill all the remaining compatible classes equally.
These are the only possibilities since both moduli are
powers of \(p'\). For the product of two terminal first
sums with different shifts the same bound follows by Cauchy.

For clarity, suppose that the ordinary shifted pure first
norm for \(l\) positions is \(O(p^{a_p l})\) and write
\(\tau_2=-.300,\ \tau_3=.135,\ H_p=2a_p\kappa_p-\tau_p>0\).
Suppose an average over gaps, in a class with prescribed
\(e_0\), has the normalized arithmetic factor
\(1+T/(e_0D_0)\) times the two ordinary terminal first
norms. This normalization results, for example, when the
gcd class has \(O(H/e_0)\) gaps up to \(H\), each quotient
has \(e_0\) lifts, and the normalization is \(H\).
Then compared with the upper twist exponent \(x=(s-m)_+\)
and the ordinary first norms, the improvement furnished by
the preceding mean, in exponents of \(X\), is at least
\[
 \min(H_p,\kappa_{p'})\min(s-d_*,\,x/\kappa_{p'})
   -O(\varepsilon)-o(1).
\]
Indeed let \(0<\ell\le s-d_*\) and \(\kappa_{p'}\ell<x\).
If \(q\le M\), or the twist factor is bounded up to the
completion loss, the raw gain \(x\) suffices. Otherwise
\(\log_X(T/D_0)=x-\kappa_{p'}\omega+O(\varepsilon)+o(1)\).
When \(\theta+\kappa_{p'}\omega\ge
\ell\min(H_p,\kappa_{p'})\) the raw twist count suffices.
In the complementary case reserve
\(l=\lfloor k(\ell-\theta/\kappa_{p'})\rfloor\) positions.
Then \(l>k\omega-O(1)\); if the difference is bounded,
the raw twist gain already suffices up to \(O(1/k)\).
The twist interval contains full \(J\)-periods by
\(\kappa_{p'}\ell<x\), with an arbitrarily small margin 
taken before the completion losses.
Relative to \(p^{2a_p l}\), the terminal square supplies
\(H_p l/k-\kappa_{p'}\omega\). The actual twist count adds
\(\theta+\kappa_{p'}\omega\) to the saving, giving
\(H_p\ell+(1-H_p/\kappa_{p'})\theta+O(\varepsilon)+o(1)\).
This is at least \(\min(H_p,\kappa_{p'})\ell\) to those
errors since \(0\le\theta\le\kappa_{p'}\ell\).
One can take \(\ell\) within \(O(\varepsilon)+o(1)\) of the
stated minimum, or use the raw bound when that minimum is
\(O(\varepsilon)+o(1)\). The gcd classes cost at most a
logarithmic factor. Fourier multipliers of absolute value
at most one may be dropped after any signed summation
requiring their constancy. If retained with modified digit
powers, gains on disjoint blocks multiply by the
concatenation inequality of Theorem~\ref{six:joint-tail-moment}.

\subsection{A finite family of cuts for the diagonal and high orders}
\label{six:adaptive}
Write $\delta=1-c$, $\kappa_p=\log_6p$, and fix
$.4229\le m\le .4801$, $n=1-m$, and $\eta\le1$. Use the moment
orders and certified upper exponents in the following table. The rows
refer respectively to the full radix grid, and to the pure grids for
the same number of digit positions. On a pure grid the base of the
power bound is its cardinality.
\[
\begin{array}{c|rrrrr}
r&1&4/3&3/2&7/4&2\\ \hline
f_r&.433890&.289650&.229470&.155800&.101756\\
b_{2,r}&.108&.052&.038&.025&.00813\\
b_{3,r}&.207&.096&.067&.040&.01215
\end{array}
\]
The bounds can also be used with a sufficiently small fixed positive
floor on the modulus of each one-digit factor. Indeed, for each entry
the unfloored estimate has a strict exponent margin, including
$\delta<.101756$ for the last entry of the first row. Choose a fixed
block length so that its shifted moment cap is strictly below the
desired power, and use uniform continuity in the floor on this
single finite block. Iteration over such blocks, including a bounded
final block, gives the desired bound. Take the minimum of the finitely
many choices for the floor. Fixing one of a block's end digits costs a
constant, since one may split off that position and estimate its
factor by one.

Put $a=f_1$, $D=.101756$, $v=(m+.10201)/2$, $w=v-.043$, and use
the common filter block $B=6^{\lfloor wk\rfloor}$ starting at the
position $V_*=6^{\lfloor vk\rfloor}$. Its normalized amplitude,
with the floor just described, is denoted by $F$. Partition according
to $F\ge B^{-.22}$, $F\le B^{-.50}$, and the intervals
$[B^{-(t+.0005)},B^{-t}]$ with left exponent
$t=.2200,.2205,\ldots,.4995$. Assign a boundary to either adjoining
set. Use the common positive convolution kernel to smooth all
indicators to multipliers summing to one. Each supported shell can be
enlarged by a factor of two at either endpoint with a uniform error
$O_A(X^{-A})$ for each prescribed $A$. Choose the cut exponent
\[
 h=
 \begin{cases}
 .101856,& F\ge B^{-.22},\\
 .077,& F\le B^{-.50},\\
 \max\{.077,\ D+.0001+
  w(\min_r(f_r+(r-2)t)-D)\},&\text{on the shell with left exponent }t .
 \end{cases}
\]
For each multiplier put $\lambda=m+h+.00003$,
$R=6^{\lfloor\lambda k\rfloor}$. In particular $v+w=m+.05901$
and $\lambda/2\le v-.000062$. Thus the common filters commute with the
protected upper steps as in the physical smoothing construction:
convolving in the $V_*$-rescaled angle at degree $O(BX^{.00005})$
translates the integer argument by multiples of $V_*$ of size
$O(X^{m+.05906})=o(R)$, with a fixed power to spare. Their
multipliers are constant on every signed lift of length at most
$6^{\lfloor\lambda k/2\rfloor}$. All $O(k^{-1})$ changes from
integer positions can be absorbed in the exponent error.

For the diagonal at $q\mid R$, split the label grid at
$D_1=(q,V_*)$, $D_2=(q/D_1,B)$, $Q_2=q/(D_1D_2)$.
After the bottom sum the frequency multiplier $V_*/D_1$ permutes
the remaining grid, and after the middle sum $B/D_2$ permutes the
terminal grid. Moreover $Q_2\mid R/(V_*B)$. Square summation on
the first and last grids uses the bound $O(D_1^\delta Q_2^\delta)$
from the shifted-grid Parseval bound (or digital occupancy on a
subgrid). On an intermediate shell or the low-amplitude terminal
shell use $F^2\ll B^{-(2-r)t}F^r$ in the middle sum, with
$t=.5$ on the latter. The resulting relative squared exponent is
at most
\[
 D-h+w\bigl(\min_r(f_r+(r-2)t)-D\bigr)+o(1).
\]
Indeed $D_1\le V_*$, $D_2\le B$, and $Q_2\le R/(V_*B)$; in this
bound one can first use $D$ throughout as an upper bound for
$\delta$. It is at most $-.0001+o(1)$ on the middle shells, and
on the terminal shell by $f_1-.5<-.066$ and $w\ge.219455$.
On the top shell unfiltered Parseval gives $D-h=-.0001$ directly.

For the off-diagonal write the difference order as
$g=p^i(p')^j$, $i\ge j$, $\{p,p'\}=\{2,3\}$. Set
$s=i/k$, $d_0=j/k$, $\xi=\kappa_p s+\kappa_{p'}d_0$,
$C=3c/4-a$ and $L_p=(1-2b_{p,1})\kappa_p$.
The general signed-lift bound, with $\eta\le1$, is
\[
 E=\delta(1-s)-C d_0-L_p(s-d_0)+(s-m)_+
                                          +\tfrac12(\xi-n)_+ .
\]
Logarithmic factors and arbitrarily small fixed powers from completion
are suppressed in these exponents. The following are lower bounds
for the improvements available in this bound. Write
$\ell=(\xi-n)_+$. If $\xi-m\ge .00001$ and
$h+\xi-m<d_0/2-.00001$, the joint intercept and ratio estimate
improves it by
\[
 J=\min\{\delta\ell+\tfrac{\delta}{2}
              \min(d_0/2,2(h+\xi-m)),\ \xi-m+\tfrac{\delta}{2}\ell\}.
\]
Use $J=0$ otherwise. The second branch here retains the
below-threshold gcd term of the joint estimate, so the reduction
does not require $\xi-m\ge .020$.

On the top shell put $t_+=.22$, and on an intermediate shell put
$t_+=t+.0005$. Set
\[
 u_1=(\min(v+w,d_0)-v)_+,\qquad
 u_2=(\min(v+w,s)-\max(v,d_0))_+ .
\]
The first expression is interpreted as zero for $d_0<v$.
They measure the parts of the common filter block that lie in
the radix and pure pieces of a coarse sum. On these shells
the multiplier is at most $O(B^{(r-1)t_+}F^{r-1})+O_A(X^{-A})$.
After the signed lift keep the factors in those two overlaps.
Splitting the coarse grid at the overlap endpoints, using the
pure-grid permutation by a power of $p'$ on each pure subblock,
and discarding the other factors of the filter give the gain
\[
 \Delta=\max_r\{(a-f_r)u_1+(b_{p,1}-b_{p,r})\kappa_p u_2
                                      -(r-1)t_+ w\}\ge0.
\]
Use $\Delta=0$ on the low-amplitude terminal shell. The two coarse
first sums give $2\Delta$ in total. The joint bound is uniform
in their two labels and precedes these sums, hence $J+2\Delta$
is an admissible gain.

Independently, the folded pure-tail estimate with its mean-square
grid exponents $-.300$ and $.135$ gives a gain of at least
\[
 T_p=\gamma_p\min\{s-d_0,(s-m)_+/\kappa_{p'}\},\qquad
                       \gamma_2=.383,\quad\gamma_3=.118 .
\]
Indeed $\gamma_p<\min(\kappa_{p'},2b_{p,1}\kappa_p-\tau_p)$,
where $\tau_2=-.300$ and $\tau_3=.135$. This estimate permits
discarding the label multipliers after the signed lift. We take
the better of the two bounds.

The following inequality holds throughout the stated range of
$m$, every shell, $p\in\{2,3\}$, and $0\le d_0\le s$ with
$.305\le s\le\lambda$:
\[
                       E-\max(J+2\Delta,T_p)\le-.0001 .
\]
Here and in the finite verification the second argument of the
maximum can equivalently be decreased by an arbitrarily small fixed
error for rounding the length of the pure tail. To check the
inequality run \texttt{adaptive/verify\_gram.py} with Python~3.
Its enclosures for $\delta$ and $\kappa_2$ are
$[.101755,.101756]$ and $[.386852,.386853]$. They follow by
raising $5$ and $2$ and comparing to powers of $6$, with integer
exponents obtained by taking a common denominator. The verifier
uses integer billionths, flooring lower products and taking ceilings
of upper products. For a box in the four variables $(m,s,d_0,t)$,
it clips $s$ to $d_0\le s\le m+h+.00003$ and $d_0$ to $d_0\le s$.
The interval bounds for $\xi=d_0+\kappa_p(s-d_0)$ use monotonicity
in both length variables on this domain. For the linear term use
$-L_ps+(L_p-C)d_0$ with lower bounds for $L_p$ and $C$, since
$s\ge d_0$ and $L_p>C$. It sets $J=0$ until its conditions hold
on the whole box. It uses lower bounds on each of the two overlap
lengths and upper bounds on the filter penalty. For the terminal
gain both arguments of the minimum are bounded below. These rules
give an upper bound on the exponent for the whole box. A box is
accepted if its upper bound is at most $-.0001$; others are
bisected, covering also shared boundaries. The top and bottom
shells are checked separately. For intermediate shells the
verifier covers all real left endpoints $.22\le t\le .5$, using
$t+.0005$ for the right endpoint. Termination and acceptance of
all boxes establish the inequality.

\subsection{Small-period removal for the lower-amplitude filters}
Let $X=6^k$, $c=\log_6 5$, $\delta=1-c$. Suppose $R,V,B$ are powers of
six, $VB\mid R/6$, $V=X^{v+o(1)}$, $B=X^{w+o(1)}$, with $v>.235$
bounded away from $.235$ and $w\ge .217-o(1)$. Let $A_R$ denote the
normalized complex Fourier polynomial of the bottom word on $[0,R)$.
One may fix its top digit and restrict the units digit to a nonempty
set of allowed units. Denote by $F_B$ the magnitude of the polynomial of
the complete unrestricted $B$-block with the same omitted digit.
Let a multiplier of magnitude at most one be negligible (to every
fixed power of $X$) off $F_B(V\alpha)\le 2B^{-.21}$. It suffices that
this localization hold for the word magnitude in which each one-digit
amplitude has been floored at a small positive constant. Then uniformly
for every $q\mid R$ with $q\le X^{.235}$,
\[
  \sum_{a\bmod q} |A_R(a/q)\omega(a/q)|^2\ll X^{-.003}.
\]
Indeed let $D=(q,V)$, $Q=q/D$, and $C=(Q,B)$, $K=Q/C$. The modulus
$Q$ is either $1$ or a pure power of $p=2$ or $3$, since a factor of
each prime in $Q$ would require $q\ge V$. For $Q=1$ one may set $p=2$.
Split $a$ first into classes modulo $Q$. The bottom $V$-block squares
sum on each shifted $D$-grid to $O(D^\delta)$. Here and below the
arbitrary-grid square bound follows by retaining a full radix block of
length $Z\asymp D$, $Z\le D$, and sampling its square on a $D$-grid.
The exact grid square is $D/|\mathcal W_Z|$, by orthogonality, also
with a fixed shift. An end restriction inside the retained block loses
at most an absolute factor. All unretained digit magnitudes are at most
one. The multiplier $V/D$ permutes the remaining residues modulo $Q$.
On the low-amplitude region, the two middle factors together with the
squared multiplier are bounded by $2 B^{-.21}F_B$ up to an arbitrarily
small power error. Split at residue classes modulo $K$. Summing the
remaining middle amplitude on the shifted $C$-grid costs
$O(C^{a_p})$, where $a_2=.151$ and $a_3=.231$: retain the first
$\log_p C$ base-six positions and use the shifted pure-power estimate
proved by \texttt{cert/verify\_mixed.py}. The top-block argument is
independent of this lift. Its squares finally cost $O(K^\delta)$,
since $B/C$ is a unit modulo $K$ and $K\mid R/(VB)$.
Thus the left hand side is at most
\[
  X^{o(1)}q^\delta B^{-.21} C^{a_p-\delta}+O(X^{-10}).
\]
With $C\le p^{\log_6 B}$ the exponent is bounded by
$.235\cdot .102+(.614(.231-.101)-.21) w+o(1)<-.004$.
The elementary logarithmic bounds used here are
$.101<\delta<.102$ and $\log_6 3<.614$.

In relative units this gives a negligible contribution for an exact
period $q$ of the lower word multiplied by the bounded upper step of
density scale $\rho_U$. Each occupied step of length $R$ is a union of
whole $q$-periods on the integer line. The sum of absolute values over
that step is $O(R\rho_L X^{-.0015})$, by Parseval and Cauchy's
inequality, with lower-word density $\rho_L$. At a box boundary one
can extend a truncated step using the same nonnegative majorant.
There are $O((X/R)\rho_U (\log X)^{O(1)})$ occupied steps. Products
of a fixed number $r$ of primes have at most $r!$ ordered prime
factorizations. The same estimate thus bounds their contribution, and
also the integer-tested continuous reference with bounded logarithmic
weight. The statement applies at every cut. On the largest cut it
permits replacement of the complementary high-amplitude multiplier by
one throughout this small-period range.

\section{Unfiltered small periods at the upper cut}
\label{six:unfiltered}

Let $X=6^k$, $c=\log_6 5$, $\delta=1-c$, and $L=\log X$.
Suppose $R\mid Y\mid X$ are powers of six with $u=\log_X(X/Y)\le .241$,
$.30\le z=\log_X(X/R)\le .479$, and $R\le Y/6$.
Write $\rho_L$ for the mean of the lower $R$-word $F$, whose last
digit is restricted to the allowed units, and $\rho_U=(Y/R)^{-\delta}$.
One may also prescribe the leading digit of $F$. Let $U$ be the upper
word on the index $\lfloor x/R\rfloor\bmod (Y/R)$, or its difference
from $5/6$ times the analogous word of period $Y/(6R)$.
A fixed displacement $O(R)$ of the upper steps is permitted. Let $E_e$
average on residue classes modulo $e$ on the $R$-circle and put
$G_q=\sum_{e\mid q}\mu(q/e)E_eF$.

\paragraph{Proposition.}
Let the $r$ prime variables have scales $P_i\ge X^\kappa$,
where $r,\kappa>0$ are fixed, and $\prod_iP_i=XL^{O(1)}$. Let $W$
be smooth with compact support on fixed-ratio intervals of these scales
and every fixed logarithmic derivative bounded by a power of $L$.
Let $\mathcal B$ be the pushforward density of
$W(\boldsymbol t)\prod_i dt_i/\log t_i$ under multiplication.
For $q\mid R$, $q\le X^{.235}$ and each fixed $J>0$,
\[
 \sum_{\boldsymbol p}W(\boldsymbol p)(G_q U)(\textstyle\prod_i p_i)
 = 3\sum_{(n,6)=1} G_q(n)U(n)\mathcal B(n)
       +O_J(X\rho_L\rho_U L^{-J}).
\]
This covers $.4195\le m\le .4801$, $.1018\le h\le .1020$, $R\asymp X^{m+h}$,
and $Y\ge X^{.760}/6$.

\paragraph{Character coefficients.}
Write $A_R=R\rho_L$ and $b_\chi=A_R^{-1}\sum_{w<R}F(w)\bar\chi(w)$
for primitive conductor $q_0$. On the units the normalized
coefficient of $G_q/\rho_L$ is
\[
 b_\chi\sum_{e\mid q,\ q_0\mid e}\mu(q/e)e/\varphi(e).
\]
Summing over the unit lifts in $E_e$ proves the formula, since each
allowed $w$ is a unit. The multiplier is bounded since there are at
most four nonzero terms and $e/\varphi(e)\le3$. Only $q_0\mid q$ enter.
We use
\[
 \sup_{\chi\ {\rm primitive}\bmod q_0}|b_\chi|\ll q_0^{-.125},
 \qquad \sum_{\chi\ {\rm primitive}\bmod q_0}|b_\chi|^2\ll q_0^\delta .
\]
For $K\le R$, the mass fraction in a residue class modulo $K$ is
$O(K^{-c})$: fix the digits above the largest radix power
$\le K$, leaving at most one choice of the remaining digits.
Parseval gives the square bound. Write $q_0=2^i3^j$,
$s=\lceil\max(i,j)/2\rceil$, $D=6^s$, $Q=q_0/(q_0,D)$.
Then $Q\mid D$, $q_0\mid D^2$, and $Q\le R/D$.
The kernel of reduction of units modulo $q_0$ to units modulo
$(q_0,D)$ is parametrized by $1+Dh$ for $h\bmod Q$, with group
law addition. A primitive character has unit slope $\ell\bmod Q$:
for each $p\mid Q$, the class $1+DQ/p$ generates the kernel at the
last $p$-power (its $p$-adic difference from $1$ has valuation
$v_p(q_0)-1$ and its other prime components are $1$).
For $w=w_0+Dh$ the upper sum, normalized by the number of its words,
has magnitude the amplitude of that word at $\ell\bar w_0/Q$.
One can retain its initial block of length the largest radix power
$B_1\le Q$. The shifted first moment with $a=.43389$ bounds
the sum on this grid by $O(Q^a)$: convolve the trigonometric
polynomial with $2K_{2B_1}-K_{B_1}$, where $K$ denotes the Fej\'er
kernel. The absolute integral is $O(B_1^{a-1})$ by the shifted grid
bound, while the kernel translates on the $Q$-grid sum in absolute
value to $O(B_1)$ pointwise. The argument also bounds a local
supremum in each cell of width $O(1/B_1)$. A prescribed end
digit costs a bounded factor by splitting into six grids.
Inversion permutes units modulo $Q$. By lower occupancy at $Q$
one thus has $|b_\chi|\ll Q^{-(c-a)}$. Apart from bounded rounding,
\[
 \log Q\ \ge\ \min_{p=2,3}
 \frac{\log p}{2\log p+\log(6/p)}\log q_0-O(1)>.27\log q_0-O(1).
\]
Indeed for fixed larger valuation $i_1$ (at $p$) the ratio before
rounding is decreasing and then increasing in the smaller valuation,
with breakpoint $i_1/2$. Since $(.898-.43389)\cdot.27>.125$,
this proves the supremum.

\subsection{Mellin transforms of translated and modulated upper words}
\label{six:upper-mellin}
Set $Z=X/R$, $B=Y/R$. In $v=\log(x/X)$ take
$V_\beta(v)=U(Xe^v)\Psi(e^v)e(-\beta Xe^v)$ with a smooth
cutoff on the product box, where $e(t)=\exp(2\pi i t)$. These
upper-word bounds hold for either cut: $1\le B\le Z$, $2\le
T=X^\tau\le Z^{1.9}$, $w=(\tau-u)_+$, with $\lambda=\log_X R$
in a fixed compact subinterval of $(0,1)$. Write $V=V_0$. Then
\[
 \|\widehat V_\beta\|_{L^2(|t|\le2T)}
      \ll \rho_U L^{O(1)}X^{\delta w/2},\qquad
 \sup_{|t|\asymp T}|\widehat V_\beta(t)|
      \ll \rho_U L^{O(1)}X^{-B_T},
 \quad B_T=\begin{cases}.566w,&\tau\le2u,\\ .232\tau,&\tau>2u .
 \end{cases}
\]
The proof of \eqref{dt:analytic-upper-square} uses only absolute log
mass, so multiplication by $e(-\beta x)$ leaves its kernel bound
unchanged. Here an interval of $H\ge1$ upper indices in a period
contains $O(H^c)$ allowed indices (cover it by two aligned blocks of
length at most $6H$). Counting periods above $Y$ and using bounded
height below $R$ gives the same local mass
$O(\rho_U T^{-1}(1+TY/X)^\delta L^{O(1)})$ on log intervals of
length $1/T$. This gives the stated square bound uniformly in $\beta$.

For \eqref{dt:analytic-upper-point} use aligned powers of six and
\eqref{six:eq-six-separated-sampling} with $a=.43389$. More precisely,
put $Z'=(\prod_i P_i)/R$, which differs from $Z$ by at most log powers.
The block length is comparable to $\min(B,Z'/\sqrt T)$ and is at
least one since $T\le Z^{1.9}$. A fixed upper prefix costs at most one.
The cell count in that proof is uniform in the translation of the
linear frequency by $-\beta R$; its two relative bounds, apart from
log powers, are $(1+TB/Z')^{a-1}$ below $T=(Z'/B)^2$ and
$(B/Z')^\delta T^{-(c-a)/2}$ above it. They agree up to constants at
the transition. Thus $1-a>.566$ and $(c-a)/2>.232$ give the displayed
estimate, retaining the stronger exponent $(c-a)\tau/2+\delta u$
in the second range. If $T\le X|\beta|L^{-C_0}$ with $C_0$ sufficiently
large, the phase derivative within each step is monotone and has size
$\gg R|\beta|$; the first derivative test gives the additional relative
bound $L^{O(1)}/(1+R|\beta|)$. Subtraction gives the signed estimates.
A short log translation changes only the cutoff and rescales $\beta$
by a bounded factor.

\subsection{Fixed-rank prime products and the small-period inequalities}
\label{six:fixed-rank-explicit}
Separate a joint smooth weight by Fourier expansion on a larger fixed
logarithmic box with individual smooth cutoffs. Truncating at $L^{O(1)}$
gives arbitrarily small inverse-log uniform error. Each integer is
an ordered product of $r$ primes in at most $r!$ ways. On an
interval $I$, $|G_q|$ has mass $O(\rho_L(|I|+q))$ by its at most
four positive periodic projections. There are
$O((X/R)\rho_U L^{O(1)})$ occupied steps. This bounds the truncation
error in the digital and reference sums. Smooth the compactly supported
upper log function at width $Z^{-1}X^{-\epsilon}$. The cost at all
jumps divided by $X\rho_L\rho_U$ is
$O((X^{-\epsilon}+q/R)L^{O(1)})$; smooth cutoff changes away from
the jumps are smaller. The transform bounds persist by translation.
Its heights above $ZX^{2\epsilon}$ give arbitrarily large power savings
by repeated differentiation of the smoothing kernel.

For the fixed primitive conductor $q_0$, write $\gamma=\log_X q_0$.
Let $N_{q_0}(1-y,T)$ count zeros of the primitive characters
at real part $\ge1-y$ and $|\Im s|\le T$, with multiplicity.
Thorner and Zaman \cite[(2.1)]{ThornerZaman} give
$N_{q_0}(1-y,T)\ll_\xi(2q_0T)^{(2.4+\xi)y}$.
Chen, Gupta and Li \cite[Theorem~1.2]{ChenGuptaLi2026} give $N_{q_0}(1-y,T)\ll_\xi(q_0T)^{(30/13)y+\xi}$
at fixed $0<y<1/2$, since $q_0$ is $T$-smooth when $T\ge3$.
At $y=1/2$ the total is $O(q_0T\log(2q_0T))$.
The Iwaniec smooth-modulus zero-free region in Thorner and Zaman (3.1)
gives $y\gg L^{-3/4}(\log(2L))^{-3/4}$ at polynomial height,
apart from exceptional real characters, whose primitive conductor here
divides 24. Their finite set of real zeros obeys the same bound for
sufficiently large $X$. For $q_0=1$ use the zeta zero-free region and
the density estimate for the principal character modulo 2.
More generally the same Thorner--Zaman log-free bound applies to
any modulus $f$, uniformly in $f$; the same Chen--Gupta--Li bound
applies whenever $f$ is $T$-smooth, regardless of its prime divisors.
For all $f$ the fixed-modulus Ingham bound
$N_f(1-y,T)\ll_\xi(fT)^{3y/(1+y)+\xi}$ for fixed $0<y<1/2$
follows from \cite[(2.5)]{FortiViola1973}; see also \cite[(1.4)]{ChenGuptaLi2026}.
In Thorner--Zaman (3.1) the general zero-free denominator is
$O(\log(2\,{\rm rad}(f))+[\log(3fT)\log\log(9fT)]^{3/4})$.
If $f\mid 6^k d$ with $d\le L^A$ and fixed $A$, an inducing
exceptional primitive real character has conductor at most $24d$,
since its odd prime exponents are at most one and its 2-exponent is
at most three. Siegel's theorem \cite[Theorem~5.28]{IwaniecKowalski2004} gives $1-\beta\gg_A L^{-1/4}$ for
these real zeros on choosing its exponent $\le 1/[4(A+1)]$.
The fixed-rank kernel estimates below also hold at any
polynomial-size conductor $f$, including after excluding its prime
divisors in the smooth prime sums: their direct cost is $O(L^{O(1)})$.

For $T\ge X^\epsilon$, the smooth prime sum twisted by $\chi(p)p^{it}$
at scale $P_i$ has the usual continuous pole term for the principal
character; on $|t|\le2T$ its remainder is bounded by $L^{O(1)}$ times
\[
 P_i^{1/2}+
 \sum_{\substack{\beta+i\omega\ {\rm zero}\\
     \beta\ge1/2,\ |\omega|\le4T}}
             P_i^\beta(1+|t+\omega|)^{-3}.
\]
Apply the smoothed von Mangoldt explicit formula with the additional
weight $1/\log x$ and integrate by parts in the logarithmic coordinate.
The Mellin weight at a zero decays to arbitrary fixed order. There
are $O(L)$ zeros per unit ordinate interval in the retained range,
higher prime powers cost $O(P_i^{1/2}L^{O(1)})$, and the omitted
high-zero tail is power-small on taking enough derivatives.
In the product of the remainders retain a maximal $\beta$ whenever
there is at least one zero. The size factor is then at most
$(\prod_i P_i)^\beta$, while the other kernels sum to $L^{O(1)}$
pointwise. If there are $N_\chi$ zeros in the layer, the sum of
their kernels has $L^1$ norm $O(N_\chi)$, supremum $O(L)$ and
squared $L^2$ norm $O(LN_\chi)$. A term with no zero has size
$(\prod P_i)^{1/2}L^{O(1)}$.

Combining the kernel estimates with the two character norms and
the two transform estimates, the exponent relative to $X\rho_L\rho_U$
is bounded by the smaller of
\[
 P(y)=-y+D(y)(\gamma+\tau)-.125\gamma-B_T,\qquad
 Q(y)=-y+\tfrac12D(y)(\gamma+\tau)+.051(\gamma+w).
\]
For $Q$ apply Cauchy--Schwarz in the Mellin integral and then over
characters. Use $D(y)=2.4y$ on $0\le y\le.04$ and $D(y)=30y/13$
above $.04$, with the arbitrarily small density losses indicated above.
At $.04$ the coefficient of $\gamma$ in $P$ is negative and
$2.4(.04)<.232,.566$. The first curvature range maximizes its
height part at $\tau=u$; in the second the height part is nonpositive.
Thus $P(.04)<-.016$ also for the smaller density slope $30/13$.
At $y_*=.215$ put $v_*=30y_*/13$. The two curvature bounds for
$P(y_*)$ before the arbitrarily small height reserve are
\[
 -y_*+(v_*-.125)(.235)+v_*(.241)<-.008,\qquad
 -y_*+(v_*-.125)(.235)+(v_*-.232)(.479)<-.0012.
\]
Also $Q(y_*)\le-y_*+(v_*/2+.051)(.235+.479)<-.0014$; its
slope for $y\ge y_*$ is at most $-1+(15/13)(.235+.479)<0$.
The height cutoff allowance $2\epsilon$ and density losses are small
compared to these margins. For fixed parameters $P$ is affine on each
of the two density ranges. Near zero, $P(0)\le0$ and
$P(y)\le-c_1 y$ for an absolute $c_1>0$ even after the arbitrarily
small log-free density slope loss. Use a real-part mesh of $1/L$ here;
moving to the farther layer endpoint costs only a bounded factor,
and the zero-free region supplies arbitrary log savings. Above $.04$,
a sufficiently fine fixed mesh yields a fixed power saving, including
the arbitrarily small density loss at its finitely many endpoints.
The no-zero term is bounded by $-1/2+.551(\gamma+\tau)<-.1$
by the square bounds. On $|t|<X^\epsilon$ truncate at $X^{2\epsilon}$
and use the trivial upper transform. Since $2.4(.235+2\epsilon)<1$,
the log-free density bound suffices for every deficit. The central
no-zero term follows from the absolute character count.

On noncentral shells every mixed pole term has power decay by integration
by parts in its smooth pole factor. On the central interval such terms
use only zeta zeros: a nonempty set of remainder slots with at least
one zero saves $X^{-\kappa y}L^{O(1)}$, and the log-free zero count
at $X^{2\epsilon}$ suffices for $(2.4+\xi)2\epsilon<\kappa/2$.
Products of baselines from remainder slots save $X^{-\kappa/2+O(\epsilon)}$.
The all-pole term integrates $U\mathcal B$ against the unit mean of
$G_q$. On an occupied step the cumulative discrepancy on replacing
this mean by $3{\bf1}_{(n,6)=1}G_q(n)$ is $O(\rho_L\,{\rm lcm}(q,6))$
by periodicity and the absolute-mass bound. Partial summation, using
$\|\mathcal B\|_\infty+X\|\mathcal B'\|_\infty\ll L^{O(1)}$,
costs $O(X\rho_L\rho_U L^{O(1)}(q+6)/R)$ summed over all steps.
The number of inducing conductors and Mellin shells is polylogarithmic.
Restoring the unsmoothed step proves the proposition.
The rational endpoint inequalities in this subsection are also
checked by \path{small_periods/verify_small_period.py}.

\section{Exceptional frequencies at large additive periods}
\subsection{Small difference orders and shortened signed lifts}
\label{six:sixmask-geometry}
Use the notation of the signed Gram estimate. In this subsection
$X=6^k$, $.4229\le m\le .4801$, $n=1-m$, $.067\le h\le .1022$,
$\lambda=m+h+O(.00003)$ with $0\le\lambda-m-h\le .00003$,
and $.76\le\eta\le1$. Write $A=\delta(\eta-\lambda)$, $\delta=1-c$,
$c=\log_6 5$, $a=.434$. Implicit rounding of digit lengths changes
the following exponents by $o(1)$. We allow an arbitrarily small
positive power loss in completion. This subsection concerns periods
$q\ge MX^{-.0001}$ and multipliers bounded by one and constant
under translation by $1/E$ for all orbits used below.
When treating the unrestricted Gram at a larger difference order,
the original amplitude multiplier must also permit the signed lifts
used in that estimate.

Write $g=p_0^i p_1^j$, $\{p_0,p_1\}=\{2,3\}$, $i\ge j$,
$s=i/k$, $d=j/k$, $\xi=\log_Xg$, $G=6^i$,
$V=6^{\lfloor j/2\rfloor}$, $Z=q/g$, $K_0=p_1^{i-j}$,
$d_2=(K_0,Z)$, $D_0=Z/d_2$ and $K=K_0/d_2$.
Then $GD_0=qK$, $(K,D_0)=1$, $D_0\mid R/G$ whenever $G\mid R$.
The ordinary signed Gram bound before the folded incidence has
coarse factor $X^{A-Cd-L_{p_0}(s-d)}\max(1,g/P)^{1/2}$, where
$C=3c/4-a$ and $L_p=(1-2a_p)\log_6 p$; the mixed-grid bounds
$a_2=.151,a_3=.231$ give $C,L_2,L_3>.2395=:\sigma$.
In particular the factor is at most $X^{A-\sigma s}$ for $s\le .305$.

Fix such an order, and write $H=X^h$, $T=(1+q/M)X^\epsilon$,
$\gamma=\log_Xq$, $\vartheta=(\gamma-m)_+$.
Fold the congruence $t+\ell z=0\bmod Z$ with
$e_0=(\ell,d_2)$, $t=e_0t'$, $\ell=e_0\ell'$.
At fixed $z_0\bmod D_0$ the fold has $e_0$ preimages whenever
$D_0\mid t'+\ell' z_0$; this cancels the change from $H$ to $H/e_0$
in the normalized gap count. Its digital square is taken at
$Kz_0/D_0$. For a given nonzero gap, the reduced fraction
$(\ell'z_0+t')/(\ell'D_0)$ has denominator $d_*\mid\ell'$.
Group by $C_*=(d_*,K)$ and dyadic projected denominator
$d_*/C_*\asymp B$. A projected center under multiplication by $K$
has at most $C_*$ preimages for fixed $C_*$.
On a nonzero displacement shell $|t'/\ell'|\asymp V_0$,
the normalized gap multiplicity per original center and grid point
is at most
\[
 (C_*B)^{-1}\min(1,T/(HV_0))
\]
up to an absolute factor. This follows by counting multiples of
$d_*$ up to $\min(H/e_0,T/(e_0V_0))$. The twist is then determined.
For zero displacement omit the minimum. Thus the factors $C_*$
cancel. There are $O((\log X)^C)$ groupings, and $BV_0\ll T$.

When $B^2\max(1,KV_0)\ll D_0$ with a sufficiently small constant,
sample the initial upper-word square on the projected centers spaced
$\gg B^{-2}$, keeping a digit block of size comparable to $B^2$.
Its summed local suprema cost $O(B^{2\delta})$. At each original
center the possible $z_0$ run through an interval of $O(1+V_0)$
integers. A disjoint terminal block at a position $6^j$ where
$D_0/(6^j,D_0)\asymp1+V_0$ costs $O((1+V_0)^\delta)$ in their
summed squares. Indeed successive reductions lose a factor at most
six; the reduced numerator multiplier is a unit, and the available
remaining digit length is at least the size of the arithmetic grid.
This holds also when $V_0<1$ by retaining no terminal digits.
The separated sampling bound follows by reproducing a polynomial of
degree $<B'$ with $2K_{2B'}-K_{B'}$, where $K_N$ is the Fej\'er
kernel, using its majorant $O(B'(1+B'\|t\|)^{-2})$. The digital
square integral of a length-$B'$ block is $O((B')^{-c})$.
These bounds apply on bounded perturbations at the respective grid
spacings. The normalized shell contribution, without the coarse
factor, is
\[
 \ll B^{2\delta-1}(1+V_0)^\delta\min(1,T/(HV_0)).
\]
At zero displacement the last two factors are one.

In the complementary case sum the complete upper square
$O(D_0^\delta)$; for fixed $z_0$ at most $O(1+B^2KV_0/D_0)$
projected centers can occur. The resulting bound is
\[
 \ll X^{.00003+O(\epsilon)+o(1)}D_0^\delta
 \left(\frac1{B\max(1,V_0)}+\frac{KT}{D_0}\right).
\]
At zero displacement only the first term applies. Since
$B\max(1,V_0)\gg(D_0/K)^{1/2}$, the first exponent including the
coarse factor is bounded by
$A-(.5-.102)\gamma+(.5-\sigma)s+.00004+O(\epsilon)+o(1)$:
use $\log_XK\le s$ and $GD_0=qK$. For the second use
$D_0\le R/G$, obtaining
$.102-(m-.0001)+(1-.102-\sigma)s+.00004+O(\epsilon)+o(1)$.
Both are strictly negative through $s=.305$.

On a separated shell put $b=\log_XB$, $v=\log_XV_0$,
$v_+=\max(v,0)$. Remove only shells satisfying
\[
 \sigma s+(1-2\delta)b+(h+v-\vartheta)_+-\delta v_+
 \le A+.00015 .
\]
For zero displacement set both terms involving $v$ to zero.
Every incident label $\alpha$ lies within $O(V_0/q)$ of a rational
with denominator dividing $G d_*/C_*$, since
$G\alpha=Kz_0/D_0\bmod1$. Zero displacement is impossible on
these shells: the indicated denominator is strictly smaller than
$q$, by the budget below, while $\alpha$ has exact order $q$.
Close each removed family under translation by $1/E$, where
$E=\max(G,6^{\lfloor.092k\rfloor})$ and $e=\log_XE$. In reduced
terms the projected center after multiplying by $E$ has denominator
$d_sd_c$, $d_s=2^l3^{l'}$, $(d_c,6)=1$. Put $b_s=\log_Xd_s$,
$b_c=\log_Xd_c$ (dyadic changes are $o(1)$), and $v_0=v_+$.
The following enlarged budget and height bound hold:
\begin{align}
 .2395s+.796(b_s+b_c)+.898v_0&\le .0523,
 \label{six:sixmask-budget}\\
 (h+v-\vartheta)_+&\le
 J:=.0523-.2395s-.796(b_s+b_c)+.102v_0 .
 \label{six:sixmask-offset}
\end{align}
Here $A\le .102(.5101)$ and $\vartheta-h\le .00003+o(1)$.
In particular $s<.219$, $b_s+b_c<.066$,
$e+4(b_s+b_c)+v_0<.356$ and $d_s\mid E^2$. The latter follows
from $.066<2(.386)(.092)$. To bound the two valuations of $q$
use $\gamma\ge m-.0001+o(1)$ and the fact that both are at most
$\lambda k+O(1)$. The smaller lower bound is
$m-.0003-1.586(.1023)>.219$, in units of $k$. Thus $6E\mid q$,
and all translations preserve primitivity. There are only
polylogarithmically many families upon grouping the same order,
$d_s$, dyadic $d_c$, displacement, $e_0$ and $C_*$. Within
any one family assign a covering center consistently along each
orbit. Order the families in decreasing $E$, removing any previously
assigned points from a later family. All assigned sets are still
unions of complete original $E$-orbits.

Here is the required correction when removing these sets from
a signed Gram. Put $W=.434(b_s+v_0)+.204b_c$ for an assigned
family. Its total of absolute \emph{orbit scalars} is
$\ll X^{W+o(1)}$; the proof by sampling is in
Section~\ref{six:sixmask-mellin}. Fix a difference order $g$. First
delete only the families with $E\ge V$, using their combined
complement multiplier on both labels. It is invariant on every
$V$-lift, since these orbit lengths are radix powers. All exceptional
shells of this particular $g$ have been removed: their original
$E\ge G\ge V$ and any prior assignment has even larger orbit length.
Thus for $s\le.305$ the separated and overflow estimates just
proved apply. For larger $s$ the unmasked estimates valid for
$V$-invariant bounded multipliers apply to this restricted Gram.
Now expand deletion of the families with $E<V$ in each of the two
labels. Each correction has either one assigned small family or
two; disjointness makes the expansion literal. For two take the
smaller of their two $E$'s, and designate its label. All remaining
family and complement multipliers in the term are constant on the
resulting $E$-lifts. The directionally reversed term is bounded
by the same argument.

For clarity, repeat the signed lift with length $E$ at this fixed
order. Write the labels as $a_0+(q/E)x,b_0+(q/E)y$. Their coarse
difference is $Zu$ with $u$ a unit modulo $g/E$. Rows $p$ lie in
one progression modulo $g/E$, since
$p u=(t+\ell b_0)/Z\bmod (g/E)$. Here $E\mid g/E$,
and $g/E\ll P$ because $e\ge.092-o(1)$ and
$\lambda-.092<1-m$ throughout the stated range. The ratio
$(p+\ell)/p\bmod E$ is fixed along the progression. The affine
intercept in the lift relation varies in unit steps modulo $E$
because each increment $g/E$ moves it by $-u/p\bmod E$.
The sum over the two low-frequency lifts is a Fourier transform
in the intercept of a sum over at most $O(E^c)$ allowed words,
with scale $O(E^{1-2c})$. Parseval and Cauchy average its absolute
value over the rows at cost
$O(E^{1-3c/2}\max(1,g/P)^{1/2})$, before the row progression
density $(g/E)^{-1}$. A forced end digit changes only constants.
For the designated coarse label the summed remaining amplitude
over \emph{one representative per orbit} costs $X^{W+o(1)}$.
For the other label its coarse difference sweeps a shifted
$(g/E)$-grid; its first moment costs $O((g/E)^a)$ by retaining
a block comparable to $g/E$ and applying separated sampling.
For fixed coarse labels and gap at most $O(1+T/Z)$ twists occur.
After summing gaps and dividing by the usual $PH\rho_U$, the
exponent is therefore bounded by
\[
 A+W-(1-a)\xi+(2-a-3c/2)e
              +(\xi-m)_++\tfrac12(\xi-n)_++O(\epsilon)+o(1).
\]
As $E<V$, $2e\le\xi\le\lambda$ up to $o(1)$. The expression is
convex in $\xi$; at either endpoint it is convex in $m,h,e$ on
their rectangular upper envelopes. The verifier checks all vertices
even with $A=.102(.5101)$ and $W=.434(.0523)/.796$, obtaining a
margin exceeding $.002$. This completes the correction uniformly
at every difference order. For $X^{.219}\le q< MX^{-.0001}$ the
unmasked zero-twist bound $X^{A-\sigma\gamma+o(1)}$ already saves
a power.

\subsection{Arithmetic comparison on the removed orbits}
\label{six:sixmask-mellin}
We prove the smooth prime-product comparison on the assigned
families of Section~\ref{six:sixmask-geometry}. Restore the lower density
and the scalar from the outer conditional expectation only after
this comparison. Write $L=\log X$, $\rho_U=(Y/R)^{-\delta}$,
$z=1-\lambda$, $u=1-\eta$, and $\tau=\log_X T$ on a Mellin
shell $|t|\asymp T$. The error required in this normalization
is $O_J(X\rho_U L^{-J})$. All test functions may carry fixed
powers of $L$ in their derivatives and scale distortions.

For an orbit choose its representative
$\alpha_0=a_0/(E d_sd_c)+\beta$, with $(a_0,d_sd_c)=1$,
$|\beta|\asymp V_0/q$. The orbit scalar $A_{a_0,\beta}$ consists
of the positional factor above $E$ at $E\alpha_0$, times any
bounded multiplier constant along this orbit. Inversion along
$\alpha_0+j/E$ gives the normalized physical expression
\[
 A_{a_0,\beta}\frac{E}{N_E}1_{\mathcal D_E}(w)
 e\!\left(-\frac{a_0}{d_sd_c}\lfloor x/E\rfloor\right)
 e(-\beta x)e(\beta w),\qquad w=x\bmod E .
\]
Here $\mathcal D_E$ includes the unit last digit and
$N_E=|\mathcal D_E|\asymp E^c$. Changing the sign convention
conjugates the expressions. The end digit at the top of the lower
word if prescribed is part of the scalar, as its positional index
is divisible by $E$. One may discard $e(\beta w)$, since
$E|\beta|\le X^{.219+.059-.4228+o(1)}$ and the total absolute
scalar mass costs at most $X^{.029+o(1)}$. On an upper step the
absolute periodic weight has mass $O(R+E)$ even against a bounded
divisor multiplicity coming from a fixed-rank prime tuple.

Write $a_0/(d_sd_c)=a_s/d_s+a_c/d_c\bmod1$ with both numerators
primitive. The complete scalar mass bounds at a fixed smooth
denominator and fixed dyadic $d_c,V_0$ are
\begin{align}
 \sum_{d_c,a_c,a_s,\beta}|A_{a_0,\beta}|
   &\ll d_s^aD_c^{2\delta}(1+V_0)^a\ll X^W,
       \qquad W=.434(b_s+v_0)+.204b_c ,
 \label{six:sixmask-lone}\\
 \sum_{d_c,a_c}\left(\sum_{a_s,\beta}|A_{a_0,\beta}|^2\right)^{1/2}
   &\ll d_s^{\delta/2}D_c^{2\delta}(1+V_0)^{\delta/2}
       \ll X^{.051(b_s+v_0)+.204b_c} .
 \label{six:sixmask-ltwo}
\end{align}
At fixed coprime numerator and denominator use first a block of
size comparable to $d_s$ for the smooth offsets, with first
moment $O(d_s^a)$ or square moment $O(d_s^\delta)$. Then move to
the first radix power divisible by $d_s$. Use two successive blocks
of sizes comparable to $D_c^2$, $d_c\asymp D_c$. In either block the
projected primitive fractions are $\gg D_c^{-2}$-separated, as a
radix power permutes the reduced fractions at each denominator.
Their summed local square suprema cost $O(D_c^{2\delta})$ per
block. Cauchy for the product of these two suprema gives
$O(D_c^{2\delta})$. In the second estimate take these suprema
before the inner square sum. Finally at each original center retain
a disjoint terminal block where $q/(q,6^j)\asymp1+V_0$. The offsets
are spaced by $1/q$ and occupy an interval of length $O(V_0/q)$.
The reduced numerator multiplier is a unit on the terminal grid,
and the block costs $O((1+V_0)^a)$ or $O((1+V_0)^\delta)$,
respectively. These blocks fit since the coprime blocks end at a
positional exponent at most $e+4(b_s+b_c)+o(1)$ (here
$1/\log_6 2<4$), and $e+4(b_s+b_c)+v_0<.356<\gamma$.
This strict inequality also permits taking local suprema at the
perturbed angles on the corresponding grids. These estimates are
uniform under deletion of arbitrary orbits: put their scalars equal
to zero before taking absolute values. Each radix rounding above
changes a physical length by a factor at most six. Consequently these
mass estimates cost at most fixed powers of $L$ even after summing all
the dyadic families. This distinction matters on zero layers approaching
real part one.

The remaining periodic profile before $e(-\beta x)$ is expanded
on unit classes modulo $Ed_sd_c$. At fixed $d_c,a_c$ denote its
character coefficients by $b_{\chi,a_s}$, excluding the scalar.
For any coefficients $f_{a_s}$,
\begin{align}
 \sum_\chi\Big|\sum_{a_s}b_{\chi,a_s}f_{a_s}\Big|^2
    &\ll E^\delta\sum_{a_s}|f_{a_s}|^2, \label{six:sixmask-charsquare}\\
 |b_{\chi,a_s}|&\ll L^C d_c^{-1/2}
             X^{-S(\gamma_s)},\nonumber\\
 S(\gamma_s)&=\max\{
 .208\gamma_s-b_s,\tfrac12 b_s+.398\min(b_s,.386e)\}.
 \label{six:sixmask-charpoint}
\end{align}
Here $X^{\gamma_s}$ is the smooth inducing conductor, hence
$\gamma_s\le e+b_s$. Parseval on the units proves the square
bound: at fixed admissible $w$ the quotient runs through all
residues modulo $d_s$ independently of the coprime component, so
its mean square over the smooth quotient produces the sum of
squares. The normalized bottom profile costs $O(E^\delta)$.
On the coprime component split the quotient exponential into
$e(-a_c\bar E x/d_c)e(a_c\bar E w/d_c)$. The first term gives
a normalized Gauss sum bounded by $O(d_c^{-1/2}L^C)$.
Indeed expanding the unit restriction for an induced character
annihilates the sum unless the complementary part of the modulus
is squarefree and coprime to the inducing modulus, when the
absolute sum equals the square root of the inducing modulus.

For the first pointwise option at a smooth conductor $q'=2^l3^{l'}$
use the linear character fibre after a radix cut at half the larger
valuation. If this largest valuation is $i$ at $p$ and the smaller
one is $j$, set $\kappa_p=\log_6p$. For $j\le i/2$ the first-moment
saving in base-six logarithms is
$(c-a_p)\kappa_p i/2+O(1)$; for $j\ge i/2$ it is
$(c-a)(j-i/2)+(c-a_p)\kappa_p(i-j)+O(1)$.
Indeed on a fixed bottom word $w$ the upper character is additive
with unit slope proportional to $w^{-1}$ on the remaining grid.
Occupancy of bottom words in a class modulo $Q$ is
$O((R_0/Q)^c)$, where $R_0$ is the half cut and $Q\mid R_0$
the slope modulus. The two bounds result by applying the shifted
mixed first moment to this upper grid. They give at least
$.208\log_6 q'+O(1)$ by checking $j=i/2,i$ for $p=2,3$.
Fixing the full smooth quotient at modulus $d_s$ translates the
word. If $d_s\mid E$, the half cut still fits. The upper grid may
require splitting into $Q/(Q,E/R_0)$ shifted subgrids;
this number divides $d_s$ since $R_0Q$ divides $E d_s$ whenever
$R_0\mid E$ and $q'\mid Ed_s$. More explicitly, prime by prime,
$v_p(R_0Q)=\max(v_p(R_0),v_p(q'))$. Each smaller mixed-grid
estimate costs no more than the original one, since its logarithmic
bound is nondecreasing in either valuation: for first moments
this uses $\kappa_p a_p\le a$. Translation and the extra additive
twist shift the labels but preserve occupancy and the shifted first
norm. If $d_s\nmid E$ then $b_s\ge .386e+o(1)$, making
$.208\gamma_s-b_s\le0$. This proves the first option.

For the second, on the kernel $1+Eh\bmod Ed_s$ the coordinate
$h-(E/2)h^2\bmod d_s$ is an additive isomorphism. Multiplication
verifies additivity since $d_s\mid E^2$; it is bijective because
$E/2$ is divisible by $2,3$ for large $k$. At fixed unit $w$ the
quotient sum is thus quadratic, with coefficient
$\pm\ell E/(2w^2)$ and linear coefficient $\pm\ell/w-a_s$.
For a prime of $d_s$ at which $\ell$ is nonunit the sum vanishes
by the primitive numerator. At the other primes its normalized
size is $O(\sqrt{(E,d_s)/d_s})$. This follows by squaring and
using additive orthogonality; periodic shifts by multiples of
$d_s/(E,d_s)$ show its support requires the linear coefficient
to vanish modulo $(E,d_s)$ up to a bounded factor at $2$.
Consequently $w$ occupies only $O(1)$ classes modulo $(E,d_s)$,
at relative cost $O((E,d_s)^{-c})$. Since
$\log_X(E,d_s)\ge\min(b_s,.386 e)+o(1)$, this proves
\eqref{six:sixmask-charpoint}.

Smooth the upper step in $\log(x/X)$ at scale
$X^{-z-W-.00004}$. The combined boundary error is negligible:
on each affected step the normalized absolute profile on the
boundary interval is $O(RX^{-W-.00004}L^C+E)$, and $e+W<\lambda$.
Apply the additive modulation after smoothing. Differentiation
discards heights above $X^{.00004}\max(X^{z+W+.00004},O(XV_0/q))$
to arbitrary power accuracy, hence
\[
 \tau\le z+\max(W,J)+.00032 .
\]
Indeed $1+v-\gamma-z=h+v-\vartheta+O(.00013)+o(1)$ by
$\gamma\ge m-.0001+o(1)$. In particular throughout
$T<(X/R)^{1.9}$, as required by the common upper-step estimate.
For fixed $d_s,d_c,a_c,a_s$, the exponentials at the distinct
$\beta$'s have spacing $1/q\ge1/R$. On each step of length $R$,
the integral of the squared sum is
$O(R\sum_\beta|A_{a_0,\beta}|^2)$ by a smooth majorant on that
step and quadratic decay of its transform. The same bound holds
after the small logarithmic translation in the smoothing.
Summing over occupied steps and taking Mellin Plancherel bounds
the squared transform norm by $O(\rho_U L^C\sum_\beta|A_{a_0,\beta}|^2)$.
More explicitly fix $d_c,a_c$, and set
$f_{a_s}(t)=\sum_\beta A_{a_0,\beta}\widehat V_\beta(t)$, where
unassigned offsets have zero coefficient. If $K_\chi(t)$ is the
sum of the retained zero kernels of the character on the layer, then
Cauchy on the product of counting measure in $\chi$ and Lebesgue
measure in $t$ bounds
\[
 \sum_{\chi}\int
  \left|\sum_{a_s} b_{\chi,a_s}f_{a_s}(t)\right|K_\chi(t)\,dt
 \ll L^C \left(E^\delta\rho_U
     \sum_{a_s,\beta}|A_{a_0,\beta}|^2\right)^{1/2}
      N_{Ed_sd_c}^{\rm all}(1-y,4T)^{1/2}.
\]
Here $N_f^{\rm all}$ sums zeros over every character modulo $f$.
For a finite-width deficit layer one uses its farther endpoint in
the count. Restriction of the sum to any subset of characters only
decreases the two nonnegative Cauchy factors. Each induced
character modulo $f$ occurs once, and the Euler factors removed
from its primitive $L$-function have zeros only on real part zero.
Thus the cited fixed-modulus estimates for all characters apply
directly. Take the density bound uniformly for $d_c\asymp D_c$
and then sum these square roots over $d_c,a_c$ using
\eqref{six:sixmask-ltwo}. No sum of zero counts over $d_c$ is involved.
For the pointwise option instead use \eqref{six:sixmask-charpoint},
\eqref{six:sixmask-lone}, and the common upper-step pointwise bound.
In this option first fix the smooth inducing conductor
$q_s=X^{\gamma_s}\mid Ed_s$; there are $O(L^2)$ choices. At each fixed
$d_c$ the corresponding characters inject into the characters modulo
$q_s d_c$, by their smooth primitive part and their coprime component.
Their zeros in the strip in use agree with those of the images.
Consequently $N_{q_sd_c}^{\rm all}$ bounds the sum of their kernel
integrals. Its bound is uniform for $d_c\asymp D_c$ and can be taken
outside \eqref{six:sixmask-lone}. This also avoids counting separately
the divisors of a growing $d_c$ near real part one.

For a zero layer of deficit $y$, suppose $D$ is an admissible
zero-count exponent at the conductor and height in question,
as in the common fixed-rank explicit formula. Up to arbitrarily
small fixed losses on layers away from zero, the normalized
pointwise and square exponents are respectively
\begin{align}
 P(y)&=-y+D(\gamma_s+b_c+\tau)-B_T+W-.5b_c-S(\gamma_s),
   \label{six:sixmask-point}\\
 Q(y)&=-y+\tfrac12D(e+b_s+b_c+\tau)
       +.051(e+b_s+v_0+\eta-\lambda)+.204b_c.
 \label{six:sixmask-square}
\end{align}
The upper pointwise savings are $B_T=(1-a)(\tau-u)_+$ for
$\tau\le2u$ and $B_T=.232\tau+.101u$ otherwise (one may
weaken the true positive coefficient of $u$). Using the second
bound in its range we may discard $.101u$. For the square
exponent use $\eta\le1$. Both dominant upper envelopes now
use $z\le .5101$. In the first curvature range maximize the
height part via
\[
 D\tau-(1-a)(\tau-u)_+\le D(.240)+
                    (D-(1-a))_+(.240).
\]
For $\tau\le.07$ one can instead disregard $B_T$ and use a
height exponent $.07$. For fixed positive deficits under
consideration $D>.208$ and the pointwise expression increases
with $\gamma_s$, so put $\gamma_s=e+b_s$ in this calculation.

Here are the exact checks and the choice of zero-density estimate.
For $y=.18$, use the log-free Thorner--Zaman slope $2.401y$
on all heights to check a strict pointwise margin. For
$.18\le y\le .24$ the first curvature regime and all $\tau\le.07$
are covered pointwise with the same slope, checking both deficit
endpoints (the envelopes are convex in $y$).
For $\tau\ge .07$ every inducing modulus divides $E d_sd_c$
and is $T$-smooth, because $d_c<X^{.066+o(1)}$.
The Chen--Gupta--Li fixed-modulus theorem therefore supplies
$D=(30/13)y$, with an arbitrarily small fixed exponent loss.
In the second curvature range use either \eqref{six:sixmask-point}
or \eqref{six:sixmask-square}. With
\[
 t=.5101+.00032+\max(J,W),\quad H=e+b_s+b_c+t,\quad
 S=S(e+b_s),
\]
their envelopes are
\[
 -y+\tfrac{30}{13}yH-.232t+W-.5b_c-S,\qquad
 -y+\tfrac{15}{13}yH+.051(e+b_s+v_0+.5101)+.204b_c .
\]
It is legitimate to maximize the pointwise height at $t$,
since $(30/13)y>.232$ on these deficits.
The file \texttt{positive/verify\_mask.py} enumerates the affine
regions of the four choices in $e=\max(s,.092)$, $\max(J,W)$,
$\min(b_s,.386e)$ and $S$. It solves the four-variable rational
facet systems of \eqref{six:sixmask-budget} and those choices and
checks the feasible intersections. There are 122 vertex records
including repeats. For each of the 600 deficit intervals of width
$.0001$ in $[.18,.24]$, the program uses the larger endpoint in
the density and the smaller in $-y$. It intersects the rational
intervals of $\theta\in[0,1]$ for which the bound
$\theta P+(1-\theta)Q\le -.0003$ holds at each vertex. The
intersections are nonempty. This proves the same inequalities in
every affine cell and allows a sufficiently small fixed real-part
mesh, height loss and density exponent loss.
At $.24$ the square exponent alone is strictly negative; it
decreases thereafter because $H<.729<13/15$. On $\tau\le.07$
this also holds by the square inequality with slope $2.401y$
and height at most $.07$. Terms without a zero use the trivial
count $O(qT L^C)$ at $y=1/2$, separately verified.

For $y\le.18$ retain the uniform $2.401y$ density slope. At $y=0$,
$S(\gamma_s)\ge b_s/2$ and the Gauss coefficient saves $b_c/2$,
leaving savings at least $.066 b_s+.296b_c$, apart from the
offset cost $.434v_0$. When $V_0\ge1$, the nonstationary upper-step
bound $O(\rho_U(1+R|\beta|)^{-1}L^C)$ pays this cost on
$T\ll X|\beta|L^{-C}$. Otherwise $\tau\ge z-o(1)$
up to logarithms: $q\le R$, $V_0\ge1$. Both
$.232z$ and $(1-a)(z-.240)$ exceed $.09$, whereas
$.434v_0<.026$. When $V_0<1$ the offset cost vanishes.
All these bounds apply at fixed shell, conductor and parameter
values; use at $y=0$ the better of the curvature and nonstationary
bounds, and at $.18$ the strict pointwise bound. The same maximum
of the two savings may be retained at both endpoints, so affine
interpolation gives an error factor
$X^{-c_0 y}$ with an additional $d_c^{-c_1}$ on $y\le .09$.
To sum the layers here use partial summation with the log-free
bound, avoiding any fixed real-part mesh next to $0$. For
$d_c>L^B$ take $B$ sufficiently large to absorb the fixed
logarithms. The only inducing divisors separated in this pointwise
estimate are the $O(L^2)$ possible smooth divisors $q_s$.
For $d_c\le L^B$ use the Iwaniec zero-free region for nonexceptional
zeros of conductors dividing $6^k d_c$: $\log\operatorname{rad}(q')
=O_B(\log L)$, so the $X^{-c_0y}$ factor beats every log power.
If a real zero is exceptional, a real primitive character has
odd-prime exponents at most one and 2-exponent at most three.
Its conductor is $O(d_c)$; Siegel's theorem then gives
$y\gg_B L^{-1/4}$ by a suitable choice of its ineffective exponent.
For the zeta function use the Vinogradov--Korobov region.
For the entire central interval $|t|<X^\epsilon$ take the zero
cutoff $X^{2\epsilon}$. When $V_0\ge1$, $X|\beta|\gg X/R$
so the nonstationary estimate gives a factor $O((1+V_0)^{-1}L^C)$;
this same bound follows trivially when $V_0<1$. Therefore the
relative pointwise exponent on this interval is at most
\[
 -y+2.401y(e+b_s+b_c+2\epsilon)
       -.066b_s-.296b_c-.566v_0 .
\]
As $e+b_s+b_c<.219$, this is $\le-.4y-.296b_c$ for sufficiently
small fixed $\epsilon$, for every $y\in[0,1/2]$. The logarithmic
layer summation just described applies at these heights as well.
The terms without retained zeros have a fixed power saving here from
the factor $X^{-1/2}$, even using the trivial character count
$O(Ed_sd_c)$ and the absolute scalar sum.

For completeness, the use of unit-group characters at the coprime
denominator allows slots occupied by a prime dividing $d_c$ to be
treated separately. Fix exactly these slots and their prime values.
There are boundedly many possibilities for each $d_c$, since
every slot has scale $\ge X^\kappa$. Let their product be $a'>1$,
let $g'=(d_c,a')$, and put $d_c'=d_c/g'$. The residual product
retains the exclusions $p\nmid d_c$ and the coprime exponential
has primitive numerator modulo $d_c'$. The upper transform
continues to use the original physical scale $X$. The pointwise
Gauss factor becomes $(d_c')^{-1/2}$, paid by the normalization
$(a')^{-(1-y)}$, since $g'\le a'$ and $y\le1/2$.
To see the latter normalization directly, the product of the remaining
scales is $(X/a')L^{O(1)}$, and the retained rightmost zero contributes
at most $(X/a')^{1-y}L^{O(1)}$. The extra Mellin twist
$(a')^{it}$ has absolute value one and is common to every $a_s,\beta$
for the fixed choices. The smooth coefficient (including the original
coprime twist in $w$) changes by a unit character phase when evaluated
at $a'x'$. Thus its square operator bound holds, with zero-count
modulus $Ed_sd_c'$; discarding the normalization only enlarges this
bound. The at most $O_{\kappa,r}(1)$ choices for each $d_c$ may be
summed before the same scalar first or mixed square summations over
$d_c,a_c$. If all slots have been fixed, the
absolute bounds $X^{.029+o(1)}E^\delta$ suffice directly.

Principal poles have zero smooth coefficient for $d_s>1$,
by averaging over the smooth quotient. For $d_s=1$ the normalized
coprime coefficient is $O(1/\varphi(d_c))$, or
$O(1/\varphi(d_c'))$ in the fixed-slot case, where
$(a')^{-1}/\varphi(d_c')\ll L^C/d_c$. At central heights the
upper integral costs $O(\rho_U(1+V_0)^{-1}L^C)$ by the
nonstationary estimate. The combined principal cost over each
family is thus $O(D_c^{2\delta-1}(1+V_0)^{a-1}L^C)$.
Mixed pole and remainder terms on the central interval use the
ordinary zeta zero bound of the common fixed-rank reduction;
on large shells a pole gives arbitrary decay. Any fixed-slot
all-pole case saves a power because $d_c\gg X^{\kappa+o(1)}$.

For the remaining all-pole case compare its periodic unit mean
with the integer-tested continuous model restricted to units of $6$.
On each step average over periods $Ed_sd_c$ against the modulated
smooth density. The two normalized absolute periodic profiles have
mass $O(|I|+Ed_sd_c)$ on an interval $I$. Partial summation thus
costs $O(Ed_sd_c(R^{-1}+|\beta|)L^C)$ per orbit, relative to
the upper occupied mass. This is power-small after aggregation,
since $e+b_s+b_c<.219$, $|\beta|\le X^{.059-.4228+o(1)}$ and
$W<.029$. Both means vanish for $d_s>1$, and they agree for
$d_s=d_c=1$. When $d_s=1<d_c$ the integer-test mean vanishes by
additive orthogonality. The remaining all-pole terms are negligible
for $d_c>L^B$ by the preceding principal cost. For $1<d_c\le L^B$
the scalar supplies an additional
$\exp(-c_2\log X/\log(2d_c))$: retain a sufficiently short fixed
positive fraction of the digits above $E$. Their argument starts
at $a_c/d_c+E\beta$ and its perturbation remains power-small over
this block. A nonzero residue modulo $d_c$ visits distance at
least $1/14$ from an integer every $O(\log(2d_c))$ radix shifts,
since below $1/14$ consecutive small distances grow sixfold.
The digit polynomial is uniformly contracted there because the
five allowed digits include consecutive ones. Long before a
possible forced top digit this supplies the asserted exponential
bound for each orbit. At fixed $d_c$ there are
$O(d_c^2(1+V_0))$ possible orbits, and the factor
$(1+V_0)^{-1}$ pays their offsets. This proves the comparison
on all assigned mask families.

\subsection{Reconstruction on the prime box}
\begin{proof}[Conclusion of Proposition~\ref{six:box}]
The initial projection is \eqref{six:eq-six-low-model}. At every succeeding
level evaluate each filtered piece at its own cut. For $q\le X^{.225}$
the small-period filter estimate bounds every lower-amplitude piece
absolutely, in both models. At the highest-amplitude cut apply that
same estimate to every complementary filter \emph{at this cut}; the
partition on the common frequency circle permits replacing the
highest multiplier by one. The proposition of Section~\ref{six:unfiltered}
gives the comparison for the resulting unfiltered term. The high cut has
$1-\log_XR<.479$ for large $k$.
For $X^{.225}<q<MX^{-.0001}$ only the zero twist contributes at
polynomial accuracy. Lemma~\ref{six:dt-gram-middle} and the diagonal
estimate of Section~\ref{six:adaptive} give a power saving.
Indeed $\rho_U^{-1}\le X^{.102(1-.4229-.077)+o(1)}$ while
$.225(.2395)>.0538$. At $q\ge MX^{-.0001}$ use
Section~\ref{six:sixmask-geometry} for the complementary sum: it
combines its small-order estimates, the large-order estimates of
Section~\ref{six:adaptive}, and the corrections on shorter lifts.
The diagonal is bounded with any further label deletions by
nonnegativity of its squared mass. Every removed family is compared
in Section~\ref{six:sixmask-mellin}.

The reference for each medium-period term and each high-period
complement is itself negligible. For its fixed upper step the lower
component $G$ is a sum of exact period-$q$ modes, multiplied by
bounded frequency weights. Its mean against the indicator of units
modulo six over the $R$-period vanishes, as the latter has Fourier
denominators dividing six, whereas $q>6$. The sum of $|G|$ over this
period is $O(R\rho_L q^{\delta/2})$ by the arbitrary-grid square
estimate and Cauchy--Schwarz. Replace $\mathcal B$ by its value at
the left endpoint of the step. In units of $X\rho_U\rho_L$ the
resulting error is $O((R/X) q^{\delta/2} L^{O(1)})$: use the smooth
derivative, including across the box boundary, then sum over the
occupied steps. It is power-small for $q\le R\ll X^{.583}$.

There are only $O(k^2)$ divisors at any cut, a fixed number of
filters and upper pieces, and $O(k)$ outer levels. Multiply by the
outer scalar $(X/Y)^{-\delta}$, enlarge the available logarithmic
saving by a fixed amount, and add the linear comparisons. The exact
partition of the ambient increment and the telescope of conditional
averages recover $f_X$ on both sides. This proves the proposition.
\end{proof}

\section{Application of the positive sieve in base six}\label{six:positive}
Apply Theorem~\ref{common:positive-sieve} with $B=6$, $s=5$,
$Q=5^k$, $L=\log X$, $u=\#(\{1,5\}\setminus\{d\})$ and
$\gamma=3u/5$. The complete word supported on units has mass
$uQ/5=\gamma Q/3$. Take $T=.55$, $\alpha=.423$, $\beta=.480$:
$T>1-\beta$, $c>.898>1-T/2$ and $1/3<\alpha<\beta<1/2$.
The interval remainder with divisor weights is \eqref{six:interval-type-i}.
For the flat weight $(\gamma Q/X)\ind_{(n,6)=1}$ on $[1,X)$ the
interval remainder for each modulus coprime to six is $O(Q/X)$ by
counting the two unit classes. Proposition~\ref{six:box} supplies the
comparison with room at both endpoints: $[.423,.480]\subset(.4229,.4801)$.
On product cutoffs supported in $[X L^{-B_0},X)$ the periodic word agrees
with its zero extension. Here $B_0$ is the fixed logarithmic cutoff
parameter, not the radix.

The $[.423,.480]$ calculation in Section~\ref{dt:sieve-certificate}
includes all ranks through seventeen. Hence the complete-word lower limit, in units of $Q/L$, is at least
$\gamma(1-.985573022923228)$. Removing permitted padding zeros
preserves admissibility. Passing from $6^k\le x<6^{k+1}$ to real
endpoints costs at most the factor $5$, so
\[
 \liminf_{x\to\infty}\frac{\pi_{6,d}(x)\log x}{x^c}
 \ge \frac{3u}{25}(1-.985573022923228)=.00173123724921264\,u.
\]
This exceeds $1/600$ even for $u=1$. The upper sieve and the maximum
of the six (possibly ineffective) thresholds finish Theorem~\ref{six:main}.

\clearpage
\part{Recovering the asymptotic in bases ten and eleven}\label{part:exact}
Here a narrower loss estimate is not the objective. To identify the
leading constant we must compare all the terms that isolate primes.
We construct a common model for the linear and bilinear estimates and
verify the hypotheses of the Ford--Maynard sieve for that model.

\section{Digital hypotheses and affine correlations}
\label{dec:section}

The asymptotic sieve needs a short-factor interval extending strictly
below $1/4$. A shifted $3/2$ moment controls the mass of a digital block;
an affine correlation controls the overlap of two such blocks when the
order of the digit splitting is reversed. Together they supply the missing
part of the interval in Section~\ref{dec:bilinear-section}. The Type I
and Type II estimates must compare with the \emph{same} nonnegative
sequence. Section~\ref{dec:prime-section} constructs it by averaging the
middle digits while keeping both ends of the suffix fixed.

Throughout these three sections, fix $b\in\{10,11\}$ and a forbidden
digit $d$. Write $\mathcal D_k=\mathcal W_k$ for the suffix words of
Section~\ref{sec:preliminaries}, and put
\[
 q=b^k,\qquad A=(b-1)^k=q^\alpha,\qquad
 \alpha=c_b=\frac{\log(b-1)}{\log b},\qquad f=f_{b,d}.
\]
We use $\widehat u(t)=\sum_nu(n)e(nt)$; for a function on
$\mathbb Z/q\mathbb Z$, the sum runs over $0\le n<q$.
Thus $F_q(t)=A^{-1}|\widehat{1_{\mathcal D_k}}(t)|$.
Write $e_q(t)=e(t/q)$. All exponents are fixed before $q\to\infty$.
In estimates with a fixed power saving, $q^{o(1)}$ may absorb divisor
bounds and $\exp(O(\sqrt{\log q}))$. In the classical rank-two estimate
the saving at the original arc can be only logarithmic: there we absorb
the representation multiplicity at the witness scale, before returning
to the original denominator.

\subsection{A criterion for the exact asymptotic}
\label{dec:inputs}

The quantitative argument uses the parameters
\begin{equation}\label{dec:constants}
 \gamma_0=\frac{27}{77},\quad
 p_0=\frac{235}{154},\quad g_0=\frac{59}{433},\quad
 \varepsilon_0=\frac{23}{500},\quad g_4=\frac3{5000},\quad
 \kappa=\frac{23}{80}.
\end{equation}
The exponents $\gamma_0=27/77$ and
$(p_0,g_0)=(235/154,59/433)$ are Maynard's decimal inputs
\cite[Lemmas 10.3--10.4]{Maynard2019}.
The additional $3/2$ moment is paired with an affine-correlation
exponent $\delta$; both will enter the reverse digit decomposition.
For every $Y=b^j$ with $j\ge0$, the moment estimates are
\begin{align}
 \sup_x\sum_{a<Y}F_Y((a+x)/Y)&\ll Y^{\gamma_0},\label{dec:moment1}\\
 \sup_x\sum_{a<Y}F_Y((a+x)/Y)^{p_0}&\ll Y^{g_0},\label{dec:momentold}\\
 \sup_x\sum_{a<Y}F_Y((a+x)/Y)^{3/2}&\ll Y^{g},\label{dec:momentnew}\\
 \sum_{a<Y}F_Y((a+x)/Y)^2&=Y^{1-\alpha},\qquad
 \sup_x\sum_{a<Y}F_Y((a+x)/Y)^4\ll Y^{g_4}.
 \label{dec:evenmoments}
\end{align}
where
\begin{equation}\label{dec:groups}
 (g,\delta)=
 \begin{cases}
 (14387/100000, 53/10000),&b=10, d\in\{4,5\},\\
 (287/2000, 1/250),&\text{otherwise}.
 \end{cases}
\end{equation}
The second and fourth moments in \eqref{dec:evenmoments} follow from
\eqref{prelim:even}. The scalar inequalities
$1/100<1-\alpha<\varepsilon_0$ and $\alpha>20/21$ hold in both bases;
the exact comparisons are recorded in Appendix~\ref{app:exact-inputs}.

For the digit laws, put
\[
 P_i(x)=\frac1{b(b-1)}
          \left|\sum_{\substack{0\le j<b\\j\ne d}}e(j(i+x)/b)\right|^2,
 \qquad \sum_{i=0}^{b-1}P_i(x)=1,
\]
and define
\begin{align*}
 H(x,E)&=\sum_{i=0}^{b-1}\sqrt{P_i(x)P_i(x+E)},
 &\rho(E)&=\sup_{0\le x\le1}H(x,E),\\
 K_t&=\sup_{\substack{x\in[0,1],\,c\in[0,b]\\a\in\{1,\ldots,b-1\}}}
 \sum_{i=0}^{b-1}P_i(x)\rho(ai+c)^t.
\end{align*}
The function $H$ is one-periodic in $x$, while $\rho$ is even and
$b$-periodic. 
The exponent $\delta$ measures the contraction of affine correlations.
The following criterion separates that input from the bilinear and sieve
arguments.

\begin{theorem}[Digital criterion for bases ten and eleven]
\label{dec:digital-criterion}\label{dec:asymptotic}
For $b\in\{10,11\}$ and a forbidden digit $d$, take the pair
$(g,\delta)$ in \eqref{dec:groups}. Assume that
\eqref{dec:moment1}--\eqref{dec:momentnew} hold uniformly in the shift
and in every power $Y$ of $b$, and that
\begin{equation}\label{dec:one-digit-hypothesis}
 K_2\le b^{-3\delta}.
\end{equation}
Then, for ordinary base-$b$ expansions and arbitrary real $X\to\infty$,
\[
 \pi_{b,d}(X)\sim\kappa_{b,d}\frac{A_{b,d}(X)}{\log X},
 \qquad
 \kappa_{b,d}=\frac{b\bigl(\varphi(b)-\mathbf1_{(d,b)=1}\bigr)}
                    {(b-1)\varphi(b)}.
\]
\end{theorem}

For the decimal central digits, the stronger affine contraction
compensates for the slightly weaker $3/2$ moment. The estimates use
each pair in \eqref{dec:groups} jointly.
In addition to the stated hypotheses, the proof uses the
digital product structure, the elementary even moments, and the
one-sided shift comparison of Lemma~\ref{prelim:shift}.
The latter is proved for these deleted-digit products in
Appendix~\ref{app:shifts}.

\begin{proposition}[Verification of the digital hypotheses]
\label{dec:verified-inputs}\label{eleven:inputs}
For every deleted digit in each of bases ten and eleven,
\eqref{dec:moment1}--\eqref{dec:momentnew} and
\eqref{dec:one-digit-hypothesis} hold with the pair
$(g,\delta)$ in \eqref{dec:groups}.
\end{proposition}

Appendix~\ref{app:exact-inputs} proves this proposition by continuous
positive-operator inequalities and one-digit correlation bounds, with
exact rational certificates. Theorem~\ref{dec:digital-criterion} is
proved in Section~\ref{dec:prime-section}; together these statements
will prove Theorem~\ref{main:exact}.

\subsection{Two-ended interpolation on probability laws}

A row-by-row Hellinger estimate produces a product along a digital path.
The next lemma converts bounds for that product under both endpoint laws
into a bound for the complete correlation.

\begin{lemma}[Two-sided probability interpolation]\label{dec:twoended}
Let $P,Q,R$ be probability laws on a finite set, and let $0\le U\le1$.
Suppose that $\sqrt{P(\omega)Q(\omega)}\le R(\omega)U(\omega)$ and,
for some $t>0$,
\[
 \mathbb E_P U^t\le\eta,\qquad \mathbb E_Q U^t\le\eta.
\]
Then
\[
 \sum_\omega\sqrt{P(\omega)Q(\omega)}
 \le\eta^{1/(1+t)}.
\]
\end{lemma}
\begin{proof}
Put $a=1/(1+t)$ and $w=\sqrt{PQ}$. The pointwise inequalities imply
\[
 w=w^a w^{1-a}
 \le(PU^t)^{a/2}(QU^t)^{a/2}R^{1-a},
\]
because $at=1-a$. H\"older's inequality with exponents
$2/a,2/a,1/(1-a)$ gives the assertion.
\end{proof}

\subsection{Affine contraction with exact carries}

An affine map transports the successive digit laws, so its correlation
can be evaluated on the two paths simultaneously. The number of levels
that contribute a gain is determined by the first digit at which the
slope differs from $1$.

For $B=b^L$, put
\[
 \mu_{B,x}(z)=((b-1)/b)^L F_B((z+x)/B)^2,
 \qquad z\in\mathbb Z/B\mathbb Z.
\]
This is a probability law whose conditional digit probabilities, read from
the least significant digit upwards, are $P_i(x_j)$, with
$x_0=x$ and $x_{j+1}=(i_j+x_j)/b$.
For a unit $\lambda$ modulo $B$ and an integer $c$, write
\[
 \mathcal H_L(\lambda,c;x,y)
 =\sum_{z\bmod B}\sqrt{\mu_{B,x}(z)\mu_{B,y}(\lambda z+c)}.
\]
We use $v_b(n)=\max\{j:b^j\mid n\}$, with $v_b(0)=\infty$;
this definition also applies when $b$ is composite.

\begin{lemma}[Affine Hellinger contraction]\label{dec:hellinger}
Let $\lambda$ be a unit modulo $B$ with $\lambda\equiv1\pmod b$, and let
$m=\min(L,v_b(\lambda-1))$. For every integer $c$, every $x,y\in[0,1]$,
and every $t>0$,
\begin{equation}\label{dec:hellingerbound}
 \sum_{z\bmod B}\sqrt{\mu_{B,x}(z)\mu_{B,y}(\lambda z+c)}
 \le K_t^{(L-m)/(1+t)}.
\end{equation}
In particular, the numerical input $K_2\le b^{-3\delta}$ gives the
upper bound $(B/b^m)^{-\delta}$.
\end{lemma}
\begin{proof}
For $m=L$, apply Cauchy's inequality. Otherwise write
$\lambda-1=b^ma$, where $b\nmid a$. Along the source path, write
\[
 r_j=z\bmod b^j,\quad r'_j=(\lambda r_j+c)\bmod b^j,
 \quad C_j=(\lambda r_j+c-r'_j)/b^j.
\]
At digit $j$, the permutation is $i\mapsto i+C_j\pmod b$.
Its Hellinger row sum is at most $\rho(E_j)$, where
\[
 E_j=C_j+\frac{r'_j+y-r_j-x}{b^j}
 =\frac{(\lambda-1)r_j+c+y-x}{b^j}.
\]
When $j=m+h$, we have
\[
 E_{m+h}\equiv ai_h+\frac{ar_h}{b^h}
                  +\frac{c+y-x}{b^{m+h}}\pmod b.
\]
Since $m\ge1$, this congruence determines $E_{m+h}$ modulo $b$
from $i_0,\ldots,i_h$. Given the preceding digits, averaging its
factor $\rho(E_{m+h})^t$ over $i_h$ costs at most $K_t$.
Conditioning in decreasing order of $h$ gives
\begin{equation}\label{dec:sourceexpect}
 \mathbb E_{\mathrm{src}}\prod_{h=0}^{L-m-1}\rho(E_{m+h})^t
 \le K_t^{L-m}.
\end{equation}
For the reverse calculation, choose integers $\lambda',c'$ with
$\lambda'\lambda\equiv1\pmod B$ and $c'\equiv-\lambda'c\pmod B$.
The first nontrivial digit of $\lambda'$ is again $m$.
At the corresponding source and destination prefixes, the reverse error
$E'_j$ satisfies
\[
 E'_j+E_j
 =(1-\lambda')C_j+
 \frac{(\lambda'\lambda-1)r_j+c'+\lambda'c}{b^j}
 \in b\mathbb Z\qquad(j<L).
\]
This identity includes the real initial states $x,y$.
Since $\rho$ is even and $b$-periodic, the reverse conditioning
argument gives \eqref{dec:sourceexpect} under the destination law
with exactly the same product.

To finish, set $P(z)=\mu_{B,x}(z)$,
$Q(z)=\mu_{B,y}(\lambda z+c)$ and
$U(z)=\prod_{h=0}^{L-m-1}\rho(E_{m+h})$.
Normalize each row of the geometric-mean transition kernel by its
row sum $H_j(z)$. This produces a probability law $\mathbb R$ on
paths with
\[
 \sqrt{P(z)Q(z)}=\mathbb R(z)\prod_{j=0}^{L-1}H_j(z)
 \le\mathbb R(z)U(z).
\]
A zero row may be completed to any probability row; the identity
remains valid because its unnormalized path weight is zero.
Apply Lemma~\ref{dec:twoended} with $\eta=K_t^{L-m}$.
\end{proof}


\section{Bilinear cancellation below one quarter}
\label{dec:bilinear-section}

We assume the digital hypotheses of Theorem~\ref{dec:digital-criterion}.
The objective of this section is the following minor-arc estimate for
arbitrary bounded coefficients. The proof combines the affine contraction
with rational-grid estimates and the classification of large bilinear
sums developed by Maynard \cite[Sections 10, 13--15]{Maynard2019}.

For $H\ge1$, let $\mathfrak M(q;H)$ be the set of indices
$a\in\mathbb Z/q\mathbb Z$ for which there are integers $c,r$ with
$1\le r\le H$ and $(c,r)=1$ such that
\[
 \left\|\frac aq-\frac cr\right\|\le\frac Hq.
\]
For a set $\mathfrak M$ of frequency indices define
\[
 (P_{\mathfrak M}u)(n)
 =\frac1q\sum_{a\in\mathfrak M}\widehat u(a/q)e_q(-an).
\]

\begin{proposition}[Minor-arc bilinear estimate]\label{dec:minorarc}
Fix an admissible positive prefix $v$ and $T>0$.
There is $D_0>0$ such that, with
$\mathfrak M=\mathfrak M(q;(\log q)^{D_0})$, one has
\begin{equation}\label{dec:minorarcII}
 \begin{split}
 &\sum_{m\in I_M,r\in I_R}\alpha_m\beta_r
 1_{\{vq\le mr<(v+1)q\}}1_{(mr,b)=1}
 ((I-P_{\mathfrak M})1_{\mathcal D_k})(mr-vq)\\
 &\hspace{35mm}\ll_{v,T}A\log^{-T}q
 \end{split}
\end{equation}
whenever $MR\asymp_vq$, $q^{0.249999}\le R\le q^{0.4244}$,
$I_M\subset[M,2M]$ and $I_R\subset[R,2R]$ are intervals, and
$|\alpha_m|,|\beta_r|\le1$.
\end{proposition}

The three ranges in \eqref{dec:rawwindow} correspond to different
bilinear estimates. For intermediate amplitudes, the positive kernel
below separates into a low-digit correlation and a high-digit
arithmetic estimate. Reversing the digit splitting provides the
correlation gain required for short factors below exponent $1/4$.

\subsection{The positive kernel and short integer products}
\label{dec:shortproduct}

For a nonnegative frequency weight $w$, define
$W=\sum w$, $V_w=\sum w^2$ and $|u|_q=q\|u/q\|$, and put
\[
 \mathcal N_q(w;R)=\sum_{r,s\in I_R}\sum_{a,c\bmod q}w(a)w(c)
 1_{\{|ra-sc|_q\le C_0R\}},\quad I_R\subset[c_1R,c_2R].
\]
The multiplier set $I_R$ may be restricted, in particular to units
modulo $b$.

For arbitrary $|\alpha_r|,|\beta_m|,|\gamma_a|\le1$ and
$M\ll_vq/R$, Cauchy's inequality in the long variable gives
\begin{align}
 \left|\sum_aw(a)\gamma_a\sum_{r,m}\alpha_r\beta_m e_q(-arm)\right|^2
 &\ll_v\frac qR\mathcal K_q(w;R),\label{dec:cauchy}\\
 \mathcal K_q(w;R)&=\sum_{a,c,r,s}w(a)w(c)
 \min(q/R,\|(ra-sc)/q\|^{-1}).\notag
\end{align}
Partition the output into cells of length $R$, and apply cyclic Cauchy
and harmonic summation to the distances between cells. This gives
\[
 \mathcal K_q(w;R)\ll(q/R)(\log q)\mathcal N_q(w;R).
\]
Thus the intermediate-amplitude bound
$\mathcal N_q\ll R^2q^{-c}$ bounds the original Fourier pairing
by $O_v(q^{1-c/2}\sqrt{\log q})$.
Returning to the digital indicator introduces the factor $A/q$.
The following observation permits decompositions of the source weight
in the complete output norm, including all cross terms.

\begin{lemma}[Positive output norms]\label{dec:outputnorm}
Let $1\le R\le q$. For nonnegative weights
$w,w_1,\ldots,w_J$ on $\mathbb Z/q\mathbb Z$,
if $w\le\sum_{j=1}^Jw_j$, then
\[
 \mathcal N_q(w;R)\ll J\sum_{j=1}^J\mathcal N_q(w_j;R).
\]
The same assertion holds for $\mathcal K_q$ in \eqref{dec:cauchy}.
The constants depend only on the fixed window constants.
Both statements concern the complete pair sum over the chosen
multiplier set $I_R$.
\end{lemma}
\begin{proof}
Use the actual output mass
$\mu_w(z)=\sum_{r,a:ra=z\bmod q}w(a)$.
Partition the circle into consecutive cells of length comparable to $R$
and at most $C_0R$. Then
\[
 E_R(w)=\sum_I\left(\sum_{z\in I}\mu_w(z)\right)^2
 \ll\mathcal N_q(w;R)\ll E_R(w).
\]
The first inequality counts pairs in the same cell. For the second,
a pair at distance at most $C_0R$ lies in one of a bounded number of
neighboring cell pairs, and Cauchy's inequality applies.
The map from $w$ to its cell masses is linear and positive. Hence
$E_R(w)\le J\sum_jE_R(w_j)$.

For $\mathcal K_q$, use the circle measure
$\widetilde\mu_w=\sum_z\mu_w(z)\delta_{z/q}$, put $Q_*=q/R$, and take
\[
 H(t)=\sum_{\substack{1\le L\le Q_*\\L\text{ a power of }2}}
 \frac1L\left|\sum_{n<L}e(nt)\right|^2.
\]
The geometric-series formula gives
$H(t)\asymp\min(Q_*,\|t\|^{-1})$: for the lower bound choose
$L\asymp\min(Q_*,\|t\|^{-1})$ small enough that $L\|t\|\le1/4$,
and use $L=1$ if necessary. Each summand has nonnegative Fourier
coefficients. Thus
\[
 E_H(\mu)=\iint H(t-u)\,d\mu(t)\,d\mu(u)
\]
is a squared Hilbert norm. Since $H\ge0$, the pointwise majorization
of the weights first gives
$E_H(\widetilde\mu_w)\le E_H(\sum_j\widetilde\mu_{w_j})$;
the Hilbert norm inequality then gives
$E_H(\widetilde\mu_w)\le J\sum_jE_H(\widetilde\mu_{w_j})$.
Pointwise comparison with $H$ proves the second assertion.
\end{proof}

In applications, pair restrictions enter only after this decomposition,
so the norm comparison always concerns the complete multiplier sum.

\begin{lemma}[Short products and congruent multipliers]\label{dec:fejer}
For $1\le V\le q$, we have
\begin{equation}\label{dec:shortbound}
 \sum_{1\le|r|,|s|\le C_1V}\sum_{a,c}w(a)w(c)
 1_{\{|ra-sc|_q\le C_2V\}}
 \ll q^{o(1)}V^2V_w+\frac{V^3W^2}{q}.
\end{equation}
If only pairs with $r\equiv s\pmod D$ are retained, where $D\le R$, then
\begin{equation}\label{dec:deep}
 \mathcal N_{q,D}(w;R)
 \ll q^{o(1)}R(1+R/D)(V_w+RW^2/q).
\end{equation}
\end{lemma}
\begin{proof}
Choose a sufficiently large fixed $C_3$.
If $V\ge q/(2C_3)$, the trivial bound $O(V^2W^2)$ is already
$O(V^3W^2/q)$. Otherwise put $L=\lfloor q/(C_3V)\rfloor\ge1$.
The normalized Fej\'er kernel
\[
 H_L(x)=L^{-2}\left|\sum_{j<L}e_q(jx)\right|^2
 =\sum_{|t|<L}\frac{L-|t|}{L^2}e_q(tx)
\]
has a fixed positive lower bound on $|x|_q\le C_2V$.
The zero frequency contributes $O(V^3W^2/q)$.
For the nonzero frequencies, apply Cauchy's inequality in $r$.
Since $0<|rt|<q/4$, each product is an actual nonzero integer with at most
$2\tau(|rt|)$ factorizations. Hence
\[
 \sum_{r,\,0<|t|<L}|\widehat w(rt/q)|^2
 \le q^{o(1)}\sum_{\xi\bmod q}|\widehat w(\xi/q)|^2=q^{1+o(1)}V_w.
\]
Multiplication by $V/L\ll V^2/q$ gives \eqref{dec:shortbound};
for $V\asymp q$, the trivial estimate suffices.
Under the restriction $r\equiv s\pmod D$, the multiplier graph has degree
$O(1+R/D)$. Replacing the preceding application of Cauchy's inequality
by the Schur bound for this graph gives \eqref{dec:deep}.
\end{proof}

When $w\le\tau$, the same lemma bounds the total mass of all
relations of small height:
\begin{equation}\label{dec:relation}
 \sum_{a,c}w(a)w(c)
 1_{\{v_1a+v_2c+v_3q+v_4=0\text{ for some }0\ne v\in[-V,V]^4\}}
 \ll q^{o(1)}(\tau WV^2+W^2V^3/q).
\end{equation}
For $v_1v_2\ne0$, apply \eqref{dec:shortbound} directly.
If one of the first two coefficients is zero, fix $(v_3,v_4)$ and use
the divisor bound for nonzero integer products. When the corresponding
source point is zero, its contribution is at most $\tau W$.
If both of the first two coefficients vanish and $V<q$, no nonzero
relation is possible. For $V\ge\sqrt q/2$, the inequality
$W^2\le q\tau W\ll\tau WV^2$ suffices.

\subsection{Forward carries and the low-digit second moment}
\label{dec:forward}

The forward splitting turns a low-digit incidence relation into a
bounded-degree graph. We work with amplitude subweights: at exponent $u$,
these are functions satisfying
\[
 0\le w(a)\le F_q(a/q)
       1_{\{q^{-u}\le F_q(a/q)<2q^{-u}\}}.
\]
Replacing the factor $2$ by another fixed constant has no effect.
This convention keeps the amplitude constraints attached to actual
nonempty digit layers, even when a smaller subweight is estimated.

Write $R=q^\nu$. Take $q=BQ$, with $Q,B$ powers of $b$,
$Q=q^{\vartheta}$, and $\vartheta=10\nu/3+O(1/\log q)$.
Write $a=z+BA$.
The exact factorization is
\begin{equation}\label{dec:forwardfactor}
 F_q(a/q)=F_B(z/B)F_Q((A+z/B)/Q).
\end{equation}
On an actual amplitude layer $F_q\asymp q^{-u}$ with $u\le\kappa$,
we have
\begin{equation}\label{dec:shift}
 F_Q((A+z/B)/Q)\le D_qF_Q(A/Q),\qquad
 D_q=\exp(O(\sqrt{\log q})).
\end{equation}
Indeed, the high factor in \eqref{dec:forwardfactor} is at least
$F_q(a/q)\gg q^{-\kappa}$, and $z/B\in[0,1)$.
Lemma~\ref{prelim:shift} therefore compares the shifted value with
the unshifted grid value at a cost $\exp(O(\sqrt{\log q}))$,
uniformly in $A,z$. Its proof in Appendix~\ref{app:shifts} uses the
unit-circle zeros of the one-digit polynomial to control the digit
positions at which such a comparison could otherwise be large. The comparison is used only on
this actual amplitude layer.

Decompose according to the low amplitude $F_B\asymp q^{-a_0}$, and put
\[
 l(z)=F_B(z/B)1_{\{q^{-a_0}\le F_B(z/B)<2q^{-a_0}\}},
 \quad \beta=\frac{q^{-u}}{C D_qq^{-a_0}},
 \quad f_Q^*(A)=\beta1_{\{F_Q(A/Q)\ge\beta\}}.
\]
The original weight is bounded by $q^{o(1)}$ times these low--high
products, and $x=u-a_0+o(1)$. The constant-height superlevel weight
ensures both pointwise bounds in \eqref{dec:highweight}.
For fixed $r,s,A,A'$, the original near-resonance condition
\[
 |B(rA-sA')+rz-sz'|_q\ll R
\]
first forces $|rA-sA'|_Q\ll R$. For each fixed $z$, the condition
$s\asymp R$ leaves only $O(1)$ possible values of $z'$, and conversely.
The low-digit bipartite graph therefore has bounded degree on both sides.
More explicitly, fixing $z$ confines $z'$ to intervals of bounded
length, separated by $q/s\gg B$, and only a bounded number meet
$[0,B)$. The same argument with $r,s$ exchanged bounds the opposite
degree. Schur's inequality gives
\begin{equation}\label{dec:forwardcarry}
 \mathcal N_q(l\otimes f_Q^*;R)
 \ll\|l\|_2^2\mathcal N_Q(f_Q^*;R).
\end{equation}
Put $\ell=\log_q\|l\|_2^2$ and $s=\log_q\sum_Af_Q^*(A)$.
Up to $o(1)$ relaxations, put $x=u-a_0$ and $\vartheta=10\nu/3$.
The full continuous parameter region needed is
\begin{align}
 &0.257\le\nu\le0.29,\quad \nu-0.07\le u\le\kappa,
 \quad 0\le a_0\le u,
 \label{dec:forwardregion}\\
 &-2a_0\le\ell\le\min\{\varepsilon_0(1-\vartheta),
        g_0(1-\vartheta)-(2-p_0)a_0,
        g_4(1-\vartheta)+2a_0\},\label{dec:lowconstraints}\\
 &-x\le s\le\min\{\gamma_0\vartheta,
        g_0\vartheta+(p_0-1)x,
        \varepsilon_0\vartheta+x,
        g_4\vartheta+3x,
        \vartheta-x\}.\label{dec:highconstraints}
\end{align}
The upper bounds come from the first, second, $p_0$th, and fourth
moments and the amplitude cap; the lower bounds for nonempty layers
come from the presence of at least one actual atom.
In the rank-two inequalities of Section~\ref{dec:ranktwo}, we will
take $p=p_0,g=g_0$.

\subsection{Reverse carries: crossing the one-quarter threshold}
\label{dec:reverse}

Now let $0.249999\le\nu\le0.258$, still with
$\vartheta=10\nu/3$ and $q=BQ$, but write $a=A+Qz$.
The product formula becomes
\begin{equation}\label{dec:reversefactor}
 F_q(a/q)=F_Q(A/Q)F_B((z+A/Q)/B).
\end{equation}
Decompose the high factor according to its actual amplitude $q^{-x}$.
The low amplitude is $q^{-a_0}$, where $a_0=u-x$.
The low profile $l_A(z)$ may depend on $A$, and the initial state
$A/Q$ is retained throughout. For each contributing $r,s,A,A'$,
write the unique centered high error as
$rA-sA'=Qk+e$, with $|e|\ll R<Q/2$.
Then the original near-resonance condition is equivalent to
\begin{equation}\label{dec:actualreversecarry}
 |e|\ll R,\qquad rz-sz'+k\equiv0\pmod B.
\end{equation}
Since the multipliers are units, the low-digit relation is an actual
affine permutation. Hence
\begin{equation}\label{dec:lowcross}
 \sum_zl_A(z)l_{A'}((r/s)z+k/s)
 \le\min\{V_{
 \rm low},B^{1-\alpha}
 \mathcal H_L(r/s,k/s;A/Q,A'/Q)\},
\end{equation}
Here the profiles obey
\[
 0\le l_A(z)\le F_B((z+A/Q)/B)
       1_{\{F_B((z+A/Q)/B)\asymp q^{-a_0}\}},
\]
and $V_{\rm low}$ is the supremum, over the actual initial states
$A/Q$, of the squared mass of the full low-amplitude band.
The first bound in \eqref{dec:lowcross} is Cauchy's inequality.
For the second, discard the amplitude restrictions and use
$F_B((z+x)/B)=B^{(1-\alpha)/2}\mu_{B,x}(z)^{1/2}$.
Thus the factor $B^{1-\alpha}$ and both initial states are exact.

First split the original output measure according to the
$\varphi(b)$ unit classes of the multiplier modulo $b$.
The Hilbert norm argument in Lemma~\ref{dec:outputnorm} costs only
this fixed number of classes. Within a class, $r/s\equiv1\pmod b$,
so Lemma~\ref{dec:hellinger} applies. Put $\ell_0=\log_qV_{\rm low}$.
We first majorize each subweight by the full nonempty amplitude band,
which justifies the constraint $\ell_0\ge-2a_0$.
Retain every constraint in \eqref{dec:lowconstraints}, with $\ell$
replaced by $\ell_0$, together with \eqref{dec:highconstraints}, and add
\begin{equation}\label{dec:newconstraints}
 \ell_0\le g(1-\vartheta)-a_0/2,\qquad s\le g\vartheta+x/2.
\end{equation}
Take $\eta=10^{-6}$, and choose a power of $b$ of the form
\begin{equation}\label{dec:deepchoice}
 D_*=q^{d_*+o(1)},\qquad d_*=\max(0,\ell_0+s-x)+\eta.
\end{equation}
When $D_*\mid r-s$, Lemma~\ref{dec:fejer} bounds the contribution
of these deep slopes by
\begin{equation}\label{dec:deepcost}
 \frac{\mathcal N_{\rm deep}}{R^2}
 \ll q^{\ell_0-d_*+o(1)}
       (q^{s-x}+q^{\nu+2s-\vartheta}).
\end{equation}
The first term is at most $q^{-\eta+o(1)}$.
The remaining slopes satisfy $v_b(r-s)<\log_bD_*$, so Hellinger
contraction acts on at least $\log_b(B/D_*)$ levels.
Equation \eqref{dec:lowcross} gives an effective low-digit load
$q^\ell$, where
\begin{equation}\label{dec:effectiveell}
 \ell\le\ell_0,\qquad
 \ell\le\varepsilon_0(1-\vartheta)-\delta(1-\vartheta-d_*).
\end{equation}
We take the minimum of these two upper bounds; the relative Hellinger
saving is not multiplied by $V_{\rm low}$.
All mass parameters and $D_*$ are chosen from the original full nonempty
amplitude bands and then held fixed. Restricting a source set decreases
the sums to be estimated; it does not change these parameters or require
a new lower bound on the mass of the restricted set.
To justify source deletion in this setting, perform the source
decomposition of Lemma~\ref{dec:outputnorm} in the original output
norm, before imposing any slope restriction.
Explicitly, a set $E$ of high frequencies lifts to the subweight
$w_E(A+Qz)=w(A+Qz)1_E(A)$.
For each resulting subweight, the deep contribution is bounded by
\eqref{dec:deepcost}, using the full masses as upper bounds.
For every other slope, \eqref{dec:lowcross} bounds the low cross term
by $q^\ell$, and enlarging the high multiplier sum gives the complete
kernel $\mathcal N_Q$ for that source subweight.
The estimates of Section~\ref{dec:ranktwo} will therefore apply to
this complete kernel.
Thus the deletion estimate follows from the positive complete kernels
$\mathcal N_q$ and $\mathcal N_Q$.

\subsection{From digital moments to arbitrary rational grids}
\label{dec:grids}

Suppose that $1\le p\le2$ and
$\sup_x\sum_{a<Y}F_Y((a+x)/Y)^p\ll Y^g$.
These are the base-$b$ sampling estimates underlying
\cite[Lemmas 10.5--10.7]{Maynard2019}; we keep the common shift explicit.
Write $F_Y=|P_Y|$ for the normalized digital trigonometric polynomial.
Integration over the shift, followed by
$\int|(|P_Y|^p)'|\le p\|P_Y\|_p^{p-1}\|P_Y'\|_p$
and Bernstein's inequality, gives
\begin{equation}\label{dec:variation}
 \int_0^1F_Y^p\ll Y^{g-1},\qquad
 \int_0^1|(F_Y^p)'|\ll Y\int_0^1F_Y^p.
\end{equation}
Consequently, a local supremum may be taken independently in each sampling
cell. Suppose that the centers have mutual distance $\gg1/Z$ and that
the integer neighborhood of each center has width $O(E)$, where $Z,E\ge1$.
Writing $\mathcal U$ for the union of these neighborhoods on the integer
grid modulo $Y$, we obtain
\begin{equation}\label{dec:separated}
 \sum_{a\in\mathcal U}F_Y(a/Y)^p\ll\min(Y^g,(ZE)^g).
\end{equation}
When $ZE<Y$, choose powers $U,V$ of $b$ with
$U\asymp Z$, $V\asymp E$, and $UV\le Y$. The product formula gives
\[
 F_Y(a/Y)^p\le F_U(a/Y)^p F_V(a/V)^p.
\]
Each neighborhood has diameter $O(1/U)$ in the first factor.
By \eqref{dec:variation}, the sum of its local suprema over the
separated centers is $O(U^g)$. In the second factor, a neighborhood
contains only $O(1)$ periods of the integer grid modulo $V$,
and contributes $O(V^g)$. This proves \eqref{dec:separated}.
For $ZE\ge Y$, use the full-grid moment.
If we sum over neighbourhoods rather than their union, a grid point
belongs to $O(1+ZE/Y)$ neighbourhoods of \emph{distinct} centres, by
spacing. For $ZE<Y$ the overlap is bounded; otherwise multiply the
full-grid moment by $O(ZE/Y)$. The resulting incidence bound is
$O((ZE)^g+ZEY^{g-1})$. It does not count different representations
of the same centre; these must be bounded separately. For the union
itself and $0\le w\le\min(F_Y,\tau)$ we have

\begin{equation}\label{dec:squarearc}
 \sum_{a\in\mathcal U}w(a)^2
 \ll\tau^{2-p}\min(Y^g,(ZE)^g).
\end{equation}

We also need two mixed estimates. The first uses two equal-sized digit
blocks. Write an arbitrary denominator as $s=S v$, where every prime
factor of $S$ divides $b$ and $(v,b)=1$. For a common shift $\beta$,
\begin{equation}\label{dec:adaptiverow}
 \sum_{j<s}\sum_{\|a/Y-\beta-j/s\|\le E/Y}w(a)
 \ll (1+E)\tau^{1-p/2}s
 \min\{v,(Y/(1+E))^{1/2}\}^{g-1}.
\end{equation}
To prove this, choose a power of $b$ with
$U\asymp\min\{v,(Y/(1+E))^{1/2}\}$ and use
$F_Y^{p/2}\le F_U(t)^{p/2}F_U(Ut)^{p/2}$.
The two transformed grids have denominators $s$ and
$s/(s,U)$, both at least $v\ge U$; repeated points are counted with
their original multiplicity. This is a statement about the entire
grids, and does not require each numerator to be coprime to $s$.
Moreover, the displaced neighborhoods have width $O(1/U)$ in both
blocks, since $U^2(1+E)\ll Y$.
By \eqref{dec:variation}, the sum of the $p$th powers of the
centerwise suprema in either block is $O(sU^{g-1})$.
Apply Cauchy's inequality, allow for the $O(1+E)$ grid points in each
neighborhood, and use $w\le\tau^{1-p/2}F_Y^{p/2}$.
For $E\gg Y$, the pointwise bound suffices.

The second is the corresponding base-$b$ proof of Maynard's mixed
rational-grid lemma \cite[Lemma 10.7]{Maynard2019}.
Fix $q_1\asymp Q_1$ with $(q_1,b)=1$, and fix $d_0\asymp D$
whose prime factors all divide $b$. If $DE\ll Y$, then
\begin{align}
 &\sum_{\substack{q_2\asymp Q_2\\(q_2,b)=1}}
 \sum_{\substack{a<d_0q_1q_2\\(a,d_0q_1q_2)=1}}
 \sum_{\substack{|\eta|\le E/Y\\Y(\beta+a/(d_0q_1q_2)+\eta)\in\mathbb Z}}
 F_Y(\beta+a/(d_0q_1q_2)+\eta)\notag\\
 &\hspace{25mm}\ll (DE)^{\gamma_0}(Q_1Q_2^2)^{1/21}
       +E^{5/6}D^{3/2}Q_1Q_2^2Y^{-10/21}.
 \label{dec:maynardmixed}
\end{align}
Here $\beta$ is fixed throughout the sum. To explain the roles of the base
and the common shift, choose powers of $b$ with
$D'\asymp D,E'\asymp E,V\asymp(Y/(D'E'))^{1/2}$, and factor
$d_0=d_1d_2d_3$ so that $d_3=(d_0,D')$ and
$d_2d_3=(d_0,VD')$. The last block sums the offsets at cost $O(E^{\gamma_0})$.
For the two middle blocks use $2uv\le u^2+v^2$: each term retains
the first block and the square of just one middle block. Multiplying
a centre by $D'$ or $D'V$ cancels respectively $g=d_3$ or $g=d_2d_3$
from its denominator (the other factors $q_1q_2$ are coprime to $b$).
The reduced image denominator is $(d_0/g)q_1q_2$, hence determines
$q_2$. Its fibre lies on a translated $g$-point grid; the first-block
suprema on this fibre cost $O(g^{\gamma_0})$ by retaining a radix
block of size comparable to $g$ and applying \eqref{dec:variation}.
Distinct image fractions share the denominator factor $(d_0/g)q_1$,
so their spacing is $\gg 1/H$ with $H=(d_0/g)Q_1Q_2^2$.
The displaced middle arguments vary by $O(1/V)$ since $D'V^2E\ll Y$.
For $H<V$ retain a block of size comparable to $H$; for $H\ge V$
sample the full block at spacing $1/H$. The second moment
\eqref{dec:evenmoments} and \eqref{dec:variation} give in the two
cases $O(H^{1-\alpha}+HV^{-\alpha})$. With
$H_1=d_1d_2Q_1Q_2^2$, $H_2=d_1Q_1Q_2^2$, the two terms total
\[
 \ll E^{\gamma_0}\{d_3^{\gamma_0}(H_1^{1-\alpha}+H_1V^{-\alpha})
 +(d_2d_3)^{\gamma_0}(H_2^{1-\alpha}+H_2V^{-\alpha})\}.
\]

Now use $1-\alpha<1/21$, $\alpha/2>10/21$,
$\gamma_0+10/21<5/6$, and $1+10/21<3/2$ to obtain
\eqref{dec:maynardmixed}. In the two blocks the common shift becomes
$D'\beta$ and $D'V\beta$; it changes neither circular spacing nor integrals
over a full period. This proves the same estimate in base eleven and
makes explicit why the shift must not vary with $q_2$.

\subsection{Uniform exponent inequalities for all rank-two denominators}
\label{dec:ranktwo}

All exponents in this subsection are measured relative to the original
$q$. Let $R=q^\nu$ and $Q=q^{\vartheta+o(1)}$, where
$\vartheta=10\nu/3$, so that $R=Q^{3/10+o(1)}$.
The high-frequency weight $f_Q^*$ satisfies
\begin{equation}\label{dec:highweight}
 0\le f_Q^*\ll\min(F_Q,q^{-x}),\qquad
 f_Q^*\ll q^{(p-1)x}F_Q^p,\qquad
 W_Q=\sum f_Q^*\ll q^s.
\end{equation}
There is also a low-digit load $q^\ell$.
After division by $R^2$, the two terms in Lemma~\ref{dec:fejer}
have exponents
\begin{equation}\label{dec:generaltargets}
 \ell+s-x,\qquad \ell+\nu+2s-\vartheta.
\end{equation}
The zero Fourier term will be controlled throughout the parameter domain.
If the general term has exponent below $-10^{-6}$, the layer is already
settled; hence we retain the constraint $\ell+s-x\ge-10^{-6}$ below.

Fix a frequency pair $(A_1,A_2)$, and let $J=q^j\ge1$ be the actual
number of near-resonant multiplier pairs. We decompose this count
dyadically; factors bounded independently of $q$ are absorbed in the
fixed width margin below. In particular, $0\le j\le2\nu+o(1)$. Any three independent integer
points in the convex body
\[
 |r|,|s|\ll R,\qquad |rA_1-sA_2+nQ|\ll R
\]
have determinant $O(R^3/Q)=o(1)$, so the rank is at most two.
If the actual points are collinear, the primitive relation has height
$O(R/J)$. Taking the smaller of \eqref{dec:relation} and the trivial
mass bound $W_Q^2$, and optimizing in $J$, bounds the normalized
contribution of this part by
\begin{equation}\label{dec:linetargets}
 q^{\ell+(3s-x)/2-\nu+o(1)}+
 q^{\ell+2s-\nu-\vartheta/3+o(1)}.
\end{equation}
Here we use $\min(t,c/t)\le c^{1/2}$ and
$\min(t,c/t^2)\le c^{1/3}$, respectively.

If there are two independent actual points, the lattice-area argument
of \cite[Lemma 14.1]{Maynard2019} gives
\begin{equation}\label{dec:rationalwitness}
 1\le q_0\ll\mathcal R:=R^2/J,\quad (b_1,b_2,q_0)=1,
 \qquad A_i/Q=b_i/q_0+O(\mathcal R/(Qq_0)).
\end{equation}
Here is the calculation in the present normalization. Write
$e(r,s,n)=rA_1-sA_2+nQ$ and transform the section lattice by
$(r,s,n)\mapsto(r/R,s/R,e(r,s,n)/R)$.
Its two successive minima satisfy $\lambda_1\lambda_2\ll J^{-1}$,
by the planar lattice-point bound, since the section contains $J$
points and two independent ones. Choose independent integer vectors
$v,v'$ attaining these minima up to fixed constants, and put
$w=v\times v'$. Then
\[
 |w_3|\ll R^2/J,\qquad
 \left|w_1-w_3A_1/Q\right|+
 \left|w_2+w_3A_2/Q\right|\ll R^2/(JQ).
\]
The last two inequalities follow by substituting
$n=(e-rA_1+sA_2)/Q$ in the cross product.
If $w_3=0$, integrality and $R^2/(JQ)=o(1)$ would force $w=0$,
a contradiction. Divide $w$ by the gcd of its coordinates and choose
its sign so that $w_3>0$; then take
$(b_1,b_2,q_0)=(w_1,-w_2,w_3)$.
This proves \eqref{dec:rationalwitness}, including the case $J=O(1)$.
Let $\lambda=2\nu-j$, so that $0\le\lambda\le2\nu+o(1)$.
Put $(b_i,q_0)=g_i$, interchange the coordinates so that $g_1\le g_2$,
and write
\[
 g_1=d_1g'_1,\quad q_0=d_0d_1q'g'_1g_2,\qquad
 g_2=d_2g'_2,\qquad(q'g'_1g'_2,b)=1,
\]
where all prime factors of $d_0,d_1,d_2$ divide $b$.
The number of their possible values is bounded by a fixed power of $\log q$.
Let $(a_1,a_2,a_3,a_4,a_5,a_6)$ be the scale exponents of
$(q',g'_1,g_2,d_0,d_1,d_2)$, respectively. To keep the matrix notation
concise, we denote them below by $(a,c,h,D,E,G)$. Thus
\begin{equation}\label{dec:geometryconstraints}
 \begin{gathered}
 a,c,h,D,E,G\ge0,\qquad 0\le\lambda\le2\nu,\\
 \rho_0=a+c+h+D+E\le\lambda,\qquad c+E\le h,\qquad G\le h.
 \end{gathered}
\end{equation}
Put $w_0=10^{-6}$ and $e_0=\lambda-\rho_0+w_0$, and define
\[
 z_1=\lambda+c-h+w_0,\quad
 z_2=\lambda+a+h-c-E+w_0,\quad
 z_{\rm arc}=\lambda+a+c-h+w_0.
\]
This fixed width margin absorbs all geometric constants.

We now define the rational-grid sums whose exponents enter the linear systems.
For an integer denominator $D_1$ and a width $E_1\ge1$, define
\[
 \mathcal S(D_1,E_1)=
 \sum_{\substack{a<D_1\\(a,D_1)=1}}
 \ \sum_{\substack{A\bmod Q\\\|A/Q-a/D_1\|\le C E_1/Q}}
 f_Q^*(A),
\]
where $C$ is fixed. Every rational representation is counted.
Within the chosen dyadic scales and fixed smooth parts, put
\[
 I(q')=\sum_{g'_1}\mathcal S(d_0d_1q'g'_1,q^{e_0}),\qquad
 O(q')=\sum_{g_2}\mathcal S(d_0q'g_2,q^{e_0}),
\]
and write $\sup_{q'}I(q')\ll q^{s_1+o(1)}$,
$\sum_{q'}O(q')\ll q^{s_2+o(1)}$.
These denominators are exactly $q_0/g_2$ and $q_0/g_1$ from the
simultaneous approximation, respectively.
Estimates \eqref{dec:separated}, \eqref{dec:adaptiverow}, and
\eqref{dec:maynardmixed} give
\begin{align}
 s_1\le\min\{&\gamma_0z_1,z_1-x,gz_1+(p-1)x,
                \max(A_1,B_1),\max(U_1,V_1)\},\label{dec:s1}\\
 s_2\le\min\{&\max(\gamma_0z_2,z_2-(1-\gamma_0)\vartheta),z_2-x,
 \notag\\[-2mm]
 & (p-1)x+\max(gz_2,z_2-(1-g)\vartheta),
   \max(A_2,B_2),\max(U_2,V_2),s+\max(0,z_2-\vartheta)\},
 \label{dec:s2}
\end{align}
where all the linear forms are as follows:
\begin{align*}
 A_1&=\gamma_0(\lambda-a-c-h+w_0)+(a+2c)/21,\\
 B_1&=5(\lambda+w_0)/6+a/6+7c/6-5h/6+2(D+E)/3-10\vartheta/21,\\
 U_1&=D+E+a+2c+e_0-(1-p/2)x-(1-g)(a+c),\\
 V_1&=D+E+a+2c+e_0-(1-p/2)x-(1-g)(\vartheta-e_0)/2,\\
 A_2&=\gamma_0(\lambda-a-c-h-E+G+w_0)+2(a+h-G)/21,\\
 B_2&=5(\lambda+w_0)/6+7a/6-5c/6+7h/6+2D/3-5E/6-G/2-10\vartheta/21,\\
 U_2&=D+2a+2h-G+e_0-(1-p/2)x-(1-g)(a+h-G),\\
 V_2&=D+2a+2h-G+e_0-(1-p/2)x-(1-g)(\vartheta-e_0)/2.
\end{align*}
For example, the separation parameter for the inner grid is
$d_0d_1q'(g'_1)^2$; multiplication by the width gives exactly the
exponent $z_1$. The orientation constraint gives
$z_1\le\lambda+w_0\le2\nu+w_0<\vartheta$ throughout the windows
used below, so the inner grid has no saturation term. In the outer grid, combine $q'g'_2$, losing only divisor
multiplicity. Thus \eqref{dec:maynardmixed} is applied with smooth part
$d_0d_2$ and coprime part $q'g'_2$, giving $A_2,B_2$.
The last term in \eqref{dec:s2} includes all outer representations:
for a fixed source point, the reduced centers are separated by at least
$\gg[d_0(Q_1G_2)^2]^{-1}$, so the number of representations is at most
$q^{o(1)}(1+q^{z_2-\vartheta})$.

For each actual frequency pair in a $J$-layer, first select one witness
from the finite set in \eqref{dec:rationalwitness}, for example the
lexicographically least $(q_0,b_1,b_2)$, with the coordinate orientation
also fixed by a deterministic rule. Group the pairs by the scales of
this chosen witness. Thus each packet $\mathcal P$ consists of ordered
source pairs, each carrying its one selected witness.
Empty packets are discarded. Its mass is
$P=\sum_{(A_1,A_2)\in\mathcal P}f_Q^*(A_1)f_Q^*(A_2)$.
For each pair, the selected witness supplies one term in
$I(q')O(q')$. Positivity therefore gives
$P\le\sum_{q'}I(q')O(q')$, while the set interpretation gives
$P\le W_Q^2$. Consequently,
\begin{equation}\label{dec:packetmass}
 P\le W_Q^2,\qquad
 P\ll \bigl(\sup_{q'} I(q')\bigr)\sum_{q'}O(q'),\qquad
 t:=\log_qP\le\min(2s,s_1+s_2)+o(1).
\end{equation}
The contribution of this packet is at most $q^\ell JP$.
Since $J/R^2=q^{-\lambda}$, its normalized target exponent is
$\ell+t-\lambda$.
The rational-grid estimate enlarges to all representations only after
the packet has been selected. This order makes it compatible with
the independent bound $P\le W_Q^2$.

We may first delete entire source arcs whose contribution is already
controlled. Fix $d_0,d_1$ and the scales above, and take the union of
the neighborhoods of width $O(q^{e_0}/Q)$ over all inner denominators
$d_0d_1q'g'_1$. Their separation parameter is $d_0d_1(Q_1G_1)^2$,
so \eqref{dec:squarearc} bounds their squared mass by
$\ll q^{gz_{\rm arc}-(2-p)x+o(1)}$. If
\[
 \ell+gz_{\rm arc}-(2-p)x<-10^{-6},
\]
let $\mathcal A$ denote this union of source arcs and apply
Lemma~\ref{dec:fejer} to $f_Q^*1_{\mathcal A}$.
Its general term is bounded by the displayed squared mass; its
zero Fourier term is bounded using the full mass $W_Q$ and the
global second target in \eqref{dec:generaltargets}.
These bounds pay for the full output norm of this source subweight.

The number of choices of scales and smooth parts is bounded by a
fixed power of $\log q$. Delete the union of their favorable source
arcs by Lemma~\ref{dec:outputnorm}, retaining its complement as one
further subweight. Restrict the previously selected packets to pairs
in that complement, without reselecting witnesses. If the inner
coordinate of a surviving pair lay in a favorable arc, the pair
would already have been deleted. Thus every remaining packet
satisfies
\begin{equation}\label{dec:arcfilter}
 \ell+gz_{\rm arc}-(2-p)x\ge-10^{-6}.
\end{equation}

\subsection{The forward and reverse bilinear windows}

We now combine the carry reductions with the rank-two estimates.
Each contribution is bounded by a power of $q$ whose exponent is a
piecewise linear function of the scale parameters. The constraints above
retain both the actual amplitude information and every rational
representation. It remains to show that all the relevant exponents are
strictly negative throughout their continuous parameter regions.

\begin{lemma}[Forward window]\label{dec:forwardlemma}
Throughout \eqref{dec:forwardregion}, every nonnegative subweight
of an actual amplitude layer satisfies
\[
 \mathcal N_q(w;R)\ll R^2q^{-c}
\]
for a fixed $c>0$.
\end{lemma}
\begin{proof}
Combine \eqref{dec:generaltargets}--\eqref{dec:arcfilter} and
\eqref{dec:forwardregion}--\eqref{dec:highconstraints} into rational
linear systems. Split each of the five binary maxima into branches,
adding the constraint that the selected term dominates the other.
In \eqref{dec:s2}, the transitions for the basic moment, the $p$th
moment, and representation multiplicity all occur at $z_2=\vartheta$.
There are therefore $32$ consistent branches.
The rational certificates in the supplementary material give
\[
\begin{array}{c|c}
 \text{target}&\text{strict upper bound}\\ \hline
 \ell+\nu+2s-\vartheta\quad(\text{unfiltered})&-6061839073/100023000000\\
 \ell+(3s-x)/2-\nu&-8736163/63218000000\\
 \ell+2s-\nu-\vartheta/3&-4100926993/142240500000\\
 \ell+t-\lambda\quad(\text{all remaining rank-two branches})&-43400641/189654000000
\end{array}
\]
For each feasible branch, a nonnegative rational linear combination
proves the stated strict upper bound; the remaining branches have
Farkas certificates. The general term has already been controlled
where $\ell+s-x<-10^{-6}$. All deleted source arcs and collinear
contributions are covered by the preceding lemmas.
Apply Lemma~\ref{dec:outputnorm} to the low-amplitude subweights and
the deleted source arcs. Within each retained subweight, add the
nonnegative contributions of the selected source-pair packets.
Both decompositions have only a logarithmic-power number of parts,
so the fixed power saving is retained.
\end{proof}

\begin{lemma}[Reverse window]\label{dec:reverselemma}
Let $w$ be a nonnegative subweight of the amplitude layer of size
$q^{-u}$ and restrict the multipliers in $\mathcal N_q$ to units
modulo $b$. For either input pair in \eqref{dec:groups}, throughout
$0.249999\le\nu\le0.258$ and $\nu-0.07\le u\le\kappa$, we have
\[
 \mathcal N_q(w;R)\ll R^2q^{-c}
\]
for a fixed $c>0$.
\end{lemma}
\begin{proof}
Take $p=3/2$ in Section~\ref{dec:ranktwo}, add
\eqref{dec:newconstraints}--\eqref{dec:effectiveell}, and split into
the two branches of the maximum in \eqref{dec:deepchoice}.
Each input pair gives $64$ rank-two branches and $12$ basic targets.
The exact rational certificates for the two parameter pairs give
\[
\begin{array}{c|c|c}
 \text{target}&(g,\delta)=(.14387,.0053)&(g,\delta)=(.1435,.004)\\ \hline
 \ell+(3s-x)/2-\nu&-568609129/16755000000000&-2061939/6275000000\\
 \ell+t-\lambda&-110754136637/2010600000000000&-1645105237/3012000000000\\
 \ell_0+\nu+2s-\vartheta&-130413037/2500000000&-52781879/1000000000
\end{array}
\]
The second collinear target is also strictly negative.
The unfiltered depth constraints give
$d_*-(1-\vartheta)\le-.093999$ and $d_*-\nu\le-.203998$.
Thus $D_*<B,R$, and the second term of \eqref{dec:deepcost}
is controlled as well. The remaining contributions are combined
as in the forward window.
\end{proof}

All these computations use only integers and rationals.
The verifier reconstructs the variables, constraints, targets, and all
branches from the fixed parameters in the text, compares them entry by
entry with the certificate matrices, and then checks the nonnegative
dual or Farkas combinations. There are $36$ forward systems and $76$
systems for each reverse input pair, for a total of $188$.
The correspondence between the mathematical variables and the entries
of each certificate matrix is fixed in the accompanying system constructors.
Since there are finitely many dual certificates, the sums of their
coefficients are bounded.
Consequently, rounding the original scales and allowing $q^{o(1)}$
losses changes each target by only $o(1)$, which the strictly negative
margins absorb.

\subsection{The classical rank-two estimate for large amplitudes and the right-hand window}
\label{dec:oldranktwo}

To give the full extension to base eleven, we record the precise form
of Maynard's rank-two proposition that is needed.
Let $\mathcal E=\{a:F_q(a/q)\ge q^{-\kappa}\}$, and let
$\mathcal F(Q_{\rm arc},E)$ be the rational-arc grid about reduced
centres $u/r$ with $(u,r)=1$, denominators
$r\asymp Q_{\rm arc}\le\sqrt q$, and offsets $\asymp E/q$.
Here $E\le100\sqrt q/Q_{\rm arc}$, and $E=0$ means that the
frequency equals its centre.
Let $\mathcal B_2(N,K,\Delta)$ be the set of frequency pairs for which
there is a rank-two integer lattice containing at least $\Delta KN^2$
points satisfying
\[
 \|\mathbf n\|_2\ll N,\qquad
 |n_1a_1+n_2a_2+n_3q|\ll\Delta q,
\]
and these points are not all collinear.

\begin{lemma}[A common form of the classical rank-two input]\label{dec:classicalrank}
If $N,K\ge1$, $\Delta\ge N/q$, and $NK\ge q^{17/40}$,
then for some fixed $c>0$,
\begin{equation}\label{dec:classicalrankbound}
 \sum_{\substack{(a_1,a_2)\in\mathcal B_2(N,K,\Delta)\\
 a_1,a_2\in\mathcal E\cap\mathcal F(Q_{\rm arc},E)}}
 F_q(a_1/q)F_q(a_2/q)
 \ll\frac{q(\log q)^{O(1)}}{NK(Q_{\rm arc}+E)^c}.
\end{equation}
This holds for both $b=10$ and $b=11$.
\end{lemma}
\begin{proof}
For base ten the model is \cite[Proposition~13.3]{Maynard2019}; we give the argument with the exceptional-set cutoff used here.
We spell out the base-dependent part of its proof.
The lattice and gcd decomposition is
\cite[Lemmas 14.1--14.3]{Maynard2019}: it uses integer determinants,
reduced fractions, and the Chinese remainder theorem.
For a different fixed base, its smooth denominators are simply those
supported on the primes dividing $b$.
The estimates below verify all the digital and exponent inputs of
\cite[Lemma 14.4]{Maynard2019}. First, \eqref{dec:momentold} gives
\[
 \sum_{a\in\mathcal E}F_q(a/q)
 \ll q^{g_0+(p_0-1)\kappa}=q^\sigma,\qquad
 \sigma=1533559/5334560<23/80.
\]
For clarity, here is the lattice-area step at thickness $\Delta q$.
Let $J\ge\Delta KN^2$ be the number of section points. Map the section
by $\mathbf n\mapsto(n_1/N,n_2/N,(n_1a_1+n_2a_2+n_3q)/(\Delta q))$.
The section points lie in a fixed box and span a plane. As in
\eqref{dec:rationalwitness}, the product of the first two minima of
the section lattice in this box is $O(1/J)$. Crossing two independent
vectors attaining these minima in the original integer coordinates gives
a nonzero integer normal $w$ with
\[
 |w_3|\ll N^2/J,\qquad |w_i-w_3a_i/q|\ll N\Delta/J\quad(i=1,2).
\]
If $w_3=0$, the latter bounds are $O(1/(NK))=o(1)$, contradicting
integrality. Make the normal primitive and orient it with $w_3=q_0>0$.
Then $q_0\ll 1/(\Delta K)\le q/(NK)$, and the offsets from $w_i/q_0$,
in units of $1/q$, are $O(q/(NKq_0))$. Thus rounding the common width
up to at least one still gives $q_0 E_0\ll q/(NK)$.
Use the gcd orientation of Section~\ref{dec:ranktwo}. Write
$Q_1,G_1,G_2$ for the scales of $q',g'_1,g_2$ and $D_0,D_1$ for
the scales of $d_0,d_1$. Then $G_1\ll G_2$, and the denominator
and error scales obey
\[
 Q_0=Q_1G_1G_2D_0D_1,\qquad E_0\ge1,\qquad
 Q_0E_0\ll q/(NK)\le q^{23/40}.
\]
Fix one witness for each source pair and group the witnesses into
packets with fixed smooth parts and dyadic size parameters. Put
$U=Q_0E_0$. In a packet let $S_1$ be the supremum of the inner grid
sum at fixed shared denominator parameter $q'$. The outer reduced
centres have denominators $d_0d_2m$ with $m=q'g_2'$, as in
Section~\ref{dec:ranktwo}. Once the combined parameter $m$ is fixed,
there are at most $\tau(m)\ll_\epsilon U^\epsilon$ choices of
$(q',g_2')$, because $m\ll Q_0$ and $E_0\ge1$. We keep this
factor outside the outer estimates.

For the first outer estimate, keep the distinct reduced centres
at combined denominators $d_0d_2m$. Write $D_2\asymp d_2$ for
the fixed factor supported on primes dividing $b$. Distinct
centres are separated by at least a constant times
$d_0d_2/(Q_0/(D_1G_1))^2$, by the integer cross determinant.
Thus we can use the separation scale
$Q_0^2/(D_0D_1^2G_1^2)$, dropping the factor $D_2^{-1}$
from this scale to make an upper bound. Counting the overlap of
the distinct centres' neighbourhoods as in \eqref{dec:separated}
gives the estimate for $S_2$ below. In its full-grid term we also drop
one factor each of $D_1$ and $G_1$ from the denominator.
Together with the inner ordinary and mixed estimates from
Section~\ref{dec:grids}, this gives
\begin{align*}
 S_1&\ll(Q_0E_0)^{\gamma_0},\\
 S_2&\ll\left(\frac{Q_0^2E_0}{D_0D_1^2G_1^2}\right)^{\gamma_0}
       +\frac{Q_0^2E_0}{q^{50/77}D_0D_1G_1},\\
 S_3&\ll q^{23/80}\left(1+\frac{Q_0}{D_0D_1NK}\right),\\
 S_1&\ll Q_0^{1/21}(D_0D_1E_0)^{\gamma_0}
       +\frac{Q_0G_1(D_0D_1)^{1/2}E_0^{5/6}}{G_2q^{10/21}}.
\end{align*}
To obtain the second outer estimate $S_3$, fix an actual frequency
$a_2$ and one of its reduced centres $u_* /(d_0d_2m_*)$ in
the packet. Choose real representatives in the common
$O(E_0/q)$ neighbourhood of $a_2/q$. For every other such centre
$u/(d_0d_2m)$ the integer $j=um_*-u_*m$ satisfies
\[
 |j|\ll \frac{E_0d_0d_2m_*^2}{q}
 \ll \frac{Q_0}{D_0D_1 NK}.
\]
Indeed $d_0d_2m\asymp d_0d_2m_*\asymp Q_0/(D_1G_1)$,
$d_0d_2\gg D_0$, and $Q_0E_0\ll q/(NK)$.
For fixed $j$ the congruence $u_*m\equiv-j\bmod m_*$ determines
the residue of $m$, because $(u_*,m_*)=1$. The dyadic size condition
leaves $O(1)$ choices of $m$, and then $u$ is determined. Summing
over $a_2\in\mathcal E$ uses the bound $q^\sigma\le q^{23/80}$.
Restoring the factor for splitting $m$ proves that the packet
mass is at most
$O_\epsilon(U^\epsilon\min\{S_1S_2,S_1S_3\})$.
The second outer estimate only counts actual centres at each
frequency, so it bounds the selected contribution separately
from the first estimate.

If the first term in the last estimate dominates, multiplication
by the bound for $S_2$ gives
\[
 S_1S_2\ll
 Q_0^{54/77+1/21}E_0^{54/77}
 +Q_0^{2+1/21}E_0^2q^{-50/77}
 \ll(Q_0E_0)^{1-c}.
\]
For the last step, use $Q_0E_0\ll q^{23/40}$ and
$\frac{22}{21}\frac{23}{40}<\frac{50}{77}$.
If the second term dominates, use $G_1\ll G_2$ and $E_0\ge1$ to obtain
\[
 S_1\ll\min\left((Q_0E_0)^{\gamma_0},
          \frac{Q_0E_0(D_0D_1)^{1/2}}{q^{10/21}}\right).
\]
For the first term of $S_2$, interpolate the two displayed
bounds for $S_1$ with weights $1/3,2/3$; for its second term,
use the second bound directly. This gives
\[
 S_1S_2\ll Q_0^{3/2}E_0^{6/5}q^{-3/10}
                +Q_0^3E_0^2q^{-9/8}.
\]
The rounded exponents are upper bounds: in particular,
$50/77+10/21>9/8$.
Since $Q_0E_0\ll q^{23/40}$, the first term is
$\ll Q_0E_0q^{-1/80}$ and hence $\ll(Q_0E_0)^{1-c}$.

Likewise, the two terms of $S_3$ give
\[
 S_1S_3\ll q^{23/80}(Q_0E_0)^{27/77}
             +\frac{Q_0^2E_0}{q^{3/16}NK}.
\]
Here $10/21-23/80>3/16$.
The second term is at most
$Q_0q^{13/16}/(NK)^2\ll Q_0q^{-3/80}$, so it too is
$\ll(Q_0E_0)^{1-c}$. We have therefore proved
\begin{align*}
 S_1S_2&\ll(Q_0E_0)^{1-c}+Q_0^3E_0^2q^{-9/8},\\
 S_1S_3&\ll(Q_0E_0)^{1-c}+q^{23/80}(Q_0E_0)^{27/77}.
\end{align*}
Taking the smaller of the last two terms and interpolating with weights
$23/100,77/100$ gives
\[
 \min\{Q_0^3E_0^2q^{-9/8}, q^{23/80}(Q_0E_0)^{27/77}\}
 \le Q_0^{96/100}E_0^{73/100}q^{-299/8000}.
\]
The preceding exponent comparisons hold with $c=1/100$ before the
representation factor is restored. Taking $\epsilon=1/1000$
therefore bounds each packet by $O(U^{1-9/1000})$.
Thus the divisor loss is absorbed in the packet scale.

Write $T=q/(NK)$. We now compare the original reduced centre
$u/r$ with the reduced centre $v/\ell$ from the witness for the
same frequency. If $T\ge q^{1/4}$, then
$Q_{\rm arc}+E\ll\sqrt q\le T^2$. For $T<q^{1/4}$ we have
$r\ll\sqrt q$, $rE\ll\sqrt q$, $\ell\ll Q_0$, and
$\ell E_0\ll Q_0E_0\ll T$. The absolute cross determinant is
at most
\[
 O\!\left(\frac{r\ell(E+E_0)}q\right)
 =O(T/\sqrt q)<1
\]
for all sufficiently large $q$. It must vanish. Since both fractions
are reduced, $r=\ell$ and the original offset is $O(E_0/q)$.
Thus $Q_{\rm arc}+E\ll T^2$ also in this case.
The packet bound is $O(T^{1-9/1000})$, which gives
$O(T(Q_{\rm arc}+E)^{-9/2000})$. The number of packets is bounded
by a fixed power of $\log q$: the size parameters lie in dyadic
ranges, and there are only polylogarithmically many fixed smooth
parts. This proves \eqref{dec:classicalrankbound}. The argument uses only
\eqref{dec:moment1}, \eqref{dec:momentold}, \eqref{dec:maynardmixed},
and comparisons of rational exponents, so it applies in base eleven
as well.
\end{proof}

Take $\zeta=.001$, and put
\[
 u_0(\nu)=\nu-\tfrac32(1-\alpha)-\zeta.
\]
On an amplitude layer with $.249999\le\nu\le.29$ and
$u\le u_0(\nu)$, the second moment gives $W\ll q^{1-\alpha+u}$.
Decompose the full Cauchy kernel according to thickness and multiplicity,
write the number of points as at least $\Delta K'R^2$, where
$\Delta\ge R/q$, and put $\mathcal R=q/(RK')$.
Fix $0<\eta<\zeta/2$. If $\mathcal R\ge q^\eta W^2$, the total
mass directly gives the bound $q^{-\eta}\mathcal R$.
Otherwise,
\begin{equation}\label{dec:largeamplitudeNK}
 RK'\gg q^{1-2\nu+(1-\alpha)+2\zeta-\eta-o(1)}.
\end{equation}
Since $1-\alpha>.01$, this exponent is still greater than $.4315$
at $\nu=.29$. Thus $RK'>q^{17/40}$, and $K'$ is bounded below
by a fixed power of $q$. This lower bound absorbs the
$O(\Delta R^2)$ points outside the plane in the lattice-point
classification of \cite[Lemma 13.2]{Maynard2019}.
Apply Lemma~\ref{dec:classicalrank} to the rank-two part.
For the collinear part, the actual nonzero multiples give height
$V\ll\mathcal R/R$. By \eqref{dec:relation}, its mass divided by
$\mathcal R$ is at most
\[
 q^{o(1)}\left(q^\eta\frac{q^{-u}W^3}{R^2}
          +q^{2\eta}\frac{W^6}{qR^3}\right).
\]
The two exponents are at most $\eta-2\zeta$ and
$2\eta+3\nu-1-3(1-\alpha)-6\zeta$, respectively, both strictly negative.
Since $1-\alpha<.046$, we have $u_0(\nu)>\nu-.07$.
This range overlaps the intermediate-amplitude bands in the two
preceding windows.

The right-hand window requires only the earlier moment estimate. Put
\[
 \nu_0=(3\sigma-\kappa)/2=3066991/10669120.
\]
On each actual amplitude layer retain
$W\ll q^{g_0+(p_0-1)u}$. Since $3p_0-4>0$, this gives
$q^{-u}W^3\ll q^{3\sigma-\kappa}$.
For $\nu_0+\varepsilon_1\le\nu\le.425-\varepsilon_1$, choose
$0<\eta<\min(\varepsilon_1,\kappa-\sigma)$.
The nontrivial multiplicity layers satisfy
$RK'>q^{1-2\sigma-\eta-o(1)}>q^{.425}$, while the two collinear
exponents are at most
\[
 \eta+2\nu_0-2\nu<0,\qquad
 2\eta+6\sigma-1-3\nu<0.
\]
At $\nu=\nu_0$ and $\eta=0$, the second expression is
$-293477/2133824$. Taking $\varepsilon_1=.0004$ and using the fixed
joining points $.2575,.288$, the windows together cover
\begin{equation}\label{dec:rawwindow}
 [.249999,.4244].
\end{equation}
The following subdivision records which estimates supply this range:
\begin{center}
\begin{tabular}{ll}
\toprule
Short-factor exponent $\nu$ & Estimates used\\
\midrule
$[.249999,.2575]$ & Reverse carries and the large-amplitude bound\\
$[.2575,.288]$ & Forward carries and the large-amplitude bound\\
$[.288,.4244]$ & The classical right-hand window\\
\bottomrule
\end{tabular}
\end{center}
In the first two rows, $u_0(\nu)>\nu-.07$ ensures that the
large-amplitude and carry estimates overlap. The ordinary spectrum
$F_q<q^{-\kappa}$ is treated separately in the completion below.
All lattice-point classifications retain the actual nonzero multiplier
points. Absorbing the fixed geometric constants into
$N\mapsto CR$, $\Delta\mapsto C\Delta$, and
$K'\mapsto cK'/C^3$ preserves the lower bound for the thickness
and the fixed power margin in $NK'$.
Moreover, a nonempty layer satisfies $\Delta K'\ll1$, so the
power lower bound for $K'$ ensures $\Delta<1$ after rescaling.

\subsection{Completion of the minor-arc estimate}
\label{dec:typeii}

\begin{proof}[Proof of Proposition~\ref{dec:minorarc}]
For the ordinary spectrum $F_q<q^{-\kappa}$, we retain the earlier
moment estimate, since
\begin{equation}\label{dec:ordinary}
 \sum_{F_q<q^{-\kappa}}F_q^2
 \ll q^{-\delta_0},\qquad
 \delta_0=\kappa(2-p_0)-g_0=127/5334560>0.
\end{equation}
The product coefficient of any bilinear convolution with bounded
coefficients is at most $\tau_2(n)$. In the physical support of a fixed
prefix, each residue class modulo $q$ contains only $O_v(1)$ integers.
Parseval's identity therefore bounds the squared frequency sum by
$\ll_vq^2\log^3q$. The ordinary spectrum consequently contributes
$O_v(Aq^{-\delta_0/2}\log^{3/2}q)$.

Choose a Dirichlet approximation with denominator at most $\sqrt q$
for every frequency and reduce it before assigning its denominator and
offset. Its Dirichlet error bound in terms of the reduced denominator
still holds. Group the frequencies by denominator
$Q_{\rm arc}$ and offset $E$. The number of packets is bounded by
a fixed power of $\log q$. If the logarithmic exponent $D_0$ in the definition of
$\mathfrak M$ is sufficiently large, every frequency outside
$\mathfrak M$ has $Q_{\rm arc}+E\gg(\log q)^{D_0}$.
The constants implicit in the packet widths are absorbed by increasing
$D_0$.
The large-amplitude and right-hand windows retain the negative power
in \eqref{dec:classicalrankbound}, while the intermediate-amplitude
range has the stronger fixed power saving.
The major arcs may subsequently be enlarged, since an arbitrary
frequency mask is already allowed in $\gamma_a$.

Lemma~\ref{dec:outputnorm}, applied to $\mathcal K_q$, combines
all frequency and amplitude packets with only a logarithmic loss.
This also handles every cross term between different packets;
no disjointness of their output supports is needed.

For a fixed admissible prefix $v$, separate the product restriction
$vq\le mr<(v+1)q$ by a truncated Perron formula with half-integer
endpoints. Take real part $1/\log q$ and height $q^4$.
The truncation error for each product is $O_v(q^{-3})$, and the
integral kernel has total absolute mass $O(\log q)$.
The integral merely multiplies the arbitrary coefficients on the two
sides by $m^{it},r^{it}$ and bounded real-part factors.
The factor $1_{(mr,b)=1}$ is absorbed separately into the two coefficient
sequences. Hence \eqref{dec:minorarcII} holds throughout the range
\eqref{dec:rawwindow}. The argument permits a further restriction of
the frequency mask, so the major arcs may be enlarged as needed in
Section~\ref{dec:prime-section}.
\end{proof}


\section{A common comparison model and the prime asymptotic}
\label{dec:prime-section}

Assume the digital hypotheses of Theorem~\ref{dec:digital-criterion}.
We now pass from Proposition~\ref{dec:minorarc} to a comparison of prime
counts. The comparison model retains the initial digits, which control
physical intervals, and the terminal digits, which control local
congruences. Averaging only the middle block makes the model nonnegative
and permits the same sequence to satisfy Type I and Type II estimates.

\subsection{The model and the common analytic estimates}
\label{dec:model}

Fix an admissible positive prefix $v$, and let $x=(v+1)q$.
Choose a sufficiently large fixed $C$ according to the final error
precision required, and put
\[
 L=\lceil C\log_b\log q\rceil,\quad B_0=b^L,\quad
 Y=q/B_0^2,\quad \rho=(b-1)^{k-2L}/Y.
\]
For sufficiently large $k$, we have $2L<k$.
For $n=vq+a(q/B_0)+B_0m+c$, with $0\le a,c<B_0$ and
$0\le m<Y$, define
\begin{equation}\label{dec:modeldef}
 f_*(n)=1_{n-vq\in\mathcal D_k}1_{(n,b)=1},\qquad
 G_C(n)=\rho\,1_{a\in\mathcal D_L}1_{c\in\mathcal D_L^*},\qquad
 W_C=f_*-G_C,
\end{equation}
and set these functions equal to zero outside this support.
Here $\mathcal D_L^*=\mathcal D_L\cap\{n:(n,b)=1\}$.
If $h=\#\{j:0\le j<b, j\ne d, (j,b)=1\}$, then
\[
 \sum_n f_*(n)=h(b-1)^{k-1}=\sum_nG_C(n).
\]
The masses also agree exactly on each cylinder obtained by fixing
the initial $L$ suffix digits. Thus the model preserves the actual
initial-digit structure at that scale.

Put $\Gamma_0=50/77$ and
\[
 \Gamma_*=\frac{\Gamma_0\alpha}{\alpha+\Gamma_0(1-\alpha)}>.629.
\]

\begin{proposition}[A common comparison sequence]\label{dec:commonmodel}
Fix an admissible positive prefix $v$, an exponent $\Gamma<\Gamma_*$,
and $D,K>0$. There is $C_0$ such that, for every fixed $C\ge C_0$,
the same sequence $W_C=f_*-G_C$ satisfies
\begin{equation}\label{dec:weightedI}
 \sum_{r\le x^\Gamma}\tau(r)^D\max_{J\ {\rm interval}}
 \left|\sum_{m\in J}W_C(rm)\right|
 \ll_{v,D,K}A\log^{-K}q
\end{equation}
and
\begin{equation}\label{dec:fullII}
 \left|\sum_{m\in I_M,r\in I_R}\alpha_m\beta_rW_C(mr)\right|
 \ll_{v,D,K}A\log^{-K}q
\end{equation}
whenever $MR\asymp_vq$, $q^{.249999}\le R\le q^{.4244}$,
$I_M\subset[M,2M]$, $I_R\subset[R,2R]$ are intervals, and
$|\alpha_m|\le\tau(m)^D$, $|\beta_r|\le\tau(r)^D$.
The latter estimate also allows a common interval restriction on $mr$.
The constants are uniform in the truncations and coefficient
sequences for fixed $v,\Gamma,D,K,C$.
\end{proposition}

The Type I estimate is proved directly from the first Fourier moment.
For Type II, the minor arcs have already been treated; the remaining
step is to compare the full major-arc projection with $G_C$.
We then extend the coefficient class using a divisor-moment bound.

\subsection{Maximal Type I with divisor weights}

\begin{lemma}[Maximal weighted Type I]\label{dec:typeI}
Given $\Gamma<\Gamma_*$ and $D,K>0$, if $C$ is sufficiently
large, then the maximal estimate \eqref{dec:weightedI} holds.
\end{lemma}
\begin{proof}
Put $Z=q/b$ and $Q=x^\Gamma$. Only $(r,b)=1$ matters: both
sequences vanish on multiples of every prime dividing $b$. By
\eqref{dec:moment1} and integration over the shift,
$\int F_Y\ll Y^{-\Gamma_0}$ at every radix scale $Y$, with
$1/2<\Gamma_0=1-\gamma_0<\alpha$. Use the coprime-modulus estimate
\eqref{arith:typeI-max} for the free word of scale $Z$. This part of
the proof needs no prime-power assumption: it uses the radix interval decomposition,
separation of reduced fractions, and contraction of the orbit $b^j a/t$
for $(a,t)=(t,b)=1$, $t>1$. These hold for $b=10$ as for $b=11$:
multiplication by $b$ preserves nonintegrality modulo $t$, and each
alphabet contains consecutive digits. Reduction of fractions and the
summation of divisor weights in that proof are independent of the base.

Writing $n=vq+bu+j$ with allowed unit last digit $j$ leaves a word
$u\in\mathcal D_{k-1}$. For an interval $I$ the condition $n\in I$
restricts $u$ to an interval; $r\mid n$ specifies one residue of $u$
modulo $r$. Thus the discrepancy from $r^{-1}\sum_{n\in I}f_*(n)$
is at most $hD_Z(r)$ in the notation of Proposition~\ref{arith:typeI},
uniformly even when $I$ depends on $r$. Since
$\Gamma<\Gamma_*<\Gamma_0$, choose $\epsilon>0$ with
$(1+\epsilon)\Gamma<\Gamma_0$. The sum of these discrepancies with
weights $\tau(r)^D$ is
\[
 \ll A\{e^{-c_1\sqrt{\log Z}}+Q^{1+\epsilon}Z^{-\Gamma_0}\}
 \ll_{v,D,K,\Gamma} A\log^{-K-2}q.
\]
This uses the interval maximum in \eqref{arith:typeI-max}, not just
distribution of the untruncated digital mass.

Let $F_*(t)=\sum_{n\le t}f_*(n)$ and
$G_0(t)=\sum_{n\le t}G_C(n)$.
The masses of the two functions agree on every complete cylinder with
its initial $L$ suffix digits fixed, so
$\sup_t|F_*(t)-G_0(t)|\ll A(b-1)^{-L}$.
For each initial interval, compare the actual progression count with
$F_*(t)/r$ and the model progression count with $G_0(t)/r$.
Their different reference masses therefore cost at most
$A(b-1)^{-L}/r$; arbitrary intervals are differences of two
initial intervals.
Once the initial and terminal words are fixed, the condition $r\mid n$
selects a unique residue class of the middle variable $m$.
The error for the model, for each modulus and arbitrary truncation,
is therefore at most
$\rho\#\mathcal D_L\#\mathcal D_L^*\ll AB_0^2/q$.
Subtracting the model from the actual function gives
\begin{align*}
 \sum_{r\le x^\Gamma}\tau(r)^D\max_J
 \left|\sum_{m\in J}W_C(rm)\right|
 \ll{}&A\log^{-K-2}q
   +Aq^{\Gamma-1}B_0^2\log^{T_D}q\\
 &+A(b-1)^{-L}\log^{T_D}q,\qquad T_D=2^{\lceil D\rceil}.
\end{align*}
Taking $C\alpha>K+T_D+3$ proves the result.
The rows with $(r,b)>1$ vanish identically.
\end{proof}

\subsection{The full major arcs and the physical model}

\begin{lemma}[Bilinear comparison of the full major arcs with the model]\label{dec:bridge}
Given $P>0$ and $\theta_0>0$, one can choose a fixed $C$ and then
the major-arc width so that, for $MR\asymp_vq$,
$\min(M,R)\ge q^{\theta_0}$, and arbitrary bounded coefficients,
\begin{equation}\label{dec:bridgeeq}
 \sum_{m,r}\alpha_m\beta_r1_{vq\le mr<(v+1)q}1_{(mr,b)=1}
 \{P_{\mathfrak M}f_{\rm all}(mr)-G_{\rm all}(mr)\}
 \ll_{v,P}A\log^{-P}q.
\end{equation}
Here $f_{\rm all}=1_{\mathcal D_k}$ is viewed as a periodic function,
and $G_{\rm all}$ is the periodic model obtained from
\eqref{dec:modeldef} by removing the unit terminal-digit restriction.
\end{lemma}
\begin{proof}
Take the major-arc denominator bound to be
$R_0=\exp(\eta\sqrt{\log q})$, with $\eta>0$ fixed and sufficiently
small, and the offset bound to be a sufficiently large fixed
logarithmic power $H$.
Write a reduced denominator $r=r_b t$, where all primes dividing
$r_b$ divide $b$ and $(t,b)=1$. If $t>1$, multiplication by
$b^{O_b(\log(2r))}$ removes $r_b$ without cancelling $t$. On the next
digits up to position $k/2$, every $O_b(\log(2t))$ consecutive orbit
points include one a fixed distance from the integers, as in the proof
of Proposition~\ref{arith:typeI}. The offset $O(H/q)$ perturbs these
points by $O(H/\sqrt q)=o(1)$. Hence
$F_q\ll\exp(-c\log q/\log(2R_0))$ there, by the one-digit product.
There are $O(R_0^2(1+H))$ such frequencies. Taking $\eta$ sufficiently
small, their total pairing is $O(Ae^{-c'\sqrt{\log q}})$ by the trivial
bilinear bound $O_v(q)$ and the Fourier inversion factor $1/q$.

Absorb the unit restriction into the coefficients and write
\[
 \mathcal B(\xi)=\sum_{m\in I_M,\,n\in I_R}\alpha_m\beta_n
                 1_{vq\le mn<(v+1)q}e(-mn\xi).
\]
Write $\mathcal B_0$ for the same sum without the product cutoff.
For a reduced denominator $r$ supported on primes dividing $b$ and
an integer offset $z$, the bilinear Cauchy estimate gives
\begin{equation}\label{dec:classicalbilinear}
 |\mathcal B(c/r+z/q)|
 \ll\frac{q\log^2q}{\sqrt{r(1+|z|)}}+q^{1-\theta_0/2+o(1)}.
\end{equation}
For a reduced fraction $a/s$ with $|\xi-a/s|\le s^{-2}$, Cauchy's inequality,
geometric sums, and grouping by separation give
$|\mathcal B_0(\xi)|^2\ll MR(MR/s+M+R+s)\log(2s)$.
For $z=0$, take $s=r$.
For $z\ne0$, apply Dirichlet approximation again with denominator
bound $\lfloor2q/(r|z|)\rfloor$.
The old fraction $c/r$ is excluded, and the lower bound $1/(rs)$
for the difference of two distinct fractions forces
$s\asymp q/(r|z|)$.
Substitution into the preceding estimate, followed by Perron's
formula to retain the product truncation, proves
\eqref{dec:classicalbilinear}.

Choose radix powers $U\asymp r$, $V\asymp Z$ with $UV\le q$.
Keeping the two end blocks bounds $F_q(c/r+z/q)$ by
$F_U(c/r+z/q)F_V(z/V)$, because $r$ divides $q/V$ for large $q$.
Variation bounds the sum over $c$ of the first factor's local suprema
by $O(r^{\gamma_0})$. The shifted first moment sums the second factor
over $z$ at cost $O(Z^{\gamma_0})$. Thus for each dyadic $Z\ge1$,
\begin{equation}\label{dec:packlone}
 \sum_{(c,r)=1}\sum_{|z|\asymp Z}
 |\widehat f_{\rm all}(c/r+z/q)|\ll A(rZ)^{\gamma_0}.
\end{equation}
The case $Z=1$ also includes $z=0$.
The chosen digit blocks are disjoint because $r,Z=q^{o(1)}$.

To control the physical boundaries of the model, let $Q_0=q/B_0$
and $T=\lfloor\log^{9C/8}q\rfloor$, and use the probability kernel
on $\mathbb Z/Q_0\mathbb Z$ given by
\[
 K_T(u)=\frac1{Q_0T}\left|\sum_{j<T}e_{Q_0}(ju)\right|^2,\qquad
 G_T(n)=\sum_uK_T(u)G_{\rm all}(n-B_0u).
\]
These translations preserve the last $L$ digits.
The Fourier support is $a=cq/B_0+z$ with $|z|<T$, and the
Fourier coefficients are multiplied by $(1-|z|/T)_+$.
The Fej\'er tail bound is
$\sum_{|u|\ge h}K_T(u)\ll\min(1,Q_0/(Th))$.
Thus the smoothing error at physical distance $D$ from an initial-word
boundary is at most $O(\rho\min(1,q/(TD)))$.
The absolute coefficient of any multiplicative test is at most
$\tau_2(n)$, and the Dirichlet hyperbola decomposition gives
\[
 \sum_{n\in J}\tau_2(n)\ll_v |J|\log q+\sqrt q
\]
for every interval $J$ in the range $n\asymp_vq$.
Decomposing by distance from the at most $B_0$ initial-word boundaries
gives
\begin{equation}\label{dec:smoothing}
 \sum_n\tau_2(n)|G_{\rm all}(n)-G_T(n)|
 \ll A B_0^{3-2\alpha}T^{-1}\log^2q
      +Aq^{-1/2}B_0^{3-2\alpha}.
\end{equation}

The function $G_T$ also satisfies \eqref{dec:packlone}. Its support
in the packet requires $r\mid B_0$, since distinct rational centres
of denominators $r,B_0$ are separated by at least $1/(rB_0)$, greater
than the sum of their polylogarithmic widths divided by $q$. For
$r\mid B_0$, the normalized magnitude of $G_{\rm all}$'s
Fourier transform is
\[
 F_{B_0}(c/r+z/q)F_{B_0}(z/B_0)
       \frac{|D_Y(z/(B_0Y))|}{Y},\qquad
 D_Y(t)=\sum_{m<Y}e(mt).
\]
The last factor is at most $\min(1,O(B_0/(1+|z|)))$, and the Fej\'er multiplier defining $G_T$ has modulus at most one.
For $Z\le B_0$, use the first moment.
For $Z>B_0$, periodic repetition of the second factor gives
$O(ZB_0^{\gamma_0-1})$; multiplication by $O(B_0/Z)$ still
leaves at most $O(Z^{\gamma_0})$.

Put $S=B_0^{1/10}$.
For $b=10$ or $11$, every smooth denominator $r\le S$ divides $B_0$.
When $r,|z|\le S$, the actual digital measure and the model differ only
in the middle block. The normalized measures both have mass one, and the phase
variation across that block is $O(|z|/B_0)$. Hence
\[
 |\widehat f_{\rm all}(c/r+z/q)-\widehat G_T(c/r+z/q)|
 \ll A|z|(B_0^{-1}+T^{-1}).
\]
The total pairing error from the small packets is at most
\[
 O(AS^4(B_0^{-1}+T^{-1})).
\]
For the remaining packets, multiply \eqref{dec:classicalbilinear}
by \eqref{dec:packlone}. The main term in each packet is
\[
 A\log^2q(rZ)^{-(1/2-\gamma_0)}.
\]
The sums over smooth denominators and dyadic offset scales converge.
The total from the second term is $Aq^{-\theta_0/2+o(1)}$.

The logarithmic exponents in the three error terms without a power
saving in $q$ are, respectively,
\[
 -C(2\alpha-15/8)+2,\qquad -3C/5,\qquad
 -C(23/154)/10+2+o(1).
\]
It is enough to take $C\ge100(P+10)$, and then to take the
major-arc offset exponent at least $9C/8+1$, increasing it further
to obtain the precision required in \eqref{dec:minorarcII}.
Once $C$ is fixed, all these logarithmic scales are eventually below
$R_0$. Adding \eqref{dec:smoothing} proves the result.
The physical unit conditions are absorbed separately into the two
coefficient sequences; the two periodic functions being compared
remain precisely those in \eqref{dec:bridgeeq}.
The argument works for every fixed $C$ above the stated threshold,
with the major-arc widths chosen afterwards. This quantifier permits
the same $C$ to be used in Type I and Type II.
\end{proof}

Combining \eqref{dec:minorarcII} with Lemma~\ref{dec:bridge} gives
an estimate for the same $W_C$:
given $v,T$, after a single choice of sufficiently large $C=C(v,T)$,
for arbitrary bounded coefficients, $MR\asymp_vq$, and
$q^{.249999}\le R\le q^{.4244}$, we obtain \eqref{dec:fullII}
with $K=T$.
An additional common product truncation is also permitted.

\subsection{Extension to divisor-bounded coefficients}

First, for fixed $J>0$, we record the growth bound
\begin{equation}\label{dec:divisormoment}
 \sum_n(f_*(n)+G_C(n))\tau(n)^J
 \ll_{v,J,C}A\log^{C_J}q,\qquad C_J=2^{\lceil5J\rceil}+1.
\end{equation}
A direct proof is as follows.
Remove the prime factors greater than $x^{1/4}$.
There are at most four, counted with multiplicity, and their
contribution to the divisor function is at most $16$.
The remaining prime factors can be grouped greedily into at most five
blocks, each with product at most $\sqrt x$. Thus
\[
 \tau(n)^J\le16^J\sum_{r\mid n, r\le\sqrt x}\tau(r)^{5J}.
\]
For $f_*$, use the weighted arithmetic-progression bound from the
proof of Lemma~\ref{dec:typeI}.
For $G_C$, use the bound $O(A/r+AB_0^2/q)$ for each modulus
from the Chinese remainder theorem.
This proves \eqref{dec:divisormoment}.
Its logarithmic exponent $C_J$ is independent of the model parameter $C$.

Suppose that $|\alpha_m|,|\beta_r|\le\tau(m)^D,\tau(r)^D$.
For a target precision $K$, take
\[
 J=3D+1,\quad L_*=C_J+K+3,\quad K_*=K+2L_*+3.
\]
Truncate both coefficient sequences at $\log^{L_*}q$.
For the small coefficients, use \eqref{dec:fullII} with precision $K_*$.
For the large coefficients, sum by the actual product $t=mr$.
There are at most $\tau(t)$ factorizations, the product of the two
weights is at most $\tau(t)^{2D}$, and the exceptional indicator is
at most $\tau(t)^D/\log^{L_*}q$.
By \eqref{dec:divisormoment}, the total is $O(A\log^{-K-3}q)$.
There are only $O_v(\log q)$ nonempty dyadic boxes.
Thus \eqref{dec:fullII} extends to coefficients bounded by any fixed
power of the divisor function. We can now conclude the common-model proposition.

\begin{proof}[Proof of Proposition~\ref{dec:commonmodel}]
Choose $J,L_*,K_*$ as above. The Type I estimate requires
$C\alpha>K+2^{\lceil D\rceil}+3$.
The bounded-coefficient Type II estimate with precision $K_*$
requires only a further lower bound for $C$.
Its extension to divisor-bounded coefficients uses a logarithmic
moment exponent $C_J$ independent of $C$, so the two choices are
compatible. Choose the common major arcs after fixing $C$.
\end{proof}

\subsection{The asymptotic sieve}
\label{dec:sieve}

We now state precisely the asymptotic sieve theorem used
\cite[Theorem 2.2(a)]{FordMaynard2024}.
Let $w_n$ be supported on $(x/2,x]$, and suppose that the parameters
$(\gamma,\theta,\nu)$ satisfy
\[
 \tfrac12\le\gamma\le1-\theta-\nu\quad\text{or}\quad
 1-\theta\le\gamma<1,\qquad
 0\le\theta<\theta+\nu\le\tfrac12.
\]
For a sufficiently large fixed $B_*$, assume
\begin{align}
 &\sum_{m\le x^\gamma}\tau(m)^{B_*}\max_I
 \left|\sum_{\substack{x/2<mn\le x\\n\in I}}w_{mn}\right|
 \le x\log^{-B_*}x,\label{dec:FM_I}\\
 &\left|\sum_{\substack{(x/2)^\theta<m\le x^{\theta+\nu}\\x/2<mn\le x}}
 \xi_m\kappa_nw_{mn}\right|\le x\log^{-B_*}x,\quad
 |\xi_m|\le\tau(m)^{B_*}, |\kappa_n|\le\tau(n)^{B_*},\label{dec:FM_II}\\
 &\sum_n|w_n|\tau(n)\le x\log^\varpi x,\qquad
 w_n\ge-x^{\nu/10}.\label{dec:FM_growth}
\end{align}
Let $M=\lfloor(1-\gamma)^{-1}\rfloor$.
Suppose that for every integer $n\ge M+1$ there is an
$a\in\mathbb N$ with $a/n\in[\theta,\theta+\nu]$, and that
there is a positive integer $h$ such that
\[
 h(1-\gamma)\in[\theta,\theta+\nu]\cup[1-\theta-\nu,1-\theta],
\]
Then, for given $A_{\rm out}>1$ and $\varpi\ge1$, taking $B_*$
sufficiently large gives $\left|\sum_pw_p\right|\ll x\log^{-A_{\rm out}}x$.

We use the exact parameters
\begin{equation}\label{dec:FMparameters}
 \begin{aligned}
 \theta&=499999/2000000,&\qquad \gamma&=1500001/2000000,\\
 \nu&=348001/2000000,& \theta+\nu&=53/125.
 \end{aligned}
\end{equation}
Use the intermediate short-factor exponents $l=.24999925$ and
$u=.4242$, strictly inside \eqref{dec:rawwindow}. If
$x^{1-u}\le m\le(x/2)^{1-l}$ and $x/2<mn\le x$, then
$(x/2)^l<n\le x^u$. Proposition~4.11(b) of \cite{FordMaynard2024}
therefore gives maximal Type I on this complementary interval: the long
coefficient carries the row phase and the divisor weight, with one
logarithm lost in taking the row-dependent interval. Summing the Type II
dyadic boxes costs one more logarithm. Since $1-u=.5758<.629$ and
$1-l=.75000075>\gamma$, this interval overlaps the direct Type I range
and reaches $x^\gamma$ for large $x$. The strict inclusions
$l<\theta<\theta+\nu<u$ also supply \eqref{dec:FM_II}; all fixed factors
from $x=(v+1)q$ and dyadic endpoints are absorbed by the exponent margins.

The parameters satisfy $\gamma=1-\theta$ and $M=4$.
For $n=5$, take $a=2$.
For $n\ge6$, we have $n\nu>1$, so take $a=\lceil n\theta\rceil$.
For the other combinatorial condition, take $h=3$, since
$3(1-\gamma)=.7499985\in(.576,.7500005)$.

Fix the parameters in the following order, so that all estimates
apply to the same model.
First take $A_{\rm out}=3$ and $\varpi=2$, and let the sieve theorem
determine $B_*$. Choose integers $D\ge B_*+4$ and $K\ge D+4$,
then determine $J,L_*,K_*$ by the extension to divisor-bounded
coefficients above. Finally make a single choice of
\[
 C\ge100(K_*+10),\qquad C\alpha>K+2^{\lceil D\rceil}+3,
\]
and fix a common major-arc width. Put
\[
 a_n=(q/A)f_*(n),\qquad b_n=(q/A)G_C(n),\qquad w_n=a_n-b_n.
\]
These sequences are supported on $(x/2,x]$.
When $v=1$, the unit condition makes them zero at the left endpoint.
Using $\tau(n)\le2\sum_{r\mid n,r\le\sqrt x}1$, the digital
arithmetic-progression estimate already proved, and the Chinese
remainder theorem for the model, we obtain
\[
 \sum_n(f_*+G_C)\tau(n)\ll_{v,C}A\log q.
\]
Thus the total growth condition in \eqref{dec:FM_growth} holds.
For the pointwise negative part, use the actual model value:
\[
 b_n\le(q/A)\rho=B_0^{2-2\alpha}=(\log q)^{O_C(1)}\le x^{\nu/10}
\]
for sufficiently large $q$.
Proposition~\ref{dec:commonmodel} applies to this fixed $C$.
After multiplication by $q/A$, its Type I and Type II bounds become
$O_v(q\log^{-K}q)$ with arbitrarily prescribed $K$.
The strict margins between the $q$-ranges and the displayed $x$-ranges
absorb the fixed factor $x/q=v+1$.
Taking the prescribed logarithmic precision gives
\eqref{dec:FM_I}--\eqref{dec:FM_II}. Hence
\begin{equation}\label{dec:primecomparison}
 \left|\sum_p(f_*(p)-G_C(p))\right|\ll_vA\log^{-3}q.
\end{equation}

\subsection{The leading constant and arbitrary real endpoints}
\label{dec:endpoints}

Fix the prefix $v$ and the model parameter $C$. An admissible top
word $a$ selects $I_a=[vq+a q/B_0,vq+(a+1)q/B_0)$; each admissible unit
bottom word selects one reduced class modulo $B_0$. By Siegel--Walfisz
at the two endpoints, the prime count in each pair is
\[
 \frac1{\varphi(B_0)}\int_{I_a}\frac{dt}{\log t}
       +O_{v,C,H}(q\log^{-H}q).
\]
Here $B_0\ll_b\log^C q$, the intervals have equal length $q/B_0$,
and $\log t=\log q+O_v(1)$ throughout. The total error after weighting
is $O_{v,C}(A\log^{-3}q)$: there are at most $B_0^2$ pairs,
$\rho B_0^2=(A/q)B_0^{4-2\alpha}$, and we take $H>C(4-2\alpha)+3$.
Counting the two end words and using $\varphi(B_0)=B_0\varphi(b)/b$ gives
\begin{align}
 \sum_pG_C(p)&\sim\rho(b-1)^Lh(b-1)^{L-1}
       \frac{q}{B_0\varphi(B_0)\log q}\notag\\
 &=\kappa_{b,d}\frac{(b-1)^k}{k\log b},\qquad
 \kappa_{b,d}=\frac{b\,h}{(b-1)\varphi(b)}.
 \label{dec:prefixasymptotic}
\end{align}
By \eqref{dec:primecomparison} the same asymptotic holds for $f_*$.
Indeed every prime in this interval is a unit for large $k$, and all
$k$ suffix positions are actual digits after the positive prefix, even
when $d=0$. This is \eqref{prelim:prefix-hypothesis}.

The fixed-prefix estimate has the same constant $\kappa_{b,d}$
for every admissible positive prefix $v$. Lemma~\ref{prelim:endpoints}
therefore gives the asserted asymptotic for ordinary expansions at every
real endpoint, completing the proof of
Theorem~\ref{dec:digital-criterion}. The allowed unit last digits give
\[
 \kappa_{10,d}=
 \begin{cases}
 5/6,&d\in\{1,3,7,9\},\\
 10/9,&d\in\{0,2,4,5,6,8\},
 \end{cases}
 \qquad
 \kappa_{11,d}=
 \begin{cases}
 11/10,&d=0,\\
 99/100,&d\in\{1,\ldots,10\}.
 \end{cases}
\]
In particular, Fourier reflection of the deleted digit preserves the
amplitude estimates but need not preserve the arithmetic constant:
in base eleven it interchanges exclusion of $0$ and exclusion of a
nonzero unit. Proposition~\ref{dec:verified-inputs} verifies the digital
hypotheses for every required digit, and hence proves
Theorem~\ref{main:exact}.

\clearpage
\section{Completion of the main results}
\label{sec:completion}

\begin{proof}[Proof of the main theorems]
For the order estimates, Sections~\ref{nine:common-tree},
\ref{sec:octal-prime-count}, \ref{sec:seven} and \ref{six:positive},
and Proposition~\ref{dt:grouped} for the decimal pairs, verify the
arithmetic hypotheses of Theorem~\ref{common:positive-sieve}.
Section~\ref{dt:sieve-certificate} gives a loss strictly below one in
each case. Proposition~\ref{common:upper} gives the matching upper
order. Thus Theorems~\ref{main:order} and~\ref{main:two-decimal}
follow for ordinary expansions, including an excluded zero.

For Theorem~\ref{main:exact} apply Theorem~\ref{dec:digital-criterion}
with the moments and correlation bounds of Appendix~\ref{app:exact-inputs}.
The fixed-prefix completion in Lemma~\ref{prelim:endpoints} retains
the digit count $A_{b,d}(x)$ itself, hence also its fluctuations.
\end{proof}

\bigskip
\noindent\textbf{AI Disclosure:} The author selected the research topic; artificial intelligence completed all mathematical derivations, proofs, and writing.

\appendix
\addtocontents{toc}{\protect\setcounter{tocdepth}{1}}
\section{Continuous bounds from finite grids}\label{app:continuous}

The finite certificates must bound a transfer operator or a correlation
profile throughout its continuous state space. We establish the
interpolation estimates that control the values between grid points and
the curvature bounds needed to apply them. For transfer operators, the
resulting supersolution inequality combines with
Lemma~\ref{prelim:supersolution} to give a moment bound for every word length.

\subsection{Interpolation and transfer supersolutions}

\begin{lemma}[Semiconvex grid interpolation]\label{prelim:grid}
If a continuous function $u$ satisfies $u''\geqslant-C$ distributionally
on an interval $[a,a+\Delta]$, then
\[
 u(t)\leqslant
 \frac{a+\Delta-t}{\Delta}u(a)+
 \frac{t-a}{\Delta}u(a+\Delta)
 +\frac C2(t-a)(a+\Delta-t).
\]
In particular, its supremum is at most the larger endpoint value plus
$C\Delta^2/8$. On a rectangular grid, if the second derivative in coordinate
$j$ is bounded below by $-C_j$, the supremum is at most the largest vertex
value plus $\sum_jC_j\Delta_j^2/8$.
\end{lemma}
\begin{proof}
The function $u(t)+Ct^2/2$ is convex and hence lies below its endpoint
chord. Subtract $Ct^2/2$ and use
$(t-a)(a+\Delta-t)\leqslant\Delta^2/4$. Applying this argument successively
in each coordinate proves the rectangular version.
\end{proof}

All transfer supersolutions used here are periodic linear interpolants
of positive rational node values. The values of a periodic $T_ph$ at
$x=0$ and $x=1$ agree, so tests at the $N$ nodes include the last cell's
right endpoint. A Lipschitz alternative is sometimes convenient:
if the residual $\Lambda h-T_ph$ has Lipschitz constant $L$ on each
cell of length $1/N$, a nodal margin $L/(2N)$ suffices.
For the sharper second-derivative criterion, Lemma~\ref{prelim:grid} gives
\begin{equation}\label{prelim:semiconvex}
 \Lambda h(j/N)-(T_ph)(j/N)\geqslant C/(8N^2)
 \quad(0\leqslant j<N)
 \quad\Longrightarrow\quad T_ph\leqslant\Lambda h,
\end{equation}
provided $T_ph$ is continuous and $1$-periodic and
$(T_ph)''\geqslant-C$ on each open cell.

\begin{lemma}[A common transfer certificate]\label{prelim:operator-cert}
Let $p\geqslant1$ and
\[
 (T_ph)(x)=\sum_{i=0}^{b-1}|Z_i(x)|^p h((i+x)/b),
 \qquad |Z_i|\leqslant1,\quad |Z_i'|\leqslant a,\quad |Z_i''|\leqslant c.
\]
Suppose that $h$ is positive, $1$-periodic, linear on the $1/N$ grid,
and $\|h\|_\infty\leqslant1$. Assume also that $T_ph$ is
$1$-periodic. If
\[
 \Lambda h(j/N)-(T_ph)(j/N)
 \geqslant\frac{bpc+2pa\operatorname{Lip}(h)}{8N^2}
 \qquad(0\leqslant j<N),
\]
then $T_ph\leqslant\Lambda h$ everywhere.
\end{lemma}
\begin{proof}
Convexity of $z\mapsto|z|^p$ gives, distributionally,
\[
 (|Z_i|^p)''\geqslant-pc,\qquad
 \bigl|(|Z_i|^p)'\bigr|\leqslant pa.
\]
For completeness, the Hessian term in the chain rule is nonnegative,
while the remaining term is bounded below by
$-p|Z_i|^{p-1}|Z_i''|\geqslant-pc$. At zeros, use the convex
regularization $(|Z_i|^2+\epsilon)^{p/2}$ and pass to the limit.
A breakpoint $(i+x)/b=m/N$ can occur only at $x=(bm-iN)/N$,
a grid node. Thus each branch interpolant is affine inside each input
cell, with slope of magnitude at most $\operatorname{Lip}(h)/b$.
The product rule therefore gives
\[
 (T_ph)''\geqslant-bpc-2pa\operatorname{Lip}(h)
\]
on that cell. Apply \eqref{prelim:semiconvex}.
\end{proof}

For the digital operators the branch magnitudes have the form
$|Z_i(x)|=f((i+x)/b)$, with $f$ and $h$ both $1$-periodic.
The change $x\mapsto x+1$ cyclically permutes these branches,
which gives the required periodicity of $T_ph$.

\subsection{Curvature of affine-correlation profiles}

The affine-correlation certificates require the same interpolation principle
in two variables. The following bounds remain valid at zeros of the
one-digit amplitudes.

\begin{lemma}[Curvature of Hellinger profiles]\label{prelim:hellinger-curvature}
Let $A(x)=(A_0(x),\ldots,A_{b-1}(x))$ be a $C^2$ complex vector with
\[
 \|A(x)\|_2=1,\qquad \|A'(x)\|_2\leqslant a,\qquad
 \|A''(x)\|_2\leqslant c.
\]
Put $P_i(x)=|A_i(x)|^2$,
\[
 H(x,E)=\sum_i|A_i(x)||A_i(x+E)|,\qquad
 \rho(E)=\sup_{x\in[0,1]}H(x,E).
\]
For arbitrary fixed shifts $s_i$, define
$G(x,t)=\sum_iP_i(x)\rho(t+s_i)^2$. Then
$0\leqslant H,\rho,G\leqslant1$, and
\[
 H_{xx}\geqslant-2c-2a^2,\qquad H_{EE}\geqslant-c,\qquad
 (\rho^2)''\geqslant-2c,\qquad
 G_{xx}\geqslant-2c,\quad G_{tt}\geqslant-2c
\]
in the sense of distributions.
\end{lemma}
\begin{proof}
Cauchy--Schwarz gives $H\leqslant1$, and the other range bounds follow.
To justify differentiation at zeros, set
$a_{i,\epsilon}=(|A_i|^2+\epsilon^2)^{1/2}$. Direct differentiation gives
\[
 |a_{i,\epsilon}'|\leqslant|A_i'|,
 \qquad a_{i,\epsilon}''\geqslant-|A_i''|,
 \qquad \|(a_{i,\epsilon})_i\|_2=(1+b\epsilon^2)^{1/2}.
\]
The product rule and Cauchy--Schwarz therefore bound the second
$x$ derivative of the regularized $H$ below by
$-2c(1+b\epsilon^2)^{1/2}-2a^2$, and its second $E$ derivative below by
$-c(1+b\epsilon^2)^{1/2}$. Let $\epsilon\downarrow0$.

For fixed $x$, squaring the nonnegative function $H(x,E)\leqslant1$
gives $(H(x,\cdot)^2)''\geqslant-2c$. Suprema preserve a common
semiconvexity bound, because adding $cE^2$ makes every such function
convex. Since $\rho^2=\sup_x H(x,\cdot)^2$, this proves the bound for
$\rho^2$, and averaging with the probabilities $P_i(x)$ gives $G_{tt}$.
Finally,
\[
 G_{xx}
 =2\sum_i\rho(t+s_i)^2
       \bigl(|A_i'|^2+\operatorname{Re}(A_i''\overline{A_i})\bigr)
 \geqslant-2\|A''\|_2\|A\|_2\geqslant-2c.
\]
\end{proof}

\subsection{Rational enclosures}

Trigonometric enclosures use Machin's identity
\[
 \pi=16\arctan(1/5)-4\arctan(1/239),
\]
alternating-series bounds, and Taylor bounds for cosine with directed
rounding. A rational-power bound $(a/D)^{u/v}\leqslant w/D$ is certified
by the integer comparison
\begin{equation}\label{prelim:integer-power}
 w^vD^{u-v}\geqslant a^u\qquad(u\geqslant v).
\end{equation}
The parameters and strict margins for the moment and affine-correlation
bounds in bases ten and eleven are given in Appendix~\ref{app:exact-inputs}.
For base eight see Lemmas~\ref{octal:uniformTypeI} and
\ref{octal:near-parseval}; for base nine see Appendix~\ref{app:nine-moments}.


\section{Unit-circle zeros and one-sided shifts}\label{app:shifts}

The shift comparison used in the carry decomposition depends on the
arithmetic location of the zeros of the digit polynomial. Their orders
are coprime to the base, so an unshifted digital argument stays a controlled
distance from every zero. We first classify these zeros and then compare
the full digital products, retaining the multiplicity of each zero.

\begin{lemma}[Zeros of a one-digit-deleted polynomial]\label{prelim:zeros}
For $b\geqslant3$ and $0\leqslant d<b$, the unit-circle zeros of
$P_{b,d}(z)=\sum_{j=0}^{b-1}z^j-z^d$ are exactly the points $z\ne1$
satisfying
\begin{equation}\label{prelim:zero-classification}
 \bigl(z^{b-1}=1,\ z^d=1\bigr)
 \quad\hbox{or}\quad
 \bigl(z^{b+1}=1,\ z^{d+1}=-1\bigr).
\end{equation}
Their orders are coprime to $b$. Every such zero is simple, except that
$z=-1$ is a double zero when $b\equiv1\pmod4$ and $d=(b-1)/2$.
\end{lemma}
\begin{proof}
Since $P_{b,d}(1)=b-1$, consider $z=e^{i\theta}\ne1$.
If $P_{b,d}(z)=0$, the geometric sum has modulus one, so
\[
 \left|\frac{\sin(b\theta/2)}{\sin(\theta/2)}\right|=1.
\]
Thus $\cos(b\theta)=\cos\theta$, which gives $z^{b-1}=1$ or
$z^{b+1}=1$. The geometric sum equals $1$ in the first case and
$-z^{-1}$ in the second; this proves both necessity and sufficiency
of \eqref{prelim:zero-classification}.

At a zero of the first type and of the second type, respectively,
\[
 zP_{b,d}'(z)=-\frac{(b-1)z}{1-z}-d,\qquad
 zP_{b,d}'(z)=z^{-1}\frac{(d-b)-(d+1)z}{1-z}.
\]
Either derivative can vanish on the unit circle only when
$d=(b-1)/2$ and $z=-1$. This point is a zero precisely when $d$ is
even, equivalently $b\equiv1\pmod4$. Writing $b=2d+1$ in that case,
\[
 P_{b,d}''(-1)=2d^2-d(d-1)=d(d+1)>0,
\]
so its multiplicity is exactly two. Finally, the orders in
\eqref{prelim:zero-classification} divide $b-1$ or $b+1$.
\end{proof}

\begin{lemma}[One-sided shift comparison]\label{prelim:shift}
Fix $b,d$ and $K>0$. Let $q$ and $Q=b^L\leqslant q$ be powers of $b$,
let $A\in\mathbb Z$, and let $|\delta|\leqslant1$. If
$F_Q((A+\delta)/Q)\geqslant q^{-K}$, then
\begin{equation}\label{prelim:shift-bound}
 F_Q((A+\delta)/Q)
 \leqslant \exp\!\bigl(O_{b,d,K}(\sqrt{\log q})\bigr)F_Q(A/Q).
\end{equation}
\end{lemma}
\begin{proof}
For $1\leqslant j\leqslant L$, write
$u_j=f_{b,d}(A/b^j)$ and $v_j=f_{b,d}((A+\delta)/b^j)$.
We first prove the uniform single-factor estimates
\begin{equation}\label{prelim:factor-comparison}
 u_j>0,\qquad v_j/u_j\ll_{b,d}1,\qquad |v_j-u_j|\ll_{b,d}b^{-j}.
\end{equation}
The last bound follows from the Lipschitz property of $f_{b,d}$.
For the first two, factor the polynomial as
\[
 P_{b,d}(z)=R(z)\prod_{\zeta\in\mathcal Z}(z-\zeta)^{m_\zeta},
\]
where $\mathcal Z$ is the set of unit-circle zeros and $R$ has none
on that circle. Thus $|R|$ is bounded above and away from zero there.
Write $\zeta=\mathrm e(a/r)$ in lowest terms. Lemma~\ref{prelim:zeros}
gives $(r,b)=1$, so
\[
 \|A/b^j-a/r\|\geqslant (rb^j)^{-1}.
\]
Consequently the denominator in
\[
 \frac{|\mathrm e((A+\delta)/b^j)-\zeta|}
      {|\mathrm e(A/b^j)-\zeta|}
 \leqslant
 1+\frac{2\pi b^{-j}}{4/(rb^j)}
\]
is nonzero and the ratio is uniformly bounded. Taking the product,
with the actual multiplicities $m_\zeta$, proves
\eqref{prelim:factor-comparison}.

At positions with $u_j>b^{-j/2}$, the ratio is
$1+O_{b,d}(b^{-j/2})$, whose product is bounded. At the other positions,
$v_j\ll_{b,d}b^{-j/2}$, and hence $v_j\leqslant b^{-j/4}$ once $j$
exceeds a fixed constant. Since $\prod_jv_j\geqslant q^{-K}$ and
$v_j\leqslant1$ at every position,
\[
 \sum_{\substack{j\text{ exceptional}\\j\text{ sufficiently large}}}
 j\leqslant4K\log_bq.
\]
There are therefore $O_{b,d,K}(\sqrt{\log q})$ exceptional positions.
Each costs a bounded ratio by \eqref{prelim:factor-comparison}, proving
\eqref{prelim:shift-bound}.
\end{proof}


\section{Verification of the exact-asymptotic digital inputs}
\label{app:exact-inputs}\label{sec:eleven}

This appendix proves Proposition~\ref{dec:verified-inputs}.
The moment exponents at orders $1$ and $235/154$ in base ten come
from Maynard \cite[Lemmas~10.3--10.4]{Maynard2019}. Karwatowski
\cite[\S\S3, 5, in particular Lemma~5.1]{Karwatowski2022} gives the corresponding
eigenvalue criteria and reports their verification in base eleven. We give rational certificates for both bases, including
the quantitative $3/2$-moments and the affine contraction used here.
Karwatowski also uses a $3/2$-moment for the base-nine endpoint digits
\cite[\S3, p.~5]{KarwatowskiBaseNine}. The continuous
certification principles are proved in Appendix~\ref{app:continuous};
the exact rational data and their verification are specified in
Appendix~\ref{app:reproduce}.

\subsection{Uniform shifted moments}

For each order $p$ we seek a positive periodic linear interpolant
$h$ on a rational grid and a number $\Lambda\le b^g$ such that
$T_ph\le\Lambda h$. By \eqref{prelim:transfer} and
Lemma~\ref{prelim:supersolution}, this single inequality gives
\begin{equation}\label{dec:transferiterate}
 H_p(b^k;x)=T_p^k1(x)\le(\max h/\min h)\Lambda^k
 \qquad(k\ge0,\ x\in\mathbb R).
\end{equation}
The branch magnitudes are $f((i+x)/b)^p$; incrementing $x$ permutes
the branches, so $T_ph$ is continuous and periodic, including across
the last grid cell. Lemma~\ref{prelim:operator-cert} therefore applies
to the nodal residuals, also when a branch amplitude vanishes.

\subsubsection{Base ten}

\begin{lemma}\label{dec:decimalcert}
For $b=10$, estimates \eqref{dec:moment1}--\eqref{dec:momentnew} hold.
\end{lemma}
\begin{proof}
Put
\[
 Z_i(x)=\frac19\sum_{\substack{0\le j<10\\j\ne d}}
                         e((j-9/2)(i+x)/10).
\]
The finite sums and the inequality $2\pi<44/7$ give
\[
 |Z_i|\le1,\quad |Z_i'|\le110/63,\quad |Z_i''|\le2662/735.
\]
Let $N$ be divisible by $10$, and let $h$ be the periodic linear
interpolant of the certified nodal vector, normalized by
$\|h\|_\infty\le1$, with Lipschitz constant $L_h$.
Lemma~\ref{prelim:operator-cert}, with
$a=110/63$ and $c=2662/735$, shows that it suffices to choose
\[
 c_0\ge10p(2662/735),\qquad c_1\ge2p(110/63),
\]
and verify, at every node $j/N$, that
\begin{equation}\label{dec:nodecert}
 \Lambda h(j/N)-T_ph(j/N)>\frac{c_0+c_1L_h}{8N^2}.
\end{equation}
The rational parameters used are
\[
\begin{array}{c|c|c|c|c}
 p&\text{exponent}&\Lambda&N&(c_0,c_1)\\ \hline
 1&27/77&1121/500&1000&(37,4)\\
 235/154&59/433&1368543/1000000&1000&(56,6)\\
 3/2, d\notin\{4,5\}&287/2000&13915537/10000000&1000&(55,6)\\
 3/2, d\in\{4,5\}&14387/100000&6963699/5000000&2000&(55,6)
\end{array}
\]
The rational nodal vectors and their exact verification are included in
the supplementary material described in Appendix~\ref{app:reproduce}.
Trigonometric values are enclosed using a rational interval for $\pi$ from
Machin's identity and Taylor polynomials with rigorous remainder bounds.
The identity
\[
 f(t)^2=\frac{9+2\sum_{r=1}^9c_r\cos(2\pi rt)}{81},\qquad
 c_r=\#\{j:j,j+r\ne d, 0\le j<j+r\le9\},
\]
allows each nodal value to be rounded upwards to a rational with denominator
$D=2^{40}$. If $p=u/v$ and $f\le a/D$, the integer condition
$w^vD^{u-v}\ge a^u$ ensures that $f^p\le w/D$.
After denominators are cleared, every margin in \eqref{dec:nodecert} is
a positive integer. The inequalities $\Lambda\le10^{\text{exponent}}$
are also checked by comparing integer powers. The five representative digits
$d=0,\ldots,4$, together with $f_{9-d}=f_d$, cover all digits.
The conclusion follows from \eqref{dec:transferiterate}.
\end{proof}

\subsubsection{Base eleven}

Put $D_d=\{0,\ldots,10\}\setminus\{d\}$ and $f_d=f_{11,d}$.
The reflection
$f_{10-d}=f_d$ reduces the verification to $d=0,\ldots,5$.

\begin{lemma}[Uniform moments in base eleven]\label{eleven:moments}
For every $d\in\{0,\ldots,10\}$, every $q=11^k$, and every real $x$,
\begin{align}
 \sum_{a<q}F_q((a+x)/q)&\ll q^{27/77},\label{eleven:l1}\\
 \sum_{a<q}F_q((a+x)/q)^{235/154}&\ll q^{59/433},\label{eleven:oldmoment}\\
 \sum_{a<q}F_q((a+x)/q)^{3/2}&\ll q^{287/2000}.\label{eleven:newmoment}
\end{align}
The implied constants are independent of $k,x$.
\end{lemma}
\begin{proof}
Each certificate uses $N=1100$ nodes and common denominator $D=2^{40}$,
specifying
\[
 h(j/N)=v_j/D,\qquad 0<v_j\leqslant D,
\]
with periodic linear interpolation between nodes. The positive vectors
$v$ and the integer arrays of nodal weights are given in the
supplementary file \path{supplement/base11/witnesses.py};
see Appendix~\ref{app:reproduce}. The three sets of parameters are
\begin{center}
\begin{tabular}{ccccc}
\toprule
$p$ & $g$ & $\Lambda$ & $c_0$ & $c_1$\\
\midrule
$1$ & $27/77$ & $23/10$ & $41$ & $4$\\
$235/154$ & $59/433$ & $137/100$ & $62$ & $6$\\
$3/2$ & $287/2000$ & $7/5$ & $62$ & $6$\\
\bottomrule
\end{tabular}
\end{center}
In each row, $\Lambda\leqslant11^g$ is verified by an integer-power
comparison.

To show that these data control the continuous state space, let
\[
 Z_i(x)=\frac1{10}\sum_{a\in D_d}
          \mathrm e((a-5)(i+x)/11).
\]
Using $\sum_{a=0}^{10}|a-5|=30$ and
$\sum_{a=0}^{10}(a-5)^2=110$, we obtain
\[
 |Z_i|\leqslant1,\qquad
 |Z_i'|\leqslant\frac{2\pi\cdot30}{110}\leqslant\frac{110}{63},
 \qquad
 |Z_i''|\leqslant\frac{4\pi^2\cdot110}{1210}
       \leqslant\frac{2662}{735}.
\]
Let $L_h=\operatorname{Lip}(h)$. The constants in the table satisfy
$c_0\geqslant11p(2662/735)$ and $c_1\geqslant2p(110/63)$.
Since $11\mid N$, Lemma~\ref{prelim:operator-cert} shows that it is enough
to check at every node that
\begin{equation}\label{eleven:node}
 \Lambda h(j/N)-(T_ph)(j/N)
 >\frac{c_0+c_1L_h}{8N^2}.
\end{equation}

Upper bounds for the nodal weights are computed directly from
\[
 f_d(t)^2=\frac{10+2\sum_{r=1}^{10}c_r\cos(2\pi rt)}{100},
 \qquad c_r=\#\{a\in D_d:a+r\in D_d\}.
\]
The program \path{exact_trig.py} uses fixed-point denominator $2^{90}$,
Machin's formula, the degree-$40$ Taylor polynomial, and remainder bound
$(22/7)^{42}/42!$ to produce directed cosine intervals. An upward-rounded
square root bounds the amplitude, and \eqref{prelim:integer-power}
bounds its rational $p$th power. These weights also have denominator $D$.
If $m=\max_j|v_{j+1}-v_j|$ and $\Lambda=u/\ell$ with coprime positive integers
$u,\ell$, clearing denominators in \eqref{eleven:node} gives the integer
inequality
\[
 \begin{split}
 \mathcal M_j={}&8N^2\left[
  11Du\,v_j-\ell\sum_{i=0}^{10}w_{iN+j}
       \bigl((11-r)v_{iN/11+\lfloor j/11\rfloor}
               +r v_{iN/11+\lfloor j/11\rfloor+1}\bigr)\right]\\
 &-11D\ell(c_0D+c_1Nm)>0,\qquad r=j\bmod11.
 \end{split}
\]
The subscripts of $v_j$ are taken modulo $N$. The index $iN+j$ of
the directed weight $w_{iN+j}$ runs over the full $11N$-point grid. For each moment, the smallest integer
margin over the six digit representatives and all nodes is
\begin{center}
\begin{tabular}{cl}
\toprule
$p$ & $\min_{d,j}\mathcal M_j$\\
\midrule
$1$ & $39943001465673152637111041290240$\\
$235/154$ & $191058315489861615435892981836800$\\
$3/2$ & $12282688731347493451960444195840$\\
\bottomrule
\end{tabular}
\end{center}
The program \path{supplement/base11/verify_moments.py} computes the trigonometric bounds
from the digit set and checks every weight, nodal interpolation, and
margin above using integer and rational arithmetic. Thus
\eqref{prelim:semiconvex} gives $T_ph\leqslant\Lambda h$, and
Lemma~\ref{prelim:supersolution} proves the bounds for every $k,x$.
\end{proof}

\subsection{One-digit affine contraction}

We use the quantities $P_i$, $H$, $\rho$ and $K_2$ defined in
Section~\ref{dec:inputs}. The path estimate itself was proved in
Lemma~\ref{dec:hellinger}; it remains to bound its one-digit coefficient.

\subsubsection{Base ten}

\begin{lemma}[One-digit affine contraction in base ten]
\label{dec:decimalhellinger}
For base ten, the following upper bounds for $K_2$ and admissible values
of $\delta$ satisfy $K_2\le10^{-3\delta}$:
\[
\begin{array}{c|ccccc}
 \text{representative }d&0&1&2&3&4\\ \hline
 \text{upper bound for }K_2&19/20&967/1000&239/250&481/500&241/250\\
 \text{admissible }\delta&7/1000&3/625&3/500&11/2000&53/10000
\end{array}
\]
The reflected digit $9-d$ has the same bounds. In particular,
\eqref{dec:groups} is valid for every decimal digit.
\end{lemma}
\begin{proof}
Use the centered amplitudes
\[
 A_i(x)=90^{-1/2}\sum_{\substack{0\le j<10\\j\ne d}}
             e((j-9/2)(i+x)/10),\qquad 0\le i<10.
\]
Discrete Parseval gives $\|A\|_2=1$ and
\[
 \|A'\|_2^2\le(2\pi/10)^2\frac{55}{6}<4,\qquad
 \|A''\|_2^2\le(2\pi/10)^4\frac{3223}{24}<25.
\]
Lemma~\ref{prelim:hellinger-curvature} therefore gives
$H_{xx}\ge-18$, $H_{EE}\ge-5$ and $(\rho^2)''\ge-10$.
For $G_a(x,t)=\sum_iP_i(x)\rho(ai+t)^2$, both coordinate second
derivatives are bounded below by $-10$.

Take $M=200$. At each $E=k/M$, maximize the directed upper bounds
for $H(j/M,E)$ and add $18/(8M^2)$. This bounds the continuous
supremum defining $\rho(E)$ by Lemma~\ref{prelim:grid}.
Evaluate $G_a$ on the product grid in $(x,t)$ and add $28/(8M^2)$;
this exceeds the required total curvature correction $20/(8M^2)$.
The same rational trigonometric enclosures used in
Lemma~\ref{dec:decimalcert}, with upward-rounded probabilities and
square roots of denominator $2^{28}$, give the displayed bounds.
The grid covers $0\le x\le1$, $0\le t\le10$ and $1\le a\le5$.
Evenness and period ten give $G_{10-a}(x,t)=G_a(x,-t)$, so the
supremum covers all nine nonzero residue classes, not just the units.
This is needed in Lemma~\ref{dec:hellinger}: the first nonzero digit
of $(\lambda-1)/10^m$ need not be coprime to ten.

All these finite operations and their directed roundings are specified
in the supplementary verification program.
Finally, if an entry is $K_2\le u/v$ and $\delta=r/s$, the comparison
with $10^{-3\delta}$ is checked by the integer inequality
$u^s10^{3r}\le v^s$. Each of the five pairs in the table satisfies it.
\end{proof}

\subsubsection{Base eleven}

\begin{lemma}[Affine contraction in base eleven]\label{eleven:hellinger}
For every one of the eleven forbidden digits,
\begin{equation}\label{eleven:K2}
 K_2\leqslant\frac{24}{25}\leqslant11^{-3/250}.
\end{equation}
Consequently, in the notation of Lemma~\ref{dec:hellinger},
\[
 \mathcal H_L\leqslant(B/11^m)^{-1/250},\qquad B=11^L,
\]
uniformly in the initial states, intercepts, and unit slopes satisfying
$\lambda\equiv1\pmod {11}$.
\end{lemma}
\begin{proof}
Let
$A_i(x)=110^{-1/2}\sum_{a\in D_d}\mathrm e((a-5)(i+x)/11)$.
Discrete Parseval and the centered exponents give
\[
 \|A\|_2=1,\qquad
 \|A'\|_2^2\leqslant(2\pi/11)^2\frac{110}{10}<4,
 \qquad
 \|A''\|_2^2\leqslant(2\pi/11)^4\frac{1958}{10}<25.
\]
Here $1958=\sum_{a=0}^{10}(a-5)^4$. Apply
Lemma~\ref{prelim:hellinger-curvature} with $a=2$, $c=5$. It gives
\[
 H_{xx}\geqslant-18,\qquad H_{EE}\geqslant-5.
\]
For $G_a(x,t)=\sum_iP_i(x)\rho(ai+t)^2$, the same lemma gives
\[
 (G_a)_{xx}\geqslant-10,\qquad (G_a)_{tt}\geqslant-10.
\]

Use a grid with $M=80$ and upward-rounded probabilities and amplitudes
with denominator $D_0=2^{28}$. First maximize over $x=j/M$ and add
$18/(8M^2)$ to bound $\rho(E)$ at each $E=k/M$. Next evaluate $G_a$ on the
full product grid in $(x,t)$ and add $20/(8M^2)$, as justified by
Lemma~\ref{prelim:grid}. Since $\rho$ is even
and $11$-periodic, it suffices to take $a=1,\ldots,5$; this reduction covers
every slope class. The program \path{supplement/base11/verify_hellinger.py}
computes these directed trigonometric intervals, square roots, and finite
suprema. The resulting upper bounds,
rounded further upward to denominator $10^4$, are
\begin{center}
\begin{tabular}{ccccccc}
\toprule
Digit representative $d$ & $0$ & $1$ & $2$ & $3$ & $4$ & $5$\\
\midrule
Upper bound for $K_2$ & $.9284$ & $.9378$ & $.9421$ & $.9476$ & $.9465$ & $.9373$\\
\bottomrule
\end{tabular}
\end{center}
Every entry is less than $24/25$. The final power comparison is exactly
$24^{250}11^3\leqslant25^{250}$. The two-ended path H\"older inequality is
proved in Lemma~\ref{dec:hellinger}; substituting $t=2$ gives the asserted
correlation contraction.
\end{proof}

\subsection{Scalar comparisons and the shift input}

For base ten, the inequalities
\[
 9^{500}>10^{477},\qquad 9^{100}<10^{99},\qquad
 9^{21}>10^{20},\qquad 730^{5000}<10^3\,729^{5000}
\]
give $1/100<1-\alpha<23/500$, $\alpha>20/21$, and
$\log_{10}(1+9^{-3})<3/5000$.
For base eleven, the corresponding integer comparisons are
\[
 10^{500}>11^{477},\qquad 10^{100}<11^{99},\qquad
 10^{21}>11^{20},\qquad 1001^{5000}<11^3\,1000^{5000}.
\]
Thus in this base as well,
\[
 1/100<1-\alpha<23/500,\qquad \alpha/2>10/21,\qquad
 \log_{11}(1+10^{-3})<3/5000.
\]
The even-moment identities \eqref{prelim:even} consequently give
\eqref{dec:evenmoments} in both bases.
The comparison needed for maximal Type I follows from
$\alpha>477/500$ and
\[
 \frac{(50/77)\alpha}{\alpha+(50/77)(1-\alpha)}
 >\frac{23850}{37879}>\frac{629}{1000}.
\]
The programs \path{supplement/base10/verify_base10.py} and
\path{supplement/base11/verify_base11.py} check these exact comparisons.

The moment and correlation bounds above are numerical; the one-sided
shift bound is not. Specializing the zero classification and product
comparison of Appendix~\ref{app:shifts} gives the following form.

\begin{corollary}[Base-eleven zeros and shift comparison]
\label{eleven:zeros}\label{eleven:shift}
The unit-circle zeros of $P_{11,d}$ are precisely
\[
 z^{10}=1,\ z^d=1,\ z\ne1;
 \qquad
 z^{12}=1,\ z^{d+1}=-1,\ z\ne1.
\]
They are all simple and have order coprime to $11$. For every fixed $K>0$,
if $Q\leqslant q$ are powers of $11$, $A\in\mathbb Z$,
$|\delta|\leqslant1$, and $F_Q((A+\delta)/Q)\geqslant q^{-K}$, then
\begin{equation}\label{eleven:distortion}
 F_Q((A+\delta)/Q)
 \leqslant\exp\!\bigl(O_{d,K}(\sqrt{\log q})\bigr)F_Q(A/Q).
\end{equation}
\end{corollary}
\begin{proof}
Apply Lemmas~\ref{prelim:zeros} and \ref{prelim:shift} with $b=11$.
The exceptional double-zero case cannot occur, since $11\equiv3\pmod4$.
\end{proof}

\begin{proof}[Proof of Proposition~\ref{dec:verified-inputs}]
Lemmas~\ref{dec:decimalcert} and \ref{eleven:moments} prove the three
uniform shifted moment estimates. Lemma~\ref{dec:decimalhellinger}
gives the required value of $\delta$ for the decimal central digits
and a value at least $1/250$ for every other decimal digit.
Lemma~\ref{eleven:hellinger} gives $\delta=1/250$ for every base-eleven
digit. These are precisely the coupled inputs in \eqref{dec:groups}.
The scalar comparisons above establish the additional elementary
inequalities used in the common argument.
\end{proof}

\section{The base-nine Fourier certificates}
\label{app:small-certificates}

The certificates provide positive supersolutions for the two transfer
operators. Interpolation extends the nodal inequalities to every
point of the circle. Appendix~\ref{app:reproduce} gives the exact
verification commands.

\subsection{Continuous Fourier moments in base nine}
\label{app:nine-moments}

Use the notation and parameters of Proposition~\ref{nine:fourier}.

\begin{proof}[Proof of Proposition~\ref{nine:fourier}]
Let
$f_d(t)=8^{-1}|\sum_{a\ne d,\,0\le a\le8}\mathrm e(at)|$,
and define a positive operator on continuous functions on the circle by
\[
 (\mathcal T_{d,s}h)(x)=\frac19\sum_{j=0}^8
 f_d\left(\frac{x+j}{9}\right)^s h\left(\frac{x+j}{9}\right).
\]
Expanding the inverse branches successively gives the exact identity
\begin{equation}\label{nine:transfer}
 \mathcal T_{d,s}^k1(x)=9^{-k}
 \sum_{a<9^k}F_{9^k}\left(\frac{a+x}{9^k}\right)^s.
\end{equation}
If $h>0$ and $\mathcal T_{d,s}h\le\lambda h$, then
$\mathcal T_{d,s}^k1\le\lambda^kh/\min h$.
Integration over $x\in[0,1]$ also gives the continuous moment.
Thus one fixed continuous supersolution controls every word length and every shift.

The base-nine certificate data in the mathematical supplement specify
ten positive rational vectors, one for each $d=0,\ldots,4$ and
$s=1,3/2$. Each defines a periodic piecewise linear function through
$8192$ nodes with denominator $10^9$.
The arrays are recorded in the file \texttt{witnesses.py}
in \texttt{supplement/base9}.
The common target constants are
\begin{equation}\label{nine:lambdas}
 \lambda_1=\frac{2477}{10000},\qquad
 \lambda_{3/2}=\frac{314229}{2000000}.
\end{equation}
The following rational estimates certify the supersolution inequality
on the whole circle.

For $1\le r\le8$, the number of ordered pairs of allowed digits with difference $r$ is
$c_r(d)=9-r-\mathbf1_{d+r\le8}-\mathbf1_{d-r\ge0}\ge0$.
Consequently,
\begin{equation}\label{nine:square-factor}
 f_d(t)^2=\frac{8+2\sum_{r=1}^8c_r(d)\cos(2\pi rt)}{64}.
\end{equation}
Take the arithmetic scale $S=10^{24}$.
Machin's identity
$\pi=16\arctan(1/5)-4\arctan(1/239)$,
with $50$ and $15$ terms of the respective alternating series, gives
\[
 \frac{3141592653589793238462643}{S}<\pi<
 \frac{3141592653589793238462644}{S}.
\]
Use periodicity and evenness to reduce the angles to $[0,\pi]$.
Evaluate the degree-$48$ Taylor polynomial for cosine with outward integer
rounding at each step, and bound the remainder by $4^{50}/50!$.
The cosine coefficients in~\eqref{nine:square-factor} are nonnegative,
so this gives an upper rational enclosure $v/S\in[0,1]$ for $f_d(t)^2$.
The corresponding first- and $3/2$-moment factors are bounded above by
\[
 S^{-1}\left\lceil\sqrt{vS}\right\rceil,\qquad
 S^{-1}\left\lceil(v^3S)^{1/4}\right\rceil
\]
respectively. All roots are computed with integer square roots and checked
again by integer powers.

Let the verification grid have $n$ nodes, with node values $h_i/G$, and put
\[
 M_h=\max_i h_i/G,\qquad
 L_h=\frac nG\max_i|h_{i+1}-h_i|.
\]
At $x=i/n$, the interpolated value on the inverse branch
$t=(i+jn)/(9n)$ is exactly
$((9-r)h_v+rh_{v+1})/(9G)$, where $i+jn=9v+r$.
Multiply by the root enclosures above and divide by the transfer average's
factor $9$. This produces an upper bound for
$\mathcal T_{d,s}h(i/n)$ with denominator $81SG$.

Centering the digit phases at $4$ gives
$\operatorname{Lip}(f_d)\le5\pi<110/7$.
Hence the circular Lipschitz constant of the residual is at most
\begin{equation}\label{nine:lip}
 L_*=\frac{s(110/7)M_h+L_h}{9}+\lambda_sL_h.
\end{equation}
Using exact fractions, the verifier checks
\begin{equation}\label{nine:residual}
 \max_i\{\mathcal T_{d,s}h(i/n)-\lambda_sh(i/n)\}
       +\frac{L_*}{2n}<0.
\end{equation}
Nine certificates use $n=8192$.
For $d=3,s=3/2$, each segment of the same linear interpolant is subdivided
exactly into sixteen equal parts, giving $n=131072$.
All ten checks give strict inequalities.
Every real point is within $1/(2n)$ of a grid point, so
\eqref{nine:residual} proves the supersolution inequality on the whole circle.
Every computation uses only fixed integer arrays and integer or rational arithmetic.

Finally, the verifier checks the integer inequalities
\[
 2477^{200}9^{127}<10000^{200},\qquad
 314229^{40000}9^{33693}<2000000^{40000}.
\]
Together with~\eqref{nine:transfer}, these give the first-moment discrete
exponent $73/200$ and the $3/2$-moment discrete exponent $6307/40000$.
The reflection identity $f_{8-d}=f_d$ covers the other four digits.
Applying H\"older's inequality on the same grid,
\[
 \sum F^{4/3}\le\left(\sum F\right)^{1/3}
                     \left(\sum F^{3/2}\right)^{2/3},
\]
gives $g_{4/3}=((73/200)+2(6307/40000))/3=13607/60000$.
\end{proof}

\section{Computational certificates}\label{app:reproduce}

\subsection{Exact data and acceptance criteria}

The electronic supplementary material contains the fixed rational data and
verification programs for the finite inequalities used in the proof.
The data consist of positive vectors, nodal upper bounds, and nonnegative
coefficients for dual and Farkas identities. Each acceptance criterion uses
only integer or rational arithmetic. Appendix~\ref{app:continuous} establishes the continuous certification
criteria. Their applications to the digital inputs are proved in
Appendices~\ref{app:exact-inputs} and~\ref{app:small-certificates}; the
exponent systems and sieve integrals are derived in the corresponding
analytic sections.

\begin{center}
\begin{tabular}{>{\raggedright\arraybackslash}p{.12\textwidth}
 >{\raggedright\arraybackslash}p{.40\textwidth}
 >{\raggedright\arraybackslash}p{.38\textwidth}}
\toprule
Base & Certificate objects & Result verified\\
\midrule
$6$ & Full- and prime-power-grid moments; mixed-period inequalities; simplex sums
 & The cut and mask bounds, and a sieve loss below $.986$\\
$7$ & Shifted moments and rational frequency inequalities
 & The single-cut comparison parameters; the shared sieve window has loss below $.934$\\
$8$ & Two shifted moment certificates and rational exponent inequalities
 & Moment and exponent inequalities for all eight digits; the triangle loss is below $.186012$\\
$9$ & Ten continuous moments; directed upper sums on full prime-factor simplices
 & Positive-tree loss less than $.834204$\\
$10$ & Fifteen continuous moments; five correlation profiles; $188$ rational systems
 & Positive continuous residuals and strictly negative exponent bounds\\
$11$ & Eighteen continuous moments; six correlation profiles; scalar parameter comparisons
 & The moments in Lemma~\ref{eleven:moments} and $K_2\leqslant24/25$\\
$10$, two omissions & Shifted decimal and prime-power moments,
signed Gram and Mellin exponents, and ordered simplex sums
 & The grouped comparison parameters and a sieve loss below $.934$\\
\bottomrule
\end{tabular}
\end{center}

\paragraph{Continuous positive operators.}
A certificate uses a positive periodic linear interpolant \(h\)
with rational node values. The node values may be stored in the
supplement or generated by iteration before the residual is checked.
Starting from the digit set, the verification program computes directed
trigonometric and rational-power enclosures. It then checks the positivity
of $h$, each nodal residual, and the continuous interpolation correction
from Appendix~\ref{app:continuous}. The comparison of the rational
supersolution constant with the claimed spectral growth is checked by
integer powers. These checks establish $T_ph\leqslant\Lambda h$
on the whole circle, and Lemma~\ref{prelim:supersolution} gives the
moment bound for every word length.

\paragraph{Rational systems of exponents.}
The forward and reverse certificates in \texttt{supplement/base10}
contain the exact systems and
their nonnegative dual or Farkas coefficients.
The verifier reconstructs every matrix from the parameters, constraints,
and branches in Sections~\ref{dec:forward}, \ref{dec:reverse},
and~\ref{dec:ranktwo} and checks
equality with the encoded system.
The six denominator coordinates \texttt{(a,b,c,d,e,f)} in the data
correspond, in that order, to $(a,c,h,D,E,G)$ in
Section~\ref{dec:ranktwo}; the coordinate order is part of the
certificate specification.
For a feasible system $Az\leqslant b$, a certificate for the target
$c^tz+c_0$ consists of $y\geqslant0$ satisfying
\[
 y^tA=c^t,\qquad y^tb+c_0<0.
\]
It follows that $c^tz+c_0\leqslant y^tb+c_0<0$ throughout the
real feasible polyhedron. An empty branch is certified by
$y\geqslant0$, $y^tA=0$, and $y^tb<0$.
Every matrix identity and sign comparison is exact.
The same $188$ systems apply in base $11$, whose digital inputs are
verified in Proposition~\ref{eleven:inputs}.

\paragraph{Sieve integrals.}
For bases six, seven and nine and for two decimal omissions, one
program enumerates the ordered simplex cells of
Proposition~\ref{nine:certificate}. Their volumes are given by the
Eulerian slice formula, with the factorial factors for equal grid
coordinates. A cell can be omitted if its upper multiplicity is zero
or if one subproduct is evaluable throughout the cell. Each remaining
cell retains the full marked
multiplicity, and every arithmetic rounding is upward.
These are the finite upper sums appearing in the sieve-loss bounds.

\subsection{Running the verification}
From the directory containing this manuscript, run
\begin{verbatim}
python3 -B -S supplement/verify.py
\end{verbatim}
Use a writable extracted copy, Python~3 with its standard library and
a C++17 compiler supporting unsigned 128-bit integers. Assertions must
be enabled. The nodal and dual witnesses are provided as readable
constants in the verification sources; no external tables or network
access are needed. The common simplex program takes the window, mesh
and ranks as inputs. For base eight the loss is checked instead by the
rational series in Section~\ref{cost:octal}.

\clearpage
\phantomsection
\addcontentsline{toc}{section}{References}
\begin{thebibliography}{99}
\small\setlength{\itemsep}{3pt}\setlength{\parsep}{0pt}
\raggedright

\bibitem{BanksShparlinski2019}
W. D. Banks and I. E. Shparlinski,
\emph{Bounds on short character sums and \(L\)-functions with characters
to a powerful modulus}, J. Anal. Math. \textbf{139} (2019), 239--263.
\href{https://doi.org/10.1007/s11854-019-0060-4}{doi:10.1007/s11854-019-0060-4}.
For (3.3) and the sentence following it see \href{https://arxiv.org/abs/1605.07553v2}{arXiv:1605.07553v2} (preprint title: \emph{Bounds on short character sums and $L$-functions for characters with a smooth modulus}).

\bibitem{ChenGuptaLi2026}
B. Chen, V. Gupta, and Y. C. Li,
\emph{Large value estimates for Dirichlet polynomials with characters and
zero density of Dirichlet $L$-functions},
\href{https://arxiv.org/abs/2507.08296v2}{arXiv:2507.08296v2}, revised 27 July 2026.

\bibitem{DartygeMauduit2000}
C. Dartyge and C. Mauduit, \emph{Nombres presque premiers dont
l'\'ecriture en base $r$ ne comporte pas certains chiffres},
J. Number Theory \textbf{81} (2000), 270--291.
\href{https://doi.org/10.1006/jnth.1999.2458}{Publisher}.

\bibitem{DartygeMauduit2001}
C. Dartyge and C. Mauduit, \emph{Ensembles de densit\'e nulle contenant
des entiers poss\'edant au plus deux facteurs premiers},
J. Number Theory \textbf{91} (2001), 230--255.
\href{https://doi.org/10.1006/jnth.2001.2681}{Publisher}.

\bibitem{ErdosMauduitSarkozy1998}
P. Erd\H{o}s, C. Mauduit and A. S\'ark\"ozy,
\emph{On arithmetic properties of integers with missing digits I:
distribution in residue classes}, J. Number Theory \textbf{70} (1998), 99--120.
\href{https://doi.org/10.1006/jnth.1998.2229}{Publisher}.

\bibitem{Ford2002ZeroFree}
K. Ford, \emph{Zero-free regions for the Riemann zeta function},
in \emph{Number Theory for the Millennium}, II,
A K Peters, Natick, MA, 2002, 25--56.
Corrected author version: \href{https://arxiv.org/abs/1910.08205v5}{arXiv:1910.08205v5}; Theorem~1 is unchanged.

\bibitem{FordMaynard2024}
K. Ford and J. Maynard,
\textit{On the theory of prime-producing sieves},
preprint (2024),
\href{https://arxiv.org/abs/2407.14368v1}{arXiv:2407.14368v1}.

\bibitem{FortiViola1973}
M. Forti and C. Viola, \emph{Density estimates for the zeros of
$L$-functions}, Acta Arith. \textbf{23} (1973), 379--391.
\href{https://doi.org/10.4064/aa-23-4-379-391}{Publisher}.

\bibitem{FriedlanderIwaniec2010}
J. Friedlander and H. Iwaniec, \emph{Opera de Cribro}, American
Mathematical Society Colloquium Publications, vol.~57, 2010.
\href{https://doi.org/10.1090/coll/057}{Publisher}.

\bibitem{Harman2007}
G. Harman,
\textit{Prime-Detecting Sieves},
London Mathematical Society Monographs, vol.~33,
Princeton University Press, Princeton, NJ, 2007.

\bibitem{Huxley}
M. N. Huxley, \emph{Large values of Dirichlet polynomials, III},
Acta Arith. \textbf{26} (1975), 435--444.
\href{https://doi.org/10.4064/aa-26-4-435-444}{doi:10.4064/aa-26-4-435-444}.

\bibitem{Iwaniec1974}
H. Iwaniec, \emph{On zeros of Dirichlet's \(L\) series},
Invent. Math. \textbf{23} (1974), 97--104.
\href{https://doi.org/10.1007/BF01405163}{doi:10.1007/BF01405163}.

\bibitem{IwaniecKowalski2004}
H. Iwaniec and E. Kowalski,
\textit{Analytic Number Theory},
American Mathematical Society Colloquium Publications, vol.~53,
American Mathematical Society, Providence, RI, 2004.
\href{https://doi.org/10.1090/coll/053}{doi:10.1090/coll/053}.

\bibitem{Karwatowski2022}
F. Karwatowski,
\textit{Primes with one excluded digit},
Acta Arith. \textbf{202} (2022), 105--121.
\href{https://doi.org/10.4064/aa191002-26-8}{doi:10.4064/aa191002-26-8}.

\bibitem{KarwatowskiBaseNine}
F. Karwatowski,
\textit{Digits of primes in base $b=9$},
preprint, 17~pp.,
\href{https://www.math.hhu.de/fileadmin/redaktion/Fakultaeten/Mathematisch-Naturwissenschaftliche_Fakultaet/Mathematik/20_Institut-Lehrstuehle/5_Algebra_und_Zahlentheorie/Karwatowski/Digits_of_primes_in_base_b_9.pdf}{available from the author's webpage}.

\bibitem{KarwatowskiDMV2019}
F. Karwatowski, \emph{Primzahlen mit einer ausgeschlossenen Ziffer},
DMV-Jahrestagung, Karlsruhe (2019), conference abstract.
\href{https://dmv2019.math.kit.edu/wp-content/uploads/2019/07/Karwatowski_PrimAusgeschZiff.pdf}{Author's abstract}.

\bibitem{Koukoulopoulos2019}
D. Koukoulopoulos,
\textit{The Distribution of Prime Numbers},
Graduate Studies in Mathematics, vol.~203,
American Mathematical Society, Providence, RI, 2019.
\href{https://doi.org/10.1090/gsm/203}{Publisher}; \href{https://dms.umontreal.ca/~koukoulo/documents/publications/primes.pdf}{author's preliminary version}.

\bibitem{Maynard2019}
J. Maynard,
\textit{Primes with restricted digits},
Invent. Math. \textbf{217} (2019), 127--218.
\href{https://doi.org/10.1007/s00222-019-00865-6}{doi:10.1007/s00222-019-00865-6}.

\bibitem{Postnikov}
A. G. Postnikov, \emph{On the sum of characters with respect to a
modulus equal to a power of a prime number}, Izv. Akad. Nauk SSSR
Ser. Mat. \textbf{19} (1955), 11--16.
\href{https://www.mathnet.ru/eng/im3519}{Math-Net.Ru}.

\bibitem{ThornerZaman}
J. Thorner and A. Zaman,
\emph{Refinements to the prime number theorem for arithmetic progressions},
Math. Z. \textbf{306} (2024), Paper No.~54.
\href{https://doi.org/10.1007/s00209-023-03414-3}{Publisher}. \href{https://arxiv.org/abs/2108.10878v2}{arXiv:2108.10878}.

\end{thebibliography}

\end{document}

